% Options for packages loaded elsewhere
\PassOptionsToPackage{unicode}{hyperref}
\PassOptionsToPackage{hyphens}{url}
\PassOptionsToPackage{dvipsnames,svgnames,x11names}{xcolor}
\documentclass[
  english,
  11pt,
  reqno]{amsart}
\usepackage{xcolor}
\usepackage{amsmath,amssymb}
\setcounter{secnumdepth}{-\maxdimen} % remove section numbering
\usepackage{iftex}
\ifPDFTeX
  \usepackage[T1]{fontenc}
  \usepackage[utf8]{inputenc}
  \usepackage{textcomp} % provide euro and other symbols
\else % if luatex or xetex
  \usepackage{unicode-math} % this also loads fontspec
  \defaultfontfeatures{Scale=MatchLowercase}
  \defaultfontfeatures[\rmfamily]{Ligatures=TeX,Scale=1}
\fi
\usepackage{lmodern}
\ifPDFTeX\else
  % xetex/luatex font selection
\fi
% Use upquote if available, for straight quotes in verbatim environments
\IfFileExists{upquote.sty}{\usepackage{upquote}}{}
\IfFileExists{microtype.sty}{% use microtype if available
  \usepackage[]{microtype}
  \UseMicrotypeSet[protrusion]{basicmath} % disable protrusion for tt fonts
}{}
\usepackage{longtable,booktabs,array}
\newcounter{none} % for unnumbered tables
\usepackage{calc} % for calculating minipage widths
% Correct order of tables after \paragraph or \subparagraph
\usepackage{etoolbox}
\makeatletter
\patchcmd\longtable{\par}{\if@noskipsec\mbox{}\fi\par}{}{}
\makeatother
% Allow footnotes in longtable head/foot
\IfFileExists{footnotehyper.sty}{\usepackage{footnotehyper}}{\usepackage{footnote}}
\makesavenoteenv{longtable}
\ifLuaTeX
\usepackage[bidi=basic,shorthands=off]{babel}
\else
\usepackage[bidi=default,shorthands=off]{babel}
\fi
\ifLuaTeX
  \usepackage{selnolig} % disable illegal ligatures
\fi
\setlength{\emergencystretch}{3em} % prevent overfull lines
\providecommand{\tightlist}{%
  \setlength{\itemsep}{0pt}\setlength{\parskip}{0pt}}
% Portable additions for the arXiv-oriented profile.
% Keep packages within the standard TeX Live distribution.
\usepackage{mathtools}
\usepackage{amssymb}
\usepackage{amsthm}
\usepackage{booktabs}
\usepackage{listings}
\usepackage{microtype}
\usepackage{centernot}
\usepackage{newunicodechar}

\ifPDFTeX
  % unicode-math supplies mathscr on Xe/LuaTeX; pdfTeX needs its own font.
  \usepackage{mathrsfs}
  % Explicit, meaning-preserving mappings only. Other unsupported input errors
  % are fatal; never discard characters or guess a different math expression.
  \newunicodechar{α}{\ensuremath{\alpha}}
  \newunicodechar{β}{\ensuremath{\beta}}
  \newunicodechar{γ}{\ensuremath{\gamma}}
  \newunicodechar{δ}{\ensuremath{\delta}}
  \newunicodechar{ε}{\ensuremath{\varepsilon}}
  \newunicodechar{θ}{\ensuremath{\theta}}
  \newunicodechar{λ}{\ensuremath{\lambda}}
  \newunicodechar{μ}{\ensuremath{\mu}}
  \newunicodechar{π}{\ensuremath{\pi}}
  \newunicodechar{σ}{\ensuremath{\sigma}}
  \newunicodechar{ω}{\ensuremath{\omega}}
  \newunicodechar{ℂ}{\ensuremath{\mathbb{C}}}
  \newunicodechar{ℕ}{\ensuremath{\mathbb{N}}}
  \newunicodechar{ℚ}{\ensuremath{\mathbb{Q}}}
  \newunicodechar{ℝ}{\ensuremath{\mathbb{R}}}
  \newunicodechar{ℤ}{\ensuremath{\mathbb{Z}}}
  \newunicodechar{∞}{\ensuremath{\infty}}
  \newunicodechar{∈}{\ensuremath{\in}}
  \newunicodechar{∅}{\ensuremath{\emptyset}}
  \newunicodechar{⊂}{\ensuremath{\subset}}
  \newunicodechar{⊆}{\ensuremath{\subseteq}}
  \newunicodechar{∩}{\ensuremath{\cap}}
  \newunicodechar{∪}{\ensuremath{\cup}}
  \newunicodechar{≤}{\ensuremath{\leq}}
  \newunicodechar{≥}{\ensuremath{\geq}}
  \newunicodechar{≠}{\ensuremath{\neq}}
  \newunicodechar{→}{\ensuremath{\rightarrow}}
  \newunicodechar{↦}{\ensuremath{\mapsto}}
  \newunicodechar{⇒}{\ensuremath{\Rightarrow}}
  \newunicodechar{⇔}{\ensuremath{\Leftrightarrow}}
  \newunicodechar{−}{\ensuremath{-}}
\fi
\newunicodechar{□}{\ensuremath{\square}}

% Portable negated hooked arrow used by several manuscripts.
\providecommand{\nhookrightarrow}{\centernot\hookrightarrow}

% U+220E END OF PROOF: render a small hollow square at the right margin.
\newcommand{\rightqedsquare}{%
  \ifhmode
    \unskip\nobreak\hfill\ensuremath{\square}%
  \else
    \noindent\hfill\ensuremath{\square}%
  \fi
}
\newunicodechar{∎}{\rightqedsquare}

% Reduce avoidable overfull lines without changing the source manuscript.
\setlength{\emergencystretch}{2em}
% Match amsart's theorem preskip while keeping Markdown statement bodies roman.
\makeatletter
\newcommand{\pandocstatementheading}[1]{%
  \par\addpenalty{\@beginparpenalty}%
  \addvspace{.5\baselineskip plus .2\baselineskip minus .2\baselineskip}%
  {\normalfont\bfseries #1\par}%
  \nobreak\@afterindenttrue\@afterheading
}
\makeatother
% Use the AMS section face at abstract-body size.
% The native abstract box, margins and body font remain owned by amsart.
\usepackage{etoolbox}
\newcommand{\pandocabstractheading}{%
  \par\begingroup
    \section*{\abstractname}%
  \endgroup
  \nobreak\noindent\ignorespaces
}
\makeatletter
\patchcmd{\abstract}
  {\item[\hskip\labelsep\scshape\abstractname.]}
  {\item[]\pandocabstractheading}
  {}%
  {\PackageError{studio-ams-abstract}
     {Unsupported amsart abstract definition}
     {Check the installed class before updating the abstract-heading patch.}}
\makeatother
% Common LaTeX macros for CCG papers
% Keep this minimal and portable across engines.

% Math shortcuts
\newcommand{\N}{\mathbb{N}}
\newcommand{\Z}{\mathbb{Z}}
\newcommand{\Q}{\mathbb{Q}}
\newcommand{\R}{\mathbb{R}}
\newcommand{\C}{\mathbb{C}}
\newcommand{\eps}{\varepsilon}

% Theorem-like environments (lightweight)
% If you need full theoremstyles, consider amsthm or thmtools here.
% \usepackage{amsthm}   % already available via TeX Live if needed
% \newtheorem{theorem}{Theorem}
% \newtheorem{lemma}{Lemma}

% Code listings default (avoid minted to keep shell-escape off)
\lstset{basicstyle=\ttfamily\small,breaklines=true,columns=fullflexible}

\usepackage{bookmark}
\IfFileExists{xurl.sty}{\usepackage{xurl}}{} % add URL line breaks if available
\urlstyle{same}
\hypersetup{
  pdftitle={A two-chirality twistor correspondence for smooth Lorentzian Ricci-flat germs},
  pdfauthor={Lior Isthmus},
  pdflang={en},
  colorlinks=true,
  linkcolor={blue},
  filecolor={Maroon},
  citecolor={blue},
  urlcolor={blue},
  pdfcreator={LaTeX via pandoc}}

\makeatletter
\let\originalsettitle\@settitle
\def\@settitle{\begingroup\let\uppercasenonmath\@gobble\originalsettitle\endgroup}
\makeatother
\title[A two-chirality twistor correspondence]{A two-chirality twistor correspondence for smooth Lorentzian Ricci-flat germs}
\author[Lior Isthmus]{Lior Isthmus\textsuperscript{1}}
\date{2026-11-19}
\setlength{\emergencystretch}{4em}

\usepackage{accsupp,needspace}
\begin{document}
\begin{abstract}
\phantomsection\label{sec:abstract}
We give a local twistor presentation of each smooth pointed, oriented and time-oriented Lorentzian Ricci-flat four-germ with spin, at its physical scale. The two inputs are ordinary cohomology classes of weights \(2\) and \(-6\) for the flat hypersurface twistor CR structure on \(\mathbb R^3\times\mathbb{CP}^1\), whose projective line is the sphere of spin directions. In the selected weighted domain of the growth-controlled coefficient complexes, the coupled reality condition gives real transverse trace-free seeds for the metric and momentum. Starting with encoded test tensors, we take three successive normalized negative-chirality contact responses, recovering a tensor at each step. Contractions of the resulting tensors give the flat scalar and York inverses for the constraints. Initialized nonlinear transport yields maximal vacuum data and both spin gradient operators. Classical Cauchy evolution develops these data. Conversely, we extend normalized maximal-disk data to weighted vacuum data on \BeginAccSupp{method=hex,unicode,ActualText=211D00B3002C}\(\mathbb R^3\),\EndAccSupp{} with exact agreement on a smaller disk, and decode the extension. The inverse gives a comparison class of presentations, with smooth local selections in finite-dimensional families and effective descent. A covariant wave kernel, determined modulo jointly smooth kernels, transports the degree-one evaluation section on each sky and its point connection between slices. Every pointed physical isometry with a specified actual spin lift has a unique full source lift modulo normalized presentation cells. On analytic germs we compare the wave source with the calibrated operator source of the analytic projective twistor correspondence. After real restriction and complete formal enrichment, the initialized comparison preserves coupled reality.
\end{abstract}
\maketitle
\footnotetext[1]{Email: \href{mailto:lioristhmus@proton.me}{\nolinkurl{lioristhmus@proton.me}}. ORCID: \href{https://orcid.org/0009-0004-7908-9876}{0009-0004-7908-9876}.}
\setcounter{footnote}{1}


{
\setcounter{tocdepth}{1}
\clearpage
\tableofcontents
}
\markboth{\MakeUppercase{Lior Isthmus}}{\MakeUppercase{A two-chirality twistor correspondence}}

\section{1. The local two-class
correspondence}\label{the-local-two-class-correspondence}\label{sec:introduction}

Penrose's nonlinear graviton construction encodes complex vacuum metrics
with one Weyl chirality in the holomorphic geometry of a deformed
twistor space (\hyperref[ref-Penrose1976]{Penrose 1976}). In his
helicity convention the construction describes left-handed gravitational
fields. The googly problem asks how the same twistor convention can also
describe interacting right-handed fields. Penrose proposed adjoining
twistor differentiation operators to holomorphic functions to form a
noncommutative twistor algebra (\hyperref[ref-Penrose2015]{Penrose
2015}).

Twistor actions also describe interactions between the two chiralities.
\hyperref[ref-AdamoMason2014]{Adamo and Mason} computed Einstein MHV
amplitudes from the conformal gravity action and proposed an Einstein
twistor action. For Euclidean gravity at \(\Lambda=0\),
\hyperref[ref-Sharma2021]{Sharma} showed that his twistor action reduces
to the Plebański action. Here the inputs are an ordinary pair in a
specified weighted domain, and the outputs are smooth local Lorentzian
Ricci-flat germs with comparisons over isometries carrying specified
spin lifts.

A second approach starts on a spacelike hypersurface. LeBrun recovered
the conformal geometry of a Riemannian three-manifold from the CR
structure on its sphere bundle (\hyperref[ref-LeBrun1984]{LeBrun 1984}).
Mason explained how hypersurface twistors also encode the trace-free
second fundamental form and distinguished these conformal data from the
scale and mean curvature (\hyperref[ref-Mason1989]{Mason 1989}). Mason
and Penrose subsequently investigated a twistorial formulation of the
full vacuum equations (\hyperref[ref-MasonPenrose1989]{Mason and Penrose
1989}).

In the analytic projective twistor correspondence
(\hyperref[ref-IsthmusLocal2026]{Isthmus 2026a}), two classes encode the
two chiral curvature records, and analytic continuation recovers an
Einstein germ. A smooth germ can have a nonzero field whose derivatives
vanish at a point. For a smooth local correspondence the data must
therefore describe a field on a neighborhood. We work at \(\Lambda=0\)
and retain the physical metric, both orientations and the actual spin
structure.

Two localization issues arise. The decoder acts on whole weighted data,
and a cutoff of vacuum data introduces a constraint error. Appendix B
uses the maximal-data extension of \hyperref[ref-Czimek2018]{Czimek} and
its explicit linear right inverses to obtain a vacuum extension that
agrees with the prescribed data on a smaller disk. Two displays of the
same germ may also use different slices.
\hyperref[ref-Radzikowski1996]{Radzikowski} characterized the Hadamard
condition by the wavefront set on a fixed spacetime, and
\hyperref[ref-gerard-wrochna2014]{Gérard and Wrochna} gave a
pseudodifferential factorization of the wave equation. Section 6 uses
the positive-frequency branch and proves uniqueness of the kernel modulo
joint smoothing in families. It transports the full degree-one
evaluation between slices, canceling the direction-dependent Doppler
factor already present for boosted flat hyperplanes.

On the fixed input space \(N_0=\mathbb R^3_x\times\mathbb{CP}^1\), each
spin direction determines a complex null line with homogeneous
representative \(z\in\mathbb C^3\). Write \(\mathscr E_m^\bullet\) for
the weight-\(m\) CR complex with growth-controlled coefficients. Its
first differential is \(d^0u=(\bar\partial u,Du)\), where
\(\bar\partial\) acts on the sphere and \(D=z\cdot\partial_x\). The
ordinary inputs lie in the first cohomology
\(H^1(\mathscr E_m^\bullet)\) at the two weights \(m=+2,-6\) fixed by
contour normalization. The coupled reality condition on the class pair
makes both decoded seeds real, so the metric and momentum seeds enter a
single vacuum development. In the diagram below the left category
records local displays of these pairs, the middle category admits
sources by independent geometric and operator equations, and the right
category records physical germs and isometries with their specified spin
lifts.

\pandocstatementheading{\leavevmode\phantomsection\label{thm:local-main-correspondence}Theorem 1.1 (Smooth two-class local correspondence).}

Let \(\mathsf{Vac}^{\mathrm{loc}}_0\) have the physical germs above as
objects and pointed physical isometries with their specified actual spin
lifts as arrows. There are equivalences

\phantomsection\label{eq:main-smooth-two-class-correspondence} \[
\mathsf{WavePair}^{\mathrm{loc}}_0
\xrightarrow{\ E\ }
\mathsf{WaveSrc}^{\mathrm{loc}}_0
\xrightarrow{\ R\ }
\mathsf{Vac}^{\mathrm{loc}}_0 .
\] The first domain consists of local displays of ordinary pairs as in
\hyperref[def:local-display-groupoids]{Definition 8.1}. All freely
specified functional data of a chart are contained in the ordinary pair
of \hyperref[def:online-effective-ordinary-cochain-domain]{Definition
3.11} up to its effective cochain arrows. The positive number \(r\),
retained neighborhood and markings specify the display. A local maximal
initial slice is part of the presentation. The middle domain is defined
by independent equations on actual geometric and spin fields and on
complete operator and module data in
\hyperref[def:online-complete-independent-wave-source]{Definition 7.7}.
The independent source uses a common causally convex, globally
hyperbolic wave neighborhood for its retained carrier hits and
connecting null segments, as specified in
\hyperref[def:online-wave-carrier-displays]{Definition 6.12}. Every
object of either of the last two categories is locally represented by a
display. For each physical arrow between recognized germs there is a
unique full source arrow over it, up to normalized presentation cells.
The inverse identifications are natural and obey the triangle
identities. The physical target includes flat and algebraically special
germs and imposes no asymptotic condition.

For smooth finite-dimensional families on common representatives with
smooth markings, a local selection has one positive normalization radius
and one buffered parent on retained compact parameter sets. This parent
is a larger common domain in which the retained patches lie a positive
distance from the boundary. For each stipulated finite derivative and
coefficient depth the selected output has a bound in finitely many input
seminorms on that parent and the positive margins. Supplied comparison
families have the corresponding estimates with their map and inverse
bounds. The local presentations and their full arrows admit effective
atlas descent over open covers of Hausdorff representatives with
effective actual spin descent.

On the analytic geometric locus, using the calibrated operator
correspondence of \hyperref[ref-IsthmusLocal2026]{Isthmus 2026a} (v1.1,
Definition H.1 and Theorem 9.10), there are structured comparisons in
both directions after real restriction and complete formal enrichment.
The analytic transport and descent results used for this clause are
specified in \hyperref[sec:analytic-wave-comparison]{Section 9}. These
comparisons carry the analytic chiral structures in their simultaneous
presentations, the optical and numerical scales, both associators,
actual root squares and base factors, the canonical finite-rank
structures, actual spin and the original coupled square cell. The
initialized framing compares the full degree-one evaluation section and
its point connection in every direction. The finite-order bounds there
apply to conditioned families on common analytic parents and their
retained real restrictions.

Different whole-space extensions of the same local data can present the
same pointed germ. \hyperref[def:local-display-groupoids]{Definition
8.1} therefore takes arrows between displays to be full comparisons of
their executed sources on a common local restriction. This makes \(E\)
fully faithful. Coverage of the physical germs follows from
\hyperref[thm:submission-vacuum-germ-coverage]{Theorem 5.7}, and full
faithfulness of \(R\) follows from the lifting theorem for physical
arrows with specified spin lifts
(\hyperref[thm:online-full-unmarked-source-hom]{Theorem 7.11}). The
proof of the combined statement concludes with
\hyperref[thm:combined-local-atlas]{Theorem 9.8}.

\textbf{Inputs and physical scale.} The two functional data are classes
in \(H^1(\mathscr E_2^\bullet)\oplus H^1(\mathscr E_{-6}^\bullet)\).
Membership uses actual primitives, the weighted linear root test and
explicit polynomial obstructions. Coupled reality is imposed on the two
inputs together. Smallness concerns the normalized whole input of the
chart. In physical variables the same display has \(h^{[r]}=r^2h\),
\(K^{[r]}=rK\) and \(\mathbb D_\pm^{[r]}=r^{-1}\mathbb D_\pm\). An
isometry between outputs preserves the metric at this restored scale,
even if the displays use different normalization radii. Classical
Einstein evolution develops the smooth vacuum data supplied by the pair.
The resulting metric fixes the wave kernel up to addition of a jointly
smooth kernel. An exact quantum state requires further data.

\textbf{Relation to the analytic correspondence.} The original
independently admitted algebroid contains an actual half-canonical
square, a framed local-system factor and factorially bounded analytic
operators (\hyperref[ref-IsthmusLocal2026]{Isthmus 2026a}, Definition
H.1). The smooth evaluation has degree one and moves with the physical
point. \hyperref[thm:analytic-structured-wave-comparison]{Theorem 9.7}
transports this evaluation while retaining the original degree-zero
calibration. Even when the metrics agree, a comparison must fix the
relative scale of the evaluation sections, respect their sign under
conjugation and compose coherently along physical arrows. In the
analytic global correspondence, complete complex skies and their point
spin maps glue, and null-path monodromy is recomputed under point
restriction (\hyperref[ref-IsthmusGlobal2026]{Isthmus 2026b}, Lemma
5.5). Here the bases are local smooth Lorentzian representatives at
\(\Lambda=0\) and the descended objects are atlases of whole-input
displays on local cores. All the comparisons are local to the retained
real ray plaques.

\Needspace{24\baselineskip}
\textbf{Organization.} In the forward construction the local transverse
trace-free encoder has Fourier order four, and each negative contact
response lowers that order by two, so three responses give
\(4-2\cdot3=-2\), the order of the flat constraint inverses
(\hyperref[lem:online-negative-response-intermediate-orders]{Lemma 4.11}
and \hyperref[thm:online-negative-response-flat-green]{Theorem 4.12}).

{\def\LTcaptype{none} % do not increment counter
\begin{longtable}[]{@{}
  >{\raggedright\arraybackslash}p{(\linewidth - 4\tabcolsep) * \real{0.1912}}
  >{\raggedright\arraybackslash}p{(\linewidth - 4\tabcolsep) * \real{0.5588}}
  >{\raggedright\arraybackslash}p{(\linewidth - 4\tabcolsep) * \real{0.2500}}@{}}
\toprule\noalign{}
\begin{minipage}[b]{\linewidth}\raggedright
Purpose
\end{minipage} & \begin{minipage}[b]{\linewidth}\raggedright
Construction
\end{minipage} & \begin{minipage}[b]{\linewidth}\raggedright
Proof
\end{minipage} \\
\midrule\noalign{}
\endhead
\bottomrule\noalign{}
\endlastfoot
Construct a spacetime germ & Read the two seeds from the classes.
Contract the three responses to obtain the flat inverses. Initialized
nonlinear transport gives vacuum initial data and hence the two spin
gradient operators. Classical Einstein Cauchy evolution develops the
data. & \hyperref[sec:online-ordinary-data-construction]{§3.6},
\hyperref[thm:online-ordinary-negative-response-initial-source]{Theorem~4.18}, \hyperref[thm:online-smooth-local-cauchy-input]{Theorem~2.6} \\
Present a given germ & Extend normalized maximal-disk data to the whole
space with exact agreement on a smaller disk. Each such extension in the
fixed gauge decodes to a unique pair of classes. The germ determines the
comparison class of its displays. &
\hyperref[thm:online-relative-extension-one-branch]{Theorem~B.1},
\hyperref[thm:online-ordinary-pair-data-bijection]{Theorems~3.13} and
\hyperref[thm:submission-vacuum-germ-coverage]{5.7} \\
Compare distinct slices & The scalar wave kernel transports the full
evaluation between carriers. A physical isometry with its specified spin
lift then has a unique full source comparison up to normalized cells. &
\hyperref[thm:online-canonical-wave-kernel]{Theorems~6.10} and
\hyperref[thm:online-full-unmarked-source-hom]{7.11} \\
\end{longtable}
}

The resulting local inverses and the descent of display atlases are
proved in \hyperref[thm:local-two-inverses]{Theorem 8.3} and
\hyperref[thm:local-display-atlas-descent]{Theorem 8.5}.

\section{2. Intrinsic initial sources and local vacuum
evolution}\label{intrinsic-initial-sources-and-local-vacuum-evolution}\label{sec:online-intrinsic-source-and-evolution}

The source has two chiral gradient operators with a common real Clifford
symbol. We first give independent local tests for these operators and
prove that they recover vacuum initial data and their actual spin
comparisons. This recognition applies without weighted smallness or a
maximality assumption. The ordinary-pair construction of Sections 3--4
produces the maximal subclass used in the local coverage argument.

An actual finite Cartan collar then carries the source off its marked
initial hypersurface. Its soldering, connection and Weyl field obey
finite equations with a closed subsidiary system. Their geometric
recognition is proved directly; smooth existence and geometric
uniqueness use the precise classical Cauchy input stated below.
Finite-dimensional families are constructed on one low-order-selected
parent, with every higher estimate referring to that same solution.

The comparison results in this section retain the initial hypersurface.
They identify all physical and actual spin comparisons in that category,
and compare exterior, gauge and scale choices for a fixed physical
initial germ. The wave enrichment and comparisons that change the
initial hypersurface require their additional structures.

\subsection{2.1. Intrinsic initial
sources}\label{intrinsic-initial-sources}\label{sec:online-intrinsic-initial-source}

The initial source constructed from an ordinary pair has a canonical
weighted presentation. Its geometric recognition uses only local
differential identities. We formulate those identities on an arbitrary
oriented smooth three-manifold \(\Sigma\), with no weighted coordinate
system or prescribed mean curvature. This separates the construction of
a display from the recognition and comparison of its actual initial
fields.

All manifolds are Hausdorff and second countable, and initial open sets
have no boundary as manifolds. A vacuum initial datum is
\(d=(\Sigma,h,K,P_\Sigma)\), where \(h\) is a positive metric, \(K\) is
a real symmetric covariant two-tensor, and \(P_\Sigma\) is an actual
\(\operatorname{Spin}(3)\) lift of the oriented \(h\)-orthonormal frame
bundle. Put \(\tau=\operatorname{tr}_hK\) and require

\phantomsection\label{eq:online-intrinsic-vacuum-constraints} \[
R(h)+\tau^2-|K|_h^2=0,\qquad
\operatorname{div}_hK-\mathrm d\tau=0.
\] A datum comparison over an orientation-preserving diffeomorphism
\(f:\Sigma\to\Sigma'\) consists of \(f^*h'=h\), \(f^*K'=K\) and an
actual spin-bundle lift of its orthonormal frame map. Denote this
groupoid by \(\mathsf{Data}_0^\infty\). The definitions apply
componentwise when the initial manifold is disconnected.

The algebraic marking of a gradient source is a real oriented and
time-oriented Lorentz rank-four bundle \((\mathbb E,\eta)\) of signature
\((+---)\), its actual \(\operatorname{Spin}^{+}(1,3)\) lift, a future
unit section \(N\), and the associated complex Dirac module
\(\mathscr S=\mathscr S_+\oplus\mathscr S_-\). It includes the real
Clifford action \(c\) and the fixed invariant nondegenerate complex
alternating form \(\mathsf B\). The normalization is

\phantomsection\label{eq:online-intrinsic-clifford-normalization} \[
c(v)c(w)+c(w)c(v)=2\eta(v,w)I,\qquad
\operatorname{tr}_{\mathscr S}(c(v)c(w))=4\eta(v,w).
\] For clarity, in a local spin frame the model is

\phantomsection\label{eq:online-intrinsic-dirac-model} \[
\begin{gathered}
c(v)=
\begin{pmatrix}0&V\\ \widetilde V&0\end{pmatrix},
\qquad
V=v^0I+\sum_a v^a\sigma_a,\qquad
\widetilde V=v^0I-\sum_a v^a\sigma_a, \\
\begin{gathered}
\mathsf B=
\begin{pmatrix}\epsilon&0\\0&\epsilon\end{pmatrix}, \\
\epsilon=\begin{pmatrix}0&1\\-1&0\end{pmatrix}, \\
\varrho(A)=\operatorname{diag}\bigl(A,(A^\dagger)^{-1}\bigr),
\quad A\in SL(2,\mathbb C).
\end{gathered}
\end{gathered}
\] Here \(\sigma_a\) are the Pauli matrices. The oriented Clifford
product fixes the chirality labels. The Clifford vectors generate the
full complex endomorphism algebra of each Dirac fiber. These are
algebraic bundle markings; a metric on \(T\Sigma\), a soldering or a
connection is not supplied.

\pandocstatementheading{\leavevmode\phantomsection\label{def:online-intrinsic-vacuum-gradient}Definition 2.1 (Intrinsic vacuum gradient source).}

Supply two actual smooth complex-linear first-order operators

\phantomsection\label{eq:online-intrinsic-gradient-operators} \[
\mathbb D_\pm:\Gamma_{\mathrm{loc}}(\mathscr S_\pm)
\longrightarrow
\Gamma_{\mathrm{loc}}(\mathbb E_{\mathbb C}\otimes\mathscr S_\pm).
\] Their scalar commutators must have a common real
orientation-preserving bundle isomorphism as symbol, and their complete
outputs must be tangent to the internal spatial plane:

\phantomsection\label{eq:online-intrinsic-symbol-tests} \[
[\mathbb D_\pm,m_a]\psi=b(\mathrm da)\otimes\psi,
\qquad
b:T^*\Sigma\xrightarrow{\sim}N^\perp,
\qquad
\eta(N,\mathbb D_\pm\psi)=0
\] for every real local scalar \(a\). Define the algebraic and
directional readouts

\phantomsection\label{eq:online-intrinsic-symbol-directional-readout} \[
\begin{aligned}
h^{-1}(\xi,\zeta)&=-\eta(b\xi,b\zeta),&
\Theta X&=b(h^\flat X),\\
P_X^\pm\psi&=-\eta(\Theta X,\mathbb D_\pm\psi),&
P_X&=P_X^+\oplus P_X^-.
\end{aligned}
\] Require, for real \(X\) and real Clifford sections \(v\),

\phantomsection\label{eq:online-intrinsic-normalization-tests} \[
\begin{aligned}
X\mathsf B(\psi,\chi)
&=\mathsf B(P_X\psi,\chi)+\mathsf B(\psi,P_X\chi),\\
[P_X,c(v)]&\in\Gamma_{\mathrm{loc}}(c(\mathbb E)).
\end{aligned}
\] The connection \(D\) specified by \(c(D_Xv)=[P_X,c(v)]\) is justified
below. The remaining admission tests are

\phantomsection\label{eq:online-intrinsic-torsion-weyl-test} \[
D\Theta=0,\qquad
\Omega_D\in\operatorname{im}\mathsf W_\Theta,
\qquad
\mathsf W_\Theta(C)(X,Y)=\mathcal C(\Theta X,\Theta Y),
\] where \(C\) ranges over real algebraic Weyl tensors on \(\mathbb E\)
and \(\eta(\mathcal C(u,v)z,w)=C(u,v,z,w)\). The curvature in this test
is computed from the recovered connection. Its Weyl preimage is a
readout, rather than an independently supplied curvature field.

The finite image test can be carried out explicitly. In an adapted
orthonormal frame \(v_0=N\), \(v_a=\Theta(e_a)\), compute

\phantomsection\label{eq:online-intrinsic-electric-magnetic-readout} \[
E_{ab}=\sum_c(\Omega_D)^c{}_b(e_c,e_a),
\qquad
B_{ab}=\frac12\epsilon_a{}^{cd}(\Omega_D)^0{}_b(e_c,e_d).
\] Require \(E,B\) to be symmetric and trace-free. Form the algebraic
tensor \(C(E,B)\) from curvature symmetries and

\phantomsection\label{eq:online-intrinsic-weyl-components} \[
C_{0a0b}=E_{ab},\qquad
C_{cd0b}=-\epsilon_{cd}{}^a B_{ab},\qquad
C_{abcd}=-\epsilon_{ab}{}^m\epsilon_{cd}{}^n E_{mn},
\] and test the equality \(\Omega_D=\mathsf W_\Theta(C(E,B))\). The
proof below shows that this is exactly the image test and that the
readout is independent of the adapted frame.

A physical source comparison over \(f:\Sigma\to\Sigma'\) is an actual
multiplication-semilinear isomorphism of the sheaves of chiral sections,
\(\widehat U=\widehat U_+\oplus\widehat U_-\). It preserves
\(\mathsf B\), the real Clifford-vector action and its orientation and
time orientation, and \(N\). Its induced vector map is denoted by \(L\).
Require the exact operator identities

\phantomsection\label{eq:online-intrinsic-gradient-intertwining} \[
\mathbb D'_\pm\widehat U_\pm
=\widehat{L\otimes U_\pm}\,\mathbb D_\pm.
\] The passage from a semilinear sheaf map to the smooth bundle maps
\(U,L\) is proved below. The map \(U\) itself is retained as a morphism.
These objects and comparisons form \(\mathsf{Grad}_0^\infty\).

Definition 2.1 tests the algebraic marking and the two operators
directly on the initial source, before either developing the data or
seeking an ordinary-pair presentation.

\pandocstatementheading{\leavevmode\phantomsection\label{lem:online-intrinsic-gradient-decoder}Lemma 2.2 (Intrinsic recovery of the directional connection).}

The symbol, tangential-output and normalization tests determine a
positive metric \(h\), an oriented spacelike soldering \(\Theta\), a
unique real Lorentz connection \(D\) and its exact spin connection
\(P\). They recover the complete supplied operators by

\phantomsection\label{eq:online-intrinsic-uncontracted-recovery} \[
\mathbb D_\pm\psi
=\sum_{a=1}^3\Theta(e_a)\otimes P^\pm_{e_a}\psi
\] in every oriented \(h\)-orthonormal frame. If \(D\Theta=0\), then

\phantomsection\label{eq:online-intrinsic-gauss-weingarten} \[
\begin{aligned}
K(X,Y)&=-\eta(D_XN,\Theta Y),\\
D_XN&=\Theta(K^\sharp X),\\
D_X(\Theta Y)&=\Theta(\nabla_X^hY)+K(X,Y)N,
\end{aligned}
\] where \(K\) is symmetric. Conversely, the last two formulas define
the unique metric connection with the given \(h,K,N,\Theta\) and zero
torsion.

\emph{Proof.} Positivity follows from the negative-definiteness of
\(\eta\) on \(N^\perp\). The readout satisfies
\(-\eta(\Theta X,\Theta Y)=h(X,Y)\) and preserves orientation. The
directional operator is function-linear in \(X\), and its scalar
commutator is \[
[P_X,m_a]\psi
=-\eta(\Theta X,b\,\mathrm da)\psi
=h^{-1}(h^\flat X,\mathrm da)\psi
=X(a)\psi.
\] It is therefore a connection. Its commutator action on real Clifford
vectors obeys the real Leibniz rule and defines \(D\). Commuting \(P_X\)
with the Clifford anticommutator gives
\(X\eta(v,w)=\eta(D_Xv,w)+\eta(v,D_Xw)\), so \(D\) is Lorentz.

The connection \(D\) has a unique lift to the supplied spin bundle. Its
induced Dirac connection and \(P\) act on Clifford vectors through the
same \(D\), so their difference commutes with every Clifford vector. The
full-matrix-algebra property makes that difference \(\lambda_X I\). Both
connections preserve \(\mathsf B\), and subtraction of their
preservation identities gives \(2\lambda_X\mathsf B(\psi,\chi)=0\). Thus
\(\lambda_X=0\). This removes the scalar connection ambiguity.

For \(v_a=\Theta(e_a)\), every vector in \(N^\perp_{\mathbb C}\) is
\(\sum_a v_a(-\eta(v_a,\cdot))\). Apply this identity to the entire
output of \(\mathbb D_\pm\). Its contractions are the \(P_{e_a}^\pm\),
and tangential output excludes a normal zero-order component. This
proves the full recovery formula and its frame independence.

Metric compatibility makes \(D_XN\) spatial and gives
\(\eta(D_X(\Theta Y),N)=K(X,Y)\). The normal component of \(D\Theta=0\)
is \(K(X,Y)-K(Y,X)=0\). The tangential component defines a metric,
torsion-free connection on \(T\Sigma\), hence \(\nabla^h\). This proves
the Gauss--Weingarten identities. Conversely they determine a connection
on the splitting \(\mathbb E=\mathbb RN\oplus\Theta(T\Sigma)\). Direct
differentiation of the three pairings of \(N\) and \(\Theta(T\Sigma)\)
proves its metric compatibility; symmetry of \(K\) and torsion-freeness
of \(\nabla^h\) prove \(D\Theta=0\). The same decomposition proves
uniqueness. \hfill\(\square\)

\subsection{2.2. Vacuum recognition and the initial unit and
counit}\label{vacuum-recognition-and-the-initial-unit-and-counit}\label{sec:online-intrinsic-vacuum-recognition}

The curvature test has a unique algebraic output even at zero Weyl
curvature. The initial recognition therefore uses no curvature-rank
inverse.

\pandocstatementheading{\leavevmode\phantomsection\label{thm:online-intrinsic-gradient-vacuum-recognition}Theorem 2.3 (Intrinsic vacuum recognition).}

An independently admitted object of \(\mathsf{Grad}_0^\infty\)
determines a datum of \(\mathsf{Data}_0^\infty\) by the metric and
second-form readouts above and the actual spin reduction at \(N\),
transported by \(\Theta\). Its unique Weyl readout is

\phantomsection\label{eq:online-intrinsic-gauss-codazzi-readout} \[
E=\operatorname{Ric}(h)+\tau K-Kh^{-1}K,\qquad
B_{ab}=\epsilon_a{}^{cd}\nabla_cK_{db},
\qquad \tau=\operatorname{tr}_hK.
\] Conversely, each datum \(d\) gives an independently admitted source.
The two operations recover all initial tensors, the actual initial spin
structure and the complete supplied operator pair under canonical
identifications. No maximality, decay, smallness or curvature-rank
condition is required for this recognition.

\emph{Proof.} The preceding lemma supplies \(h,K,D\) with symmetric
\(K\). In an adapted frame put \(\kappa_a=K_{ab}\theta^b\) and let
\(\omega\) be the Levi--Civita connection of \(h\). The Lorentz
connection has blocks \(A^0{}_a=A^a{}_0=\kappa_a\) and
\(A^a{}_b=\omega^a{}_b\). Consequently

\phantomsection\label{eq:online-intrinsic-gauss-codazzi-blocks} \[
\begin{aligned}
\Omega^a{}_b&=(\Omega^h)^a{}_b+\kappa_a\wedge\kappa_b,\\
\Omega^0{}_b(e_c,e_d)&=\nabla_cK_{db}-\nabla_dK_{cb}.
\end{aligned}
\] The positive sign in the quadratic term follows from the two equal
mixed blocks in signature \((+---)\). Contracting gives exactly the
displayed \(E,B\). The matrix \(E\) is symmetric. The trace of \(B\)
vanishes by symmetry of \(K\), and

\phantomsection\label{eq:online-intrinsic-four-constraints} \[
\operatorname{tr}E=R(h)+\tau^2-|K|_h^2,\qquad
\epsilon_a{}^{bc}B_{bc}
=\nabla^jK_{ja}-\nabla_a\tau.
\] The stated component formula for \(C(E,B)\) gives a real algebraic
Weyl tensor whenever \(E,B\) are symmetric and trace-free: alternation
and pair symmetry define the other components. In an oriented
orthonormal frame the cyclic sum of
\(C_{0abc}=-\epsilon_{bc}{}^d B_{da}\) over \(a,b,c\) is
\(-\epsilon_{abc}\operatorname{tr}B\). The mixed Ricci contraction is,
up to the contraction sign, \(\sum_{b,d}\epsilon_{ab}{}^d B_{db}\),
which vanishes exactly when \(B\) is symmetric. The purely spatial
cyclic identity follows from the spatial component formula with
symmetric \(E\). The time-time Ricci contraction is a signed trace of
\(E\); in the spatial Ricci contraction the two \(E_{ab}\) terms cancel
and only a multiple of \(\delta_{ab}\operatorname{tr}E\) remains.
Conversely, the contractions in the definition recover \(E,B\) from
every algebraic Weyl tensor. These matrices determine all its components
by the same formula. In particular the restriction \(\mathsf W_\Theta\)
is injective and has the asserted finite inverse on its image.
Equivalently, if a Weyl tensor vanishes on every spatial first pair,
pair symmetry removes every component except possibly
\(C(N,v_a,N,v_b)\); its Ricci trace then removes those components as
well. The uniqueness makes the reconstruction independent of the adapted
frame.

Admission therefore gives symmetric trace-free \(E,B\). The two
identities above imply the vacuum constraints, with their stated signs.
Reduction of the actual Lorentz spin bundle at \(N\) gives a
\(\operatorname{Spin}(3)\) lift of the spatial orthonormal frame bundle
in \(N^\perp\). Pulling it back through the oriented isometry
\(\Theta:(T\Sigma,h)\to(N^\perp,-\eta)\) gives \(P_\Sigma\).

For the converse, given \(d=(\Sigma,h,K,P_\Sigma)\), set

\phantomsection\label{eq:online-intrinsic-constructor-connection} \[
\begin{aligned}
\mathbb E_d&=\mathbb R\oplus T\Sigma, \\
\eta_d((a,X),(b,Y))&=ab-h(X,Y), \\
N_d&=(1,0), \\
\Theta_dX&=(0,X), \\
D_X^d(a,Y)&=\bigl(Xa+K(X,Y),\,
                    \nabla_X^hY+aK^\sharp X\bigr).
\end{aligned}
\] Extend \(P_\Sigma\) through
\(\operatorname{Spin}(3)\subset\operatorname{Spin}^{+}(1,3)\) and take
its associated chiral modules and connections \(\nabla^{d,\pm}\). Define

\phantomsection\label{eq:online-intrinsic-constructor-gradients} \[
\mathbb D_\pm^d\psi
=\sum_a\Theta_d(e_a)\otimes\nabla^{d,\pm}_{e_a}\psi.
\] The expression is invariant under an oriented orthonormal change of
frame. Its symbol is \(b_d(\xi)=\Theta_d(\xi^{\sharp_h})\), and its full
output is spatial. Its directional readout is the actual spin connection
of \(D^d\), so it passes the alternating and real Clifford tests. The
previous lemma gives torsion-freeness.

The constraints make the displayed \(E,B\) symmetric and trace-free. The
mixed curvature block is the epsilon reconstruction of \(B\). The
spatial block is an algebraic three-dimensional curvature tensor with
Ricci contraction \(E\). Such a tensor is determined by that
contraction: writing its components in the order
\(G_{cdba}=\epsilon_{cdm}\epsilon_{ban}S_{mn}\) identifies it with a
symmetric matrix \(S\), and the contraction is
\(S-(\operatorname{tr}S)I\), an invertible linear map since its trace is
\(-2\operatorname{tr}S\). The spatial Weyl restriction of \(C(E,B)\) has
the same contraction, so it equals the Gauss block. Metric skewness
determines the other mixed block. Thus
\(\Omega_{D^d}=\mathsf W_{\Theta_d}(C(E,B))\), establishing the
independent image test.

Decoding this source recovers \(h,K\) exactly. Reduction of its extended
spin bundle at \(N_d\) canonically recovers \(P_\Sigma\). For a source
followed by decoding and reconstruction, the vector comparison is

\phantomsection\label{eq:online-intrinsic-source-counit} \[
L_d:\mathbb R\oplus T\Sigma\longrightarrow\mathbb E,
\qquad L_d(a,X)=aN+\Theta X.
\] It is an oriented and time-oriented isometry, and the
Gauss--Weingarten formula identifies its transported connection with
\(D\). The source spin bundle is canonically the extension of its
reduction at \(N\): the map on associated principal bundles is
\([p,g]\mapsto p g\). The corresponding Dirac counit is the evaluation
map \([p,\psi]\mapsto p(\psi)\). It has its original actual spin label.
The exact uncontracted recovery lemma now identifies the reconstructed
pair with the supplied \(\mathbb D_\pm\), including every zero-order
coefficient. These are the inverse object identifications.
\hfill\(\square\)

The symbol metric retains physical scale. In particular a physical
comparison must preserve \(h\) itself. On the maximal weighted chart,
the source of
\hyperref[thm:online-ordinary-negative-response-initial-source]{Theorem
4.18} satisfies these intrinsic tests by its four commutator equations,
and its intrinsic readout agrees with its canonical one. Restricting
that constructed source to an initial open set restricts the fields and
these local readouts. The whole weighted decoder of Section 3 is used
before this restriction.

\subsection{2.3. Physical comparisons on a fixed initial
hypersurface}\label{physical-comparisons-on-a-fixed-initial-hypersurface}\label{sec:online-initial-physical-comparisons}

The comparison maps retain the actual spin action on spinor sections. We
prove that the source intertwining equations determine precisely the
physical comparisons of the decoded data.

\pandocstatementheading{\leavevmode\phantomsection\label{thm:online-initial-gradient-full-hom}Theorem 2.4 (Full physical Hom on the initial manifold).}

The intrinsic constructor and decoder extend to an equivalence

\phantomsection\label{eq:online-intrinsic-initial-equivalence} \[
\mathsf{Grad}_0^\infty\simeq\mathsf{Data}_0^\infty
\] which is bijective on Hom sets over each specified initial
diffeomorphism. Its unit and counit are natural exact comparisons and
satisfy the triangle identities. The two central spin lifts on a
connected initial domain remain distinct morphisms.

\emph{Proof.} First, a multiplication-semilinear isomorphism of local
section sheaves over \(f\) is a unique smooth bundle isomorphism over
\(f\). Indeed, identify the target with the pullback bundle on the
source and write a section in a local frame. Semilinearity says that its
image is the same linear combination of the images of those frame
sections. These images are smooth, so they form the columns of a smooth
matrix. The inverse sheaf map gives its inverse matrix. Frame changes
give the bundle transformation law. This argument is local and requires
no functional-analytic assumption on the section map.

The induced conjugation on real Clifford vectors defines a smooth bundle
map \(L\). The Clifford relation makes it a Lorentz isometry. The
required algebraic orientations and preservation of \(N\) make it proper
and orthochronous. In local spin frames choose one spin lift of \(L\),
with Dirac action \(U_0\). Any other Clifford intertwiner \(U\) differs
by a matrix commuting with the entire complex Clifford algebra; hence
\(U_0^{-1}U=\lambda I\). Preservation of \(\mathsf B\) gives
\(\lambda^2=1\). Thus the permitted intertwiners are precisely the two
actual spin actions. Their local lifts glue because their Dirac actions
are the given smooth \(U\) and the Dirac representation is faithful. The
central element \(-1\) acts as \(-I\), so its two actions have not been
identified.

Take scalar commutators of the exact gradient intertwining equation. At
each \(x\) they give

\phantomsection\label{eq:online-initial-comparison-symbol-law} \[
b'_{f(x)}
 \bigl(((\mathrm df_x)^{-1})^*\xi\bigr)
=L_xb_x(\xi),\qquad \xi\in T_x^*\Sigma.
\] The metric readout then gives \(f^*h'=h\) and
\(L\Theta X=\Theta'f_*X\). Contracting the full gradient identity with
these soldered vectors gives

\phantomsection\label{eq:online-initial-comparison-connection-law} \[
P'_{f_*X}\widehat U=\widehat U P_X,
\qquad
D'_{f_*X}(Lv)=L(D_Xv).
\] The second equality follows by commuting the first with each Clifford
vector. Preservation of \(N\) now gives \(f^*K'=K\). The spin map
restricts at \(N\) and, through \(\Theta\), becomes precisely a datum
spin comparison. The unique Weyl readout is carried to its counterpart
because the connection and its curvature are carried exactly.

Conversely, a datum comparison induces the oriented Lorentz isometry
between the split bundles in the constructor and the actual extension of
its supplied spin lift. The formula for \(D^d\) shows that it
intertwines the full connections. Its uncontracted gradient formula then
gives the exact source intertwining equations. Thus every datum
comparison has a source comparison, and the two assignments on arrows
are inverse.

The unit identifies the recovered \(h,K,P_\Sigma\) with the original
datum. The counit is the vector map \(L_d\) and the associated
principal-spin and spinor evaluation maps in the recognition theorem.
Every physical comparison commutes with these evaluations, proving
naturality. Reconstructing and then evaluating acts as the identity on
each vector, spin frame and spinor; this gives the two triangle
identities. It also verifies that the Hom bijection retains the supplied
actual spin map on both passes. \hfill\(\square\)

This Hom statement compares sources over a specified map of their
initial hypersurfaces. Comparisons in which the initial hypersurface
itself moves inside a spacetime require the additional local evolution
and point-family structures.

\subsection{2.4. Initial descent and finite local
estimates}\label{initial-descent-and-finite-local-estimates}\label{sec:online-initial-descent-and-estimates}

\pandocstatementheading{\leavevmode\phantomsection\label{prop:online-initial-gradient-descent-families}Proposition 2.5 (Initial descent and local smooth dependence).}

Intrinsic vacuum gradient sources and their physical comparisons admit
effective descent over open covers of a fixed initial manifold. The
constructor, decoder, Weyl readout and their unit and counit commute
with restriction and preserve smooth finite-dimensional families with
supplied local algebraic and spin identifications.

More precisely, in retained coordinate and spin charts let
\(K\Subset\operatorname{int}K'\) be compact spatial sets and
\(\mathcal K\Subset\operatorname{int}\mathcal K'\) compact parameter
sets. Impose a positive lower bound for the symbol metric and finite
bounds for the inverse marking matrices on \(K'\times\mathcal K'\).
Every prescribed finite spatial and parameter norm of the readouts and
comparison maps on \(K\times\mathcal K\) is controlled by finitely many
such norms of the supplied fields and comparison maps on
\(K'\times\mathcal K'\). For the object constructor and decoder, spatial
input order \(m+2\) suffices for output order \(m\), with constants
depending on the stated finite coefficient bounds and inverse margins.

In particular, locally smooth convergence of admitted sources and their
markings to a nondegenerate spacelike-symbol limit gives an admitted
limit and locally smooth convergence of the decoded data and Weyl
readout, including zero-curvature limits.

\emph{Proof.} Let sources be supplied on \(\{O_i\}\) with actual overlap
module maps \(U_{ij}\) covering the identity and satisfying

\phantomsection\label{eq:online-initial-actual-spin-cocycle} \[
U_{ij}U_{jk}=U_{ik}
\] on every triple overlap. The preceding theorem makes these smooth
bundle maps and actual spin maps. The ordinary bundle cocycle glues the
Dirac modules, Clifford-vector bundles, Lorentz and alternating forms,
future section and actual spin reductions. The scalar commutator
identity glues the symbol and soldering. Exact operator intertwining
glues the two first-order operators as actions on the sheaves of
sections.

All admission conditions are local. Their finite readouts agree under
the same transitions, so the torsion and Weyl-image conditions hold for
the glued object. Its Weyl field glues by the uniqueness of the
algebraic preimage. Compatible morphisms glue as actual bundle maps and
retain the exact intertwining equations. If two gluings have the
prescribed local identifications, those identifications give compatible
isomorphisms on every \(O_i\) and hence a unique global comparison. This
is effective descent over the given manifold, including the exact spin
signs on triples.

For the local estimates, work in finitely many retained spin frames. The
inverse metric and \(\Theta\) are finite matrix operations on the symbol
and markings. The directional connection is obtained by contraction of
the full supplied operator coefficients. Its vector connection follows
by finite Clifford commutators, and the second form uses this connection
and one derivative of \(N\). Its curvature and the algebraic Weyl
readout use one further spatial derivative. The positive inverse-matrix
margins bound every denominator in these operations.

In the converse construction, an oriented local frame for \(h\) is
obtained from a fixed coordinate frame by its positive square root. The
connection uses one derivative of \(h\) and the value of \(K\); its
curvature uses two derivatives of \(h\) and one of \(K\). The matrix
square root and inverse are smooth on the retained positive set.
Differentiating these finite formulas proves the stated spatial bound.
Repeated passive-parameter differentiation obeys the finite product and
inverse-matrix rules and therefore uses finitely many input norms, with
no additional symbol-smallness condition.

For a supplied comparison, semilinearity reads its local bundle matrix
directly from its values on a frame. The vector map follows by Clifford
conjugation; the actual spin map is its locally initialized lift through
the smooth double covering. On a connected lift chart its sign is fixed.
The unit and counit are the finite associated-bundle evaluations already
described. The supplied comparison and the bounds on each lift chart give the same finite parameter estimates, with the sign fixed only locally.

Every local formula commutes with restriction. Their compositions with
overlap maps agree by the Hom theorem, which proves compatibility of
recognition and descent. For a convergent sequence, a retained compact
set has a uniform positive symbol margin for all sufficiently late
members. All finite operations then converge in each finite derivative
norm. Torsion and curvature admission pass to the limit because they are
finite differential identities and the Weyl-image test uses the explicit
continuous algebraic reconstruction. It follows that the limit is
admitted and that all the stated readouts converge, even where the
curvature evaluation loses rank. \hfill\(\square\)

The estimates in this proposition are estimates for finite initial
differential constructions. The common spacetime domain and evolution
estimates are supplied by the local Cauchy and collar results below.

\subsection{2.5. The smooth Cauchy input and its local
use}\label{the-smooth-cauchy-input-and-its-local-use}\label{sec:online-smooth-cauchy-input}

The geometric construction uses the smooth Einstein Cauchy theorem. The
finite Cartan equations below identify the resulting source and
propagate its constraints; existence is supplied by the hyperbolic
Cauchy theory.

\pandocstatementheading{\leavevmode\phantomsection\label{thm:online-smooth-local-cauchy-input}Theorem 2.6 (Local smooth Einstein Cauchy input).}

Let \((h,K)\) be smooth vacuum initial data on an open three-manifold,
with \(h\) positive definite and a chosen future direction. There is a
smooth Ricci-flat development germ along the initial hypersurface. Two
developments of the same data have a unique germ of
time-orientation-preserving isometry extending a fixed initial
identification.

The external geometric input is the existence and common-development
statement of \hyperref[ref-Sbierski2016]{Sbierski} (2016, Theorem 2.6).
His Lemma 3.1 and Corollary 3.2 give uniqueness of the comparison with
fixed initial identification. We use the local germs of these
conclusions.

For the reduced hyperbolic construction we use
\hyperref[ref-HughesKatoMarsden1977]{Hughes, Kato and Marsden} (1977,
Theorem III and its following regularity remark; Theorem V): the
harmonic metric equation has a unique solution on one time interval for
data in a fixed \(H^{s+1}\times H^s\) neighborhood; the solution depends
continuously on the data there, and additional smooth regularity
persists on that interval. We choose an integer base index \(s_0\ge5\)
and retain a positive hyperbolicity margin.

\emph{Proof.} We justify the local application. The vacuum constraints
are exactly the smooth initial-data conditions in the cited geometric
theorem. Multiplying a metric of signature \((-+++)\) by \(-1\) leaves
its Levi--Civita connection and its Ricci tensor unchanged and gives our
signature \((+---)\). With \(K=\frac12\partial_t h_t\) in Gaussian
coordinates, the initial tensors have the stated convention. Restricting
the common development gives local isometries. A time-oriented isometry
fixing the initial hypersurface also fixes its normal, so the cited
comparison uniqueness applies. On overlapping initial opens the unique
germs agree. This proves the germ assertion; it does not require an open
product collar itself to be a globally hyperbolic development.

Here is the localization used for estimates. In a coordinate ball choose
\(V\Subset W\Subset W_2\Subset U\), where the initial data are supplied
on \(U\). Use a fixed cutoff to extend \(h\) from a neighborhood of
\(\overline{W_2}\) to a smooth positive metric on \(\mathbb R^3\), equal
to a fixed Euclidean metric outside a compact set. Convex interpolation
of positive matrices gives such an extension with a positive lower
bound. Extend \(K\) smoothly by the same fixed cutoff. For a family near
one parameter, these extensions depend smoothly on that parameter and
have finite estimates in every specified order.

For the reduced metric equation take, at coordinate time \(t=0\),

\phantomsection\label{eq:online-harmonic-initial-tensors} \[
g_{00}=1,\qquad g_{0i}=0,\qquad g_{ij}=-h_{ij},
\qquad
\partial_tg_{ij}=-2K_{ij}.
\] Choose the four remaining time derivatives to make
\(\mathcal H^\mu=g^{\alpha\beta}\Gamma^\mu_{\alpha\beta}=0\) initially.
Explicitly, with \(\tau=\operatorname{tr}_hK\) and
\(\Gamma_h^k=h^{ij}\Gamma(h)^k{}_{ij}\), they are

\phantomsection\label{eq:online-harmonic-initial-gauge} \[
\partial_tg_{00}=2\tau,\qquad
\partial_tg_{0i}=-h_{ik}\Gamma_h^k.
\] These are finite differential operations on \((h,K)\). The reduced
equations \(\operatorname{Ric}(g)-\nabla_{(\mu}\mathcal H_{\nu)}=0\)
have principal part
\(-\frac12g^{\alpha\beta}\partial_\alpha\partial_\beta g_{\mu\nu}\) and
lower terms smooth in \(g^{-1},g,\partial g\). Subtract a fixed extended
initial metric. The new unknown and its initial time derivative have the
Sobolev behavior required by the hyperbolic input. The coefficients are
smooth, their spatial derivatives are bounded, and the spatial principal
matrices and time coefficient are uniformly definite after the sign
convention is adjusted. At \(t=0\) this follows from the positive bound
for \(h\) and the zero shift; it persists in one sufficiently small
solution neighborhood. Thus the hypotheses of the cited reduced theorem
hold.

The cutoff extension need not satisfy the vacuum constraints away from
\(W_2\). To see why this is sufficient, the contracted Bianchi identity
for a reduced solution gives the homogeneous gauge equation

\phantomsection\label{eq:online-harmonic-gauge-propagation} \[
\nabla^\alpha\nabla_\alpha\mathcal H_\nu
+R_\nu{}^\mu\mathcal H_\mu=0.
\] Initially \(\mathcal H=0\). On \(W_2\) its normal derivative is zero
as well: the constraints give the normal components of the Einstein
tensor equal to zero, while the reduced equation expresses these
components, at \(\mathcal H=0\), as one half of the corresponding normal
derivatives of \(\mathcal H\). Tangential derivatives of its zero
initial value vanish. Uniqueness for the linear gauge equation on the
domain of dependence of \(W_2\) therefore gives \(\mathcal H=0\) there,
and the reduced metric is Ricci-flat on that domain.

Finite propagation follows from the wave energy identity on truncated
cones whose lateral speed exceeds the characteristic speeds; their
outward flux has the required sign. Uniform low-order metric bounds give
one such speed bound. Hence a common smaller time interval and the fixed
spatial buffers put the retained reduced-coordinate region, and
subsequently the normal geodesics from \(W\), inside this domain of
dependence. This proves the required local applicability without an
exterior extension of the vacuum constraints. \hfill\(\square\)

\subsection{2.6. Actual finite Cartan
collars}\label{actual-finite-cartan-collars}\label{sec:online-finite-cartan-collars}

The intrinsic initial source determines a soldering \(\Theta\), a real
metric connection \(D\) and a unique algebraic Weyl tensor \(C_0\). We
now give its evolution an independent finite-field presentation. Its
unknowns are a spacetime soldering, a Lorentz connection and an
algebraic Weyl tensor. The frame formulation by these fields is
classical (\hyperref[ref-FriedrichNagy1999]{Friedrich and Nagy 1999},
§2, equations (2.2)--(2.5)). The temporal gauge and subsidiary
identities used here are proved below. Existence is supplied by
\hyperref[thm:online-smooth-local-cauchy-input]{the smooth local Cauchy
theorem}.

A collar is a neighborhood of the marked zero slice in a product with a
time interval. We work with its germ along a retained initial patch, and
may shrink that patch or the neighborhood. A product \((-T,T)\times V\)
with spatial boundary is used as a coordinate domain; Cauchy uniqueness
is applied on smaller domains of dependence inside a buffered initial
patch. A common positive width is asserted on retained compact patches,
not on an arbitrary noncompact initial manifold.

Let \(\mathbb G\) be an independently admitted intrinsic vacuum gradient
source on an initial open set \(U\), as in
\hyperref[def:online-intrinsic-vacuum-gradient]{the intrinsic
definition}. Write its readout as \((\mathbb E,\eta,N,\Theta,D,C_0)\),
with its specified actual spin bundle. Extend these algebraic bundles by
product in time and put \(n(t,x)=N(x)\). In an adapted initial spin
frame extended without time dependence, \(n\) has the constant
components \(e_0\). Changes between such frames are independent of time.

For an algebraic curvature tensor \(C\), write \(\mathcal C\) for its
endomorphism form, with

\phantomsection\label{eq:online-cartan-curvature-endomorphism} \[
\eta\bigl(\mathcal C(u,v)z,w\bigr)=C(u,v,z,w).
\] This uses the same last-slot convention as the intrinsic Weyl
readout.

\pandocstatementheading{\leavevmode\phantomsection\label{def:online-finite-cartan-collar}Definition 2.7 (An initialized finite Cartan collar).}

An initialized finite Cartan collar over \(\mathbb G\) consists of
actual smooth fields

\phantomsection\label{eq:online-cartan-temporal-fields} \[
\theta=n\,\mathrm dt+\theta_i\,\mathrm dx^i,
\qquad
A=A_i\,\mathrm dx^i,
\qquad C.
\] The spatial \(\theta_i\) are an oriented basis of \(n^\perp\), the
\(A_i\) are the matrices of a real \(\eta\)-metric connection, and \(C\)
is a real algebraic Weyl tensor. The connection on the product bundle is
\(\widehat D_0=\partial_t\) and \(\widehat D_i=\partial_i+A_i\), with
its induced actual spin connections. Thus its temporal connection
coefficient vanishes in the initialized product gauge. The following
identities are required on the actual collar:

\phantomsection\label{eq:online-cartan-temporal-equations} \[
\partial_t\theta_i=A_i n,
\qquad
\partial_t A_i=\mathcal C(n,\theta_i),
\] \phantomsection\label{eq:online-cartan-mixed-bianchi} \[
\mathcal B_{0ij}
:=
(\widehat D_0\mathcal C)(\theta_i,\theta_j)
+(\widehat D_i\mathcal C)(\theta_j,n)
+(\widehat D_j\mathcal C)(n,\theta_i)
=0.
\] The connection acts on every tensor slot. In particular, for bundle
sections \(u,v\),

\phantomsection\label{eq:online-cartan-evaluated-tensor-derivative} \[
\widehat D_i^{\operatorname{End}}\bigl(\mathcal C(u,v)\bigr)
=(\widehat D_i\mathcal C)(u,v)
+\mathcal C(\widehat D_i u,v)
+\mathcal C(u,\widehat D_i v).
\] The initial values are prescribed by the source:

\phantomsection\label{eq:online-cartan-exact-initial-values} \[
\theta_i(0)=\Theta_i,
\qquad
A_i(0)=D_i-\partial_i,
\qquad
C(0)=C_0.
\] All equations are finite first-order identities on supplied fields.
Their definition uses the independent initial source tests and no
spacetime metric. The Weyl field has ten real algebraic components,
equivalently two Lorentz-related chiral Weyl spinors. Higher curvature
derivatives may subsequently be read from the solution, but are not
additional evolution variables in this definition.

The equations are invariant under time-independent initial coordinate
and adapted spin-frame changes, so this definition patches on initial
coordinate opens. The local collar germs are always taken with their
initialized vector and actual spin identifications.

\pandocstatementheading{\leavevmode\phantomsection\label{prop:online-cartan-subsidiary-system}Proposition 2.8 (The exact finite subsidiary system).}

Suppose fields of these algebraic types satisfy the two temporal
equations. Define their spatial defects by

\phantomsection\label{eq:online-cartan-spatial-defects} \[
\begin{aligned}
\mathcal T_{ij}
&=\widehat D_i\theta_j-\widehat D_j\theta_i,\\
\Omega_{ij}
&=\partial_i A_j-\partial_j A_i+[A_i,A_j],\\
\Delta_{ij}
&=\Omega_{ij}-\mathcal C(\theta_i,\theta_j).
\end{aligned}
\] Then

\phantomsection\label{eq:online-cartan-closed-subsidiary-system} \[
\partial_t\mathcal T_{ij}=\Delta_{ij}n,
\qquad
\partial_t\Delta_{ij}
=-\mathcal B_{0ij}+\mathcal C(n,\mathcal T_{ij}).
\] At each fixed spatial point this is a linear system for thirty real
defect components, with forcing \((0,-\mathcal B)\). If \(L\) bounds its
coefficient matrix in a fixed positive component norm on a retained
compact cylinder, then

\phantomsection\label{eq:online-cartan-forced-defect-estimate} \[
|(\mathcal T,\Delta)(t,x)|
\le e^{L|t|}
\left(
|(\mathcal T,\Delta)(0,x)|
+\int_{\min(0,t)}^{\max(0,t)}
|\mathcal B(u,x)|\,\mathrm du
\right).
\] In particular, the actual mixed Bianchi identities and zero initial
spatial defects imply zero spatial defects throughout the retained
interval.

\emph{Proof.} Differentiate the torsion defect and use the temporal
equations: \[
\begin{aligned}
\partial_t\mathcal T_{ij}
={}&\mathcal C(n,\theta_i)\theta_j
-\mathcal C(n,\theta_j)\theta_i\\
&+\widehat D_i(A_jn)-\widehat D_j(A_in).
\end{aligned}
\] The algebraic first Bianchi identity makes the first line
\(-\mathcal C(\theta_i,\theta_j)n\). Because the components of \(n\) are
constant in the adapted product frame, the second line is
\(\Omega_{ij}n\). This proves the first identity.

For the connection curvature, direct differentiation gives \[
\partial_t\Omega_{ij}
=\widehat D_i^{\operatorname{End}}\bigl(\mathcal C(n,\theta_j)\bigr)
-\widehat D_j^{\operatorname{End}}\bigl(\mathcal C(n,\theta_i)\bigr).
\] Here \(\widehat D_i n=A_i n=\partial_t\theta_i\). Expand each
evaluated tensor using
\hyperref[eq:online-cartan-evaluated-tensor-derivative]{the derivative
rule}, and subtract \[
\partial_t\bigl(\mathcal C(\theta_i,\theta_j)\bigr)
=(\widehat D_0\mathcal C)(\theta_i,\theta_j)
+\mathcal C(\partial_t\theta_i,\theta_j)
+\mathcal C(\theta_i,\partial_t\theta_j).
\] The terms involving \(\partial_t\theta\) cancel, by alternation of
the curvature arguments. The remaining expression is \[
\begin{aligned}
\partial_t\Delta_{ij}
={}&(\widehat D_i\mathcal C)(n,\theta_j)
-(\widehat D_j\mathcal C)(n,\theta_i)
-(\widehat D_0\mathcal C)(\theta_i,\theta_j)\\
&+\mathcal C(n,\widehat D_i\theta_j-\widehat D_j\theta_i),
\end{aligned}
\] which is the second asserted identity.

There are three spatial pairs. Each has a four-component torsion defect
and a six-component Lorentz-skew curvature defect, giving thirty
components in total. The coefficients of their homogeneous system are
finite algebraic functions of the supplied \(C\) and fixed \(n\). The
integral equation and Gronwall's inequality give the displayed estimate
for positive time. Reversing time gives its stated form for negative
time. With zero forcing and zero initial value, uniqueness of this
finite ODE gives the zero solution. \hfill\(\square\)

The mixed Bianchi defect is the forcing term throughout the collar. The
definition requires its vanishing there, a condition strictly stronger
for smooth fields than vanishing of its normal jets on the initial
slice.

\subsection{2.7. Recognition and initialized
comparison}\label{recognition-and-initialized-comparison}\label{sec:online-cartan-recognition-comparison}

\pandocstatementheading{\leavevmode\phantomsection\label{thm:online-cartan-collar-recognition}Theorem 2.9 (Recognition and existence of the initialized collar).}

An initialized finite Cartan collar reconstructs a smooth Lorentzian
Ricci-flat metric with its specified initial data and actual spin
marking. Under its soldering, the supplied connection is the full
Levi--Civita connection and \(C\) is the actual Weyl curvature. The full
differential Bianchi identity and its vacuum divergence equation hold.

Conversely, a smooth vacuum development of the initial data recognized
from \(\mathbb G\) has a Gaussian collar germ with exactly this
initialized finite Cartan presentation. Its germ is unique after the
initial vector and actual spin identifications are fixed. The existence
and geometric uniqueness used here are those of
\hyperref[thm:online-smooth-local-cauchy-input]{the external smooth
Cauchy input}.

The initial value \(C_0\) is determined by the initial source
connection. It contributes no additional freely prescribed
normal-curvature datum.

\emph{Proof.} The intrinsic source recognition gives

\phantomsection\label{eq:online-cartan-initial-spatial-identities} \[
D\Theta=0,\qquad
\Omega_D(X,Y)=\mathcal C_0(\Theta X,\Theta Y).
\] Thus the spatial defects in the preceding proposition vanish
initially. The temporal equations already give \[
(\widehat D\theta)_{0i}
=\partial_t\theta_i-A_i n=0,
\qquad
\Omega_{0i}
=\partial_t A_i=\mathcal C(n,\theta_i).
\] The actual mixed Bianchi equation and the subsidiary proposition now
imply all the remaining components:

\phantomsection\label{eq:online-cartan-full-spacetime-identities} \[
\widehat D\theta=0,
\qquad
\Omega_{\widehat D}=\mathcal C(\theta,\theta).
\] The soldering is an isomorphism because \(\theta(\partial_t)=n\) and
the three \(\theta_i\) are a basis of \(n^\perp\). Define

\phantomsection\label{eq:online-cartan-reconstructed-metric-connection}
\[
g(X,Y)=\eta(\theta X,\theta Y),
\qquad
\nabla_XY=\theta^{-1}\widehat D_X(\theta Y).
\] These are actual smooth fields. Metric compatibility of
\(\widehat D\) gives \(\nabla g=0\), and its torsion is
\(\theta^{-1}\widehat D\theta=0\). Hence \(\nabla\) is the Levi--Civita
connection. The full curvature identity identifies its curvature with
\(C\) under the soldering. The algebraic Weyl conditions imply
\(\operatorname{Ric}(g)=0\). Consequently \(C\) is the actual Weyl
tensor, and the differential Bianchi identity for the Levi--Civita
connection, followed by its Ricci-flat contraction, gives both remaining
assertions about Bianchi.

The metric has the Gaussian form

\phantomsection\label{eq:online-cartan-gaussian-metric} \[
g=\mathrm dt^2-h_t,
\qquad
(h_t)_{ij}=-\eta(\theta_i,\theta_j).
\] Its future normal along the zero slice is \(\partial_t\). Our
second-form convention is \(K(X,Y)=-g(\nabla_X\partial_t,Y)\) there. The
temporal equation gives

\phantomsection\label{eq:online-cartan-initial-second-form} \[
\left.\frac12\partial_t(h_t)_{ij}\right|_{t=0}
=-\frac12\left(
\eta(D_iN,\Theta_j)+\eta(\Theta_i,D_jN)
\right)
=K_{ij}.
\] The last equality uses the symmetry of \(K\) already obtained from
initial torsion. Thus the initial metric, full second form and
orientations are exactly the source readouts. The soldering transports
the product actual spin bundle to the oriented and time-oriented
orthonormal frame bundle of \(g\), with the prescribed initial
reduction. On a smaller domain of dependence this is the required Cauchy
development germ.

For the assertion about \(C_0\), suppose an algebraic Weyl tensor has
zero endomorphism restriction when its first two arguments are spatial.
Pair symmetry then also annihilates components with spatial last pair.
Only components with index pattern \(0a0b\) could remain. Their Ricci
trace-free equations express them as spatial contractions, which already
vanish. Therefore the tensor is zero. The initial spatial curvature and
\(\Theta\) determine at most one Weyl tensor; the intrinsic admission
test supplies its existence. This proves the claimed uniqueness of
\(C_0\) directly.

For existence, decode \(\mathbb G\) by
\hyperref[thm:online-intrinsic-gradient-vacuum-recognition]{intrinsic
initial recognition} and apply the smooth Cauchy input to its vacuum
data. Near a retained initial patch, the normal exponential map gives
Gaussian coordinates after shrinking to avoid a focal point and to
retain an injective coordinate map. The initial Lorentz identification
is

\phantomsection\label{eq:online-cartan-initial-geometric-identification}
\[
J_0(aN+\Theta X)=a\nu+\mathrm d\iota(X),
\] where \(\iota\) is the marked initial embedding and \(\nu\) its
future unit normal. Extend \(J_0\) by Levi--Civita parallel transport
along the normal geodesics. The pullback soldering and connection then
have \(\theta(\partial_t)=n\) and vanishing temporal connection
coefficient. Transport the product actual spin bundle through this
vector identification: its frame projection is the original product
frame projection followed by \(J\). This equips the collar with the
required actual spin structure and its initial reduction. If the
development already carries an actual spin extension with the prescribed
initial identification, path lifting along the normal geodesics gives
the unique compatible spin-bundle identification. In either description
the induced spin connection has vanishing temporal coefficient.

Put \(C\) equal to the actual Weyl tensor in this identification.
Torsion-freeness gives the first temporal equation, and the curvature
identity gives the second. The actual differential Bianchi identity
gives the mixed equation. At zero, the full initial tangent connection
agrees with \(D\) by the Gauss--Weingarten formula and equality of
\(h,K\). Its curvature therefore has the same spatial restriction, so
the unique Weyl readout just proved gives \(C(0)=C_0\). Every initial
value is exactly the stipulated one.

For uniqueness, take two initialized Cartan germs over \(\mathbb G\).
Their recognized metrics induce the same vacuum data on a buffered
initial open set. Restrict to common Cauchy neighborhoods of a smaller
retained patch and apply geometric Cauchy uniqueness. The resulting
isometry fixes the initial embedding and carries its future normal to
the future normal. It therefore carries each normally parametrized
geodesic to its counterpart, with the same proper-time parameter. In the
Gaussian coordinates its point map is the identity on a smaller collar
germ. The two metrics agree there.

Their initialized parallel identifications solve the same finite
parallel-transport equation for this metric, so the solderings agree.
Levi--Civita uniqueness then identifies the connections, and curvature
identifies \(C\). The two actual spin identifications agree as lifts of
the same initialized frame identification: their difference is a locally
constant central sign, fixed on each initial component. This proves
uniqueness of the entire initialized source germ. \hfill\(\square\)

The smooth Einstein Cauchy theorem supplies a collar realizing the
initial data. Conversely, the Cartan identities recognize the vacuum
metric of a supplied collar. The homogeneous subsidiary system
propagates vanishing of the spatial torsion and curvature defects from
the initial slice when the mixed Bianchi identity holds throughout the
collar.

A comparison between initialized Cartan collars consists of a local
point diffeomorphism \(F\) carrying the marked initial hypersurface to
its counterpart, an oriented and time-oriented Lorentz bundle isometry
\(L\) over \(F\), and an actual compatible spin lift \(\widetilde L\).
Require the full identities

\phantomsection\label{eq:online-cartan-physical-comparison} \[
\begin{aligned}
\theta'\circ\mathrm dF&=L\circ\theta,\\
\widehat D'_{F_*X}(Lv)&=L(\widehat D_Xv),\\
\mathcal C'(Lu,Lv)Lw&=L\bigl(\mathcal C(u,v)w\bigr).
\end{aligned}
\] The fields on the right and left are evaluated at corresponding
points. These conditions are imposed directly on supplied fields. Their
restriction is an intrinsic initial gradient comparison, with its actual
spin map. Germs and comparisons are restricted to a common collar
whenever they are composed.

\pandocstatementheading{\leavevmode\phantomsection\label{thm:online-marked-collar-full-hom}Theorem 2.10 (Full comparison for a marked initial hypersurface).}

Restriction gives a bijection between actual Cartan collar comparisons
and intrinsic initial gradient comparisons over each specified initial
diffeomorphism. Every initial comparison extends uniquely. These
assignments preserve identities, inverse maps, composition and
restriction.

Reconstruction also identifies these comparisons with all
orientation-preserving and time-orientation-preserving spin isometries
of the recognized vacuum germs that carry their marked initial
hypersurfaces to one another. Thus initialized finite Cartan collar
germs, intrinsic initial gradient sources and their marked vacuum
development germs form equivalent groupoids.

\emph{Proof.} A Cartan comparison preserves \(g=\eta(\theta,\theta)\) by
the soldering identity. It therefore induces an isometry of the
recognized metrics. On the initial slice its connection and soldering
identities recover the initial gradient intertwining relation by
\hyperref[lem:online-intrinsic-gradient-decoder]{the exact intrinsic
decoder}. The stipulated spin lift restricts to the actual initial lift,
so restriction is well defined.

Conversely, let an intrinsic source comparison cover \(f:U\to U'\), with
induced Lorentz map \(L_0\) and actual spin map \(\widetilde L_0\). By
\hyperref[thm:online-initial-gradient-full-hom]{initial full Hom}, it is
exactly a physical comparison of the decoded \(h,K\) and initial spin
structures. Regard the target metric as a development of the source data
using the initial embedding \(\iota'f\). Geometric Cauchy uniqueness, on
the buffered domains just used, supplies an isometry germ \(F\)
extending this identification. It preserves the future normal.
Normal-geodesic uniqueness gives

\phantomsection\label{eq:online-cartan-comparison-gaussian-point-map} \[
F(t,x)=(t,f(x))
\] in the two Gaussian presentations. Parallel transport of \(L_0\) and
\(\widetilde L_0\) along these geodesics gives the full vector and spin
comparisons.

In the initialized product frames the transported maps are independent
of time. Their complete coefficient identities are

\phantomsection\label{eq:online-cartan-comparison-full-coefficients} \[
\begin{aligned}
\theta'_j(t,f(x))\,\partial_i f^j
&=L_0(x)\theta_i(t,x),\\
A'_j(t,f(x))\,\partial_i f^j
&=L_0A_i(t,x)L_0^{-1}
-(\partial_iL_0)L_0^{-1}.
\end{aligned}
\] The first follows from the soldering and normally parallel
identifications; the second is the full connection transformation law,
including the derivative of the comparison. Curvature gives the
displayed \(\mathcal C\) identity. The induced actual spin connection
transforms by the corresponding principal spin law, with the same
prescribed \(\widetilde L_0\). These formulas prove that the extension
satisfies every source comparison condition.

An isometry carrying the marked initial hypersurfaces and preserving
time orientation must carry their future normal geodesics with their
proper-time parametrizations. Hence its point germ is already determined
by its initial map. The full soldering identity then determines its
Lorentz bundle map. Two compatible spin lifts differ by a locally
constant central sign, and their initial lift fixes that sign on every
component. This proves uniqueness and faithfulness. It also shows that
every spin isometry of marked vacuum germs has precisely the source
comparison just constructed.

The extension of an identity is the identity. The composite of two
extensions has the prescribed composite initial comparison, so
uniqueness makes it the extension of that composite; the same argument
gives inverses and compatibility with restriction. The source-to-metric
and metric-to-source object constructions are inverse up to the
initialized identifications in the recognition theorem. The
data-to-source unit is the intrinsic unit, and the collar counit is the
initialized normal-geodesic and parallel-spin identification. They
commute with every comparison by uniqueness; their composites are the
identity on the corresponding vector and spinor bundles. This proves the
equivalences and their triangle identities. Distinct initial central
spin lifts remain distinct throughout. \hfill\(\square\)

The marked hypersurface is part of this Hom statement. Comparison
between displays based on different hypersurfaces will use the later
unmarked source comparison theorem.

\subsection{2.8. Ordinary displays and physical
scale}\label{ordinary-displays-and-physical-scale}\label{sec:online-ordinary-collars-and-scale}

\pandocstatementheading{\leavevmode\phantomsection\label{prop:online-cartan-physical-scale}Proposition 2.11 (Physical scaling of the actual collar).}

Let \(r>0\) be constant on a given collar and put
\(\psi_r(t',x)=(t'/r,x)\). On the corresponding rescaled time domain
define

\phantomsection\label{eq:online-cartan-physical-scale-fields} \[
\theta^{(r)}=r\,\psi_r^*\theta,
\qquad
A^{(r)}=\psi_r^*A,
\qquad
C^{(r)}=r^{-2}\psi_r^*C,
\] with the pulled-back actual spin bundle and its initialized marking.
These fields satisfy the same finite Cartan equations. Their recognized
metric and initial data are

\phantomsection\label{eq:online-cartan-physical-scale-metric-data} \[
g^{(r)}=r^2\psi_r^*g,
\qquad
h^{(r)}=r^2h,
\qquad
K^{(r)}=rK.
\] With the fixed internal Clifford metric and unit normal, the initial
chiral gradients obey

\phantomsection\label{eq:online-cartan-physical-scale-gradients} \[
\mathbb D_\pm^{(r)}=r^{-1}\mathbb D_\pm.
\] For an ordinary display, first restrict the already constructed
complete initial contact presentation to its retained patch. In the
fixed cotangent-ray coordinates of
\hyperref[prop:online-contact-complete-kinematic-codec]{Proposition
4.15}, its physical-scale coefficients satisfy

\phantomsection\label{eq:online-cartan-physical-scale-contact} \[
\beta^{(r)}=\beta,\qquad
\eta_c^{(r)}=\eta_c,\qquad
\omega^{(r)}=\omega+3\log r,
\qquad
\nu^{(r)}=r^3\nu.
\] The local coefficient laws retain the maximal marking \(\tau=0\). The
whole-space normalization relative to \(\delta\) applies before this
change of scale and need not hold for \(r^2h\).

The construction is compatible with restriction, initialized actual spin
comparisons and successive positive scale changes.

\emph{Proof.} In the new clock,
\(\theta^{(r)}=n\,\mathrm dt'+r\theta_i(t'/r,x)\,\mathrm dx^i\), and
\(A_i^{(r)}=A_i(t'/r,x)\). Therefore \[
\partial_{t'}\theta_i^{(r)}=A_i^{(r)}n,
\qquad
\partial_{t'}A_i^{(r)}
=r^{-1}\mathcal C(n,\theta_i)
=\mathcal C^{(r)}(n,\theta_i^{(r)}).
\] In the mixed Bianchi equation the time derivative of \(C^{(r)}\)
contributes \(r^{-3}\) and its two spatial soldering arguments
contribute \(r^2\). Each spatial derivative term has factor \(r^{-2}\)
from \(C\) and factor \(r\) from its spatial argument. Thus the entire
mixed defect is \(r^{-1}\psi_r^*\mathcal B\) and vanishes. Weyl
symmetries, reality and the initial spatial identities are preserved.

The metric formula follows by applying \(\eta\) to the rescaled
soldering. Differentiating \(h^{(r)}_{t'}=r^2h_{t'/r}\) gives the
displayed initial \(K\) law. At zero, the connection itself is unchanged
in the fixed internal frame, while an orthonormal tangent vector changes
from \(e_a\) to \(r^{-1}e_a\) and the soldering changes from \(\Theta\)
to \(r\Theta\). Consequently \[
\sum_a\Theta^{(r)}(r^{-1}e_a)
\otimes P^\pm_{r^{-1}e_a}
=r^{-1}\sum_a\Theta(e_a)\otimes P^\pm_{e_a},
\] which proves the gradient law.

For the local contact formulas, \(A=h(z,z)\), \(B=h(z,\bar z)\) and
\(\Delta_c\) each acquire a factor \(r^2\), so \(\beta\) and
\(\ell=z-\beta\bar z\) are unchanged. Also \(\rho_n\) acquires
\(r^{-1}\), whereas \(K\) acquires \(r\) and the Christoffel symbols are
unchanged. The complete shear formula therefore gives
\(\eta_c^{(r)}=\eta_c\). Finally \(\det(r^2h)=r^6\det h\) gives the
stated \(\omega\) and density laws. The initial mean curvature remains
zero.

Pullback carries the supplied spin identification in the same internal
frames. The even tensor coefficients alone cannot distinguish its
central sign. The formulas compose under multiplication of positive
scales and commute with the initialized comparisons, proving the
remaining claims. \hfill\(\square\)

\pandocstatementheading{\leavevmode\phantomsection\label{cor:online-ordinary-pair-cartan-collar}Corollary 2.12 (The actual collar of an ordinary-pair display).}

Let \(c\) belong to the ordinary domain in
\hyperref[thm:online-ordinary-negative-response-initial-source]{Theorem
4.18}, with its fixed physical and actual spin markings. Its constructed
initial gradient source has an initialized finite Cartan collar germ.
All initial collar values are read from that source, and the recognized
metric is a smooth Ricci-flat development of its exactly recovered data.

Conversely, suppose a supplied initialized Cartan collar has an initial
source satisfying the independent canonical tests and retained-domain
conditions of Theorem 4.18. The ordinary class

\phantomsection\label{eq:online-cartan-canonical-ordinary-decoder} \[
c=\sigma_\gamma\mathscr P_\gamma^{-1}\pi\mathbb G
\] reconstructs that initial source exactly. Its constructed collar is
identified with the supplied collar by the unique comparison inducing
the identity on the initial source and its actual spin bundle.

Effective ordinary cochain comparisons give these initialized collar
comparisons. The physical scale \(r\) is retained by the preceding
proposition. On a smaller initial open set, the collar construction uses
restriction of the already constructed source; it does not apply the
whole weighted ordinary decoder to a spatially truncated field.

\emph{Proof.} Theorem 4.18 gives an actual initial source with
independent kinematic and vacuum admission, exact ordinary-class
recovery, and the specified spin marking. Its positive-frame
presentation is a special case of the intrinsic source definition. The
intrinsic recognition theorem provides \(\Theta,D,C_0\) and the vacuum
data. The collar recognition theorem then gives an actual initialized
Cartan germ. Its only initial values are those finite readouts; smooth
Cauchy existence supplies the geometric evolution from them.

For the converse, Theorem 4.18 applies to the initial source and the
displayed decoder. It reconstructs the same complete initial operator
pair, or equivalently its complete canonical contact presentation with
scale and spin. The uniqueness part of collar recognition identifies the
entire initialized finite-field germ. Initial full Hom and its extension
theorem give the same conclusion for the effective ordinary comparisons.
In particular the initial ordinary-class unit persists after passing to
the collar.

The scale formulas apply directly to these actual fields. Restriction of
the finite Cartan equations and their initialized comparison is local,
and uniqueness identifies the restricted construction. The inverse in
the displayed ordinary formula remains the whole weighted inverse on its
stated canonical domain. This proves all assertions. \hfill\(\square\)

This construction supplies the actual geometric collar used by a local
display. The full wave section and its point connection, and the
comparison of displays with different initial hypersurfaces, retain
their separate admission and comparison statements.
\subsection{2.9. Higher estimates for the same reduced
solution}\label{higher-estimates-for-the-same-reduced-solution}\label{sec:online-reduced-family-estimates}

Write \(\mathcal P\subset\mathbb R^d\) for a parameter domain and use
fixed local coefficient trivializations. A smooth family means joint
smoothness of its supplied coefficients in parameter and spatial
variables. All compact parameter estimates below use
\(\mathcal K\Subset\operatorname{int}\mathcal K'\Subset\mathcal P\). The
input seminorm \(\mathcal N_N\) is the maximum of a finite collection of
spatial and parameter derivatives up to order \(N\) on the stated larger
compact set, including the supplied coordinate and frame markings.

\pandocstatementheading{\leavevmode\phantomsection\label{lem:online-reduced-cauchy-finite-family-estimates}Lemma 2.13 (One reduced solution family at every finite order).}

For a smooth local family of initial data near any fixed parameter
value, the localized reduced construction in
\hyperref[thm:online-smooth-local-cauchy-input]{the Cauchy input} has
one parameter neighborhood and one positive time interval chosen at the
finite base order. On every fixed buffered compact part of that
interval, the selected solutions are jointly smooth. Every finite
spacetime and parameter seminorm is bounded by a finite locally bounded
function of finitely many input seminorms and the inverses of the
retained hyperbolicity and domain-buffer margins.

For two such families in the same fixed gauge and cutoff construction,
write \(d,d'\) for their full reduced Cauchy data. On a common retained
compact domain, for each \(m,q\) there are a finite \(N=N(m,q)\) and a
finite locally bounded function \(G_{m,q}\) such that

\phantomsection\label{eq:online-reduced-family-difference-bound} \[
\|g-g'\|_{C^q_\lambda C^m_{t,x}}
\le
G_{m,q}\bigl(\mathcal N_N(d),\mathcal N_N(d'),
\mathfrak m^{-1}\bigr)\,
\mathcal N_N(d-d').
\] Here all input norms use the larger compact parameter and spatial
sets, and \(\mathfrak m\) records the retained positive margins. At each
output order the estimate uses a finite, possibly higher, input order.

\emph{Proof.} Fix the cutoff and the reference extension. The initial
reduced fields depend continuously on the parameter in the base
\(H^{s_0+1}\times H^{s_0}\) norm. Shrink its neighborhood once so that
they lie in one existence neighborhood of the external theorem. Its
common time interval and a smaller retained interval are then fixed. The
base solution remains in a set with uniform positive hyperbolicity and
bounded low-order coefficients.

For completeness, the estimates behind the higher and parameter
assertions can be obtained on this same solution. For a linear equation
with the reduced wave principal part, let

\phantomsection\label{eq:online-linearized-reduced-wave} \[
\mathcal L_g v
=g^{\alpha\beta}\partial_\alpha\partial_\beta v
+B^\alpha\partial_\alpha v+Dv=f,
\] where the matrices \(B^\alpha,D\) have the required finite
coefficient bounds. For an integer \(n\ge0\), a spatially differentiated
energy is equivalent to \(\|v(t)\|_{H^{n+1}}+\|\partial_tv(t)\|_{H^n}\).
One may use the square root of

\phantomsection\label{eq:online-reduced-wave-energy} \[
\sum_{|\beta|\le n}\int_{\mathbb R^3}
\left(
g^{00}|\partial_t\partial_x^\beta v|^2
-g^{ij}\partial_i\partial_x^\beta v\cdot
          \partial_j\partial_x^\beta v
+|\partial_x^\beta v|^2
\right)\mathrm dx.
\] In the retained coordinates both quadratic forms are positive.
Multiplication of the commuted equation by
\(\partial_t\partial_x^\beta v\), followed by integration by parts,
gives

\phantomsection\label{eq:online-reduced-linear-energy-bound} \[
E_n(v;t)\le
C_n\left(
E_n(v;0)+\int_0^{|t|}
\|f(\pm u)\|_{H^n}\,\mathrm du
\right).
\] The constant uses the fixed time interval, positive margins and
finitely many coefficient norms. A spatial cutoff or cone version gives
the same estimate on retained nested domains. The terms involving
\(g^{0i}\) are flux terms and terms involving its first derivatives. The
remaining commutators are bounded by the product rule, Sobolev
multiplication and interpolation. Smooth approximation justifies the
estimate for solutions at the stated finite regularity.

For the nonlinear reduced equation put \(u=g-g_{\mathrm{ref}}\), where
\(g_{\mathrm{ref}}\) is the fixed extended initial reference. Its
principal coefficients depend on \(g\) and not on \(\partial g\). In the
commuted equation the term with the largest derivative of a coefficient
is a derivative of \(g\) multiplied by a bounded low-order derivative of
\(g\); the other products admit the same tame bound. Thus, on the fixed
low-order neighborhood, the high energy satisfies an inequality of the
form

\phantomsection\label{eq:online-reduced-higher-energy-bound} \[
E_n(u;t)\le
C_n\left(E_n(u;0)+
\int_0^{|t|}\bigl(1+E_n(u;\pm u')\bigr)\,\mathrm du'\right),
\] with fixed-reference terms included in \(C_n\). Applying the estimate
to regularized initial data, taking the base limit by uniqueness, and
using weak compactness and the estimate proves the bound for the already
selected solution. The usual difference estimate then gives strong
continuity at each lower retained order. Equivalently, the persistence
clause of the external input identifies the high-regularity solution
with this same base solution. Gronwall's inequality supplies each finite
bound on the same time interval. No high-order smallness is introduced.

To justify parameter differentiation, differentiate the reduced equation
formally once and move all terms linear in the highest parameter
derivative to the left. For \(v=\partial_{\lambda^a}g\) its principal
part is \(\mathcal L_g\); the variation of \(g^{\alpha\beta}\) times
\(\partial_{\alpha\beta}g\) is a zeroth-order matrix acting on \(v\).
The variation of the nonlinear first-derivative terms contributes first-
and zeroth-order matrices. Their finite norms are controlled by higher
spatial bounds just proved.

This formal equation identifies an actual derivative. Difference
quotients
\(v_\varepsilon=(g_{\lambda+\varepsilon e_a}-g_\lambda)/\varepsilon\)
satisfy linear equations with averaged coefficients and uniformly
bounded finite data. The linear energy bound controls them at any fixed
desired order after finitely many stronger input norms have been
retained. Base continuity and uniform higher spatial bounds imply
coefficient convergence in every required finite lower norm by Sobolev
interpolation. Subtract the limiting linear equation from the quotient
equation. The remaining initial data and forcing tend to zero, so the
same energy estimate makes the difference tend to zero. Existence for
the linear equation follows by regularization and this energy estimate,
or by the linear case of the stated hyperbolic input. Thus \(v\) is the
actual parameter derivative.

For a multi-index \(\alpha\), the equation for
\(\partial_\lambda^\alpha g\) has the same linear principal operator.
Its forcing is a finite sum of products involving only lower parameter
derivatives, each with at most two additional spacetime derivatives.
Induction on \(|\alpha|\), allowing these finite spatial derivative
costs, proves existence, continuity and finite estimates for every
specified mixed derivative. Its difference-quotient justification is the
same. Compact containment of \(\mathcal K\) in the interior of
\(\mathcal K'\) keeps all quotients inside the supplied parameter set.

Time derivatives are recovered successively from the equation because
\(g^{00}\) has a positive lower bound. Sobolev embedding then gives
every specified mixed \(C^m\) estimate on a retained compact region. For
two families, subtraction of their equations and repeated parameter
differentiation give linear difference equations with coefficients
controlled by both families and forcing proportional to finitely many
input differences. The same argument proves the final assertion. Every
step consumes finitely many stronger norms; the base interval was fixed
before their orders. \hfill\(\square\)

\subsection{2.10. Common Gaussian collars and their finite-family
bounds}\label{common-gaussian-collars-and-their-finite-family-bounds}\label{sec:online-common-collar-families}

\pandocstatementheading{\leavevmode\phantomsection\label{thm:online-common-collar-finite-families}Theorem 2.14 (Common collars for smooth finite-dimensional families).}

Let a smooth finite-dimensional family of independently admitted
intrinsic initial gradient sources be supplied on common coordinate and
spin representatives. Near each parameter value and over each retained
compact initial patch, there are one parameter neighborhood, a positive
\(T\), and buffered spatial parents such that their recognized finite
Cartan collars are defined on a common parent \((-2T,2T)\times W\).

The metric, soldering, connection, actual Weyl field, initialized
normal-geodesic map and initialized actual spin parallel transport are
jointly smooth there. For every \(m,q<\infty\), their coefficient
seminorms on \(\mathcal K\times[-T,T]\times\overline V\), with
\(V\Subset W\), satisfy

\phantomsection\label{eq:online-collar-finite-family-bound} \[
\begin{multlined}
\|\mathcal Y\|_{m,q}
+\|\Phi\|_{m,q}
+\|\Phi^{-1}\|_{m,q}
\\
\le
F_{m,q}\left(
\mathcal N_N(\text{initial source and markings};
\mathcal K'\times\overline U),\,
\mathfrak m^{-1}
\right)
\end{multlined}
\] for some finite \(N\) and a finite locally bounded function
\(F_{m,q}\). The map and inverse norms use their respective retained
source and image coordinate buffers. Here \(\mathcal Y\) is any fixed
finite list of the preceding coefficient fields and transports. The
margins \(\mathfrak m\) include coordinate and symbol nondegeneracy,
hyperbolicity, spacelikeness, domain buffers and the retained Gaussian
inverse margin. The parameter neighborhood and parent are chosen before
\(m,q\).

A supplied smooth family of initial physical comparisons with actual
spin lifts extends to a smooth family of the marked collar comparisons,
after one common buffered restriction. It has the same kind of finite
estimates, with the supplied comparison, inverse and spin-marking
seminorms included. A supplied smooth positive scale is allowed, with
its finite derivatives and a positive lower bound included.

\emph{Proof.}
\hyperref[thm:online-intrinsic-gradient-vacuum-recognition]{Intrinsic
recognition} reads \((h,K)\) and the actual initial spin structure by
finite differential and matrix operations. On each retained compact set,
symbol nondegeneracy gives one positive margin. These operations
therefore preserve the stipulated family and supply the finite bounds
for its local Cauchy data. Apply the preceding lemma to the fixed
localized reduced gauge. On the smaller domain of dependence its metric
is Ricci-flat by the gauge-propagation argument.

Solve the normal geodesic equation for this actual reduced metric, with
initial position \((0,x)\) and the future unit initial normal. On the
chosen initial gauge this velocity is \(\partial_t\). Its coefficients
use the first metric derivatives, and its derivatives obey finite
variational ODEs. Low-order bounds and the spatial buffer give a common
interval for all retained \(x\) and parameters. At time zero the
spacetime differential of the map

\phantomsection\label{eq:online-initialized-gaussian-family-map} \[
\Phi_\lambda(t,x)=
\exp^{g_\lambda}_{(0,x)}
\bigl(t\,n_\lambda(0,x)\bigr)
\] is invertible. Shrinking the time interval once gives a uniform
inverse margin. In sufficiently small convex coordinate patches its
differential remains uniformly close to the initialized one; the
mean-value estimate then gives injectivity. For a general retained
compact patch, take a finite cover by such smaller patches with buffered
parents. By its Lebesgue number, sufficiently close initial base points
belong to one injective parent. Pairs of base points separated by at
least that number remain separated for small time, since
\(\Phi_\lambda(t,x)\) tends uniformly to \((0,x)\) as \(t\to0\).
Compactness makes this choice uniform over the retained parameters. Thus
the common restricted collar map is injective. On overlaps the maps
describe the same initialized normal geodesics, so they give the same
collar germ.

Retain a larger parent than the final compact observation set before
differentiating. The ODE variation equations give finite mixed bounds
for \(\Phi\). The inverse margin and the spatial buffers also give
common image buffers. Differentiating
\(\Phi^{-1}\circ\Phi=\operatorname{id}\) gives the inverse bounds from
that margin. Pullback gives the Gaussian metric. Parallel transport
along these normal geodesics identifies the tangent Lorentz bundle with
its product extension and gives the temporal gauge \(A_0=0\). Principal
spin parallel transport with the supplied initial identification gives
its actual spin lift. These finite-dimensional ODEs have unique
solutions, and their variations obey linear ODEs with forcing involving
finitely many earlier derivatives. Their finite estimates follow
successively from Gronwall's inequality. The initialized spin lift
retains its initial central sign.

In that identification, \(\theta\) is the metric soldering, \(A\) the
Levi--Civita connection and \(C\) its Weyl curvature. Reading them uses
finitely many derivatives of the already constructed fields.
\hyperref[thm:online-cartan-collar-recognition]{Cartan recognition}
identifies this presentation with the independently admitted finite
Cartan collar. This proves joint smoothness and the displayed bound. All
higher estimates use the same low-order-selected reduced solution,
geodesic domain and parallel-transport parent.

For a supplied initial comparison \((f_\lambda,\widetilde f_\lambda)\),
\hyperref[thm:online-marked-collar-full-hom]{marked collar comparison}
has Gaussian point map \[
(t,x)\longmapsto(t,f_\lambda(x)).
\] In the reduced-coordinate representatives it is
\(\Phi'_\lambda\circ(\operatorname{id},f_\lambda)\circ\Phi_\lambda^{-1}\).
If \(P_\lambda\) and \(P'_\lambda\) denote the initialized principal
spin transports, its actual spin map is

\phantomsection\label{eq:online-collar-comparison-spin-family} \[
\widetilde F_\lambda(t,x)
=P'_\lambda(t,f_\lambda(x))\,
\widetilde f_\lambda(x)\,
P_\lambda(t,x)^{-1}.
\] The supplied initial images and inverse margins permit one common
restriction of the parents. Composition, inversion on this retained
nondegenerate domain and the finite ODE estimates give the asserted
comparison bounds. Uniqueness gives their identity and composition laws.
The comparison is chosen by the supplied initial physical and spin data.

Finally apply \hyperref[prop:online-cartan-physical-scale]{physical
scale transport}. For a parameter-dependent \(r>0\), its physical clock
is \(t'=rt\). On a retained compact parameter set, a positive lower
bound for \(r\) gives one common physical time parent. The chain rule
for \(t'/r\), the factors \(r\) in the soldering and \(r^{-2}\) in the
internal Weyl tensor use only finitely many derivatives of \(r\) and its
reciprocal at each order. The connection and initialized spin transport
obey the same pullback law. This proves the scaled family assertion with
the stated finite dependence. \hfill\(\square\)

The theorem applies in particular to the ordinary-pair initial sources
of
\hyperref[thm:online-ordinary-negative-response-initial-source]{Theorem
4.18}, by composition with that theorem's finite-family estimates. Its
present output is the finite geometric collar and its marked
comparisons. The full wave kernel, evaluation section and all-direction
point connection have their own subsequent constructions. On a
noncompact initial manifold these conclusions are local over retained
compact patches; their time widths need not have a positive global lower
bound.

\subsection{2.11. Exterior, gauge and scale choices on one initial
germ}\label{exterior-gauge-and-scale-choices-on-one-initial-germ}\label{sec:online-fixed-slice-geometric-choices}

The ordinary input is a whole pair, while its physical display retains
only a local source. We now compare the initial and finite Cartan parts
of two such displays directly. Their comparison conditions are the
source identities already defined, and do not require equality of the
whole pairs.

\pandocstatementheading{\leavevmode\phantomsection\label{cor:online-fixed-slice-geometric-choice-independence}Corollary 2.15 (Geometric choice independence on a fixed initial germ).}

Let \(d=(O,h,K,P_O)\) be a physical maximal vacuum initial germ. For
\(a=0,1\), let \(c_a\) be an ordinary pair in the selected domain of
\hyperref[thm:online-ordinary-negative-response-initial-source]{Theorem
4.18}. Execute the pair on its whole domain and restrict its actual
initial source to \(V_a\). Supply a positive scale \(r_a\), constant on
that spacetime display, and a physical embedding \(\iota_a:V_a\to O\),
with an actual spin identification \(\Psi_a\) over it, such that

\phantomsection\label{eq:online-fixed-slice-physical-embeddings} \[
\iota_a^*h=r_a^2h(c_a)|_{V_a},\qquad
\iota_a^*K=r_aK(c_a)|_{V_a}.
\] On a common germ in the two images, the scaled initial gradients and
initialized Cartan collars have a unique physical comparison inducing
the common initial identification and its specified actual spin map. Its
initial point and spin maps are

\phantomsection\label{eq:online-fixed-slice-comparison-from-markings} \[
f_{01}=\iota_1^{-1}\iota_0,\qquad
\widetilde U_{01}=\Psi_1^{-1}\Psi_0.
\] The initial complete contact graph and physical density are also
carried by the induced positive cotangent-ray map. These geometric
comparisons have exact restriction, identity, inverse and cocycle laws.
They include changes of exterior extension, based spatial gauge,
coordinate or frame normalization, cutoff and scale that preserve the
displayed physical initial germ.

For smoothly supplied finite-dimensional families of such displays and
embeddings, the comparisons are smooth locally in the parameter on a
common retained parent. Each prescribed finite derivative norm is
bounded by finitely many norms of the supplied constructions and
embeddings, with positive inverse and buffer margins. The positive
scales may depend on the passive parameters. This assertion uses the
initialized local family choices; it does not select isometries or spin
signs over an arbitrary parameter space.

\emph{Proof.} Put \(\mathbb G_a^{(r_a)}=r_a^{-1}\mathbb G_a\) for the
two uncontracted gradient operators, keeping their algebraic Clifford
markings and unit normals.
\hyperref[prop:online-cartan-physical-scale]{Physical scale transport}
gives their readouts as \[
h_a^{\mathrm{phys}}=r_a^2h(c_a),\qquad
K_a^{\mathrm{phys}}=r_aK(c_a).
\] The supplied embeddings imply
\(f_{01}^*h_1^{\mathrm{phys}}=h_0^{\mathrm{phys}}\) and
\(f_{01}^*K_1^{\mathrm{phys}}=K_0^{\mathrm{phys}}\). Together with
\(\widetilde U_{01}\) these are a physical initial-data comparison.
\hyperref[thm:online-initial-gradient-full-hom]{Initial full Hom} gives
its exact source intertwiner. In particular, if \(U_{01,\pm}\) are its
chiral actions and \(L_{01}\) its Lorentz vector map, then

\phantomsection\label{eq:online-fixed-slice-scaled-source-intertwining}
\[
(r_1^{-1}\mathbb D_{1,\pm})\widehat U_{01,\pm}
=\widehat{L_{01}\otimes U_{01,\pm}}
  (r_0^{-1}\mathbb D_{0,\pm}).
\] This is precisely the independently defined operator comparison,
including its zero-order coefficients. The marking determines the actual
spin action; passage to even tensor coefficients has not identified its
two central lifts.

Apply \hyperref[thm:online-marked-collar-full-hom]{marked collar full
Hom} after restoring the physical clocks by the scale proposition. It
extends the initial comparison uniquely to the finite Cartan fields. In
their physical Gaussian coordinates its point map is

\phantomsection\label{eq:online-fixed-slice-physical-clock-comparison}
\[
F_{01}(t,x)=(t,f_{01}(x)).
\] The vector and spin maps are the initialized normal parallel
transports of the initial ones. Thus the extension preserves the
physical clock, soldering, full connection and actual Weyl field. This
applies on buffered Cauchy neighborhoods and hence to the retained
collar germs.

For the initial contact assertion, the induced map is

\phantomsection\label{eq:online-fixed-slice-cotangent-ray-comparison} \[
\widehat f_{01}(x,[\xi])
=\left(f_{01}(x),
  [((\mathrm df_{01,x})^{-1})^*\xi]\right).
\] The cotangent lift preserves the tautological form before positive
projectivization, so it preserves the projective contact distribution.
In the intrinsic description of
\hyperref[prop:online-contact-complete-kinematic-codec]{Proposition
4.15}, the graph is defined from the oriented metric screen, its cross
product, the Levi--Civita horizontal splitting and the second form. An
oriented \(h\)-isometry preserves the screen and cross product, and its
derivative carries the Levi--Civita splitting to the target splitting.
Preservation of \(K\) then carries the shear term as well. It therefore
intertwines the whole graph, including the terms arising from the
derivative of the coordinate transformation. The density comparison is
\(f_{01}^*\operatorname{vol}_{h_1^{\mathrm{phys}}}
=\operatorname{vol}_{h_0^{\mathrm{phys}}}\). The contact coefficient
formulas are used after restriction and physical scaling; no claim is
made that the scaled local tensors remain in the original whole weighted
chart.

For an explicit exterior example, let two whole extensions \((H_a,L_a)\)
coincide with the same physical data on \(O\) and let their based gauges
be \(\phi_a\). On the common local images take
\(f_{01}=\phi_1^{-1}\phi_0\). Then

\phantomsection\label{eq:online-fixed-slice-exterior-gauge-identity} \[
f_{01}^*(\phi_1^*H_1,\phi_1^*L_1)
=(\phi_0^*H_0,\phi_0^*L_0).
\] Initialized actual spin lifts give the stipulated spin comparison.
Additional local coordinate, frame and scale choices are already
included by the embeddings and scale factors above. This proves the
stated applicability without comparing whole Green kernels or
recomputing a whole decoder after truncation.

For three displays, the actual formulas give \(f_{12}f_{01}=f_{02}\) and
\(\widetilde U_{12}\widetilde U_{01}=\widetilde U_{02}\). Initial full
Hom preserves these compositions, and uniqueness of the initialized
Cartan extension gives the same equality on their common collars. The
cotangent-ray maps compose by the ordinary chain rule. The same argument
proves identities, inverses and compatibility with restriction.

For families, first retain compact spatial and parameter sets inside the
supplied larger parents. The inverse-function bounds for \(\iota_a\),
the positive scale bounds and the supplied spin identifications control
each finite derivative of \(f_{01}\) and \(\widetilde U_{01}\).
\hyperref[prop:online-initial-gradient-descent-families]{The finite
initial estimates} and
\hyperref[thm:online-common-collar-finite-families]{the common-collar
family theorem} then control the induced fields and comparisons on one
retained parent. The contact comparison is the displayed cotangent
formula and the preceding finite tensor construction, so it has the same
finite-order dependence. This proves the family assertion.
\hfill\(\square\)

For the initial and finite Cartan parts, the local inverse retains the
source-comparison class of an ordinary display. The two whole pairs can
differ. Also, changing the radius used to select a maximal disk can
change the initial hypersurface itself. Such a change is outside this
fixed-hypersurface result unless its physical initial germs actually
agree. Full wave comparisons and changes of initial hypersurface will be
treated with their additional source structures.

\section{3. Ordinary classes and the initial-data
construction}\label{ordinary-classes-and-the-initial-data-construction}\label{sec:online-ordinary-input-construction}

Each local display starts from two ordinary classes in the smooth
coefficient complexes on the fixed flat Cauchy source. Their root values
give linear coordinates for the free transverse-traceless seeds. A
selected nonlinear constraint correction then gives whole maximal vacuum
data, with an exact inverse on an independently specified tensor domain.

Three distinctions determine the construction. The linear root complex
is a polynomial quotient, whereas nonlinear coefficient operations use
actual representatives. The actual cohomology has polynomial classes
that the root map does not see; explicit obstruction tests select the
ordinary domain. Finally, the smallness condition is imposed at one
finite base order, and all higher weighted estimates concern that same
selected solution.

The weighted analytic proofs are given in
\hyperref[app:online-weighted-construction]{Appendix A}.
\hyperref[thm:online-ordinary-pair-data-bijection]{Theorem 3.13}
supplies the ordinary input and tensor-decoder foundation.
\hyperref[sec:online-actual-source-execution]{Section 4} constructs the
actual initial source from those classes by normalized negative
responses and current input transport. Section 5 uses this initial
source in the local coverage construction, with its stated Cauchy and
wave-completion inputs.

The classical linear Penrose transform relates twistor cohomology to
massless fields (\hyperref[ref-EastwoodPenroseWells1981]{Eastwood,
Penrose and Wells}). The normalization of the two contour residues is
that of the original analytic correspondence
(\hyperref[ref-IsthmusLocal2026]{Isthmus 2026a}, Theorem 3.1 and
Appendix C.2).

\subsection{3.1. Root coefficients on the flat Cauchy
source}\label{root-coefficients-on-the-flat-cauchy-source}\label{sec:online-root-complex}

The root calculation supplies linear coordinates for the weighted free
seeds. We first compute it in a coefficient complex modulo spatial
polynomials. This quotient is used only for the linear transform. Actual
ordinary classes and their permitted primitives are defined in the
separate actual coefficient complex; products of arbitrary coefficients
do not descend through the polynomial quotient.

Fix the Euclidean orientation, scale and spin marking on
\(\mathbb R^3\). Fourier transformation uses

\phantomsection\label{eq:online-root-fourier-convention} \[
\widehat f(\xi)=(2\pi)^{-3/2}\int e^{-ix\cdot\xi}f(x)\,\mathrm dx,
\qquad
f(x)=(2\pi)^{-3/2}\int e^{ix\cdot\xi}\widehat f(\xi)\,\mathrm d\xi.
\] On \(N_0=\mathbb R^3_x\times\mathbb{CP}^1\), write the homogeneous
conic as

\phantomsection\label{eq:online-root-conic} \[
z(s,t)=(s^2-t^2,\ i(s^2+t^2),\ 2st),
\qquad z\cdot z=0.
\] In the affine chart \(s=1\), the antiholomorphic distribution is
generated by \(\partial_{\bar\zeta}\) and
\(W=z(1,\zeta)\cdot\partial_x\). The conic line is \(\mathcal O(-2)\),
with its retained spin square root \(\mathcal O(-1)\). Evaluation of a
weight-\(m\) one-form on \(W\) therefore has weight \(m+2\).

For \(\xi\ne0\), put \(r=|\xi|\), \(n=\xi/r\) and

\phantomsection\label{eq:online-root-frequency-quadratic} \[
q_\xi(s,t)=\xi\cdot z(s,t),
\qquad L_\xi=iq_\xi.
\] The two zeros of \(q_\xi\) are simple. Indeed, \(q_\xi=0\) means that
\(\xi\) is perpendicular to both the real and imaginary parts of the
null conic vector. The corresponding unit normal to this oriented
two-plane is \(n\) or \(-n\). Thus its divisor \(D_\xi\) consists of the
antipodal roots \(n_\varepsilon=\varepsilon n\), \(\varepsilon=\pm1\).
Let

\phantomsection\label{eq:online-root-bundle} \[
(\mathcal R_m)_\xi=\mathcal O(m+2)|_{D_\xi}
=\mathcal O(m+2)|_{n_+}\oplus\mathcal O(m+2)|_{n_-}.
\] This is a radially pulled-back rank-two bundle over punctured
frequency space. Its individual root lines are not trivialized by a
polarization frame.

\pandocstatementheading{\leavevmode\phantomsection\label{def:online-moderate-root-coefficients}Definition 3.1 (Moderate coefficients and the polynomial quotient).}

For a smooth bundle \(V\) on \(\mathbb{CP}^1\), let \(\mathscr A_M(V)\),
\(M\ge0\), consist of smooth sections on
\((\mathbb R^3\setminus\{0\})\times\mathbb{CP}^1\) satisfying, in
finitely many bundle charts,

\phantomsection\label{eq:online-root-moderate-amplitude-estimates} \[
\begin{aligned}
|\partial_\xi^\alpha\nabla_\zeta^\nu a(\xi,\zeta)|
&\le C_{\alpha,\nu}|\xi|^{-M-|\alpha|},
&&0<|\xi|\le1,\\
|\partial_\xi^\alpha\nabla_\zeta^\nu a(\xi,\zeta)|
&\le C_{\alpha,\nu,N}(1+|\xi|)^{-N},
&&|\xi|\ge1,
\end{aligned}
\] for every finite derivative order and every \(N\). One exponent \(M\)
works for all fiber derivatives. Write
\(\mathscr A(V)=\bigcup_{M\ge0}\mathscr A_M(V)\). Each fixed-\(M\) space
has its displayed seminorms; continuity between these spaces is allowed
to replace \(M\) by one larger finite exponent. For bundles pulled
radially from the frequency sphere, the same definition uses finite
angular trivializations.

For a bundle \(V\) over \(\mathbb{CP}^1\), let
\(\widetilde{\mathscr A}(V)\) consist of actual smooth sections
\(f(x,\zeta)\) with polynomial growth in \(x\) together with all
derivatives, whose distributional Fourier restriction away from zero
belongs to \(\mathscr A(V)\). At each fixed spatial differential order,
a single polynomial growth exponent must work for all fiber derivatives.
Let \(\mathscr P(V)\) be the finite-degree polynomials in \(x\) with
smooth \(V\)-valued coefficients, and set

\phantomsection\label{eq:online-root-spatial-polynomial-quotient} \[
\mathscr C(V)=\widetilde{\mathscr A}(V)/\mathscr P(V).
\] The weight-\(m\) complex has terms \[
\begin{aligned}
\mathscr C_m^0&=\mathscr C(\mathcal O(m)),\\
\mathscr C_m^1&=\mathscr C(\Lambda^{0,1}\otimes\mathcal O(m))
\oplus\mathscr C(\mathcal O(m+2)),\\
\mathscr C_m^2&=\mathscr C(\Lambda^{0,1}\otimes\mathcal O(m+2)),
\end{aligned}
\] and differential

\phantomsection\label{eq:online-root-spatial-differential} \[
d_m^0u=(\bar\partial u,Wu),
\qquad
d_m^1(\alpha,\beta)=\bar\partial\beta-W\alpha.
\] Here \(\bar\partial\) differentiates only in the projective fiber.
The two generators commute because \(z\) is holomorphic. The operators
preserve spatial polynomials, so the differential is defined on the
quotient. A lift of a closed quotient cochain has polynomial
differential. Closedness of the lift requires that polynomial to vanish.

\pandocstatementheading{\leavevmode\phantomsection\label{lem:online-fourier-polynomial-quotient}Lemma 3.2 (Fourier restriction and the polynomial ambiguity).}

Fourier restriction induces an isomorphism

\phantomsection\label{eq:online-root-fourier-quotient-isomorphism} \[
\mathscr C(V)\longrightarrow\mathscr A(V).
\] It identifies the differential with

\phantomsection\label{eq:online-root-frequency-differential} \[
d_m^0u=(\bar\partial u,L_\xi u),
\qquad
d_m^1(\alpha,\beta)=\bar\partial\beta-L_\xi\alpha.
\] Every punctured coefficient has an actual smooth spatial
representative of polynomial growth. The ambiguity between its tempered
extensions is precisely a polynomial spatial section. On a fixed
amplitude order these constructions have bounds in finitely many
seminorms, with a finite choice of distributional order.

\emph{Proof.} Choose a scalar \(\chi\in C_c^\infty(\mathbb R^3)\) equal
to one near zero and an integer \(L\ge0\) with \(L+4>M\). For
\(a\in\mathscr A_M(V)\) define, in a bundle chart,

\phantomsection\label{eq:online-root-finite-part-extension} \[
\langle E_La,\varphi\rangle
=\int a(\xi)\left[
\varphi(\xi)-\chi(\xi)
\sum_{|\mu|\le L}\frac{\partial^\mu\varphi(0)}{\mu!}\xi^\mu
\right]\,\mathrm d\xi.
\] The bracket is \(O(|\xi|^{L+1})\) near zero. Including the radial
measure, its absolute integral is bounded there by a constant times
\(\int_0^1r^{L+3-M}\,\mathrm dr\), which converges. At infinity rapid
decrease suffices. This gives a tempered extension with order bounded in
terms of \(M,L\) and finitely many amplitude seminorms. The same
argument applies after every fiber derivative without increasing \(L\),
because the fixed exponent \(M\) controls them all. Since the cutoff and
Taylor operation act only in \(\xi\), the construction commutes with the
bundle transition maps and gives a smooth family of distributions over
the fiber.

Split this distribution with a second frequency cutoff. The part away
from zero is Schwartz, with smooth fiber dependence. The part near zero
has compact support and finite distributional order. Its inverse Fourier
transform is an actual smooth function: differentiating in \(x\) inserts
powers of \(\xi\) in the compactly supported distribution's action on
the exponential. Each such derivative is bounded by a polynomial in
\(x\). The same order works for fiber derivatives at each fixed spatial
order. Thus the inverse transform belongs to
\(\widetilde{\mathscr A}(V)\).

Two tempered extensions differ by a distribution supported at zero. If
its order is at most \(L'\), it depends only on the Taylor polynomial of
a test function through degree \(L'\). To verify this assertion,
subtract that Taylor polynomial near zero. The remainder and all
derivatives through order \(L'\) vanish there; multiplication by a
cutoff with shrinking support and the distributional estimate make its
value tend to zero. Consequently the distribution is a finite linear
combination of \(\partial_\xi^\mu\delta_0\). Smooth fiber dependence
gives smooth coefficients, and the uniform distributional order gives
one finite degree. Its inverse Fourier transform is a polynomial
section.

Conversely, every polynomial has Fourier transform supported at zero. If
a section of \(\widetilde{\mathscr A}(V)\) has zero punctured Fourier
restriction, the same point-support argument makes it a polynomial. This
proves injectivity and surjectivity. Fourier transformation changes
\(\partial_{x^j}\) to \(i\xi_j\) and commutes with fiber
differentiation. The differential and its polynomial ambiguity therefore
have exactly the stated form. The estimates in this construction use
finitely many seminorms at each requested order. \hfill\(\square\)

\pandocstatementheading{\leavevmode\phantomsection\label{thm:online-moderate-root-cohomology}Theorem 3.3 (The two-root cohomology calculation).}

For \(m=2,-6\), evaluation of the horizontal component at \(D_\xi\)
gives

\phantomsection\label{eq:online-root-cohomology-isomorphism} \[
\mathscr H_m^1:=H^1(\mathscr C_m^\bullet)
\xrightarrow{\ \sim\ }\mathscr A(\mathcal R_m),
\qquad[(\alpha,\beta)]\longmapsto\beta|_{D_\xi},
\qquad
\mathscr H_m^0=\mathscr H_m^2=0.
\] There are linear normal representatives and primitives continuous
between fixed coefficient orders, allowing a finite shift of \(M\). They
commute with smooth finite-dimensional parameter differentiation on such
fixed orders. A Fourier-flat Schwartz amplitude is a Schwartz section
all of whose frequency derivatives vanish at zero. The restriction to
those amplitudes agrees with their root cohomology and embeds it
injectively.

\emph{Proof.} Work in punctured frequency space using the preceding
lemma. Fix the Fubini--Study metric and invariant Hermitian metrics on
\(\mathcal O(k)\). Let \(H_k^j\) denote harmonic projection in Dolbeault
degree \(j\), and let \(T_k=\bar\partial^*G_k\) on one-forms. The fiber
identities are

\phantomsection\label{eq:online-root-fiber-homotopy} \[
\bar\partial T_k=1-H_k^1,
\qquad
T_k\bar\partial=1-H_k^0.
\] These are fixed continuous maps on smooth sections of the compact
fiber. Their estimates use finitely many stronger fiber seminorms. They
commute with spatial spin rotations for the chosen invariant metrics.
The usual cohomology of \(\mathcal O(k)\) on \(\mathbb{CP}^1\) gives
\(H^1(\mathcal O(2))=H^1(\mathcal O(4))=0\),
\(H^0(\mathcal O(-6))=H^0(\mathcal O(-4))=0\), and dimensions \(3,5\)
for \(H^0(\mathcal O(2)),H^0(\mathcal O(4))\) and \(5,3\) for
\(H^1(\mathcal O(-6)),H^1(\mathcal O(-4))\).

For \(m=2\), subtract \(d^0(T_2\alpha)\) from a closed pair. Its first
component becomes zero and its second becomes the holomorphic section
\(\beta_0=\beta-L_\xi T_2\alpha\). Multiplication \[
L_\xi:H^0(\mathcal O(2))\longrightarrow H^0(\mathcal O(4))
\] is injective because \(q_\xi\) is nonzero. Put
\(L_\xi^+=(L_\xi^*L_\xi)^{-1}L_\xi^*\). Subtracting the remaining
holomorphic coboundary gives the unique normal representative

\phantomsection\label{eq:online-root-positive-normal-form} \[
(0,\beta_+),
\qquad
\beta_+=(1-L_\xi L_\xi^+)\bigl(\beta-L_\xi T_2\alpha\bigr)
\in\ker L_\xi^*.
\] If this expression vanishes, an explicit primitive is
\(T_2\alpha+L_\xi^+(\beta-L_\xi T_2\alpha)\).

For \(m=-6\), put \[
U_6=H^1(\mathcal O(-6)),\quad
U_4=H^1(\mathcal O(-4)),\quad
A_\xi=H_{-4}^1L_\xi|_{U_6}.
\] After subtracting \(d^0(T_{-6}\alpha)\), the first component is
\(\alpha_0=H_{-6}^1\alpha\). Closedness gives \(A_\xi\alpha_0=0\). Since
\(H^0(\mathcal O(-4))=0\), the second component is uniquely
\(T_{-4}L_\xi\alpha_0\). The normal representative is therefore

\phantomsection\label{eq:online-root-negative-normal-form} \[
(\alpha_0,T_{-4}L_\xi\alpha_0),
\qquad \alpha_0\in\ker A_\xi.
\] It is exact only when \(\alpha_0=0\), in which case \(T_{-6}\alpha\)
is a primitive of the original pair. The dual of \(A_\xi\) is, up to the
fixed nonzero pairing factor, multiplication by \(q_\xi\) from
\(H^0(\mathcal O(2))\) to \(H^0(\mathcal O(4))\). That map is injective,
so \(A_\xi\) is surjective and has a two-dimensional kernel.

All finite-dimensional inverses used here have the required estimates.
At \(\xi=rn\), \(L_\xi=rL_n\) and \(A_\xi=rA_n\). Their nonzero singular
values have a positive lower bound for \(r=1\), by constant rank and
compactness of the unit sphere. Thus \(L_\xi^+\) and
\(A_\xi^*(A_\xi A_\xi^*)^{-1}\) are smooth homogeneous multipliers of
degree \(-1\); the kernel and complement projections have degree zero. A
smooth multiplier of any fixed degree changes a small-frequency exponent
by only a finite amount, and each Cartesian derivative has the
corresponding additional power of \(r^{-1}\). At infinity it preserves
rapid decrease. The normal forms and primitives consequently preserve
\(\mathscr A\), with the claimed fixed-order continuity.

The divisor sequence

\phantomsection\label{eq:online-root-divisor-sequence} \[
0\longrightarrow\mathcal O(m)
\xrightarrow{\,L_\xi\,}\mathcal O(m+2)
\longrightarrow\mathcal O(m+2)|_{D_\xi}\longrightarrow0
\] identifies the positive cokernel and the negative kernel above with
the two root fibers. For the negative weight the connecting class is
represented by \(L_\xi^{-1}\bar\partial\beta\), which is \(\alpha\) for
a closed pair. Thus in the total complex the common root map is exactly
\(\beta|_{D_\xi}\). The root-to-positive-normal-form map has degree
zero. In the negative normal form the map to \(\alpha_0\) has degree
\(-1\), and the map to its horizontal component \(T_{-4}L_\xi\alpha_0\)
has degree zero. These maps are smooth in the direction and satisfy the
same multiplier estimates.

One can see surjectivity without an implicit choice of a sheaf
representative. Extend prescribed root values to a smooth section
\(\beta_b\) of \(\mathcal O(m+2)\), holomorphic near each root. Choose
uniformly separated caps around the antipodal roots in finitely many
angular charts, use holomorphic bundle frames within the caps, and patch
by a partition of unity in the frequency direction. Put
\(\alpha_b=L_\xi^{-1}\bar\partial\beta_b\). The numerator is zero near
the roots, so this division is smooth on the whole projective line.
Division costs one power of \(r^{-1}\); all direction derivatives have
finite uniform bounds. The pair is closed and has the stipulated root
values. Its inverse Fourier transform is an actual smooth cochain
representative; closedness at this stage is closedness in the polynomial
quotient.

A degree-zero cocycle is fiber holomorphic and annihilated by the
nonzero \(q_\xi\), hence zero. For degree two, the positive-weight right
inverse is \(b\mapsto(0,T_4b)\). For the negative weight, given a
degree-two form \(b\), choose \[
\alpha=-A_\xi^*(A_\xi A_\xi^*)^{-1}H_{-4}^1b,
\qquad
\beta=T_{-4}(b+L_\xi\alpha).
\] Then \(b+L_\xi\alpha\) has zero harmonic part, and
\(d^1(\alpha,\beta)=b\). These operations have the already verified
estimates, proving the remaining cohomology assertions.

All displayed operators act continuously on each specified coefficient
order, with a finite order shift. They therefore commute with
finite-dimensional parameter differentiation and give finite seminorm
bounds for each specified derivative. On Fourier-flat amplitudes the
same homogeneous operations preserve flatness at zero and Schwartz
decay. A Fourier-flat class whose root values vanish has the primitive
just constructed in that Fourier-flat space itself. This proves the
final injectivity and agreement assertion. \hfill\(\square\)

\subsection{3.2. Decay and the zero-frequency
normalization}\label{decay-and-the-zero-frequency-normalization}\label{sec:online-root-weighted-fourier}

For a real exponent \(\beta<0\), let \(\mathscr S^\beta\) consist of
actual smooth fields with seminorms \[
\sup_x\langle x\rangle^{-\beta+|\alpha|}
|\partial_x^\alpha f(x)|<\infty
\] at every finite order. Let \(\mathscr B_\beta\) consist of smooth
functions on punctured frequency space, rapidly decreasing at infinity
together with all derivatives, and satisfying

\phantomsection\label{eq:online-root-decay-amplitude-test} \[
|\partial_\xi^\alpha a(\xi)|
\le C_\alpha|\xi|^{-\beta-3-|\alpha|},
\qquad 0<|\xi|\le1.
\] These definitions apply componentwise to tensors. The amplitude is
locally integrable at zero because \(\beta<0\). Its specified locally
integrable extension is used whenever an inverse Fourier integral is
taken.

\pandocstatementheading{\leavevmode\phantomsection\label{lem:online-decay-fourier-isomorphism}Lemma 3.4 (Weighted Fourier transformation and its inverse estimate).}

For every \(\beta<0\), the ordinary inverse Fourier integral defines a
continuous map \(\mathscr B_\beta\to\mathscr S^\beta\). For
\(-3<\beta<0\), Fourier transformation and this inverse give continuous
inverse maps

\phantomsection\label{eq:online-root-weighted-fourier-isomorphism} \[
\mathscr S^\beta\xrightarrow{\ \sim\ }\mathscr B_\beta.
\] In the two-sided assertion, the distributional Fourier transform is
precisely the locally integrable extension of its punctured restriction,
with no additional origin-supported distribution. Reality, tensor
symmetries and the transverse-traceless conditions are preserved. For
any fixed \(0<\alpha<1\),
\(\mathscr S^\beta=\bigcap_{j\ge0}C^{j,\alpha}_\beta\) with equivalent
families of seminorms, using the weighted Hölder norms of
\hyperref[def:online-weighted-holder-chart-spaces]{Definition A.1}. The
one-way inverse estimate for \(\beta\le-3\) does not assert this
two-sided Fourier isomorphism outside \((-3,0)\).

\emph{Proof.} First take \(-3<\beta<0\) and \(f\in\mathscr S^\beta\). A
dyadic spatial partition decomposes its exterior part into
\(f_j(x)=2^{j\beta}g_j(2^{-j}x)\), \(j\ge0\), where \(g_j\) is supported
in one fixed annulus and has uniformly bounded derivatives of every
prescribed finite order. Integration by parts gives

\phantomsection\label{eq:online-root-forward-annular-estimate} \[
|\partial_\xi^\mu\widehat f_j(\xi)|
\le C_{\mu,N}2^{j(\beta+3+|\mu|)}
(1+2^j|\xi|)^{-N}.
\] For \(0<r=|\xi|\le1\), split at \(2^jr=1\). The sum with \(2^jr\le1\)
is bounded by \(Cr^{-\beta-3-|\mu|}\) because \(\beta+3+|\mu|>0\).
Choose \(N>\beta+3+|\mu|\) for the other geometric tail; it has the same
bound. If \(r\ge1\), this second estimate applies to every annulus and
gives arbitrary rapid decrease. The central compactly supported term is
Schwartz. The differentiated series converges uniformly on compact
subsets away from zero, so it gives a smooth punctured coefficient in
\(\mathscr B_\beta\). The original dyadic series converges to \(f\) as a
tempered distribution, identifying that coefficient with its
distributional Fourier restriction.

For the inverse direction, now let any \(\beta<0\) and
\(a\in\mathscr B_\beta\). Separate a Schwartz high-frequency part and
annular pieces supported where \(|\xi|\asymp2^{-j}\). On a fixed annulus
in \(\eta=2^j\xi\), the rescaled amplitudes
\(2^{-j(\beta+3)}a(2^{-j}\eta)\), including their cutoffs, have bounded
derivatives of every prescribed finite order. The inverse Fourier
contribution of such a piece satisfies

\phantomsection\label{eq:online-root-inverse-annular-estimate} \[
|\partial_x^\mu f_j(x)|
\le C_{\mu,N}2^{j(\beta-|\mu|)}
(1+2^{-j}|x|)^{-N}.
\] For \(|x|\ge1\), split at \(2^{-j}|x|=1\). In the small-frequency
tail, \(\beta-|\mu|<0\) gives a sum bounded by \(C|x|^{\beta-|\mu|}\).
For the other part choose \(N>|\mu|-\beta\); its finite geometric sum
has the same bound. On bounded sets the tail converges with every
spatial derivative. The high-frequency part is Schwartz. This proves
membership in \(\mathscr S^\beta\). Each target seminorm uses finitely
many amplitude seminorms, and the same is true for the preceding forward
estimate.

All \(\xi^\mu a\) are integrable near zero and at infinity. Hence the
inverse integral and its derivatives are actual functions and the
inverse tends to zero at infinity. For an original
\(f\in\mathscr S^\beta\) with \(-3<\beta<0\), the difference between its
distributional Fourier transform and the integrable extension just
obtained is supported at zero. The Fourier restriction lemma makes its
inverse a polynomial. Both \(f\) and the inverse integral tend to zero,
so this polynomial is zero. This proves the two inverse identities and
the absence of an added zero-frequency mode.

Conjugation and reversal of \(\xi\) give the real Fourier relation. The
tensor symmetry conditions commute with Fourier transformation. If
\(\xi_i a_{ij}=0\) off zero, the same identity holds for the locally
integrable extension, hence its inverse satisfies
\(\partial_i f_{ij}=0\). Trace is handled in the same way. Conversely,
transforming those actual identities gives the punctured tensor
conditions.

On each rescaled spatial annulus, the Hölder seminorm of a derivative is
bounded by its supremum and the supremum of its next derivative on a
slightly larger annulus. The symbol estimates thus give every weighted
Hölder seminorm using finitely many stronger symbol seminorms. The
pointwise part of the weighted Hölder norm gives the converse. This
proves the final seminorm equivalence. \hfill\(\square\)

The lower bound \(-3<\beta\) was used in the forward spatial-annulus
sum. It was not used in the inverse-frequency sum. In particular, the
one-way estimate applies to a specifically constructed amplitude of
order \(\beta=\gamma-3\); it does not assign that Fourier behavior to
every field with the same faster spatial decay.

\subsection{3.3. Calibrated root tensors and the weighted seed
transform}\label{calibrated-root-tensors-and-the-weighted-seed-transform}\label{sec:online-root-calibration-and-seeds}

For \(\xi\ne0\), let \(P_n=1-n\otimes n\) and \[
\begin{gathered}
\mathcal T_\xi=
\{Q\in\operatorname{Sym}^2(\mathbb C^3):Qn=0,\ \operatorname{tr}Q=0\}, \\
\Pi_nQ=P_nQP_n-\tfrac12\operatorname{tr}(P_nQP_n)P_n.
\end{gathered}
\] Put \(j_n v=n\times v\) and \(\mathcal J_nQ=j_nQ\) on
\(\mathcal T_\xi\). In an oriented orthonormal basis of \(n^\perp\), a
tracefree symmetric matrix anticommutes with \(j_n\). Hence \(j_nQ\) is
again symmetric and tracefree, \(\mathcal J_n^2=-1\), and its two
eigenspaces are lines. Define

\phantomsection\label{eq:online-root-helicity-projectors} \[
\mathcal L_\sigma(\xi)=\ker(\mathcal J_n-\sigma i),
\qquad
\Pi^{\mathrm{hel}}_\sigma(\xi)
=\tfrac12(1-\sigma i\mathcal J_n)\Pi_n,
\quad \sigma=\pm1.
\] These are smooth homogeneous degree-zero projectors. Their parity is

\phantomsection\label{eq:online-root-helicity-parity} \[
\Pi^{\mathrm{hel}}_\sigma(-\xi)\overline Q
=\overline{\Pi^{\mathrm{hel}}_\sigma(\xi)Q},
\] because \(j_{-n}=-j_n\). The projector formulas define the helicity
lines intrinsically, without a choice of polarization vectors.

At \(n_\varepsilon\), the conic vector has \(j_n\) eigenvalue
\(-\varepsilon i\). Its tensor square identifies
\(\mathcal O(-4)|_{n_\varepsilon}\) with
\(\mathcal L_{-\varepsilon}(\xi)\). The complex bilinear pairing
\(\operatorname{tr}(QU)\) is nondegenerate between opposite helicity
lines and zero on each line with itself. It therefore identifies the
dual root line \(\mathcal O(4)|_{n_\varepsilon}\) with
\(\mathcal L_\varepsilon(\xi)\). Denote these degree-zero maps by
\(\iota_{-6,\varepsilon}\) and \(\iota_{2,\varepsilon}\).

For a paired root class
\(b=(b_+,b_-)\in\mathscr H_2^1\oplus\mathscr H_{-6}^1\), let

\phantomsection\label{eq:online-root-raw-tensors} \[
U^+_\varepsilon=\iota_{2,\varepsilon}(b_+|_{n_\varepsilon}),
\qquad
U^-_\varepsilon=\iota_{-6,\varepsilon}(b_-|_{n_\varepsilon}).
\] The superscript labels the projective weight, the subscript labels
the root branch, and \(\sigma\) labels helicity. These labels have
different roles. The ordinary pair in the actual coefficient cohomology
is denoted by \(c\); its root readout will be \(\rho(c)=b\).

We retain the original residue calibration. With these line
identifications its metric-amplitude factors are

\phantomsection\label{eq:online-root-calibration-factors} \[
k^+_\varepsilon(r)=4i\varepsilon r,
\qquad
k^-_\varepsilon(r)=\frac{i\varepsilon}{4r^3}.
\] Their normalization can be checked by one nonzero-frequency plane,
using the original contour convention
(\hyperref[ref-IsthmusLocal2026]{Isthmus 2026a}, §3, Theorem 3.1). Use
the Hermitian soldering \[
X^{AA'}=\frac1{\sqrt2}
\begin{pmatrix}t-x_3&x_1-ix_2\\x_1+ix_2&t+x_3\end{pmatrix},
\qquad
u=X^{00'},\ v=X^{01'},\ q=X^{10'}.
\] For \(F_m=Ae^{\lambda w^0}/\zeta\), the ordered north-south Čech
differential and a north cutoff \(\rho_N\) give the smooth CR
representative \[
\alpha=Ae^{\lambda u}\zeta^{-1}\bar\partial\rho_N,
\qquad
\beta=-\sqrt2\lambda Ae^{\lambda u}\rho_N.
\] Indeed subtract the differential of \(\rho_NG\), where
\(G=Ae^{\lambda u}(e^{\lambda v\zeta}-1)/\zeta\), from the usual cutoff
representative. The identities \(Wu=-\sqrt2\zeta\), \(Wv=\sqrt2\), and
\(WG=\sqrt2\lambda Ae^{\lambda u}\) give the displayed pair. At the
north root its coefficient is \(a=-\sqrt2\lambda A\). Put
\(\lambda=-i\sqrt2\varepsilon r\); the plane has spatial frequency
\(\xi=\varepsilon re_3\) and time factor \(e^{-i\varepsilon rt}\), with
\(a=2i\varepsilon rA\).

The positive and negative contours have highest curvature components
\(A\lambda^4e^{\lambda u}\) and \(Ae^{\lambda u}\). In the original
last-slot curvature convention they are realized by metric variations
\(-2A\lambda^2e^{\lambda u}(\mathrm dv)^2\) and
\(-2A\lambda^{-2}e^{\lambda u}(\mathrm dq)^2\), since
\(\Psi^+_{0000}=-C_{uvuv}\) and \(\Psi^-_{0'0'0'0'}=-C_{uquq}\)
(\hyperref[ref-IsthmusLocal2026]{Isthmus 2026a}, §2 and Appendix C.2).
At the north root \(z=(1,i,0)\), and \[
\iota_{-6}(a)=a\,z\otimes z,
\qquad
\iota_2(a)=\tfrac a4\,\overline z\otimes\overline z,
\] because \((z\cdot\overline z)^2=4\). The corresponding spatial metric
amplitudes are \(A\lambda^2\overline z\otimes\overline z\) and
\(A\lambda^{-2}z\otimes z\). Substitution gives precisely the two
factors in \hyperref[eq:online-root-calibration-factors]{the calibration
formula}. Spin rotation equivariance carries this computation to every
root, with the same scalar factors. The plane fixes the local
normalization; no plane wave is asserted to belong to the decaying
weighted coefficient domain. The structured nonlinear analytic
comparison has its own hypotheses and is separate from this flat
normalization.

\pandocstatementheading{\leavevmode\phantomsection\label{def:online-weighted-root-domain}Definition 3.5 (The admitted weighted root pair).}

Fix \(-1<\gamma<-1/2\). For the root tensors of \(b\), put

\phantomsection\label{eq:online-root-cauchy-amplitudes} \[
\begin{aligned}
A_\varepsilon(\xi)
 =4i\varepsilon r\,U^+_\varepsilon(\xi)
 +\frac{i\varepsilon}{4r^3}U^-_\varepsilon(\xi),\\
Q_b=A_++A_-,\qquad
P_b=-\frac{ir}{2}(A_+-A_-).
\end{aligned}
\] The real space \(\mathscr W_\gamma\) consists of paired root classes
satisfying

\phantomsection\label{eq:online-root-residue-reality} \[
U^+_-(\xi)=\frac1{16r^4}\overline{U^-_+(-\xi)},
\qquad
U^-_-(\xi)=16r^4\overline{U^+_+(-\xi)},
\] and

\phantomsection\label{eq:online-root-weighted-admission} \[
Q_b\in\mathscr B_\gamma,
\qquad
P_b\in\mathscr B_{\gamma-1}.
\] Its seminorms are those of these two amplitude spaces. The tests use
only the linear root map, fixed multipliers, conjugation and amplitude
estimates, before the nonlinear constraint chart is applied. Let
\(\mathcal X^\infty_\gamma\) be the real transverse-traceless pairs
\((q,p)\) with \(q\in\mathscr S^\gamma\) and
\(p\in\mathscr S^{\gamma-1}\), equipped with all their weighted
seminorms. In particular these fields belong to every finite weighted
Hölder space required by the nonlinear chart.

\pandocstatementheading{\leavevmode\phantomsection\label{thm:online-root-seed-isomorphism}Theorem 3.6 (The exact weighted root-to-seed transform).}

The map

\phantomsection\label{eq:online-root-seed-transform} \[
\mathscr P_\gamma:\mathscr W_\gamma\longrightarrow\mathcal X^\infty_\gamma,
\qquad
b\longmapsto\bigl(\mathcal F^{-1}Q_b,\mathcal F^{-1}P_b\bigr)
\] is a continuous real-linear isomorphism. The inverse transforms are
convergent integrals with the specified zero-frequency normalization.
The exact inverse is obtained by setting \[
A_\varepsilon=\tfrac12\widehat q+\frac{i\varepsilon}{r}\widehat p
\] and then

\phantomsection\label{eq:online-root-seed-exact-inverse} \[
U^+_\varepsilon
=\frac1{4i\varepsilon r}\Pi^{\mathrm{hel}}_\varepsilon A_\varepsilon,
\qquad
U^-_\varepsilon
=\frac{4r^3}{i\varepsilon}\Pi^{\mathrm{hel}}_{-\varepsilon}A_\varepsilon,
\] followed by the inverse root-line maps and the root cohomology
theorem. On the Fourier-flat subspace it is the calibrated residue
Cauchy transform with the same constants. On smooth finite-dimensional
families it commutes with every parameter derivative, and each specified
finite output seminorm is bounded by finitely many input seminorms on
the same retained parameter domain.

\emph{Proof.} The projective root lines have precisely the helicity
types displayed above, so \(Q_b\) and \(P_b\) are symmetric, tracefree
and transverse away from zero. The coupled reality equations imply \[
A_-(\xi)=\overline{A_+(-\xi)}.
\] For example, the negative-helicity part of the left side is
\(-4irU^+_-(\xi)=-i\overline{U^-_+(-\xi)}/(4r^3)\), which is the
corresponding component of \(\overline{A_+(-\xi)}\). The other line is
identical with the factors exchanged. It follows that
\(Q_b(-\xi)=\overline{Q_b(\xi)}\) and \(P_b(-\xi)=\overline{P_b(\xi)}\).
Both weights \(\gamma\) and \(\gamma-1\) lie in \((-3,0)\). The weighted
Fourier lemma therefore puts the inverse transforms in the stated real
smooth weighted TT spaces, with no polynomial addition.

Conversely, take \((q,p)\in\mathcal X^\infty_\gamma\) and its unique
Fourier amplitudes \(Q,P\). The first inverse formula solves
\(Q=A_++A_-\) and \(P=-ir(A_+-A_-)/2\). The two helicity projectors
split each \(A_\varepsilon\) into its two line components, and the
second formula divides by the nonzero calibrated factors on those
components. All line maps and projectors are smooth homogeneous
degree-zero maps; the remaining factors have finite homogeneous orders.
Thus the four raw root tensors belong to the moderate amplitude spaces.
The root cohomology theorem gives a unique pair \(b\) with these values.

Reality of \(Q,P\) gives \(A_-(\xi)=\overline{A_+(-\xi)}\). Apply the
parity of the projectors and divide by the calibrated factors. This
yields exactly \hyperref[eq:online-root-residue-reality]{the coupled
residue equations}, so \(b\in\mathscr W_\gamma\). The recovered
amplitudes are the original \(Q,P\); hence the admission estimates are
unchanged. The two compositions are identities because the root
evaluation is an isomorphism and each displayed amplitude operation has
the explicit inverse just used.

The amplitude seminorms and the weighted Fourier lemma give continuity
in both directions. The inverse root formulas and the Hodge normal forms
also bound each requested representative seminorm by finitely many seed
seminorms, with a fixed finite small-frequency exponent shift. All these
maps are linear and continuous on the stated spaces. Difference
quotients therefore show that they commute with finite-dimensional
parameter derivatives; applying the same finite seminorm estimate to
each derivative proves the family assertion. No smallness is needed for
this linear theorem. On the Fourier-flat subspace the root evaluation,
line maps and calibrated factors are the same formulas, proving exact
agreement there. \hfill\(\square\)

The coupled equations are also the fixed-point equations of the
antilinear involution \[
(\mathfrak c_{\mathrm{res}}U)^+_\varepsilon(\xi)
=\frac1{16r^4}\overline{U^-_{-\varepsilon}(-\xi)},
\qquad
(\mathfrak c_{\mathrm{res}}U)^-_\varepsilon(\xi)
=16r^4\overline{U^+_{-\varepsilon}(-\xi)}.
\] Its square is the identity: reflection preserves \(r\), the
reciprocal factors cancel, and conjugation squares to one. The actual
spin marking remains supplied. These even projective weights do not
detect the two central spin lifts by their action alone.

The linear isomorphism covers all weighted smooth TT seeds in the
specified orders, including permitted low-frequency singularities. It
introduces no independent polynomial mode and imposes no
Fourier-flatness or analytic-background condition on those seeds. To
obtain the ordinary inputs used for source execution, one must next lift
these root classes to actual closed cochains and impose the actual
polynomial-obstruction tests. That next step is not a nonlinear
operation on the polynomial quotient used here.

\subsection{3.4. Actual coefficient complexes and ordinary
classes}\label{actual-coefficient-complexes-and-ordinary-classes}\label{sec:online-actual-coefficient-complexes}

The root complex is a linear model for the two free fields. To perform
nonlinear operations on smooth representatives, we retain the
coefficient space before taking its polynomial quotient. We now
determine the additional cohomology in that space.

For a smooth bundle \(V\) over the sphere, write \(\mathscr S^\mu(V)\)
for the space of actual smooth coefficients with seminorms

\phantomsection\label{eq:online-actual-symbol-seminorms} \[
\sup_{x,\zeta}\langle x\rangle^{-\mu+|\alpha|}
\left|\partial_x^\alpha\nabla_\zeta^\beta f(x,\zeta)\right|,
\qquad \alpha,\beta\ge0,
\] in the fixed finite bundle charts. One spatial order \(\mu\) applies
to every derivative. Put
\(\mathscr E(V)=\bigcup_{\mu\in\mathbb R}\mathscr S^\mu(V)\), and let
\(\mathscr P(V)\) denote the finite-degree polynomials in \(x\) with
smooth fiber coefficients. Statements of continuity on this union will
specify an input order, an output order and, when relevant, a
polynomial-degree bound.

With \(D=z\cdot\partial_x\), the actual weight-\(m\) complex is

\phantomsection\label{eq:online-actual-cr-complex} \[
\begin{aligned}
\mathscr E_m^0&=\mathscr E(\mathcal O(m)),\\
\mathscr E_m^1&=\mathscr E(\Lambda^{0,1}\otimes\mathcal O(m))
 \oplus\mathscr E(\mathcal O(m+2)),\\
\mathscr E_m^2&=\mathscr E(\Lambda^{0,1}\otimes\mathcal O(m+2)),\\
d^0u&=(\bar\partial u,Du),\qquad
d^1(\alpha,\beta)=\bar\partial\beta-D\alpha.
\end{aligned}
\] The differential squares to zero because \(z\) is holomorphic in the
sphere variable and independent of \(x\). An ordinary class is
represented by a closed smooth cochain in this complex, modulo
differentials of smooth zero-cochains in the same coefficient space.

\pandocstatementheading{\leavevmode\phantomsection\label{prop:online-actual-coefficient-exact-sequence}Proposition 3.7 (Actual coefficients and their polynomial kernel).}

The actual moderate coefficient space used to define the root complex
equals \(\mathscr E(V)\). Punctured Fourier restriction gives an exact
sequence

\phantomsection\label{eq:online-actual-symbol-exact-sequence} \[
0\longrightarrow\mathscr P(V)
\longrightarrow\mathscr E(V)
\xrightarrow{\mathcal F|_{\xi\ne0}}\mathscr A(V)
\longrightarrow0.
\] For every fixed \(M\), a continuous linear extension
\(\mathscr A_M(V)\to\mathscr S^\mu(V)\) exists when \(\mu>0\) is
nonintegral and \(\mu>M-3\). Its choice may depend on the spatial
anchor. Actual multiplication and differentiation satisfy

\phantomsection\label{eq:online-actual-symbol-algebra} \[
\mathscr S^\mu(V)\mathscr S^\nu(W)
\subset\mathscr S^{\mu+\nu}(V\otimes W),
\qquad
\partial_{x^i}\mathscr S^\mu(V)\subset\mathscr S^{\mu-1}(V).
\] The collection of actual CR complexes over all integer weights is a
differential graded algebra.

\emph{Proof.} If \(f\in\mathscr S^\mu(V)\), increase \(\mu\) to a
nonnegative number for the Fourier estimate. A spatial dyadic
decomposition, rescaled to fixed annuli, gives \[
\left|\partial_\xi^\alpha\nabla_\zeta^\beta\widehat f_j(\xi,\zeta)\right|
\le C_{\alpha,\beta,N}
2^{j(\mu+3+|\alpha|)}(1+2^j|\xi|)^{-N}.
\] Splitting the sum at \(2^j|\xi|=1\) gives the punctured estimate
\(C|\xi|^{-\mu-3-|\alpha|}\) near zero and rapid decrease, with all
derivatives, at infinity. This is the same annular estimate as in
\hyperref[lem:online-decay-fourier-isomorphism]{the decay Fourier
lemma}; here it is used only for the restriction away from zero and does
not require decay of \(f\). Thus \(\mathscr E(V)\) lies in the actual
moderate space.

Conversely, let \(a\in\mathscr A_M(V)\). Separate a Schwartz coefficient
away from zero and frequency annuli \(a_j\) supported where
\(|\xi|\asymp2^{-j}\), \(j\ge0\). Their inverse transforms can be
written \[
f_j(x,\zeta)=2^{j(M-3)}g_j(2^{-j}x,\zeta),
\] where the \(g_j\) are bounded in every Schwartz and fiber seminorm by
finitely many seminorms of \(a\). Let \(L=\lfloor\mu\rfloor\). The
strict inequality \(\mu>M-3\) implies \(M-4-L<0\). If \(T_Lf_j\) is the
spatial Taylor polynomial at zero through degree \(L\), define

\phantomsection\label{eq:online-actual-dyadic-extension} \[
Ea=f_{\mathrm{high}}+\sum_{j\ge0}(f_j-T_Lf_j).
\] For \(|\alpha|\le L\), Taylor's formula bounds a differentiated
summand on a compact set by \(C2^{j(M-4-L)}|x|^{L+1-|\alpha|}\). For
\(|\alpha|>L\), its polynomial part is absent and the bound is
\(C2^{j(M-3-|\alpha|)}\). Both tails are summable. The same statements
hold after each specified fiber derivative, so \(Ea\) is actual smooth.

To estimate it at \(R=|x|\ge2\), split the sum at \(2^j=R\). The tail is
bounded by \(CR^{M-3-|\alpha|}\). The nonpolynomial part of the
remaining finite sum has the same bound by Schwartz decay. A Taylor term
of degree \(d\) contributes at most \[
C R^{d-|\alpha|}
\sum_{2^j\le R}2^{j(M-3-d)},
\qquad |\alpha|\le d\le L.
\] According to the sign of \(M-3-d\), this is bounded by a multiple of
\(R^{M-3-|\alpha|}\), \(R^{d-|\alpha|}\), or \(R^{d-|\alpha|}\log R\).
Since \(\mu>M-3\) and \(\mu>L\), each expression is
\(O(R^{\mu-|\alpha|})\). These estimates prove
\(Ea\in\mathscr S^\mu(V)\) and continuity in finitely many source
seminorms for each target seminorm. They also control partial sums as
tempered distributions. The subtracted polynomials have Fourier support
at zero, so the punctured Fourier restriction is \(a\).

The point-support statement in
\hyperref[lem:online-fourier-polynomial-quotient]{the Fourier quotient
lemma} identifies the kernel with \(\mathscr P(V)\). For any coefficient
in the original actual moderate space, subtract the extension just
constructed from its Fourier restriction. The difference is polynomial,
hence belongs to \(\mathscr E(V)\). This proves the equality and
exactness.

The multiplication and derivative bounds are direct consequences of the
Leibniz rule. Fiber derivatives preserve the order. The operators
\(\bar\partial\) and \(D\) commute and obey the graded Leibniz rule. In
degree one, the product is \[
(\alpha,\beta)\wedge(\alpha',\beta')
=\alpha\beta'-\beta\alpha',
\] with the line-bundle weights added. This proves the differential
graded algebra assertion. \hfill\(\square\)

The polynomial coefficients form a differential subcomplex, so the
quotient supports the linear root calculation. Since the polynomials
contain \(1\), the algebra ideal they generate is the entire coefficient
algebra. We therefore form nonlinear products in the coefficient algebra
\(\mathscr E\).

\pandocstatementheading{\leavevmode\phantomsection\label{prop:online-polynomial-cr-cohomology}Proposition 3.8 (The polynomial cohomology and actual root map).}

Let \(\mathcal P\) be the scalar polynomials in \(x\). Normalize the
fiber identifications by \[
v\longmapsto v_i z^i\in H^0(\mathcal O(2)),
\qquad
S\longmapsto S_{ij}z^iz^j\in H^0(\mathcal O(4)),
\] and use their Serre-dual pairings to identify \(H^1(\mathcal O(-6))\)
with constant complex symmetric tracefree tensors and
\(H^1(\mathcal O(-4))\) with complex vectors. Put

\phantomsection\label{eq:online-polynomial-killing-divergence} \[
(\mathcal Kv)_{ij}
=\frac12(\partial_i v_j+\partial_jv_i)
-\frac13\delta_{ij}\partial_kv_k,
\qquad
(\operatorname{div}S)_j=\partial_iS_{ij}.
\] Then

\phantomsection\label{eq:online-polynomial-first-cohomology} \[
\begin{aligned}
H^1(\mathscr P_2^\bullet)
&=(\mathcal P\otimes\operatorname{STF}^2\mathbb C^3)
  /\mathcal K(\mathcal P\otimes\mathbb C^3),\\
H^1(\mathscr P_{-6}^\bullet)
&=\ker\bigl(\operatorname{div}:
\mathcal P\otimes\operatorname{STF}^2\mathbb C^3
\longrightarrow\mathcal P\otimes\mathbb C^3\bigr).
\end{aligned}
\] For both weights \(H^2(\mathscr P_m^\bullet)=0\). The negative
degree-zero cohomology vanishes, while
\(H^0(\mathscr P_2^\bullet)=\ker\mathcal K\) is the ten-dimensional
complex space of polynomial conformal Killing fields.

For \(m=2,-6\), the actual root map fits into the exact sequence

\phantomsection\label{eq:online-actual-cohomology-root-extension} \[
0\longrightarrow H^1(\mathscr P_m^\bullet)
\longrightarrow H^1(\mathscr E_m^\bullet)
\xrightarrow{\rho_m}\mathscr H_m^1
\longrightarrow0.
\] Moreover \(H^0(\mathscr E_m^\bullet)=H^0(\mathscr P_m^\bullet)\) and
\(H^2(\mathscr E_m^\bullet)=0\). Every root class has an actual closed
lift, obtained by finite polynomial corrections after a coefficient
extension. These constructions have finite estimates on fixed
coefficient orders and bounded polynomial degrees.

\emph{Proof.} The fixed fiber Hodge operators preserve polynomial
spatial coefficients. In positive weight they reduce the polynomial
complex to \[
\mathcal P\otimes H^0(\mathcal O(2))
\xrightarrow{D}\mathcal P\otimes H^0(\mathcal O(4)).
\] Since \(z\cdot z=0\), this is precisely \(\mathcal K\) in the
displayed tensor identifications. The relevant higher fiber cohomology
vanishes. This proves the positive formulas and its degree-two
vanishing.

In negative weight the reduced map is \[
\mathcal P\otimes H^1(\mathcal O(-6))
\xrightarrow{H_{-4}^1D}\mathcal P\otimes H^1(\mathcal O(-4)).
\] It is exactly \(\operatorname{div}\). Indeed, the dual of
multiplication by \(z^i\) pairs \(S\) with the tracefree part of
\(\tfrac12(e_i\otimes v+v\otimes e_i)\), giving \(S_{ij}v_j\).
Differentiation in \(x^i\) gives \(\partial_iS_{ij}\) with no additional
scalar.

This polynomial divergence map is onto. Set \(\mathcal L=2\mathcal K\)
and \[
\mathcal N v=\operatorname{div}\mathcal Lv
=\Delta v+\frac13\nabla\operatorname{div}v.
\] The coefficient of \(\partial_1^2\) in \(\mathcal N\) is
\(M=\operatorname{diag}(4/3,1,1)\). Write its other terms as
\(B^A\partial_1\partial_A+C^{AB}\partial_A\partial_B\),
\(A,B\in\{2,3\}\). For a polynomial vector \(w\) of degree at most
\(d\), prescribe \(v_0=v_1=0\) and recursively define

\phantomsection\label{eq:online-polynomial-navier-recursion} \[
v_{n+2}=M^{-1}\left(
\partial_1^n w|_{x^1=0}
-B^A\partial_Av_{n+1}
-C^{AB}\partial_A\partial_Bv_n
\right).
\] Induction gives \(\deg v_n\le d+2-n\). The recursion therefore
terminates, and \(v=\sum_{n=0}^{d+2}(x^1)^nv_n/n!\) satisfies
\(\mathcal Nv=w\). Thus \(\mathcal Lv\) is a polynomial tracefree tensor
with divergence \(w\). This proves negative degree-two vanishing. Its
degree-zero group vanishes because \(H^0(\mathcal O(-6))=0\).

For the positive kernel, writing \(\sigma=\operatorname{div}v/3\) in
\(\mathcal Kv=0\) and differentiating gives \[
\partial_i\partial_jv_k
=\delta_{jk}\partial_i\sigma
+\delta_{ik}\partial_j\sigma
-\delta_{ij}\partial_k\sigma.
\] Commuting one more derivative and contracting gives
\(\partial_j\partial_l\sigma+\delta_{jl}\Delta\sigma=0\). Taking the
trace yields \(\Delta\sigma=0\), hence \(\nabla^2\sigma=0\).
Substitution gives precisely \[
v(x)=a+\Omega x+\lambda x+2(c\cdot x)x-|x|^2c,
\qquad \Omega^{\mathsf T}=-\Omega.
\] The dimensions of these summands are \(3,3,1,3\).

Apply the long exact cohomology sequence to the coefficient exact
sequence above. \hyperref[thm:online-moderate-root-cohomology]{The
root-cohomology theorem} gives \(\mathscr H_m^0=\mathscr H_m^2=0\). The
polynomial degree-two vanishing just proved therefore gives the
displayed exact sequence and the assertions in degrees zero and two. In
particular its first map is injective; a nonzero polynomial
first-cohomology class remains nonzero in actual cohomology.

Actual closed lifts can also be obtained directly. Extend a root cocycle
coefficientwise to an actual cochain \(e\). Its residual \(r=d^1e\) is a
polynomial degree-two cochain. For positive weight, \(p=(0,T_4r)\)
satisfies \(d^1p=r\). For negative weight, let \(R=H_{-4}^1r\), solve
\(\mathcal Nv=R\) by the finite recursion, and let \(a\) be the harmonic
representative corresponding to \(-\mathcal Lv\). Then
\(H_{-4}^1(r+Da)=0\), and \[
p=\bigl(a,T_{-4}(r+Da)\bigr),\qquad d^1p=r.
\] Thus \(e-p\) is an actual closed lift. All operations are fixed fiber
maps, coefficient extensions with their stated order bounds, and finite
polynomial linear operations. Their estimates have the asserted finite
dependence. \hfill\(\square\)

The map \(\rho_m\) is therefore not injective on the full actual
cohomology. We next specify the lift used for the admitted weighted
pairs and give a direct test for its actual cohomology class.

\pandocstatementheading{\leavevmode\phantomsection\label{thm:online-canonical-actual-pair-lift}Theorem 3.9 (The canonical actual lift of the weighted pair).}

Let \(-1<\gamma<-1/2\) and put \(\vartheta=\gamma+1\in(0,1/2)\). There
is a continuous real-linear section

\phantomsection\label{eq:online-canonical-ordinary-section} \[
\sigma_\gamma:\mathscr W_\gamma
\longrightarrow
H^1(\mathscr E_2^\bullet)\oplus
H^1(\mathscr E_{-6}^\bullet),
\qquad \rho\sigma_\gamma=I.
\] For a paired root class \(b\in\mathscr W_\gamma\), let
\(\beta_+(\xi)\) and \(\alpha_-(\xi)\) be its positive harmonic normal
coefficient and negative harmonic first coefficient in the
root-cohomology theorem. A fixed actual representative of
\(\sigma_\gamma(b)\) is

\phantomsection\label{eq:online-canonical-actual-representatives} \[
\begin{aligned}
e_+(b)&=(0,B_+(b)),\\
B_+(b)(x)&=(2\pi)^{-3/2}\int_{
\mathbb R^3}(e^{ix\cdot\xi}-1)\beta_+(\xi)\,\mathrm d\xi,\\
e_-(b)&=(A_-(b),B_-(b)),\\
A_-(b)&=\mathcal F^{-1}\alpha_-,
\qquad B_-(b)=T_{-4}D A_-(b).
\end{aligned}
\] These cochains are exactly closed and satisfy

\phantomsection\label{eq:online-canonical-actual-orders} \[
\begin{aligned}
B_+&\in\mathscr S^{\vartheta}(\mathcal O(4)),\qquad B_+(0)=0,\\
A_-&\in\mathscr S^{\vartheta-3}
(\Lambda^{0,1}\otimes\mathcal O(-6)),\\
B_-&\in\mathscr S^{\vartheta-4}(\mathcal O(-4)).
\end{aligned}
\] The positive representative is unique with its punctured harmonic
normal coefficient, the displayed growth, and the anchor. The negative
representative is unique in its harmonic normal form with the displayed
decay. Changing the positive anchor changes the actual representative by
an explicit ordinary coboundary, so its actual class is unchanged. The
section is translation-equivariant in actual cohomology, preserves the
calibrated coupled root reality, and preserves smooth passive families
in the stated seminorms.

\emph{Proof.} For \((q,p)=\mathscr P_\gamma b\), let \(Q=\widehat q\)
and \(P=\widehat p\). The weighted root-domain bounds give
\(Q=O(r^{-\gamma-3})\) and \(P=O(r^{-\gamma-2})\), \(r=|\xi|\), with
their differentiated versions. The inverse root-seed formulas give
\(U^+_\varepsilon=O(r^{-\gamma-4})\) and
\(U^-_\varepsilon=O(r^{-\gamma})\). The positive root-to-normal map has
degree zero and the negative map has degree \(-1\). Consequently

\phantomsection\label{eq:online-canonical-normal-low-frequency-orders}
\[
\begin{aligned}
|\partial_\xi^\alpha\nabla_\zeta^\beta\beta_+|
&\le C_{\alpha,\beta}r^{-\vartheta-3-|\alpha|},\\
|\partial_\xi^\alpha\nabla_\zeta^\beta\alpha_-|
&\le C_{\alpha,\beta}r^{-\vartheta-|\alpha|},
\qquad 0<r\le1.
\end{aligned}
\] Both coefficients decrease rapidly at infinity with all derivatives.

For \(B_+\), the single bracket \(e^{ix\cdot\xi}-1\) supplies one power
of \(r\) near zero on compact spatial sets. Its two summands are not
integrated separately. Including the three-dimensional volume factor
leaves the integrable radial bound \(Cr^{-\vartheta}\,\mathrm dr\).
Every positive-order spatial derivative is also integrable. This proves
actual smoothness, the fiber holomorphic property, and \(B_+(0)=0\).

The one-way inverse assertion in
\hyperref[lem:online-decay-fourier-isomorphism]{the decay Fourier lemma}
applies to every prescribed punctured amplitude with a negative spatial
exponent. Its two-sided characterization is used only in its stated
interval \(-3<s<0\). In particular the one-way estimate can be applied
below to the constructed coefficient of \(B_-\), whose spatial order is
below \(-3\).

Apply it to the first derivatives of \(B_+\), whose Fourier coefficients
have exponent \(s=\vartheta-1<0\). Their higher spatial derivatives have
the corresponding lower orders. Integrating their bound along the
segment from \(0\) to \(x\) gives
\(|B_+(x)|\le C\langle x\rangle^{\vartheta}\), also after fiber
differentiation. This proves the first order statement. The
differentiated integral gives
\(\partial_jB_+=\mathcal F^{-1}(i\xi_j\beta_+)\). On punctured frequency
space these three identities imply \(\widehat B_+=\beta_+\), because at
least one coordinate of \(\xi\) is nonzero. The anchor fixes the
constant left undetermined by the spatial derivatives. Since \(B_+\) is
a holomorphic quartic in the fiber, \(d^1e_+=0\) exactly.

The coefficient \(\alpha_-\) and every spatial Fourier multiple are
locally integrable since \(\vartheta<3\). The reverse estimate with
\(s=\vartheta-3\) gives the stated order for \(A_-\). The coefficient of
\(B_-\) is \(T_{-4}L_\xi\alpha_-\), with low-frequency order
\(r^{1-\vartheta}\); the same reverse estimate, now with
\(s=\vartheta-4\), gives its stated order. The identity
\(H_{-4}^1L_\xi\alpha_-=0\) holds for these locally integrable
extensions as an actual distributional identity. Multiplication by
\(\xi\) produces no point-supported term. The fiber homotopy formula
therefore gives \(\bar\partial B_-=D A_-\), proving exact closedness.

Two positive normal representatives with the same punctured coefficient
differ by a polynomial. Their order \(\vartheta<1\) makes it constant,
and the anchor makes it zero. For the negative normal forms, the
difference of first components is a decaying polynomial and hence zero.
The second component is then fixed because \(H^0(\mathcal O(-4))=0\).

An anchor \(a\) replaces \(B_+(x)\) by \(B_+(x)-B_+(a)\). Write
\(B_+(a)=S_{ij}z^iz^j\) with \(S\) tracefree, and set
\(u_a=S_{ij}x^iz^j\). Then

\phantomsection\label{eq:online-positive-anchor-coboundary} \[
d^0u_a=(0,B_+(a)),
\] so the change is the coboundary \(-d^0u_a\). These comparisons add
exactly for successive anchor changes. The anchored representative of
the translated Fourier coefficient \(e^{ia\cdot\xi}\beta_+\) is
\(B_+(x+a)-B_+(a)\), proving translation equivariance of its ordinary
class. The negative inverse transform already commutes with translation
on representatives.

Every output seminorm has been estimated using finitely many input
seminorms. Differentiation of a supplied smooth finite-dimensional
family under the integrals therefore gives the same estimates for all
specified parameter derivatives. Root readout recovers the initial
normal representatives. Thus \(\rho\sigma_\gamma=I\) and the section is
injective on \(\mathscr W_\gamma\). It preserves the projective weights
and the specified coupled root reality. This assertion does not impose
an extra pointwise conjugation law on arbitrary off-root
representatives. \hfill\(\square\)

\pandocstatementheading{\leavevmode\phantomsection\label{thm:online-actual-cochain-obstruction-test}Theorem 3.10 (Independent normalization tests for actual cochains).}

Let \(e=(e_+,e_-)\) be an actual closed cochain pair with root readout
\(b\in\mathscr W_\gamma\). There are polynomial obstructions

\phantomsection\label{eq:online-actual-polynomial-obstructions} \[
\begin{aligned}
\kappa_+(e)&\in
(\mathcal P\otimes\operatorname{STF}^2\mathbb C^3)
/\mathcal K(\mathcal P\otimes\mathbb C^3),\\
\kappa_-(e)&\in
\ker\bigl(\operatorname{div}:
\mathcal P\otimes\operatorname{STF}^2\mathbb C^3
\longrightarrow\mathcal P\otimes\mathbb C^3\bigr)
\end{aligned}
\] such that

\phantomsection\label{eq:online-actual-normalization-test} \[
e-e(b)=d^0u
\quad\Longleftrightarrow\quad
\kappa_+(e)=\kappa_-(e)=0.
\] The obstructions and a primitive when they vanish are computed by
fixed fiber Hodge maps, punctured Fourier multipliers, coefficient
extension and finite polynomial linear algebra. The negative primitive
is unique. The positive primitive is unique modulo the ten-dimensional
polynomial conformal Killing kernel. The tests are invariant under
actual coboundaries and have finite estimates on fixed symbol orders and
bounded polynomial degrees, allowing a finite increase of order. They
preserve smooth passive families and use no spacetime metric or Einstein
evolution.

\emph{Proof.} Write \(e_+=(\alpha,\beta)\) and \(e_+(b)=(0,B_b)\). Fiber
normalization gives \[
\beta_0=\beta-DT_2\alpha,
\qquad \bar\partial\beta_0=0.
\] At nonzero frequency put \(v=L_\xi^+\widehat\beta_0\). Its
homogeneous order is \(-1\), so \(v\) belongs to a fixed moderate
coefficient space. Extend it to an actual coefficient \(Ev\) by the
coefficient proposition, keeping its values in the finite-dimensional
space of holomorphic quadratic sections. Then

\phantomsection\label{eq:online-positive-polynomial-residual} \[
r=\beta_0-B_b-D(Ev)
\] has zero punctured Fourier transform, because
\((I-L_\xi L_\xi^+)\widehat\beta_0\) is exactly the positive normal
coefficient defining \(B_b\). Thus \(r\) is a polynomial holomorphic
quartic. Define \(\kappa_+(e)=[r]\) in the stated quotient. Another
extension differs by a polynomial quadratic and changes \(r\) by an
element of \(\operatorname{im}\mathcal K\), so this class is independent
of that choice.

If \([r]=0\), solve \(Dp=r\) among polynomial holomorphic quadratic
sections. For a nonzero residual of degree \(D\), a solution can be
chosen of degree at most \(D+1\): every higher homogeneous component of
any polynomial solution lies separately in \(\ker\mathcal K\) and can be
removed. For \(r=0\), take \(p=0\). On the resulting bounded-degree
space this is finite-dimensional linear algebra on the image; choose a
complement to the conformal Killing kernel to fix a solution. Then \[
u_+=T_2\alpha+Ev+p
\] satisfies \(d^0u_+=e_+-e_+(b)\). Conversely, if that difference
equals \(d^0u_+\), the identity \(T_2\bar\partial u_+=u_+-H_2^0u_+\)
gives \(\beta_0=B_b+D(H_2^0u_+)\). Since the canonical positive normal
coefficient lies in \(\ker L_\xi^*\), \(v\) is the punctured Fourier
restriction of \(H_2^0u_+\). Hence \(H_2^0u_+-Ev\) is polynomial and
\(r\) lies in \(\operatorname{im}\mathcal K\). This proves the positive
equivalence.

For \(e_-=(\alpha,\beta)\), let \(\alpha_0=H_{-6}^1\alpha\), and write
\(e_-(b)=(A_b,B_b)\). The two harmonic first coefficients have the same
punctured Fourier restriction. Thus \[
S=\alpha_0-A_b
\] is a polynomial harmonic section. Closedness gives \(H_{-4}^1DS=0\),
which is exactly \(\operatorname{div}S=0\) in the normalized tensor
identification. Set \(\kappa_-(e)=S\). If \(S=0\), subtract
\(d^0(T_{-6}\alpha)\) from \(e_--e_-(b)\). Its first component is zero
and its second component is holomorphic in \(\mathcal O(-4)\), hence
zero. This gives \(u_-=T_{-6}\alpha\). Necessity follows by taking
\(H_{-6}^1\) of the first component of an exact cochain.

Both obstructions are precisely the polynomial cohomology class of
\(e-e(b)\) in the injective map of the actual cohomology sequence.
Adding a coboundary in the coefficient complex preserves both
obstructions. Two primitives differ by a closed zero-cochain, so the
degree-zero groups computed above give the stated ambiguity. The
coefficient extensions and multipliers have finite seminorm estimates,
and the remaining inverses are fixed linear maps on bounded-degree
polynomial spaces. This proves continuity and the family assertions at
the specified stages. \hfill\(\square\)

\subsection{3.5. The selected ordinary
domain}\label{the-selected-ordinary-domain}\label{sec:online-selected-ordinary-domain}

Let \(\mathcal B_\gamma\) be a fixed open finite-base neighborhood of
zero in the smooth TT seed space, contained in the retained nonlinear
chart domain and, when source execution is invoked, in its stated
domain. Put

\phantomsection\label{eq:online-selected-ordinary-class-domain} \[
\mathscr D_{\gamma,\mathcal B}=\mathscr P_\gamma^{-1}(\mathcal B_\gamma),
\qquad
\mathscr C^{\mathrm{ord}}_{\gamma,\mathcal B}
=\sigma_\gamma(\mathscr D_{\gamma,\mathcal B}).
\] For an ordinary pair
\(c\in\mathscr C^{\mathrm{ord}}_{\gamma,\mathcal B}\), use the distinct
notations

\phantomsection\label{eq:online-ordinary-root-seed-types} \[
b=\rho(c)\in\mathscr D_{\gamma,\mathcal B},
\qquad
s(c)=\mathscr P_\gamma\rho(c)\in\mathcal B_\gamma.
\] The first is a paired root class, and the second is the actual whole
TT seed. No quotient by spatial polynomials is applied to the ordinary
pair itself.

\pandocstatementheading{\leavevmode\phantomsection\label{def:online-effective-ordinary-cochain-domain}Definition 3.11 (Admitted actual cochains and effective ordinary
arrows).}

An admitted cochain pair is an element
\(e\in\mathscr E_2^1\oplus\mathscr E_{-6}^1\) satisfying

\phantomsection\label{eq:online-effective-ordinary-admission} \[
d^1e_+=d^1e_-=0,
\qquad \rho(e)\in\mathscr D_{\gamma,\mathcal B},
\qquad \kappa_+(e)=\kappa_-(e)=0.
\] The positive obstruction is tested in its polynomial quotient, rather
than by setting a chosen residual polynomial equal to zero. An effective
ordinary arrow \(e\to e'\) is an actual zero-cochain pair \(u\) with
\(e'-e=d^0u\), modulo

\phantomsection\label{eq:online-effective-ordinary-arrows} \[
u\sim v\quad\Longleftrightarrow\quad d^0(u-v)=0.
\] Composition is addition of primitives. Coefficient topology and
family estimates are taken at fixed symbol orders and bounded polynomial
degrees, together with the \(\mathscr W_\gamma\) topology of the root
readout.

\pandocstatementheading{\leavevmode\phantomsection\label{cor:online-effective-ordinary-class-identification}Corollary 3.12 (The actual ordinary input represented by the tests).}

The actual cohomology classes represented by the admitted cochains are
precisely \(\mathscr C^{\mathrm{ord}}_{\gamma,\mathcal B}\). On this
selected ordinary domain, root readout and \(\sigma_\gamma\) are
inverse. Two admitted cochains admit an effective ordinary arrow exactly
when their roots agree, and that effective arrow is then unique. The
representatives and comparison primitives have the finite seminorm and
parameter estimates stated above.

\emph{Proof.} If \(e\) is admitted, the normalization test gives
\(e-e(\rho(e))=d^0u\), so \([e]=\sigma_\gamma\rho(e)\). Conversely every
\(e(b)\) with \(b\in\mathscr D_{\gamma,\mathcal B}\) is closed and has
zero obstruction, as does every actual cohomologous representative. The
identity \(\rho\sigma_\gamma=I\) gives the inverse assertion on classes.

Every coboundary in the coefficient complex has zero root readout, as
required for the arrow statement. If two roots agree, subtract their two
explicit comparisons with the same canonical representative to obtain an
actual primitive for their difference. Any two such primitives differ by
a closed zero-cochain and hence define the same effective arrow. The
displayed formulas give addition under composition and the asserted
estimates. \hfill\(\square\)

These effective arrows compare cochain representatives of the ordinary
input. Comparisons of completed sources also carry a physical map and
its actual spin lift, as specified in §7.

\subsection{3.6. The ordinary-pair initial-data
constructor}\label{the-ordinary-pair-initial-data-constructor}\label{sec:online-ordinary-data-construction}

The root transform, actual lift and nonlinear chart now have compatible
domains. We state their exact composition before using it to cover local
spacetime germs. All fields in this construction remain defined on the
whole marked \(\mathbb R^3\); a local restriction is taken only after
construction and decoding.

Let \(\mathcal U\subset\mathcal X_{k_0}\) and
\(\mathcal V\subset\mathcal Y_{k_0}\) be the convex base neighborhoods
of \hyperref[thm:online-chart-base-constraint-bijection]{the constraint
chart}. Choose the finite-base neighborhood \(\mathcal B_\gamma\) of the
preceding subsection inside \(\mathcal U\cap\mathcal X^\infty_\gamma\).
In applications requiring the initial source execution, it is also
contained in that operation's retained domain, as specified in
\hyperref[sec:submission-coverage-conventions]{§5.1}. This extra use of
a smaller seed ball does not change the coefficient complexes, transform
or decoder.

For a whole tensor pair in the ambient gauge domain, write the already
defined algebraic and differential readout as

\phantomsection\label{eq:online-ordinary-whole-decoder-parts} \[
\operatorname{Dec}(h,K)
=\bigl(\operatorname{Dec}_{s}(h,K),
       \operatorname{Dec}_{y}(h,K)\bigr)
=\bigl((q,p),(u,W)\bigr).
\] Define \(\mathsf{Data}^{\mathrm g}_{\gamma,\mathcal B}\) by the
following tests on the tensors themselves:

\phantomsection\label{eq:online-ordinary-independent-data-domain} \[
\begin{gathered}
h-\delta\in\mathscr S^\gamma,
\qquad K\in\mathscr S^{\gamma-1},
\qquad \tfrac12\delta<h<\tfrac32\delta,\\
\operatorname{div}_\delta\operatorname{TF}_\delta h=0,
\qquad \operatorname{tr}_hK=0,
\qquad R(h)=|K|_h^2,
\qquad\operatorname{div}_hK=0,\\
\operatorname{Dec}_{s}(h,K)\in\mathcal B_\gamma,
\qquad\operatorname{Dec}_{y}(h,K)\in\mathcal V.
\end{gathered}
\] The fields are real covariant tensors, and the coordinate, scale and
actual spin markings are retained. The decoder is defined on the ambient
gauge domain before the two constraint equations are imposed. Thus these
tests require no supplied generating pair and no membership in a
constructor image. They select the part of the independent tensor target
in the fixed uniqueness neighborhood.

\pandocstatementheading{\leavevmode\phantomsection\label{thm:online-ordinary-pair-data-bijection}Theorem 3.13 (Ordinary classes and the independent maximal-data domain).}

The maps

\phantomsection\label{eq:online-ordinary-data-constructor-and-inverse}
\[
\begin{aligned}
\mathfrak D(c)
&=\mathcal T\bigl(\mathscr P_\gamma\rho(c)\bigr),
&&c\in\mathscr C^{\mathrm{ord}}_{\gamma,\mathcal B},\\
\mathfrak C(h,K)
&=\sigma_\gamma\mathscr P_\gamma^{-1}
       \operatorname{Dec}_{s}(h,K),
&&(h,K)\in\mathsf{Data}^{\mathrm g}_{\gamma,\mathcal B},
\end{aligned}
\] are inverse bijections. In particular,

\phantomsection\label{eq:online-ordinary-data-two-inverse-identities} \[
\mathfrak C\mathfrak D(c)=c,
\qquad
\mathfrak D\mathfrak C(h,K)=(h,K).
\] The forward operation uses the ordinary pair, the fixed coefficient
normalizations and the retained flat markings. Its nonlinear correction
is the selected solution of the constraint chart. Both maps are
continuous in their full weighted seminorm topologies. They preserve
every supplied smooth finite-dimensional family on its retained domain.
For retained compact parameter sets
\(\mathcal K\Subset\operatorname{int}\mathcal K'\), each specified
finite output estimate on \(\mathcal K\) is controlled by finitely many
input spatial and parameter seminorms on \(\mathcal K'\) and the fixed
base margins.

If the input is given by an admitted actual cochain pair, the result is
independent of its effective ordinary presentation. Its comparison with
the canonical representative uses a primitive in the coefficient complex
satisfying the finite estimates in
\hyperref[thm:online-actual-cochain-obstruction-test]{the obstruction
theorem}.

\emph{Proof.} For an admitted ordinary pair, set
\(s=\mathscr P_\gamma\rho(c)\). The selected-domain definition gives
\(s\in\mathcal B_\gamma\). The base chart gives
\(\mathcal T(s)=\mathcal A(s,\mathcal C(s))\) satisfying the metric
gauge, maximality and the two vacuum constraints.
\hyperref[thm:online-chart-one-branch-weighted-smoothness]{The all-order
theorem} puts its tensors in the required symbol spaces. The exact
ambient decoder returns \((s,\mathcal C(s))\), so the two neighborhood
tests also hold. Hence the forward map lands in the independently
defined target.

Conversely, take an independently admitted tensor pair and decode it to
\((s,y)\). By definition \(s\in\mathcal B_\gamma\subset\mathcal U\) and
\(y\in\mathcal V\). The ambient identities give
\((h,K)=\mathcal A(s,y)\). Its vacuum constraints are exactly
\(\mathcal F(s,y)=0\). Uniqueness of the base correction therefore
forces \(y=\mathcal C(s)\). This proves that its inverse-class formula
constructs an ordinary input whose output is exactly the initial tensor
pair.

On the ordinary domain,
\hyperref[cor:online-effective-ordinary-class-identification]{the actual
normalization theorem and its corollary} give
\(\sigma_\gamma\rho(c)=c\). Together with the two inverse identities for
\(\mathscr P_\gamma\) and the seed decoder, this proves
\(\mathfrak C\mathfrak D=I\). The identity \(\rho\sigma_\gamma=I\) and
the preceding constraint-uniqueness argument prove
\(\mathfrak D\mathfrak C=I\).

Continuity and the finite-family estimates follow in the order in which
the maps are applied: the exact root transform and actual section, the
same-branch nonlinear chart, and the exact ambient decoder. For any
admitted cochain presentation, the polynomial-obstruction theorem
constructs the comparison primitive at fixed coefficient orders and
bounded polynomial degrees. Its effective-arrow assertion makes the
resulting ordinary class independent of that representative. Every
finite estimate is a finite composition of the displayed estimates. The
one base neighborhood is fixed before choosing any higher output order.
\hfill\(\square\)

\pandocstatementheading{\leavevmode\phantomsection\label{cor:online-ordinary-pair-ungauged-access}Corollary 3.14 (Whole maximal data with the based comparison retained).}

Let \((H,L)\) be whole smooth maximal vacuum data with
\(H-\delta\in\mathscr S^\gamma\) and \(L\in\mathscr S^{\gamma-1}\),
sufficiently small in the finite weighted norms of
\hyperref[cor:online-gauge-whole-data-decoding]{Corollary A.11}. After
one restriction of that threshold, there are the initialized based gauge
\(\phi_H\) and a unique ordinary pair
\(c_H\in\mathscr C^{\mathrm{ord}}_{\gamma,\mathcal B}\) such that

\phantomsection\label{eq:online-ordinary-based-whole-data-recovery} \[
\phi_H^*(H,L)=\mathfrak D(c_H),
\qquad
c_H=\sigma_\gamma\mathscr P_\gamma^{-1}
       \operatorname{Dec}_{s}\bigl(\phi_H^*(H,L)\bigr).
\] The comparison retains the actual spin lift initialized along the
based gauge isotopy. Its gauge, ordinary representatives and decoded
data have all specified finite spatial and parameter estimates on the
same retained family domain, with the compact-set and finite-input
quantifiers of the preceding theorem.

\emph{Proof.} Corollary A.11 supplies the actual gauge and an exact seed
in the selected constraint chart. The gauge estimate and continuity of
its decoder at the base order allow the input threshold to be reduced
once so that this seed lies in the prescribed neighborhood
\(\mathcal B_\gamma\). Apply the preceding theorem. Its uniqueness is
uniqueness of the ordinary pair for the normalized whole datum in the
retained chart; the based gauge itself is unique in its displacement
ball. The spin and family assertions are those of
\hyperref[prop:online-gauge-rotation-and-spin]{the gauge comparison}
followed by the finite estimates of the preceding theorem.
\hfill\(\square\)

\hyperref[thm:online-ordinary-pair-data-bijection]{Theorem 3.13}
recovers the normalized whole datum from its ordinary pair. Section 4
constructs the chiral gradient and contact operators of the same pair.
\hyperref[thm:online-ordinary-negative-response-initial-source]{Theorem
4.18} proves that the operators recover the whole datum. The local
coverage argument in Section 5 uses the following initial-data results.

{\def\LTcaptype{none} % do not increment counter
\begin{longtable}[]{@{}
  >{\raggedright\arraybackslash}p{(\linewidth - 2\tabcolsep) * \real{0.5000}}
  >{\raggedright\arraybackslash}p{(\linewidth - 2\tabcolsep) * \real{0.5000}}@{}}
\toprule\noalign{}
\begin{minipage}[b]{\linewidth}\raggedright
Input to local coverage
\end{minipage} & \begin{minipage}[b]{\linewidth}\raggedright
Proof in this article
\end{minipage} \\
\midrule\noalign{}
\endhead
\bottomrule\noalign{}
\endlastfoot
Weighted root pair and exact TT seed &
\hyperref[thm:online-root-seed-isomorphism]{Theorem 3.6} \\
Actual ordinary lift and independent cochain tests &
\hyperref[thm:online-canonical-actual-pair-lift]{Theorems 3.9--3.10} and
\hyperref[cor:online-effective-ordinary-class-identification]{Corollary
3.12} \\
Nonlinear constraints and exact whole decoder &
\hyperref[thm:online-ordinary-pair-data-bijection]{Theorem 3.13},
\hyperref[prop:online-chart-exact-ambient-decoder]{Proposition A.4} and
\hyperref[thm:online-chart-base-constraint-bijection]{Theorems
A.6--A.7} \\
Based coordinate gauge with actual spin &
\hyperref[thm:online-gauge-based-normalization]{Theorem A.9},
\hyperref[prop:online-gauge-rotation-and-spin]{Proposition A.10} and
\hyperref[cor:online-ordinary-pair-ungauged-access]{Corollary 3.14} \\
\end{longtable}
}

\section{4. Actual source execution from ordinary
classes}\label{actual-source-execution-from-ordinary-classes}\label{sec:online-actual-source-execution}

Starting from the fixed flat source, we construct the initial gradient
and contact operators of a selected ordinary pair. The weighted
negative-to-dual operation acts on closed cochain representatives and
includes the primitives that compare them. Contact normalization over
the complete sphere then gives the negative response. A finite
composition of those responses realizes the scalar and York inverses
required by the nonlinear current correction.

The source curvature supplies four independently defined defects. At
each current source, its vertical response agrees with the transverse
differential of these defects. The resulting horizontal field changes
the prescribed ordinary seed and preserves the defects. Its endpoint is
selected on one finite base neighborhood, and its weighted smoothness
and family estimates concern that same endpoint. The argument in
\hyperref[sec:online-current-flat-input-transport]{Appendix A} supplies
actual Picard convergence, path independence and a finite-execution
error bound.

The canonical initial target is specified by tests on supplied operators
or complete contact coefficients. The resulting unit and counit identify
this target with the ordinary-pair domain. Full spacetime admission,
wave enrichment and changes of initial carrier use the additional
structures identified in Section 5.

\subsection{4.1. The actual weighted negative-to-dual
operation}\label{the-actual-weighted-negative-to-dual-operation}\label{sec:online-actual-negative-dual}

The negative source determines a positive class for the conjugate
projective sphere. This is a complex-linear operation on the original
negative class. Its conjugate complex structure changes the angular
differential; the operation itself preserves the real spatial Fourier
frequency. We first fix the original scalar calibration, then prove the
operation on actual weighted cochains and exhibit the primitives
required to compare their representatives.

Retain the flat conic \(z\), Fourier normalization and physical
root-line identifications of
\hyperref[sec:online-root-calibration-and-seeds]{the root transform}.
Write \(D=z\cdot\partial_x\) and \(D^\vee=\bar z\cdot\partial_x\). The
conjugate positive line is \(\overline{\mathcal O(2)}\): it has degree
\(2\) for the conjugate sphere complex structure and Chern number \(-2\)
for the original sphere orientation. Its differential and horizontal
frequency multiplier are

\phantomsection\label{eq:online-dual-positive-differential} \[
\begin{aligned}
d_\vee^0v&=(\partial v,D^\vee v),&
d_\vee^1(\alpha,\beta)&=\partial\beta-D^\vee\alpha,\\
L_\xi^\vee&=i\xi\cdot\bar z,&
L_\xi&=i\xi\cdot z .
\end{aligned}
\] Here \(\partial\) acts only in the projective fiber. In particular
\(L_\xi^\vee\) is not \(\overline{L_\xi}\) at fixed \(\xi\). Conjugation
of an actual spatial coefficient also reflects its frequency.

The roots remain \(n_\varepsilon=\varepsilon\xi/r\), where \(r=|\xi|\)
and \(\varepsilon=\pm1\). Write \(U^-_\varepsilon\) for the original
negative root tensor and \(V^\vee_\varepsilon\) for the dual positive
root tensor. Both take values in the line generated by
\(z(n_\varepsilon)\otimes z(n_\varepsilon)\). At the north root the
scalar-to-tensor maps are

\phantomsection\label{eq:online-dual-root-line-normalizations} \[
c\longmapsto c\,z\otimes z
\quad\text{and}\quad
c\longmapsto \frac c4\,z\otimes z
\] for the negative and dual positive weights, respectively. The second
is the conjugate of the positive root-line normalization already used.

\pandocstatementheading{\leavevmode\phantomsection\label{prop:online-original-negative-dual-calibration}Proposition 4.1 (The original derived-dual multiplier).}

On smooth Fourier amplitudes with compact support away from zero, the
original analytic derived-dual class has root tensors

\phantomsection\label{eq:online-negative-dual-calibrated-multiplier} \[
V^\vee_\varepsilon(\xi)=\frac1{16r^4}U^-_\varepsilon(\xi).
\] The multiplier is complex-linear, has positive sign, and uses the
same frequency and root label.

\emph{Proof.} The original dual coefficient formula
(\hyperref[ref-IsthmusLocal2026]{Isthmus 2026a}, §5, equation for the
derived dual class) sends the negative plane cocycle
\(F_-=A e^{\lambda w^0}/\zeta\), \(\lambda\ne0\), to

\phantomsection\label{eq:online-negative-plane-derived-dual} \[
G(\eta,\varrho)=\frac{A\lambda^{-4}}{\varrho}
\left(e^{-\lambda\eta^0}
-\sum_{j=0}^{3}\frac{(-\lambda\eta^0)^j}{j!}\right).
\] Indeed its original coefficients are \(c^-_{j,0,0}=A\lambda^j/j!\),
and the dual formula multiplies this coefficient by
\((-1)^j j!/(j+4)!\). Summation gives the displayed expression. Each
subtracted term is a positive-weight coboundary: \((\eta^0)^j/\varrho\)
becomes \((\tilde\eta^0)^j\tilde\varrho^{\,3-j}\) in the weight-two
south frame, with \(\tilde\eta=\eta/\varrho\) and
\(\tilde\varrho=1/\varrho\), and is holomorphic there when \(j\le3\).

The derived class is consequently represented by
\(A\lambda^{-4}e^{-\lambda\eta^0}/\varrho\). Lorentz incidence is
\(\eta=-\bar w\), \(\varrho=\bar\zeta\). Its restriction is
\(A\lambda^{-4}e^{\lambda\bar w^0}/\bar\zeta\). Apply the
conjugate-angular version of the plane-to-root calculation preceding
\hyperref[def:online-weighted-root-domain]{the weighted root
definition}, keeping \(\lambda\) and \(A\) unchanged. The relevant
identities are \(D^\vee u=-\sqrt2\bar\zeta\) and
\(D^\vee\bar v=\sqrt2\). Subtracting the same smooth cutoff primitive
gives

\phantomsection\label{eq:online-negative-dual-plane-root-coefficients}
\[
c^\vee=-\sqrt2\lambda A\lambda^{-4}
=\lambda^{-4}c^-,
\qquad c^-=-\sqrt2\lambda A.
\] The pole exponent is still \(e^{\lambda u}\), so the real spatial
frequency is unchanged. At \(\xi=\varepsilon r e_3\), the fixed
soldering convention gives \(\lambda=-i\sqrt2\varepsilon r\), hence
\(\lambda^4=4r^4\). Combining this with the additional tensor-line
factor \(1/4\) proves the multiplier and its sign.

The conic, incidence and original dual spin pairing commute with spatial
spin rotations, so the axial calculation applies to both root branches.
Let the frequency support lie in \(r_0\le |\xi|\le R_0\) with \(r_0>0\).
For each branch take a finite partition of unity in direction, supported
in small caps with smooth spin frames. On each cap fix the angular chart
at its central direction, so the corresponding pole has coordinate
\(a(\xi)\) and \(a=0\) at the cap centre. Fix a bounded annulus
\(\rho<|\zeta|<R\) in this chart. The rotation at the centre is the
identity in the chosen frame. By shrinking the cap its Möbius
denominator stays bounded away from zero on the closed annulus and
\(|a(\xi)|\le\rho/2\). Thus \(|\zeta-a(\xi)|\ge\rho/2\) throughout the
overlap. Use the south coordinate at infinity and the conjugate
coordinate for the dual cocycle.

Restrict the remaining complex coordinates to a common bounded open set
with closure in the coordinate domains, taking the minimum of the
available radii over the finite cover. For either the plane or its
derived cocycle, write \(F_\xi(Z)\) for its coefficient in these fixed
complex coordinates. If \(S\) is the supported frequency set in the cap
and \(K\) is relatively compact in the resulting overlap,
\(\sup_{(\xi,Z)\in S\times K}|\partial_\xi^\beta D_Z^\nu F_\xi(Z)|\) is
finite for every specified pair of multi-indices. Indeed the frame
factors have bounded derivatives there, \(|\lambda|\) is bounded above
and away from zero, and the exponential series and all its finite
derivatives converge normally on a slightly larger compact set. These
bounds justify integration against the smooth amplitude term by term,
differentiation under the frequency integral and the contour computation
of the root coefficient on a fixed circle enclosing the moving pole. In
the real cochain calculation, subtract the cutoff primitive before
reading the root, as in §3.3. On the real initial slice the corrected
coefficients have the spatial phase \(e^{ix\cdot\xi}\) times smooth
compactly supported frequency amplitudes, also after fiber
differentiation. Integration by parts in \(\xi\) gives rapid spatial
decay of their superposition. \hfill\(\square\)

\pandocstatementheading{\leavevmode\phantomsection\label{def:online-negative-dual-weighted-complexes}Definition 4.2 (Weighted actual negative and dual complexes).}

Fix \(-1<\gamma<-1/2\) and put \(\vartheta=\gamma+1\in(0,1/2)\). Use the
actual symbol spaces, including every fiber derivative, of
\hyperref[sec:online-actual-coefficient-complexes]{the coefficient
algebra}. The negative complex is

\phantomsection\label{eq:online-negative-actual-weighted-complex} \[
\begin{aligned}
\mathscr C_{-,\vartheta}^{0}
&=\mathscr S^{\vartheta-3}(\mathcal O(-6)),\\
\mathscr C_{-,\vartheta}^{1}
&=\mathscr S^{\vartheta-3}
  (\Lambda^{0,1}\otimes\mathcal O(-6))
\oplus\mathscr S^{\vartheta-4}(\mathcal O(-4)),\\
\mathscr C_{-,\vartheta}^{2}
&=\mathscr S^{\vartheta-4}
  (\Lambda^{0,1}\otimes\mathcal O(-4)),
\end{aligned}
\] with \(d_-^0u=(\bar\partial u,Du)\) and
\(d_-^1(\alpha,\beta)=\bar\partial\beta-D\alpha\). The dual positive
complex is

\phantomsection\label{eq:online-dual-positive-actual-weighted-complex}
\[
\begin{aligned}
\mathscr C_{\vee,\vartheta}^{0}
&=\mathscr S^{\vartheta+1}(\overline{\mathcal O(2)}),\\
\mathscr C_{\vee,\vartheta}^{1}
&=\mathscr S^{\vartheta+1}
  (\Lambda^{1,0}\otimes\overline{\mathcal O(2)})
\oplus\mathscr S^{\vartheta}(\overline{\mathcal O(4)}),\\
\mathscr C_{\vee,\vartheta}^{2}
&=\mathscr S^{\vartheta}
  (\Lambda^{1,0}\otimes\overline{\mathcal O(4)}),
\end{aligned}
\] with the differential in
\hyperref[eq:online-dual-positive-differential]{the dual formula}. These
are actual complexes: a horizontal derivative lowers the spatial order
by one and commutes with the fiber differential.

Their root spaces consist of tensor sections on the specified two root
lines with rapid decrease at infinity, with all frequency derivatives,
and bounds

\phantomsection\label{eq:online-negative-dual-weighted-root-bounds} \[
\begin{aligned}
|\partial_\xi^\mu U^-_\varepsilon(\xi)|
&\le C_\mu r^{1-\vartheta-|\mu|},\\
|\partial_\xi^\mu V^\vee_\varepsilon(\xi)|
&\le C_\mu r^{-\vartheta-3-|\mu|},
\qquad 0<r\le1.
\end{aligned}
\] Derivatives are taken in the tensor realization or finitely many
angular root frames. The complexes carry their full symbol seminorms,
their cohomologies the quotient topologies, and the root spaces the
displayed amplitude seminorms. These are marginal weighted domains; an
auxiliary negative input is not required by itself to be a Lorentz-fixed
two-class pair.

\pandocstatementheading{\leavevmode\phantomsection\label{thm:online-actual-weighted-negative-dual}Theorem 4.3 (Actual weighted dualization and its comparison primitives).}

Root readout identifies the degree-one cohomology of each displayed
actual complex with its root space. The two maps are continuous
isomorphisms, with continuous choices of actual closed representatives
and actual primitives for closed cochains of zero root. Consequently

\phantomsection\label{eq:online-actual-negative-dual-isomorphism} \[
\mathcal D_-^{\mathrm{dual}}:
H^1(\mathscr C_{-,\vartheta}^{\bullet})
\xrightarrow{\ \sim\ }
H^1(\mathscr C_{\vee,\vartheta}^{\bullet}),
\qquad
V^\vee_\varepsilon=\frac1{16r^4}U^-_\varepsilon
\] is a continuous complex-linear isomorphism, with inverse multiplier
\(16r^4\). It extends the original derived-dual operation in the
preceding proposition.

For the fixed real anchor \(x=0\), its selected dual representative is

\phantomsection\label{eq:online-actual-dual-anchored-representative} \[
\begin{aligned}
e_\vee&=(0,B_\vee),\\
T_\vee(x)
&=(2\pi)^{-3/2}\int_{\mathbb R^3}
 (e^{ix\cdot\xi}-1)
 \sum_{\varepsilon=\pm1}V^\vee_\varepsilon(\xi)
 \,\mathrm d\xi,\\
B_\vee(x,\zeta)&=T_\vee(x)(\bar z(\zeta),\bar z(\zeta)).
\end{aligned}
\] The tensor \(T_\vee\) is smooth, transverse-traceless, of order
\(\vartheta\), and vanishes at the anchor. Every closed cochain in
either weighted complex is recovered up to an actual coboundary within
that same complex. The change of dual anchor is the explicit coboundary
in that complex. Inclusion of either weighted complex into the
corresponding unrestricted actual coefficient complex is injective on
degree-one cohomology.

Each specified output seminorm, also after specified finite parameter
differentiation, is bounded by finitely many input seminorms. The
weight, anchor and coefficient domains are fixed before those orders are
chosen.

\emph{Proof.} We first classify the negative cochains. For
\(e_-=(\alpha,\beta)\), subtract \(d_-^0(T_{-6}\alpha)\), using
\hyperref[eq:online-root-fiber-homotopy]{the fixed fiber Hodge
operators}. The first component becomes
\(\alpha_0=H^1_{-6}\alpha\in\mathscr S^{\vartheta-3}\). Closedness and
\(H^0(\mathcal O(-4))=0\) give the actual identity

\phantomsection\label{eq:online-negative-weighted-hodge-normalization}
\[
e_--d_-^0(T_{-6}\alpha)
=(\alpha_0,T_{-4}D\alpha_0),
\qquad H^1_{-4}D\alpha_0=0 .
\] Since \(\vartheta-3\in(-3,0)\),
\hyperref[lem:online-decay-fourier-isomorphism]{the two-sided Fourier
estimate} applies to \(\alpha_0\). Its punctured amplitude has order
\(r^{-\vartheta}\) near zero, with the corresponding derivative bounds.
Root evaluation of \(T_{-4}L_\xi\widehat\alpha_0\) has degree one and
gives the asserted \(U^-\) bounds.

Conversely, the negative root inverse in
\hyperref[thm:online-moderate-root-cohomology]{the root-cohomology
theorem} has degree \(-1\) for its harmonic first coefficient. Applied
to the stipulated \(U^-\), it gives \(\alpha_-(\xi)\) with bounds
\(r^{-\vartheta-|\mu|}\). This coefficient and every spatial Fourier
multiple are locally integrable. The inverse integral gives an actual
\(A_-\in\mathscr S^{\vartheta-3}\), and

\phantomsection\label{eq:online-negative-dual-selected-negative-representative}
\[
e_-(U^-)
=(A_-,T_{-4}D A_-)
\in\mathscr C_{-,\vartheta}^1 .
\] For its second component, the known amplitude is
\(T_{-4}L_\xi\alpha_-\), of low-frequency order \(r^{1-\vartheta}\). The
one-way inverse estimate, valid for every negative spatial exponent,
gives order \(\vartheta-4\). This does not use a two-sided Fourier
characterization below \(-3\). The harmonic identity holds for locally
integrable amplitudes; multiplication by \(\xi\) introduces no
point-supported term. The fiber homotopy identity consequently proves
exact closedness.

For a supplied closed cochain with these roots, \(\alpha_0-A_-\) has
zero punctured Fourier transform.
\hyperref[prop:online-actual-coefficient-exact-sequence]{The actual
coefficient exact sequence} makes it polynomial, and its negative order
makes it zero. The normalized second components then agree by the
absence of holomorphic sections of \(\mathcal O(-4)\). Thus the actual
comparison primitive is \(T_{-6}\alpha\) and has the required order
\(\vartheta-3\). This proves both directions on the negative side.

For the dual side, use the conjugate fiber Hodge operator
\(T_2^\vee=\overline{T_2\,\overline{(\,\cdot\,)}}\), which satisfies the
analogous \(\partial\) identities. Set

\phantomsection\label{eq:online-dual-positive-hodge-normalization} \[
v_0=T_2^\vee\alpha,\qquad
\beta_0=\beta-D^\vee v_0.
\] Here \(v_0\in\mathscr S^{\vartheta+1}\),
\(\beta_0\in\mathscr S^\vartheta\), and closedness makes \(\beta_0\) an
anti-holomorphic quartic on the whole sphere. The positive-order annular
Fourier estimate in the proof of
\hyperref[prop:online-actual-coefficient-exact-sequence]{Proposition
3.7} gives \(\partial_\xi^\mu\widehat\beta_0=O(r^{-\vartheta-3-|\mu|})\)
on punctured frequency space. Its root values therefore obey the
required dual bounds.

Given such root values, the tensor sum
\(\widehat T_\vee=\sum_\varepsilon V^\vee_\varepsilon\) is
transverse-traceless, and its quartic

\phantomsection\label{eq:online-dual-canonical-frequency-quartic} \[
\beta_\vee(\xi,\zeta)
=\left(\sum_\varepsilon V^\vee_\varepsilon(\xi)\right)
  (\bar z(\zeta),\bar z(\zeta))
\] has exactly those roots. The complementary tensor line evaluates to
zero at either root, while the retained line evaluates with the
normalization in \hyperref[eq:online-dual-root-line-normalizations]{the
root-line formula}. Thus this is a right inverse to root readout.

The integral defining \(T_\vee\) is one bracket integral; its two terms
are not integrated separately. Near zero the bracket contributes one
power of \(r\), giving the radial majorant
\(Cr^{-\vartheta}\,\mathrm dr\) on each compact spatial set. Every
positive spatial derivative is also integrable. The derivative
amplitudes have spatial exponent \(\vartheta-1<0\), so the one-way
inverse estimate gives \(\partial T_\vee\in\mathscr S^{\vartheta-1}\)
with all higher derivatives. Integration from the anchor yields
\(T_\vee\in\mathscr S^\vartheta\). The differentiated integrals identify
its punctured Fourier transform, and the anchor fixes the constant left
by differentiation. These estimates hold also after fiber
differentiation.

Its trace vanishes under the convergent bracket integral. Its divergence
vanishes under the convergent first-derivative integral because
\(\xi_iV^\vee_{\varepsilon,ij}=0\). The coefficient \(B_\vee\) is
anti-holomorphic in the fiber, so \(e_\vee\) is exactly closed. This
proves surjectivity and the claimed representative estimates on the dual
side.

It remains to compare a general \(\beta_0\) to that representative. At
nonzero frequency their difference has zero roots. Multiplication by
\(L_\xi^\vee\) maps anti-holomorphic quadratics injectively to
anti-holomorphic quartics; the root divisor sequence identifies its
image with the zero-root quartics. Choose its invariant Hermitian left
inverse \((L_\xi^\vee)^+\) and put

\phantomsection\label{eq:online-dual-comparison-primitive-amplitude} \[
\widehat w=(L_\xi^\vee)^+
               (\widehat\beta_0-\beta_\vee).
\] This is an anti-holomorphic quadratic coefficient with low-frequency
bounds \(|\partial_\xi^\mu\widehat w|
\le C_\mu r^{-\vartheta-4-|\mu|}\) and rapid decrease at infinity. The
left inverse is smooth homogeneous of degree \(-1\): at \(r=1\) its
nonzero singular values have a positive minimum on the compact direction
sphere, as in the root-cohomology proof.

Use the twice-anchored integral

\phantomsection\label{eq:online-dual-twice-anchored-primitive} \[
w(x)=(2\pi)^{-3/2}\int_{\mathbb R^3}
\bigl(e^{ix\cdot\xi}-1-ix\cdot\xi\bigr)
\widehat w(\xi)\,\mathrm d\xi .
\] Its bracket gains two powers of \(r\) near zero, again giving the
integrable radial majorant \(Cr^{-\vartheta}\,\mathrm dr\). Rapid
decrease controls infinity, including the subtracted terms. The second
derivatives are ordinary inverse integrals with spatial order
\(\vartheta-1\). Consequently \(D_x^2w\in\mathscr S^{\vartheta-1}\) at
all orders. Since \(w(0)=0\) and \(D_xw(0)=0\), integration along
segments gives \(D_xw\in\mathscr S^\vartheta\) and
\(w\in\mathscr S^{\vartheta+1}\). The same argument controls every fiber
derivative. This proves the required actual primitive order directly.

Its punctured Fourier transform is \(\widehat w\), as follows by taking
two spatial derivatives. Hence

\phantomsection\label{eq:online-dual-constant-polynomial-residual} \[
R=\beta_0-B_\vee-D^\vee w
\] has zero punctured transform and spatial order \(\vartheta<1\). It is
a constant anti-holomorphic quartic, specifically \(R=\beta_0(0)\).
Write \(R=C_{ij}\bar z^i\bar z^j\) with \(C\) tracefree symmetric, and
set \(w_C=C_{ij}x^i\bar z^j\). Then \(\partial w_C=0\) and
\(D^\vee w_C=R\). The actual identity is

\phantomsection\label{eq:online-dual-actual-comparison-primitive} \[
(\alpha,\beta)-(0,B_\vee)
=d_\vee^0(v_0+w+w_C).
\] Every term on the right has degree-zero order \(\vartheta+1\). This
proves classification within the displayed complex, without discarding a
polynomial first-cohomology class. A primitive is determined modulo a
closed zero-cochain in that same complex; the formula provides one
continuous choice.

All operations used are fixed compact-fiber Hodge maps, homogeneous
multipliers with finite order shifts, convergent anchored integrals and
finite tensor readouts. Each output seminorm therefore uses finitely
many input seminorms. Root readout is continuous by the forward
estimates; the representative maps are continuous by the inverse
estimates. The explicit comparisons identify the induced inverse maps on
the quotient topologies. Multiplication by \(r^{-4}/16\) and its inverse
interchange the displayed root bounds, proving the actual cohomology
isomorphism and both inverse identities.

If a weighted closed cochain is exact in the unrestricted actual
complex, its root readout is zero. The just-proved weighted
classification then supplies a primitive within the weighted complex
itself. This proves the two asserted injections into unrestricted actual
cohomology. It also identifies the negative component of
\hyperref[thm:online-canonical-actual-pair-lift]{the selected ordinary
pair} with a unique class in the present weighted negative group.

An anchor \(a\) changes \(B_\vee\) to \(B_\vee(x)-B_\vee(a)\). If
\(B_\vee(a)=C_a(\bar z,\bar z)\), the change is the actual coboundary of
\(-C_{a,ij}x^i\bar z^j\), of allowed order \(\vartheta+1\). Successive
anchor comparisons add exactly.

Finally the constructions are continuous linear maps between the stated
spaces. They commute with differentiation of any supplied smooth
finite-dimensional family in those seminorms. On retained compact
parameter sets, applying the finite seminorm estimates to each
prescribed derivative gives the asserted bounds. No choice of a
primitive by an additional smooth-selection theorem is required.
\hfill\(\square\)

\pandocstatementheading{\leavevmode\phantomsection\label{prop:online-negative-dual-reality-and-markings}Proposition 4.4 (Coupled reality and the retained markings).}

For an admitted ordinary pair \(c=(c_+,c_-)\), the calibrated residue
reality is equivalent to

\phantomsection\label{eq:online-negative-dual-coupled-reality} \[
\mathcal D_-^{\mathrm{dual}}c_-=\overline{c_+},
\qquad
V^\vee_\varepsilon(\xi)
=\overline{U^+_{-\varepsilon}(-\xi)} .
\] Here \(\overline{c_+}\) is formed by conjugating an actual positive
representative and using the conjugate positive complex. With the common
real anchor and the transverse-traceless quartic representatives, the
equality is literal conjugation of actual closed cochains. Other
retained representatives have the actual comparison primitives proved
above.

The class transformation commutes with translations and the retained
spatial spin rotations. A translation of its anchored representative is
accompanied by the explicit anchor coboundary. All these operations
preserve the stipulated smooth finite-dimensional families and their
finite seminorm estimates.

\emph{Proof.} Conjugation reflects the Fourier frequency, and
\(n_\varepsilon(\xi)=n_{-\varepsilon}(-\xi)\). The original coupled
reality is \[
U^-_\varepsilon(\xi)
=16r^4\,\overline{U^+_{-\varepsilon}(-\xi)}.
\] The multiplier proved above gives precisely the displayed dual root
equality, with its converse obtained by multiplying by \(16r^4\). Root
classification and injectivity into actual cohomology turn this equality
into the claimed class identity.

To see the representative assertion, choose the positive representative
\(e_+=(0,S_+(z,z))\), with

\phantomsection\label{eq:online-positive-tt-anchored-representative} \[
S_+(x)=(2\pi)^{-3/2}\int
(e^{ix\cdot\xi}-1)\sum_\varepsilon U^+_\varepsilon(\xi)
\,\mathrm d\xi .
\] It is the conjugate-angular counterpart of the dual representative
construction, has order \(\vartheta\), and represents the original
positive class by the conjugate of the theorem's actual comparison. This
also compares it with any retained harmonic normal representative from
Theorem 3.9. In the bracket integral, conjugation followed by
\(\xi\mapsto-\xi\) returns the same bracket at a real anchor. The root
identity therefore gives \(T_\vee=\overline{S_+}\) and
\(B_\vee=\overline{S_+(z,z)}\). Both vertical components vanish, proving
literal cochain equality for these representatives.

The multiplier is radial on geometrically identified root lines. It
commutes with the translation factor \(e^{ia\cdot\xi}\) and with spatial
spin rotations. The prescribed representative after translation is
\(T_\vee(x+a)-T_\vee(a)\); its comparison with the literal translate is
exactly the anchor coboundary in the theorem. Rotation equivariance
follows from the invariant Hodge choices and fixed conic
identifications. The continuous linear formulas give the asserted
parameter statements. \hfill\(\square\)

The dual class is thus derived from the negative ordinary input; it is
not an additional independent class. We retain the spin marking
separately because the even projective weights cannot distinguish the
two central lifts. On unrestricted actual cohomology, the surviving
polynomial kernel prevents defining this operation by root values alone.
The construction above is a whole-space weighted operation with actual
primitives; local source restrictions are taken after it has been
performed.

\subsection{4.2. Actual contact cochains and the normalized linear
response}\label{actual-contact-cochains-and-the-normalized-linear-response}\label{sec:online-actual-contact-response}

The fixed inverse needed by the nonlinear correction will be constructed
from actual negative cochains. We first define the contact operation
used in each auxiliary call. This identifies its complete linear initial
graph before any Fourier multiplier is read from that graph.

On the marked flat carrier \(N_0=\mathbb R^3\times\mathbb{CP}^1\),
retain the conic and Fourier conventions of
\hyperref[sec:online-root-complex]{the root calculation}. In a north
spin chart put

\phantomsection\label{eq:online-contact-response-flat-frame} \[
\begin{gathered}
d=1+|\zeta|^2,\qquad f=d^{-2},\qquad
a_0=\frac{2\bar\zeta}{d},\qquad b_0=-\frac{2\zeta}{d},\\
n=\frac{i\,z\times\bar z}{z\cdot\bar z},\qquad
\alpha_0=n\cdot\mathrm dx,\qquad R=n\cdot\partial_x,\\
V=\partial_{\bar\zeta},\qquad U=\partial_\zeta,\qquad
H=z\cdot\partial_x,\qquad \bar H=\bar z\cdot\partial_x.
\end{gathered}
\] Here \(\mathcal A_0=\operatorname{span}_{\mathbb C}(V,H)\) is
involutive,
\(\overline{\mathcal A}_0=\operatorname{span}_{\mathbb C}(U,\bar H)\),
and \(R\) is the Reeb field of \(\alpha_0\). For a complex contact
Hamiltonian \(h\), fix \[
\alpha_0(X_h)=h,\qquad
\iota_{X_h}\mathrm d\alpha_0=(Rh)\alpha_0-\mathrm dh.
\] The Bott connection on \(Q=T_{\mathbb C}N_0/\mathcal A_0\) is
\(\nabla^B_Y[Z]=[Y,Z]\bmod \mathcal A_0\) for \(Y\in \mathcal A_0\). Let
\(d_B\) be its differential and set \(\mathcal G_Bh=d_B[X_h]\). A
horizontal Bott cocycle is contact-compatible when
\(\mathrm d\alpha_0(\Phi Y,Z)=\mathrm d\alpha_0(\Phi Z,Y)\) for
\(Y,Z\in \mathcal A_0\).

The positive and dual positive coefficient lifts are the global scalar
operators

\phantomsection\label{eq:online-actual-chiral-contact-lifts} \[
\mathcal D_+u=-f\bar H u,\qquad
\mathcal D_-v=-fHv.
\] Their domains are, respectively, \(\mathcal O(2)\) and its conjugate
line on the real sphere. On a north-south overlap,
\(u_S=\zeta^{-2}u_N\), \(v_S=\bar\zeta^{-2}v_N\), \(z_S=\zeta^{-2}z_N\),
and \(d_S=d_N/|\zeta|^2\). The factors cancel in both formulas, since
the spatial derivative does not differentiate an angular transition. The
lifts therefore act on actual smooth sections, with one spatial
derivative and the corresponding symbol-order decrease.

These factors retain the original contact normalization. With the
soldering \(X\) of §3.3, put \(w=X(0,x)\pi\), \(\pi=(1,\zeta)\),
\(\kappa=(1,\bar\zeta)\) and \(\eta=-\kappa X(0,x)\). The original
contact form \(\theta=W\,\mathrm dZ\) restricts to
\(-(d/\sqrt2)\alpha_0\), and direct differentiation gives

\phantomsection\label{eq:online-contact-response-incidence-derivatives}
\[
\begin{aligned}
\bar H w^0&=-\sqrt2d\,\bar\zeta,&
\bar H w^1&=\sqrt2d,\\
H\eta^{0'}&=\sqrt2d\,\zeta,&
H\eta^{1'}&=-\sqrt2d.
\end{aligned}
\] Consequently the original lifts \(F_{w^1}-\bar\zeta F_{w^0}\) and
\(-G_{\eta^{1'}}+\zeta G_{\eta^{0'}}\), after contact-line conversion by
\(-\sqrt2/d\), give precisely \(\mathcal D_+\) and \(\mathcal D_-\) on
actual holomorphic restrictions (\hyperref[ref-IsthmusLocal2026]{Isthmus
2026a}, §5 and Appendix F.1). The displayed differential formulas define
their actual smooth extensions here. The negative lift acts after the
derived dual-class transformation; its domain is not the original
weight-minus-six line.

\pandocstatementheading{\leavevmode\phantomsection\label{prop:online-actual-contact-cochain-map}Proposition 4.5 (The actual cochain-to-Bott map).}

For an actual closed positive cochain \(e=(\alpha,\beta)\), whose local
equation is \(V\beta=H\alpha\), define

\phantomsection\label{eq:online-contact-positive-cochain-map} \[
\begin{aligned}
\Gamma_+(\alpha,\beta)(V)
&=\frac f2\bigl[\bar H(V\alpha-b_0\alpha)+4R\alpha-2f\beta\bigr]\bar H
 -\frac f2\bigl[\bar H(V\beta)+2R\beta\bigr]U,\\
\Gamma_+(\alpha,\beta)(H)
&=\frac f2\bigl[\bar H(H\alpha)+2R\beta\bigr]\bar H
 -\frac f2\bar H(H\beta)U.
\end{aligned}
\] This is a global linear differential operation of order at most two
on closed cochains. Its output is a horizontal contact-compatible Bott
cocycle, and

\phantomsection\label{eq:online-contact-cochain-gauge-identity} \[
\Gamma_+d^0u=\mathcal G_B\mathcal D_+u.
\] It commutes with spatial restriction and smooth passive parameter
differentiation. On each finite collection of retained charts, every
finite output norm is bounded by finitely many stronger input norms. The
conjugate construction

\phantomsection\label{eq:online-contact-dual-cochain-map} \[
\Gamma_\vee(e^\vee)=\overline{\Gamma_+(\overline{e^\vee})}
\] is complex-linear from actual closed dual positive cochains to the
corresponding Bott cocycles over \(\overline{\mathcal A}_0\), with
\(\Gamma_\vee\bar d^0=\bar{\mathcal G}_B\mathcal D_-\), where
\(\bar{\mathcal G}_Bh=\overline{\mathcal G_B\bar h}\).

\emph{Proof.} The conic formulas give

\phantomsection\label{eq:online-contact-response-commutators} \[
\begin{gathered}
[V,\bar H]=-b_0\bar H+2R,\qquad
[H,U]=-a_0H-2R,\qquad [V,R]=-fH,\\
[V,H]=[H,\bar H]=[H,R]=0,\qquad
Vf=b_0f,\qquad Hf=0,\\
\mathrm d\alpha_0(U,H)=-2,\qquad
\mathrm d\alpha_0(\bar H,V)=2.
\end{gathered}
\] The other pairings between these two horizontal frames vanish.
Solving the contact-field equations modulo \(\mathcal A_0\) gives
\([X_h]=hR+\tfrac12Hh\,U-\tfrac12Vh\,\bar H\). Taking its Bott
derivative yields

\phantomsection\label{eq:online-contact-gauge-components} \[
\begin{aligned}
\mathcal G_Bh(V)
&=-\frac12(V^2h-b_0Vh)\bar H+\frac12VHh\,U,\\
\mathcal G_Bh(H)
&=-\frac12HVh\,\bar H+\frac12H^2h\,U.
\end{aligned}
\] For \(h=\mathcal D_+u\), \(\alpha=Vu\) and \(\beta=Hu\), the
commutators give \(Vh=-f(\bar H\alpha+2Ru)\) and \(Hh=-f\bar H\beta\).
Substitution, using \(VRu=R\alpha-f\beta\), gives the formula for
\(\Gamma_+\) and the exact gauge identity. If \(d^0u=0\), the displayed
expression is zero.

For a general closed cochain, no actual local primitive is required. At
a point choose \(j\) with \(z_j\ne0\). In the complexified real Taylor
ring the analytic coordinates \[
s=\frac{x^j-x_0^j}{z_j(\zeta)},\qquad
y^k=x^k-x_0^k-\frac{z_k(\zeta)}{z_j(\zeta)}(x^j-x_0^j),
\quad k\ne j,
\] together with \(\zeta,\bar\zeta\), make \(H=\partial_s\) and
\(V=\partial_{\bar\zeta}\). With \(\sigma=\bar\zeta-\bar\zeta_0\), a
formally closed pair has the formal primitive

\phantomsection\label{eq:online-contact-finite-jet-primitive} \[
\widehat u(s,\sigma)
=\int_0^\sigma\widehat\alpha(0,r)\,\mathrm dr
 +\int_0^s\widehat\beta(t,\sigma)\,\mathrm dt.
\] The integrals are coefficientwise. The primitive through order
\(k+1\) depends linearly on the cochain through order \(k\), with smooth
coefficients in the center. A formally CR section is independent of
\(s,\sigma\); truncating in the remaining analytic coordinates realizes
each of its finite jets by a local analytic CR section.

Define the comparison at the point by the finite jet of
\(d_B[X_{\mathcal D_+\widehat u}]\). Two choices differ by a formally CR
section and contribute zero by the preceding differential identity. This
remains true at the second projective pole: the required finite CR jet
is realized analytically, the global lift has zero Bott derivative on
the north-chart overlap, and continuity gives zero at the pole. The
construction is therefore independent of the formal primitive and
angular frame. It agrees with the north formula and patches globally.
Only the finite jet of the formal primitive enters the comparison at a
point. Neither convergence of the series nor a smooth local solution of
the primitive equation is needed.

The closedness equation gives \[
\begin{aligned}
\mathrm d\alpha_0(\Gamma_+e(V),H)
&=f\bigl(\bar H(V\beta)+2R\beta\bigr)\\
&=f\bigl(\bar H(H\alpha)+2R\beta\bigr)
=\mathrm d\alpha_0(\Gamma_+e(H),V).
\end{aligned}
\] Thus the cocycle is contact-compatible. Its horizontality is
explicit. To verify \(d_B\Gamma_+e=0\), choose the finite formal
primitive one order farther and apply \(d_B^2=0\) for the involutive
flat distribution. This proves an identity of actual finite jets at
every point.

The formula is differential with smooth coefficients on a finite angular
cover, so the asserted restriction, parameter and finite-order estimates
follow. Complex conjugation identifies the positive and dual complexes
and interchanges the two distributions. Conjugating the proved
identities gives the dual assertions. \hfill\(\square\)

For symmetric complex tensors \(q,p\), write their linear initial-sky
graph as

\phantomsection\label{eq:online-linear-tensor-bott-graph} \[
\begin{aligned}
\dot\Phi_{q,p}(V)&=\beta_1U,\qquad
\dot\Phi_{q,p}(H)=-\beta_1\bar H+\eta_1U,\\
\beta_1&=\frac f4q(z,z),\\
\eta_1&=\frac12\bigl(p(z,z)-n_a\Gamma(q)^a{}_{ij}z^iz^j\bigr),\\
\Gamma(q)^a{}_{ij}
&=\frac12(\partial_iq_{aj}+\partial_jq_{ai}-\partial_aq_{ij}).
\end{aligned}
\] Traces of \(q,p\) do not change this graph, so we choose their
trace-free representatives. The symbol \(\eta_1\) is a graph
coefficient, separate from the decay exponent \(\gamma\). The graph over
\(\overline{\mathcal A}_0\) obtained by replacing \((V,H,U,z)\) by
\((U,\bar H,V,\bar z)\), while keeping \(q,p\) fixed, is denoted by
\(\dot{\bar\Phi}_{q,p}\).

Define

\phantomsection\label{eq:online-contact-curl-and-momentum-operators} \[
\begin{aligned}
(\mathsf CS)_{ij}
&=\frac12\left(\epsilon_i{}^{k\ell}\partial_kS_{\ell j}
             +\epsilon_j{}^{k\ell}\partial_kS_{\ell i}\right),\\
\mathcal P(S)
&=2\Delta S-4\operatorname{TF}_\delta
               (\operatorname{sym}\partial\operatorname{div}_\delta S).
\end{aligned}
\] Here \(\operatorname{sym}\) means one half the sum with the two free
indices exchanged. On TT tensors, \(\mathsf C^2=-\Delta\) and
\(\mathcal P(S)=2\Delta S\). The first identity follows from the two
helicity eigenvalues \(\pm|\xi|\) of the Fourier curl symbol;
equivalently it is the constant-coefficient curl identity after setting
trace and divergence to zero.

\pandocstatementheading{\leavevmode\phantomsection\label{prop:online-positive-marked-sky-response}Proposition 4.6 (Complete-sky normalization of the positive cochain).}

Let \(e=(\alpha,\beta)\) be an actual closed positive cochain on a
complete sphere bundle. With the fixed fiber Hodge operator, put
\(u=T_2\alpha\) and write \[
e-d^0u=(0,S(z,z)),\qquad S^{\mathsf T}=S,\quad\operatorname{tr}_\delta S=0.
\] The tensor \(S\) is unique and actual smooth. Set \(k=-S(n,n)\). Then

\phantomsection\label{eq:online-positive-normalized-sky-response} \[
\Psi_e=\Gamma_+(e)+\mathcal G_B(k-\mathcal D_+u)
      =\dot\Phi_{4i\mathsf CS,\,\mathcal P(S)}.
\] This is a global contact-compatible Bott cocycle preserving the
vertical sky tangent bundle. All operations have finite estimates on the
actual coefficient stages and commute with smooth parameters.

\emph{Proof.} Since \(H^1(\mathbb{CP}^1,\mathcal O(2))=0\), fiber Hodge
decomposition gives \(Vu=\alpha\). Closedness then makes \(\beta-Hu\) a
holomorphic quartic. The quadratic restriction map from trace-free
symmetric tensors to \(H^0(\mathcal O(4))\) is an isomorphism, giving
the unique \(S\). These are fixed continuous fiber maps, so they
preserve actual smoothness and the indicated families.

The cochain-gauge identity reduces \(\Psi_e\) to
\(\Gamma_+(0,S(z,z))+\mathcal G_Bk\). Since \(Vn=-fz\) and \(Vf=b_0f\),
\[
Vk=2fS(n,z),\qquad
(V^2-b_0V)k=-2f^2S(z,z).
\] The \(\bar H\) component of \(\mathcal G_Bk(V)\) cancels that of
\(\Gamma_+(0,S(z,z))(V)\). The remaining graph coefficients are

\phantomsection\label{eq:online-positive-normalized-graph-coefficients}
\[
\begin{aligned}
\beta_1&=f\bigl[(\partial_zS)(n,z)-(\partial_nS)(z,z)\bigr],\\
\eta_1&=-\frac f2(\partial_{\bar z}\partial_zS)(z,z)
        -\frac12(\partial_z^2S)(n,n).
\end{aligned}
\] In spatial differentiation the angular vectors are held fixed. The
identity \(n\times z=-iz\) gives \(\beta_1=if(\mathsf CS)(z,z)\), the
required conic coefficient.

To verify the momentum, rotate the constant frame at one angular point
to \(n=e_3\), \(z=e_1+ie_2\) and \(f=1\). Put
\(D=\partial_1+i\partial_2\), \(\bar D=\partial_1-i\partial_2\),
\(N=\partial_3\), \(A=S(z,z)\), \(B=S(n,n)\) and \(C=S(n,z)\). For
\(q=4i\mathsf CS\), direct expansion gives \[
\begin{aligned}
q(z,z)&=4(DC-NA),&
q(n,z)&=\bar DA+3DB-2NC,\\
(\operatorname{div}S)(z)&=\tfrac12\bar DA-\tfrac12DB+NC.
\end{aligned}
\] The coefficient formula gives \(2\eta_1=-\bar DDA-D^2B\). Hence \[
\begin{aligned}
2\eta_1+n_a\Gamma(q)^a{}_{ij}z^iz^j
&=2\eta_1+Dq(n,z)-\tfrac12Nq(z,z)\\
&=2N^2A+2D^2B-4DNC\\
&=2\Delta A-4D\bigl((\operatorname{div}S)(z)\bigr)
=\mathcal P(S)(z,z).
\end{aligned}
\] These tensor identities are rotation-covariant and respect the conic
transitions, so the calculation proves the formula on the entire sphere.
Globality and Bott closedness also follow from the preceding proposition
and the global Hamiltonian \(k\). Each operation is a fixed fiber map or
a finite spatial differential operator. This proves the estimates and
parameter assertions. \hfill\(\square\)

\pandocstatementheading{\leavevmode\phantomsection\label{thm:online-actual-linear-contact-graph}Theorem 4.7 (The actual normalized negative graph).}

For an actual closed dual positive cochain
\(e^\vee=(\alpha^\vee,\beta^\vee)\), define \[
v=\overline{T_2\overline{\alpha^\vee}},\qquad
e^\vee-\bar d^0v=(0,T(\bar z,\bar z)),\qquad
T^{\mathsf T}=T,\quad\operatorname{tr}_\delta T=0.
\] Put

\phantomsection\label{eq:online-dual-complete-sky-normalization} \[
\Psi^\vee_{e^\vee}
=\Gamma_\vee(e^\vee)
 +\bar{\mathcal G}_B\bigl(-T(n,n)-\mathcal D_-v\bigr).
\] There is a complex-linear conversion \(\mathfrak C\) between the
global vertical-preserving contact-compatible Bott cocycles over
\(\overline{\mathcal A}_0\) and those over \(\mathcal A_0\),
characterized by \(\mathfrak C\dot{\bar\Phi}_{q,p}=\dot\Phi_{q,p}\). The
calibrated negative graph is

\phantomsection\label{eq:online-negative-normalized-graph} \[
\mathcal B_-(e^\vee)
=-\mathfrak C\Psi^\vee_{e^\vee}
=\dot\Phi_{q_-,p_-},\qquad
q_-=4i\mathsf CT,\qquad p_-=-\mathcal P(T).
\] These are equalities of actual complete-sphere linear graph fields.
If \(T\in\mathscr S^\vartheta\), then \(q_-\in\mathscr S^{\vartheta-1}\)
and \(p_-\in\mathscr S^{\vartheta-2}\), with finite seminorm estimates
and the same statements for all supplied finite parameter derivatives.

\emph{Proof.} A vertical-preserving contact-compatible cocycle over
\(\overline{\mathcal A}_0\) has the form \[
\Xi(U)=b_1V,\qquad \Xi(\bar H)=-b_1H+c_1V.
\] Its Bott equations are \(Ub_1=-a_0b_1\) and \(Uc_1=\bar H b_1\).
Since \(Uf=-a_0f\), the first makes \(4f^{-1}b_1\) an antiholomorphic
quartic on the complete sphere. Thus \(b_1=fq(\bar z,\bar z)/4\) for a
unique trace-free tensor \(q\). Using \(Un=-f\bar z\) and \(U\bar z=0\)
gives \[
U\left[-\tfrac12n_a\Gamma(q)^a{}_{ij}\bar z^i\bar z^j\right]
=\bar H b_1.
\] Consequently \(2c_1+n_a\Gamma(q)^a{}_{ij}\bar z^i\bar z^j\) is
another antiholomorphic quartic, giving the unique \(p\). Reversing
these calculations proves the converse. The same reasoning in the
opposite polarization proves the asserted conversion and its global
finite angular readout.

It is useful to distinguish this conversion from the antipodal map. At
\(a\zeta=-1/\bar\zeta\), \(n(a\zeta)=-n(\zeta)\),
\(\bar z(a\zeta)=-\zeta^{-2}z(\zeta)\) and
\(f(a\zeta)=|\zeta|^4f(\zeta)\). The converted coefficients are

\phantomsection\label{eq:online-opposite-polarization-conversion} \[
\beta_1=\frac{\zeta^2}{\bar\zeta^2}b_1(a\zeta),\qquad
\eta_1=\zeta^4c_1(a\zeta)
       -n_a\Gamma(q)^a{}_{ij}z^iz^j.
\] The connection term retains the same extrinsic tensor \(p\). Literal
antipodal pullback instead yields \(\dot\Phi_{q,-p}\). The tensor
construction proves extension across both chart poles.

Conjugating the positive normalization, with \(T\) held as the resulting
actual coefficient, gives
\(\Psi^\vee_{e^\vee}=\dot{\bar\Phi}_{-4i\mathsf CT,\,\mathcal P(T)}\).
The factor \(i\) changes sign; the real differential operator
\(\mathcal P\) does not. Conversion preserves the two tensors. The final
minus sign is the original coordinate-operator convention
\(E=(\Phi^{-1})^*\): the point-map infinitesimal generator is the
negative of the coordinate-operator generator. This proves the
calibrated graph formula.

The calibration can be checked directly against §3.3. For the derived
dual TT representative write
\(\widehat T=\sum_{\varepsilon}U^-_\varepsilon/(16r^4)\), using the
actual anchored integral when necessary. At the root
\(n_\varepsilon=\varepsilon\xi/r\), the negative line has curl
eigenvalue \(\varepsilon r\). Thus

\phantomsection\label{eq:online-negative-graph-residue-calibration} \[
\widehat q_-=\sum_\varepsilon\frac{i\varepsilon}{4r^3}U^-_\varepsilon,
\qquad
\widehat p_-=\sum_\varepsilon\frac1{8r^2}U^-_\varepsilon.
\] These are the original negative metric and momentum amplitudes,
including their signs. For the selected positive TT quartic
representative, with \(\widehat S=\sum_\varepsilon U^+_\varepsilon\) up
to its constant anchor, the calibrated graph \(-\Psi_e\) likewise has
\(\widehat q_+=\sum_\varepsilon4i\varepsilon r\,U^+_\varepsilon\) and
\(\widehat p_+=2r^2\sum_\varepsilon U^+_\varepsilon\), because the
positive root line has curl eigenvalue \(-\varepsilon r\). Hence the sum
for these selected representatives agrees with the exact two-class seed
transform already proved. A retained longitudinal quartic is instead
governed by the full tensor formula and its contact gauge; its graph is
not determined by root values alone. Fourier reflection enters only when
conjugating spatial fields for Lorentz reality; it is not part of the
complex-linear conversion.

For the weighted dual representative the single anchored integral is
differentiated before the conic tensor is read. Its constant anchor
vanishes under these derivatives.
\hyperref[thm:online-actual-weighted-negative-dual]{The actual weighted
dual theorem} supplies the finite bounds, and the operators above use at
most two additional spatial derivatives. Fixed fiber maps and finite
angular coefficient readouts preserve all stated parameter estimates.
\hfill\(\square\)

For use with retained representatives, the contact gauge has the
elementary tensor readout

\phantomsection\label{eq:online-contact-scalar-vector-gauge} \[
\mathcal G_B(n\cdot W+F)
=\dot\Phi_{-\mathcal L_W\delta,\,\nabla^2F}.
\] Indeed \(V(n\cdot W)=-fz\cdot W\) gives the conic coefficient in the
formula for \(\mathcal G_B\). Its \(H\) coefficient is
\(\tfrac12\partial_z^2(n\cdot W+F)\), and
\(\Gamma(-\mathcal L_W\delta)^a{}_{ij}=-\partial_i\partial_jW^a\)
supplies the same connection term. Traces are invisible to \(\dot\Phi\).
This verifies the full graph identity, not only its conic part.

\subsection{4.3. Actual representative changes in the linear contact
response}\label{actual-representative-changes-in-the-linear-contact-response}\label{sec:online-dual-representative-contact-comparison}

The actual primitives of
\hyperref[thm:online-actual-weighted-negative-dual]{the weighted
dualization theorem} also determine the comparison of retained linear
contact responses. We use the fixed-polarization graph
\(\dot\Phi_{q,p}\), its actual cochain realization, and the contact
gauge \(\mathcal G_B\) defined in
\hyperref[thm:online-actual-linear-contact-graph]{the linear contact
construction}. This paragraph concerns that linearized response.

\pandocstatementheading{\leavevmode\phantomsection\label{prop:online-dual-primitive-linear-contact-gauge}Proposition 4.8 (The contact gauge of an actual dual primitive).}

Let \(e^{\vee\prime}-e^\vee=d_\vee^0w\) be an actual dual cochain change
in the weighted complex above. Write the fiber anti-holomorphic
projection of \(w\) as

\phantomsection\label{eq:online-dual-primitive-holomorphic-projection}
\[
w_0=w-T_2^\vee\partial w=Y_i(x)\bar z^i.
\] Then the calibrated negative graph satisfies

\phantomsection\label{eq:online-dual-primitive-contact-hamiltonian} \[
\mathcal B_-(e^{\vee\prime})-\mathcal B_-(e^\vee)
=\mathcal G_B h_-,
\qquad
h_-=-i\,n\cdot\operatorname{curl}Y+\frac23\operatorname{div}Y.
\] For a positive representative change whose holomorphic projection is
\(Y_+\cdot z\), the calibrated positive graph has

\phantomsection\label{eq:online-positive-primitive-contact-hamiltonian}
\[
h_+=i\,n\cdot\operatorname{curl}Y_+
                   +\frac23\operatorname{div}Y_+ .
\] Thus \(Y=\overline{Y_+}\) gives a real paired Hamiltonian
\(h_++h_-\). Strict additive equations for supplied degree-zero cochain
comparisons give strict additive equations for these Hamiltonians. All
formulas are global on the complete sphere and obey finite symbol and
parameter estimates.

\emph{Proof.} Fiber normalization of the change gives

\phantomsection\label{eq:online-dual-primitive-quartic-change} \[
\begin{aligned}
\delta\beta_0
&=D^\vee w-D^\vee T_2^\vee\partial w
=D^\vee w_0,\\
\delta T&=\mathcal K Y
=\operatorname{TF}_\delta(\operatorname{sym}\partial Y).
\end{aligned}
\] Here \(\mathcal K\) is the tracefree symmetric gradient already
defined in the polynomial cohomology calculation, with
\(\operatorname{sym}\) containing the factor \(1/2\). To compute the
calibrated graph, retain the operators
\(\mathsf C=\operatorname{curl}_{\mathrm{sym}}\) and
\(\mathcal P(T)=2\Delta T-4\operatorname{TF}_\delta
(\operatorname{sym}\partial\operatorname{div}T)\) from the linear
contact theorem. Direct differentiation gives

\phantomsection\label{eq:online-dual-primitive-tensor-identities} \[
\begin{aligned}
\mathsf C(\mathcal K Y)
&=\frac12\operatorname{sym}\partial(\operatorname{curl}Y),\\
\operatorname{div}(\mathcal K Y)
&=\frac12\Delta Y+\frac16\,\mathrm d(\operatorname{div}Y),\\
\mathcal P(\mathcal K Y)
&=-\frac23\operatorname{TF}_\delta
                   \nabla^2(\operatorname{div}Y).
\end{aligned}
\] For the first identity, insert
\((\mathcal K Y)_{ij}=(\partial_iY_j+\partial_jY_i)/2
-\delta_{ij}\partial_aY_a/3\). Terms with two commuting derivatives
contracted against the alternating tensor vanish, and the symmetrized
trace terms cancel. The surviving terms are
\(\frac14(\partial_i(\operatorname{curl}Y)_j+
\partial_j(\operatorname{curl}Y)_i)\). The second identity follows by
taking the divergence of the same formula. Substituting it into
\(\mathcal P\) cancels the two
\(\operatorname{TF}_\delta(\operatorname{sym}\partial\Delta Y)\) terms
and gives the third identity.

The calibrated negative graph has \(q_-=4i\mathsf C T\) and
\(p_-=-\mathcal P(T)\), so its change is

\phantomsection\label{eq:online-dual-primitive-graph-tensor-change} \[
\delta q_-=2i\operatorname{sym}\partial(\operatorname{curl}Y),
\qquad
\delta p_-=\frac23\operatorname{TF}_\delta
                       \nabla^2(\operatorname{div}Y).
\] Use \hyperref[eq:online-contact-scalar-vector-gauge]{the
scalar-vector contact identity} \[
\mathcal G_B(n\cdot W+F)
=\dot\Phi_{-\mathcal L_W\delta,\,\nabla^2F}
\] with \(W=-i\operatorname{curl}Y\) and
\(F=\frac23\operatorname{div}Y\). The metric tensor is exactly
\(\delta q_-\) because
\(\mathcal L_W\delta=2\operatorname{sym}\partial W\). The momentum
differs from \(\delta p_-\) only by a trace, which vanishes upon
contraction with \(z\otimes z\) in the graph. This proves the negative
formula. The calibrated positive graph uses \(q_+=-4i\mathsf C S_+\) and
the same \(p_+=-\mathcal P(S_+)\); the identical calculation uses
\(W=i\operatorname{curl}Y_+\) and gives the positive formula.

The vector \(n\) is globally defined on the marked sphere, and the
remaining tensors are actual fields on the initial space. The
Hamiltonians therefore patch globally. Their expressions are linear in
the primitive through its fixed Hodge projection and tensor readout.
This proves strict additivity. Complex conjugation gives
\(h_-=\overline{h_+}\) for the stated paired change, hence a real sum.

If \(w\in\mathscr S^{\vartheta+1}\), then \(Y\) has that order and
\(h_-\in\mathscr S^\vartheta\). The metric and momentum changes have
orders \(\vartheta-1\) and \(\vartheta-2\), respectively. Their finite
seminorms use at most the displayed finite number of spatial derivatives
and the compact-fiber projection. The same linear formulas commute with
every supplied finite parameter derivative. \hfill\(\square\)

When the constructor first takes its fixed canonical cohomology
representative, an ordinary original coboundary leaves that canonical
input and its response unchanged; the comparison is then the identity.
The proposition instead applies to actual representatives that are
retained, including a supplied dual cochain and the selected
representative with the same class. It proves equality of actual
linearized graph fields and their infinitesimal contact gauges. The
nonlinear original point maps and their finite compositions require the
separate nonlinear comparison laws.

\subsection{4.4. A local encoder for actual negative
cochains}\label{a-local-encoder-for-actual-negative-cochains}\label{sec:online-negative-harmonic-input}

Let \(\mathscr J_{-6}\) be the harmonic Serre-dual identification from
trace-free symmetric tensors to \(H^1(\mathcal O(-6))\), normalized by
\[
\langle\mathscr J_{-6}(B),Q(z,z)\rangle_{\mathrm{Serre}}=B:Q.
\] The Serre integral is
\((2\pi i)^{-1}\int p(\zeta)\,\mathrm d\zeta\wedge\alpha\), with the
contour orientation of §3.3. For an actual weighted TT field
\(B\in\mathscr S^{\gamma-2}_{\mathrm{TT}}\), with the fixed
\(-1<\gamma<-1/2\), define

\phantomsection\label{eq:online-negative-harmonic-input-cochain} \[
e_-(B)=\left(\mathscr J_{-6}(B),
T_{-4}(z\cdot\partial_x)\mathscr J_{-6}(B)\right).
\] \pandocstatementheading{\leavevmode\phantomsection\label{prop:online-negative-harmonic-root-encoder}Proposition 4.9 (The harmonic negative input and its exact roots).}

The cochain \(e_-(B)\) is exactly closed and preserves the spatial
support of \(B\). Its original negative root tensors obey

\phantomsection\label{eq:online-negative-harmonic-exact-root-map} \[
U^-_+(\xi)+U^-_-(\xi)=2i\mathsf C(\xi)\widehat B(\xi).
\] With \(\vartheta=\gamma+1\), its coefficients have orders
\(\vartheta-3\) and \(\vartheta-4\), so it lies in the actual weighted
negative complex used by the dual transformation. The differential
closedness and support assertions alone also hold for arbitrary smooth
TT fields.

\emph{Proof.} Under the normalized Serre identifications, the fiber
harmonic projection \(H_{-4}^1(z\cdot\partial_x)\mathscr J_{-6}(B)\) is
exactly \(\operatorname{div}B\), as in
\hyperref[prop:online-polynomial-cr-cohomology]{the polynomial
cohomology calculation}. This is an identity of constant-coefficient
spatial differential operators and fixed fiber maps, so it applies to
all actual smooth tensors. It vanishes for TT \(B\). The fiber Hodge
identity therefore gives
\(\bar\partial T_{-4}(z\cdot\partial_x)\mathscr J_{-6}(B)
 =(z\cdot\partial_x)\mathscr J_{-6}(B)\), which proves closedness. Both
operations are local in the spatial variable and have the displayed
derivative orders.

On the two-dimensional TT plane at \(\xi\ne0\), the sum of roots defines
a rotation-equivariant smooth homogeneous multiplier of degree one. The
root sum is odd under \(\xi\mapsto-\xi\): its horizontal coefficient
contains \(i\xi\cdot z\), and root reversal just interchanges the two
summands. The stabilizer of a unit covector has distinct characters on
the two helicity lines. Hence the multiplier is
\(a|\xi|I+b\mathsf C(\xi)\) for constants \(a,b\); oddness forces
\(a=0\).

The remaining constant is fixed by the plane used in §3.3. For
\(Ae^{\lambda w^0}/\zeta\), \(\xi=\varepsilon r e_3\) and
\(\lambda=-i\sqrt2\,\varepsilon r\), the corrected first CR coefficient
is \(Ae^{\lambda u}\zeta^{-1}\bar\partial\rho_N\). The normalized
contour pairing against \(Q(z,z)\) gives \(Ae^{\lambda u}Q(z_0,z_0)\),
where \(z_0=(1,i,0)\). Thus its harmonic tensor is
\(Ae^{\lambda u}z_0\otimes z_0\). Its original root tensor is
\(2i\varepsilon r\) times that tensor, while the curl eigenvalue on this
line is \(\varepsilon r\). Therefore \(b=2i\). This proves the root
formula with its normalization. The actual symbol estimates follow
directly from the finite differential and fiber operations.
\hfill\(\square\)

\pandocstatementheading{\leavevmode\phantomsection\label{prop:online-normalized-negative-half-newton}Proposition 4.10 (The actual normalized negative response).}

For \(B\in\mathscr S^{\gamma-2}_{\mathrm{TT}}\), form \(e_-(B)\), apply
\hyperref[thm:online-actual-weighted-negative-dual]{the actual
negative-to-dual operation} with its prescribed anchored TT
representative, apply the normalized contact graph above, and read its
conic tensor. Denote this operation by \(\mathcal N_-(B)\). It is a
continuous complex-linear map

\phantomsection\label{eq:online-normalized-negative-half-newton} \[
\mathcal N_-:\mathscr S^{\gamma-2}_{\mathrm{TT}}
\longrightarrow\mathscr S^\gamma_{\mathrm{TT}},
\qquad
\widehat{\mathcal N_-(B)}(\xi)
=-\frac1{2|\xi|^2}\widehat B(\xi),\quad \xi\ne0.
\] The equality is between actual decaying tensors. The response is
unchanged by ordinary coboundaries before the canonical input
projection, and commutes with translations, spin rotations, conjugation
and every specified finite parameter derivative. Retained noncanonical
dual representatives have the actual comparison primitives of the dual
theorem and
\hyperref[prop:online-dual-primitive-linear-contact-gauge]{their linear
contact gauges}.

\emph{Proof.} The preceding proposition proves that the actual input is
closed and belongs to the stated negative domain. Its root sum and
\hyperref[eq:online-negative-dual-calibrated-multiplier]{the calibrated
multiplier} \(1/(16r^4)\) give the dual coefficient tensor \[
\widehat T(\xi)=\frac{i}{8r^4}\mathsf C(\xi)\widehat B(\xi).
\] The normalized actual graph has conic tensor \(q_-=4i\mathsf CT\).
Since \(\mathsf C(\xi)^2=r^2I\) on TT tensors, its multiplier is exactly
\(-1/(2r^2)\). Thus the multiplier is computed from the specified
cochain and contact operations.

Write \(\vartheta=\gamma+1\in(0,1/2)\). Since
\(B\in\mathscr S^{\vartheta-3}\) with \(\vartheta-3\in(-3,0)\),
\hyperref[lem:online-decay-fourier-isomorphism]{the weighted Fourier
lemma} gives differentiated low-frequency bounds
\(r^{-\vartheta-|\alpha|}\) and rapid decrease at infinity. The conic
amplitude has order \(r^{-2-\vartheta}\), which is locally integrable in
three dimensions. The curl of the anchored dual integral is precisely
this inverse integral, because it annihilates the constant subtraction.
The inverse Fourier estimate gives the actual order
\(\vartheta-1=\gamma\) with finite seminorm bounds. The TT identities
hold for the locally integrable extensions. There is no added polynomial
because the resulting field decays.

Canonical projection makes ordinary input coboundaries identical before
the graph is constructed. For retained representatives, the dual theorem
supplies actual primitives and the stated contact comparison. The final
scalar multiplier is real, even and rotation-invariant, so translations,
conjugation and spin rotations commute with the actual response. In
particular the intermediate anchor has disappeared. Applying the same
continuous linear maps to each finite parameter derivative proves the
family assertion. \hfill\(\square\)

The auxiliary calls are made in the complex source domain. A
negative-only call is not required to be a Lorentz-fixed pair. Its conic
response is nevertheless real on real tensor inputs by the proved
multiplier, and the final contractions preserve the real residual
fields. The actual spin marking is retained separately throughout.

\subsection{4.5. The fixed inverse from finitely many negative
responses}\label{the-fixed-inverse-from-finitely-many-negative-responses}\label{sec:online-negative-response-green-word}

Define the local fourth-order operator on symmetric tensors by

\phantomsection\label{eq:online-negative-response-local-encoder} \[
D_{ij}=\delta_{ij}\Delta-\partial_i\partial_j,\qquad
(\mathcal R_4T)_{ij}
=\left(D_{ik}D_{j\ell}-\frac12D_{ij}D_{k\ell}\right)T_{k\ell}.
\] The identities \(\partial_iD_{ij}=0\), \(\operatorname{tr}D=2\Delta\)
and \(D^2=\Delta D\) show directly that its output is symmetric,
trace-free and divergence-free. At \(\xi\ne0\), put

\phantomsection\label{eq:online-negative-response-tt-projector} \[
P=I-\frac{\xi\otimes\xi}{r^2},\qquad
\Pi_{ij,k\ell}
=\frac12(P_{ik}P_{j\ell}+P_{i\ell}P_{jk})
 -\frac12P_{ij}P_{k\ell}.
\] Then \(\mathcal R_4\) has polynomial symbol \(r^4\Pi\). The displayed
differential formula implements this symbol without first applying a
nonlocal TT projection.

Choose a real orthonormal basis \(E_A\), \(1\le A\le5\), of trace-free
symmetric tensors. For a scalar \(g\in\mathscr S^{\gamma-2}\), form

\phantomsection\label{eq:online-three-negative-response-word} \[
F_A^{(0)}(g)=\mathcal R_4(E_Ag),\qquad
F_A^{(j+1)}(g)=\mathcal N_-\bigl(F_A^{(j)}(g)\bigr),
\quad j=0,1,2.
\] After each response its conic tensor is read before the next negative
cochain is formed. The auxiliary graph also has a momentum tensor, which
is not asserted to vanish.

\pandocstatementheading{\leavevmode\phantomsection\label{lem:online-negative-response-intermediate-orders}Lemma 4.11 (Actual domains of the three successive responses).}

All four tensors are TT and obey

\phantomsection\label{eq:online-negative-response-intermediate-orders}
\[
F_A^{(j)}(g)\in\mathscr S^{\gamma-6+2j},\qquad
\widehat F_A^{(j)}(\xi)
=\left(-\frac12\right)^j r^{4-2j}\Pi(\xi)E_A\widehat g(\xi),
\quad 0\le j\le3.
\] Every prescribed finite output seminorm is bounded by finitely many
input seminorms, uniformly in this finite basis and for specified finite
parameter derivatives. The actual coefficient domains in the three calls
are as follows. The last column gives the orders of both tensors in the
auxiliary contact graph.

{\def\LTcaptype{none} % do not increment counter
\begin{longtable}[]{@{}
  >{\raggedright\arraybackslash}p{(\linewidth - 6\tabcolsep) * \real{0.2500}}
  >{\raggedright\arraybackslash}p{(\linewidth - 6\tabcolsep) * \real{0.2500}}
  >{\raggedright\arraybackslash}p{(\linewidth - 6\tabcolsep) * \real{0.2500}}
  >{\raggedright\arraybackslash}p{(\linewidth - 6\tabcolsep) * \real{0.2500}}@{}}
\toprule\noalign{}
\begin{minipage}[b]{\linewidth}\raggedright
Call
\end{minipage} & \begin{minipage}[b]{\linewidth}\raggedright
TT input
\end{minipage} & \begin{minipage}[b]{\linewidth}\raggedright
Negative first and second coefficients
\end{minipage} & \begin{minipage}[b]{\linewidth}\raggedright
Conic and momentum tensors
\end{minipage} \\
\midrule\noalign{}
\endhead
\bottomrule\noalign{}
\endlastfoot
\(0\) & \(F_A^{(0)}\in\mathscr S^{\gamma-6}\) &
\(\mathscr S^{\gamma-6}\oplus\mathscr S^{\gamma-7}\) &
\(\mathscr S^{\gamma-4}\oplus\mathscr S^{\gamma-5}\) \\
\(1\) & \(F_A^{(1)}\in\mathscr S^{\gamma-4}\) &
\(\mathscr S^{\gamma-4}\oplus\mathscr S^{\gamma-5}\) &
\(\mathscr S^{\gamma-2}\oplus\mathscr S^{\gamma-3}\) \\
\(2\) & \(F_A^{(2)}\in\mathscr S^{\gamma-2}\) &
\(\mathscr S^{\gamma-2}\oplus\mathscr S^{\gamma-3}\) &
\(\mathscr S^{\gamma}\oplus\mathscr S^{\gamma-1}\) \\
\end{longtable}
}

Each negative cochain lies in \(\mathscr C_{-,\vartheta}^1\) for the
same \(\vartheta=\gamma+1\). Each selected dual cochain lies in
\(\mathscr C_{\vee,\vartheta}^1\), with actual tensor coefficient of
order \(\vartheta\), and the dual comparison primitives lie in
\(\mathscr C_{\vee,\vartheta}^0\) of order \(\vartheta+1\). The selected
coefficients, their anchor-comparison primitives, and the displayed
graph tensors have finite seminorms bounded by finitely many seminorms
of \(g\). For a separately supplied representative, the comparison
bounds also use finitely many seminorms of that representative.

\emph{Proof.} The differential encoder gives the first assertion and its
symbol at \(j=0\). Its order is at most \(\gamma-2\), so the normalized
response applies. Inductively the Fourier multiplier is the stated one.
It remains to prove the stronger orders needed to apply the next
response without losing the domain.

Put \(\vartheta=\gamma+1\). The scalar amplitude has low-frequency
estimates \(r^{-\vartheta-|\alpha|}\), since its input order
\(\vartheta-3\) lies in \((-3,0)\). With \(\mu=4-2j\), the explicit
amplitude \(r^\mu\Pi E_A\widehat g\) therefore has differentiated bounds
\(r^{\mu-\vartheta-|\alpha|}\) and rapid decrease at infinity. It is in
the amplitude space of
\hyperref[lem:online-decay-fourier-isomorphism]{the one-way inverse
Fourier estimate} with exponent \(\beta=\vartheta-3-\mu=\gamma-6+2j<0\).
That estimate was proved for every negative \(\beta\), so it applies
also to the first two exponents below \(-3\). It gives the stated
spatial orders and finite seminorm bounds. Equivalently, a low-frequency
annulus \(r\asymp2^{-k}\) contributes at most \[
C_{\alpha,N}\,
2^{-k(3+\mu-\vartheta+|\alpha|)}
(1+2^{-k}|x|)^{-N}
\] to the absolute value of its \(\alpha\)th spatial derivative;
splitting at \(2^k\asymp\langle x\rangle\) gives the same bound. For all
four values of \(j\), \(3+\mu-\vartheta>0\), so the amplitudes and their
spatial Fourier multiples are integrable at zero.

The improved orders follow from the explicit factors in the encoded
amplitude at low frequency. The gain thus applies to these encoded
inputs. Faster decay alone would not justify it for a general Newton
potential. At each stage the source response and the displayed inverse
transform are decaying actual fields with the same punctured Fourier
transform. Their difference is a polynomial and hence zero. This proves
the exact identities and the zero-frequency normalization of every conic
output.

For the other intermediate objects, the local harmonic encoder preserves
the first coefficient's spatial order and lowers the second by one. All
its inputs have order at most \(\gamma-2\), so the fixed weighted dual
theorem applies to every call. It gives the dual representatives and
comparison primitives in the common stated spaces, with finite seminorm
bounds. In calls \(0\) and \(1\), the dual tensor is more precisely
\(\tau-\tau(0)\), where the same inverse-frequency estimate gives
\(\tau\in\mathscr S^{\gamma-3+2j}\). The constant subtraction is
retained; these two anchored tensors are not claimed to have negative
spatial order. Their first and second derivatives do have the improved
orders, and the graph formulas give the table. In call \(2\), the
once-anchored integral has order \(\vartheta\) and its differentiated
integral gives the same conclusion. The degree-zero primitive of any
anchor change is the explicit linear cochain of the dual theorem and has
allowed order \(\vartheta+1\).

All these operations are continuous linear maps on the displayed symbol
stages. For each finite list of output, fiber and parameter derivatives,
one first chooses the finite input seminorms required by the dual
theorem and the finite differential operations, and then those required
by the inverse-frequency estimate. For retained representative
comparisons, that list includes the supplied representative. The
composition proves the finite primitive and coefficient bounds for each
of the three calls. The list is finite and the weight and anchor are
unchanged. This proves the parameter assertion. \hfill\(\square\)

Set \(\mathcal A_A(g)=F_A^{(3)}(g)\). Define scalar and vector
operations by

\phantomsection\label{eq:online-negative-response-invariant-contractions}
\[
\begin{aligned}
\mathcal N^{\mathrm{src}}g
&=4\sum_AE_A:\mathcal A_A(g),\\
(\mathcal MF)_i
&=\sum_{A,j,k}(E_A)_{kj}[\mathcal A_A(F_k)]_{ij},\\
\mathcal V^{\mathrm{src}}F
&=\frac34\mathcal N^{\mathrm{src}}F+2\mathcal MF,
\end{aligned}
\] where the scalar operation in the last line acts componentwise.

\pandocstatementheading{\leavevmode\phantomsection\label{thm:online-negative-response-flat-green}Theorem 4.12 (The complete flat inverse from negative source responses).}

On actual weighted smooth scalar and vector fields of order
\(\gamma-2\), \[
\mathcal N^{\mathrm{src}}=\Delta^{-1},\qquad
\mathcal V^{\mathrm{src}}=\Delta_L^{-1},
\] with the decaying normalization of
\hyperref[lem:online-chart-decaying-flat-inverses]{Lemma A.2}. Thus the
complete inverse of the flat constraint correction operator is

\phantomsection\label{eq:online-negative-response-full-green} \[
\mathsf G_0^{\mathrm{src}}
=\operatorname{diag}\left(-\frac12\mathcal N^{\mathrm{src}},
                          \mathcal V^{\mathrm{src}}\right)
=\mathsf A_0^{-1},\qquad
\mathsf A_0=\operatorname{diag}(-2\Delta,\Delta_L).
\] The operations are independent of the auxiliary orthonormal basis,
preserve real inputs, and commute with Euclidean translations and
rotations together with the supplied spatial spin lifts. Each smooth
operation uses a local differential encoder, finitely many actual
normalized negative responses, and finite coefficient contractions. A
spin lift is retained as part of the marking; it is not recovered from
the tensor contraction.

\emph{Proof.} The preceding lemma gives the multiplier
\(-\Pi E_A/(8r^2)\) for \(\mathcal A_A\). Since \(\Pi\) is the rank-two
orthogonal projector on the TT plane in the five-dimensional trace-free
symmetric space, \[
\sum_AE_A:\Pi E_A=\operatorname{tr}\Pi=2.
\] Therefore \(\mathcal N^{\mathrm{src}}\) has multiplier \(-r^{-2}\).
Its actual output has order \(\gamma\) and decays. Its difference from
the decaying inverse of \(\Delta\) has zero punctured Fourier transform,
hence is a decaying polynomial and is zero.

Completeness of the basis gives
\(\sum_A[\Pi E_A]_{ij}(E_A)_{k\ell}=\Pi_{ij,k\ell}\). The partial trace
is

\phantomsection\label{eq:online-negative-response-partial-trace} \[
\sum_j\Pi_{ij,kj}
=\frac12\bigl(P_{ik}\operatorname{tr}P+(P^2)_{ik}\bigr)
 -\frac12(P^2)_{ik}
=P_{ik}.
\] Hence \(\mathcal M\) has multiplier \(-P/(8r^2)\) and \[
\widehat{\mathcal V^{\mathrm{src}}}(\xi)
=-\frac1{r^2}\left(\frac34I+\frac14P\right)
=-\frac1{r^2}\left(I-\frac14\frac{\xi\otimes\xi}{r^2}\right).
\] This is the inverse of the Fourier symbol
\(-r^2(I+\tfrac13\xi\otimes\xi/r^2)\) of \(\Delta_L\). Decay and the
same polynomial-kernel argument identify the actual output with its
uniquely normalized inverse. The scalar factor \(-1/2\) gives the
inverse of \(-2\Delta\).

The basis sums are invariant traces and partial traces, so the result is
basis-independent. The encoder and final contractions are real and
rotation-equivariant, and the normalized response has the required
conjugation and covariance laws. There are five chains of three
responses for one scalar component, and finitely many such chains for
the vector operation. The preceding lemma has verified the actual
cochain domain and finite estimates of every call. No extra Green or
biharmonic primitive is used by the execution. \hfill\(\square\)

The encoder takes four derivatives, so the individual stages need not
act boundedly on the base residual Hölder space. After contraction the
exact identity identifies the composite with the bounded inverse
\(\mathcal Z_m\to\mathcal Y_m\), with its two-derivative estimate. The
identity specifies the realization at finite base order directly,
without a density argument in the large weighted Hölder space.

At a smooth current state \(Z\), the fixed inverse in the current
Neumann correction can therefore be replaced by this equal source
operation:

\phantomsection\label{eq:online-negative-response-current-inverse-word}
\[
\mathsf E_Z=\mathsf G_0^{\mathrm{src}}(\mathsf A_Z-\mathsf A_0),
\qquad
\mathsf G_Z^{\mathrm{src}}
=\sum_{j\ge0}(-\mathsf E_Z)^j\mathsf G_0^{\mathrm{src}}.
\] The current inverse identity uses the single base margin in
\hyperref[prop:online-current-inverse-horizontal-flatness]{Proposition
A.13}. At smooth coefficients, the preceding formulas provide cochain
representatives for every finite word. The series converges in the base
operator norm. At higher orders one estimates the same inverse solution
through its equation, without taking a limit of the partial sums in
those norms.

On the analytic domain the comparison identifies the calibrated linear negative response. For smooth input the anchored cochains and the graph formulas on the complete sphere define each auxiliary response chain directly, without summing a holomorphic Hamiltonian series. The construction therefore needs no stagewise identification with the nonlinear analytic constructor.

\subsection{4.6. Actual initial source
coefficients}\label{actual-initial-source-coefficients}\label{sec:online-actual-initial-source-carriers}

The ordinary classes will generate two uncontracted chiral gradient
operators and a complete contact graph on the initial sphere bundle. We
first specify these fields and their independent kinematic tests, before
imposing the vacuum constraints. The physical scale is retained by the
gradient symbol and, in the contact presentation, by the volume density.
The initial spin bundle and its actual marking are part of both
presentations.

Use the ambient weighted coordinates \(Z=(s,y)\) of
\hyperref[prop:online-chart-exact-ambient-decoder]{Proposition A.4},
restricted to the fixed product \(D\) of
\hyperref[prop:online-current-inverse-horizontal-flatness]{the current
inverse proposition}. Thus \[
\begin{gathered}
h=\delta+q+u\delta, \\
\kappa=p+L_0W, \\
K=\operatorname{TF}_h\kappa, \\
s=(q,p),\quad y=(u,W).
\end{gathered}
\] These tensors are auxiliary coefficients of a source operation. Its
output will consist of actual operators and a contact graph, with
explicit inverse readouts. This distinction is needed to verify that
tensor correction and source correction have the same meaning.

The smooth carriers below use the intersection of the stated weighted
spaces. At a finite order, a carrier and its tangents mean their
corresponding coefficient completions, with the same exact identities
wherever defined. The base order is high enough for all displayed
curvature tests. Thus a residual direction of finite regularity produces
an actual tangent of that regularity; weighted smooth residuals produce
weighted smooth tangents.

Fix the trivial real Lorentz bundle
\(\mathbb E=\mathbb R^3\times\mathbb R^{1,3}\), with signature
\((+---)\), future unit vector \(N=E_0\) and oriented spatial
orthonormal vectors \(E_a\). Fix its actual spin reduction and complex
Dirac module \(\mathscr S=\mathscr S_+\oplus\mathscr S_-\), with the
invariant nondegenerate complex alternating form \(\mathsf B\). The
Clifford conventions are

\phantomsection\label{eq:online-source-clifford-marking} \[
c(v)c(w)+c(w)c(v)=2\eta(v,w)I,
\qquad
\operatorname{tr}_{\mathscr S}\bigl(c(v)c(w)\bigr)=4\eta(v,w).
\] The Clifford vectors generate the full complex endomorphism algebra.
For \(\Omega\in\mathfrak{so}(1,3)\), let \(\varrho_\pm(\Omega)\) be the
chiral restrictions of \(\frac14\Omega^{\mu\nu}c(E_\mu)c(E_\nu)\), where
\(\Omega^{\mu\nu}=\Omega^\mu{}_{\lambda}\eta^{\lambda\nu}\). Thus
\([\varrho(\Omega),c(v)]=c(\Omega v)\). The choice of actual spin bundle
is retained even when its central sign acts trivially on a vector or
contact coefficient.

\pandocstatementheading{\leavevmode\phantomsection\label{def:online-kinematic-gradient-carrier}Definition 4.13 (Canonical kinematic gradient carrier).}

Supply two actual first-order complex-linear operators
\(\mathbb D_\pm:\Gamma(\mathscr S_\pm)\to\Gamma(\mathbb E_{\mathbb C}\otimes\mathscr S_\pm)\).
Their commutators with real scalar multiplication must have a common
symbol given by a real orientation-preserving bundle isomorphism
\(b:T^*\mathbb R^3\xrightarrow{\sim}N^\perp\), and their complete
outputs must be tangent to \(N^\perp\):

\phantomsection\label{eq:online-source-symbol-and-tangential-tests} \[
[\mathbb D_\pm,m_f]\psi=b(\mathrm df)\otimes\psi,
\qquad \eta(N,\mathbb D_\pm\psi)=0.
\] Define

\phantomsection\label{eq:online-source-directional-decoder} \[
\begin{aligned}
h^{-1}(\xi,\zeta)&=-\eta(b\xi,b\zeta),
&\Theta X&=b(h^\flat X),\\
P_X^\pm\psi&=-\eta(\Theta X,\mathbb D_\pm\psi),
&P_X&=P_X^+\oplus P_X^-.
\end{aligned}
\] Require, for real \(X\) and real Clifford sections \(v\),

\phantomsection\label{eq:online-source-bilinear-reality-tests} \[
\begin{aligned}
X\mathsf B(\psi,\chi)
&=\mathsf B(P_X\psi,\chi)+\mathsf B(\psi,P_X\chi),\\
[P_X,c(v)]&\in\Gamma(c(\mathbb E)).
\end{aligned}
\] As proved below, \(c(D_Xv)=[P_X,c(v)]\) then defines a real Lorentz
connection. Require the torsion test \(D\Theta=0\), and set

\phantomsection\label{eq:online-source-second-form-readout} \[
K(X,Y)=-\eta(D_XN,\Theta Y).
\] The canonical frame condition is that the matrix of \(b\) in the
fixed coordinate covectors and the fixed spatial vectors is the positive
matrix \(h^{-1/2}\). Require the decoded \((h,K)\) to satisfy the
weighted ambient gauge and maximality tests of Definition A.3, and its
exact ambient decoder to belong to \(D\). No curvature-image test is
imposed in this kinematic carrier. Its seed readout \(\pi\) is the \(s\)
component of that ambient decoder.

These tests use only the supplied fields and therefore apply before any
choice of a generating ordinary pair.

The ambient seed readout is evaluated on the whole weighted field. A
later local restriction restricts the constructed operators and graph;
it does not apply this whole inverse to a spatially truncated field.

\pandocstatementheading{\leavevmode\phantomsection\label{lem:online-gradient-exact-kinematic-codec}Lemma 4.14 (Exact gradient construction and readout).}

Write \(F=h^{1/2}\), \(e=F^{-1}\), \(\Gamma_i=(\Gamma^a{}_{ib})\), and
\(K_i=(K_{ij})_j\) as a row. Define

\phantomsection\label{eq:online-source-positive-frame-encoding} \[
\begin{aligned}
\Theta X&=(0,FX),\\
\Omega_i&=
\begin{pmatrix}
0&K_i e\\
eK_i^{\mathsf T}&F\Gamma_i e-(\partial_iF)e
\end{pmatrix},\\
\mathbb D_\pm(Z)
&=\sum_{a,i}E_a\otimes e_a^{\,i}
\bigl(\partial_i+\varrho_\pm(\Omega_i)\bigr).
\end{aligned}
\] Here \(e_a^{\,i}\) are the components of the oriented
\(h\)-orthonormal frame \(e_a\). These formulas define a map
\(\operatorname{Enc}_{\mathrm{gr}}\) from \(D\) to the independently
tested kinematic carrier. Its exact inverse is the source readout
followed by Proposition A.4. At every finite weighted order the forward
and inverse maps are smooth, with finite estimates on retained bounded
sets.

More precisely, measure a coefficient difference by the
\(C^{m,\alpha}_\gamma\) norm of its symbol and the
\(C^{m-1,\alpha}_{\gamma-1}\) norms of both zero-order coefficients.
Denote the resulting norm by \(\|\cdot\|_{\mathrm{op},m}\). On such
sets,

\phantomsection\label{eq:online-source-balanced-operator-norm} \[
\|\delta\operatorname{Enc}_{\mathrm{gr}}\|_{\mathrm{op},m}
\asymp
\|\delta h\|_{C^{m,\alpha}_\gamma}
+\|\delta K\|_{C^{m-1,\alpha}_{\gamma-1}}
\asymp \|\delta s\|_{\mathcal X_m}+\|\delta y\|_{\mathcal Y_m}.
\] The assertion holds for tangents and finite differences, with
constants depending on the fixed positive margin and finitely many
coefficient bounds.

\emph{Proof.} The scalar commutator in the definition gives the Leibniz
rule for \(P_X\), since \(-\eta(\Theta X,b\,\mathrm df)=Xf\). It also
gives function-linearity in \(X\), so \(P\) is a connection. Its
commutators with real Clifford vectors define a real connection \(D\).
Differentiate the Clifford relation by \(P\) to see that \(D\eta=0\).
The difference between \(P\) and the spin lift of \(D\) commutes with
every Clifford vector, hence is a scalar one-form because those vectors
generate the full matrix algebra. Preservation of \(\mathsf B\) makes
this scalar zero. Thus \(P\) is the exact spin connection, with no
undetected scalar connection term.

For any oriented \(h\)-orthonormal frame, tangential output gives the
full recovery identity

\phantomsection\label{eq:online-source-uncontracted-exact-recovery} \[
\mathbb D_\pm\psi
=\sum_a\Theta(e_a)\otimes P_{e_a}^\pm\psi.
\] Indeed, contraction with each \(\Theta(e_a)\) recovers its
corresponding component, and the output has no component along \(N\).

The connection is Lorentz, so \(D_XN\) is spatial and defines \(K\) by
the displayed readout. The normal component of \(D\Theta=0\) says that
\(K\) is symmetric. Its tangential component says that the metric
connection induced on \(T\mathbb R^3\) is torsion-free, hence is
\(\nabla^h\). Consequently

\phantomsection\label{eq:online-source-gauss-weingarten} \[
D_XN=\Theta(K^\sharp X),\qquad
D_X\Theta Y=\Theta(\nabla_X^hY)+K(X,Y)N.
\] In the positive frame these identities are exactly the displayed
matrices \(\Omega_i\). Conversely those matrices obey the same
identities; their spatial block is skew because \(\nabla^h h=0\). They
therefore define a Lorentz connection whose spin lift preserves
\(\mathsf B\), the real Clifford action and the torsion test. The
constructed symbol is \(h^{-1/2}\), so every canonical-frame condition
holds. The two constructions are inverse by the full recovery identity
and the unique Gauss--Weingarten formula.

The forward formulas use finite matrix operations, one derivative of
\(h\), and the value of \(K\). The inverse reads \(h\) algebraically
from the symbol, and \(K\) from the directional connection and the fixed
\(N\). Their weighted product estimates give the first equivalence.
Proposition A.4, including its bounded decaying York inverse, gives the
second. The same argument applied to differences or derivatives gives
the stated smoothness and estimates. \hfill\(\square\)

\pandocstatementheading{\leavevmode\phantomsection\label{prop:online-contact-complete-kinematic-codec}Proposition 4.15 (Complete contact coefficients and their exact
inverse).}

On \(N_0=P_+T^*\mathbb R^3=\mathbb R^3\times S^2\), retain
\(\alpha_0=n\cdot\mathrm dx\), its contact distribution and its
projection to \(\mathbb R^3\). In an affine angular chart put \[
\begin{gathered}
z=(1-\zeta^2,i(1+\zeta^2),2\zeta),\qquad
d=1+|\zeta|^2,\\
n=d^{-1}(-2\operatorname{Re}\zeta,-2\operatorname{Im}\zeta,1-|\zeta|^2),
\qquad Z_0=\partial_{\bar\zeta},\quad W_0=z\cdot\partial_x.
\end{gathered}
\] Thus \(\partial_{\bar\zeta}n=-z/d^2\), \(z\cdot\bar z=2d^2\) and
\(n\times z=-iz\). For an ambient pair define

\phantomsection\label{eq:online-source-complete-contact-coefficients} \[
\begin{aligned}
A&=h(z,z),& B&=h(z,\bar z),&
\Delta_c&=(B^2-|A|^2)^{1/2},\\
\beta&=\frac{A}{B+\Delta_c},&
\ell&=z-\beta\bar z,&
\rho_n&=(n^{\mathsf T}h^{-1}n)^{1/2},\\
\eta_c&=\frac12\left(\rho_nK(\ell,\ell)
-n_a\Gamma^a{}_{ij}\ell^i\ell^j\right),
&\omega&=\frac12\log\det h.
\end{aligned}
\] The graph spanned by

\phantomsection\label{eq:online-source-contact-graph} \[
Z_0+\beta\bar Z_0,\qquad
W_0-\beta\bar W_0+\eta_c\bar Z_0
\] is a real transverse, contact-compatible, integrable CR distribution.
The formulas agree on a finite rotated angular cover of the whole
sphere. Together with the density \(e^\omega\mathrm d^3x\), the fixed
value \(\tau=0\) of the mean curvature, and the actual spin marking,
they define \(\operatorname{Enc}_{\mathrm{CR}}(Z)\).

There is an independent test and exact decoder for these complete
coefficients. Fix a sufficiently small uniform graph neighborhood of the
flat graph, reducing \(D\) once at the base order if needed so that its
encodings lie there. Require the supplied graph to lie in this
neighborhood, to have \(|\beta|<1\) and the displayed triangular form on
every chart, and to obey angular compatibility and real transversality.
At three fixed sphere directions with flat null vectors proportional to
\((1,i,0)\), \((1,0,i)\) and \((0,1,i)\), use the supplied \(\beta\) to
form \(\ell\). Solve five independent real equations among
\(h^\sharp(\ell,\ell)=0\), together with
\(\operatorname{tr}_\delta h^\sharp=3\). Require the candidate
\(h^\sharp\) to be positive and to reproduce the supplied conic at every
angular direction. Set

\phantomsection\label{eq:online-source-contact-density-scale} \[
h=\left(\frac{e^{2\omega}}{\det h^\sharp}\right)^{1/3}h^\sharp.
\] Using this metric, solve five independent real null-square equations

\phantomsection\label{eq:online-source-contact-shear-decoder} \[
K(\ell,\ell)
=\rho_n^{-1}\bigl(2\eta_c+n_a\Gamma^a{}_{ij}\ell^i\ell^j\bigr)
\] together with \(\operatorname{tr}_hK=0\). Require this real symmetric
candidate to reproduce the supplied shear at every angular direction.
Finally apply the ambient gauge, decay and \(D\) tests to the decoded
tensors. These are the complete kinematic contact tests. They require no
vacuum equation and no prior membership in an ordinary-pair image. They
are equivalent to exact reconstruction by
\(\operatorname{Enc}_{\mathrm{CR}}\).

For each fixed finite angular order \(l\), give coefficient differences
the balanced norm

\phantomsection\label{eq:online-source-balanced-contact-norm} \[
\|\delta Q\|_{\mathcal B_{m,l}}
=\|\delta\beta\|_{C^{m,\alpha}_\gamma(C^l_{\mathrm{ang}})}
+\|\delta\eta_c\|_{C^{m-1,\alpha}_{\gamma-1}(C^l_{\mathrm{ang}})}
+\|\delta\omega\|_{C^{m,\alpha}_\gamma}.
\] On retained bounded sets this norm is equivalent, in both
finite-difference and tangent estimates, to
\(\|\delta h\|_{C^{m,\alpha}_\gamma}+\|\delta K\|_{C^{m-1,\alpha}_{\gamma-1}}\).
The forward and inverse maps preserve smooth finite-dimensional families
with finite seminorm estimates.

\emph{Proof.} Positivity of \(h\) on the real plane spanned by
\(\operatorname{Re}z,\operatorname{Im}z\) gives \(B>0\) and
\(B^2-|A|^2>0\). The chosen root has \(|\beta|<1\), and
\(h(\ell,\ell)=A-2B\beta+\bar A\beta^2=0\). Its line is the oriented
\(-i\) screen eigenspace, since it equals the flat line at \(h=\delta\)
and positive metrics form a connected set.

For an intrinsic verification of the entire graph, let
\(u=h^{-1}n/\rho_n\) and \(j_hX=u\times_hX\) on \(E_n=\ker n\). Define
\[
\begin{aligned}
M_hw&=\rho_n^{-1}h^{-1}w
-\rho_n^{-3}h^{-1}n\,(n^{\mathsf T}h^{-1}w),\\
C_hX&=P_n\bigl(\Gamma^a{}_{ji}X^jn_a\bigr)_i,
\qquad P_n=I-n\otimes n.
\end{aligned}
\] Differentiation of the ray map \(n\mapsto h^{-1}n/\rho_n\) gives
\(M_h\), and substitution gives \(M_h^{-1}v=\rho_nP_nhv\). A parallel
covector has coordinate derivative \(\Gamma^a{}_{ji}X^jn_a\), so the
fixed contact splitting converts to the metric splitting by
\((X,w)\mapsto(X,M_h(w-C_hX))\).

Let \(B_K\) be the tracefree part of \(h^{-1}K\) on the \(h\)-screen
\(u^\perp=E_n\). It anticommutes with \(j_h\), and the intrinsic complex
structure has matrix
\(\left(\begin{smallmatrix}j_h&0\\2j_hB_K&j_h\end{smallmatrix}\right)\).
For screen vectors, \(h(X,M_hw)=\rho_n^{-1}X\cdot w\). This shows that
\(Z_0+\beta\bar Z_0\) has eigenvalue \(-i\). For the second graph vector
its base component is \(\ell\) and its fixed vertical component is
\(-\eta_c\bar z/d^2\). Its required \(+i\) vertical projection is
\(-B_K\ell\). Pairing this condition with \(\ell\) gives \[
-\frac{2\eta_c}{\rho_n}
-\frac{n_a\Gamma^a{}_{ij}\ell^i\ell^j}{\rho_n}
+K(\ell,\ell)=0,
\] which is the defining formula. Thus the graph is the entire intrinsic
distribution; this proves its angular compatibility and real
transversality. Its contact compatibility also follows directly from
\((z+\beta\bar z)\cdot(z-\beta\bar z)=0\).

To prove integrability, fix a base point and choose an oriented
\(h\)-orthonormal frame normal there. In its angular coordinate the
antiholomorphic vectors are \(\partial_{\bar\zeta_h}\) and
\(z_h^iH_i^h+\tfrac12K(z_h,z_h)\partial_{\zeta_h}\), where \(H_i^h\) are
the Levi-Civita horizontal lifts. At the chosen point the vertical
coefficients of the lifts vanish along the entire fiber; their vertical
derivatives vanish there as well. The remaining coefficients are
holomorphic in \(\zeta_h\). The bracket of these two vectors therefore
vanishes on that fiber, in each angular chart. The obstruction to
involutivity is tensorial and the point was arbitrary, so the
distribution is integrable. No constraint or spacetime evolution has
been used.

At the flat metric the three imaginary null-square equations read off
the three off-diagonal entries of a real symmetric form, while two real
equations read off its diagonal differences. Adding its trace gives an
invertible real six-by-six system. The same selected equations remain
invertible in a fixed small graph neighborhood. They therefore give the
claimed candidate \(h^\sharp\). The complete conic test makes that
candidate conformally equal to the supplied conic, and the density gives
its unique physical scale by the determinant formula. The corresponding
five null-square equations determine a symmetric \(K\) modulo the
metric; its prescribed zero trace determines the missing component.
Their coefficient matrix is again invertible near flat data. The
complete shear test is precisely the remaining equality with the
supplied graph. This proves both inverse identities and the independent
nature of the tests.

All these formulas are smooth finite-dimensional operations with
uniformly nonzero denominators in the retained angular charts. The
forward map uses one spatial derivative of \(h\) and no derivative of
\(K\). The inverse obtains \(h\) by fixed angular evaluations and uses
one derivative of this recovered metric when obtaining \(K\). Weighted
multiplication gives the displayed balanced estimates, in both
directions. Taking a fixed number of angular derivatives costs only
further finite derivatives of these angular formulas, with no additional
spatial derivative. Parameter differentiation obeys the same finite
rules. \hfill\(\square\)

\subsection{4.7. The commutator residual and its normalized
response}\label{the-commutator-residual-and-its-normalized-response}\label{sec:online-source-curvature-residual}

The kinematic source already determines its connection. Its curvature
therefore gives an independent vacuum test before a corrected source is
constructed. The test has four components; no inversion of a curvature
record is involved.

\pandocstatementheading{\leavevmode\phantomsection\label{prop:online-four-source-curvature-defects}Proposition 4.16 (The four source curvature defects).}

For a kinematic gradient source form the actual multiplication operator
\[
\mathcal R^{\mathscr S}(X,Y)=[P_X,P_Y]-P_{[X,Y]}.
\] Let \(e_a\) be an oriented \(h\)-orthonormal frame,
\(v_a=\Theta(e_a)\) and \(\theta^a\) its dual coframe. The following are
finite differential readouts of the supplied operators:

\phantomsection\label{eq:online-source-four-commutator-readouts} \[
\begin{aligned}
E_{ab}
&=-\frac14\sum_c\operatorname{tr}_{\mathscr S}
\left(c(v_c)[\mathcal R^{\mathscr S}(e_c,e_a),c(v_b)]\right),\\
B_{ab}
&=\frac18\epsilon_a{}^{cd}\operatorname{tr}_{\mathscr S}
\left(c(N)[\mathcal R^{\mathscr S}(e_c,e_d),c(v_b)]\right),\\
H&=\operatorname{tr}E,
\qquad M_a=\epsilon_a{}^{bc}B_{bc},
\qquad M=M_a\theta^a,
\qquad\mathfrak d=(H,M).
\end{aligned}
\] The matrix \(E\) is symmetric, \(B\) is tracefree, and on the maximal
weighted carrier

\phantomsection\label{eq:online-source-defect-equals-constraint} \[
H=R(h)-|K|_h^2,\qquad
M=\operatorname{div}_hK,\qquad
\mathfrak d(\operatorname{Enc}_{\mathrm{gr}}(Z))=\mathcal F(Z).
\] For any real symmetric tracefree matrices \(E_0,B_0\), define
\(W(E_0,B_0)\) by the algebraic curvature symmetries and

\phantomsection\label{eq:online-source-weyl-component-convention} \[
C_{0a0b}=(E_0)_{ab},\qquad
C_{cd0b}=-\epsilon_{cd}{}^a(B_0)_{ab},\qquad
C_{abcd}=-\epsilon_{ab}{}^m\epsilon_{cd}{}^n(E_0)_{mn}.
\] These are all real algebraic Weyl tensors. Put
\(E^\circ=E-H\delta/3\) and \(B^{\mathrm s}=(B+B^{\mathsf T})/2\). With
\(\eta(\mathcal C(u,v)z,w)=C(u,v,z,w)\),

\phantomsection\label{eq:online-source-four-defect-complement} \[
\begin{aligned}
\Omega_D-W(E^\circ,B^{\mathrm s})(\Theta,\Theta)
&=\mathcal C_\Theta(H,M),\\
\mathcal C_\Theta(H,M)(e_a,e_b)N
&=\frac12(M_bv_a-M_av_b),\\
\mathcal C_\Theta(H,M)(e_a,e_b)v_c
&=\frac12(\delta_{ac}M_b-\delta_{bc}M_a)N
+\frac H6(\delta_{bc}v_a-\delta_{ac}v_b).
\end{aligned}
\] Consequently \(\mathfrak d=0\) is equivalent to the finite Weyl-image
test
\(\Omega_D\in\operatorname{im}(C\mapsto\mathcal C(\Theta,\Theta))\). The
Weyl preimage is unique. In the adapted orthonormal frame, using the
positive Frobenius norm of the endomorphism matrices,

\phantomsection\label{eq:online-source-four-defect-norm} \[
\sum_{a<b}\|\mathcal C_\Theta(H,M)(e_a,e_b)\|_{\mathrm{Frob}}^2
=\frac{H^2}{6}+|M|^2.
\] The contact presentation has the same independent test by its exact
decoder and the gradient construction. Its complete kinematic tests and
\(\mathfrak d=0\) are therefore equivalent to the maximal vacuum
constraints, before referring to any generating ordinary pair.

\emph{Proof.} The connection Leibniz rule cancels the first-order part
of the commutator, and
\([\mathcal R^{\mathscr S}(X,Y),c(v)]=c(\Omega_D(X,Y)v)\). The Clifford
trace normalization therefore makes \(E\) the spatial Ricci contraction
of \(\Omega_D\) and \(B\) its magnetic contraction.

The Gauss--Weingarten identities give, without a constraint hypothesis,
\[
E=\operatorname{Ric}(h)+\tau K-Kh^{-1}K,
\qquad B_{ab}=\epsilon_a{}^{cd}\nabla_cK_{db},
\qquad \tau=\operatorname{tr}_hK.
\] The first matrix is symmetric. Symmetry of \(K\) gives
\(\operatorname{tr}B=0\). Contraction yields \(H=R(h)+\tau^2-|K|_h^2\)
and \(\epsilon_a{}^{bc}B_{bc}=\nabla^jK_{ja}-\nabla_a\tau\). Since the
carrier imposes \(\tau=0\), these are the two rows of \(\mathcal F\). In
particular the second row is the actual covariant one-form \(M\), not
just its moving frame components.

The component formula for \(W\) verifies the algebraic curvature
symmetries, the cyclic identity and all vanishing traces. Conversely the
electric and magnetic contractions of an algebraic Weyl tensor are
symmetric tracefree matrices and recover all its components by that
formula. Restriction to spatial first arguments retains these
contractions and is injective.

The spatial curvature block of a torsion-free Gauss--Weingarten
connection is a three-dimensional algebraic curvature tensor. It is
determined by its symmetric Ricci contraction: in dimension three its
components are the metric wedge with the Ricci tensor, minus half its
scalar trace times the metric wedge with itself. Subtracting the
displayed Weyl restriction leaves Ricci contraction \(H\delta/3\), hence
the spatial term in the complement formula. The remaining magnetic
matrix is \((B-B^{\mathsf T})/2=\epsilon_{ab}{}^cM_c/2\). Epsilon
inversion gives the normal component in that formula, and Lorentz
skewness gives the value on \(N\). This proves the entire complement and
its equivalence with \(\mathfrak d=0\). Summing the squares of the
displayed matrices gives the norm identity: the spatial terms contribute
\(H^2/6\) and the normal-spatial terms contribute \(|M|^2\). All
operations used are contractions, matrix operations and finitely many
derivatives of the supplied coefficients. \hfill\(\square\)

Write \(\mathsf A_Z=D_y\mathcal F_Z\) and
\(\mathsf B_Z=D_s\mathcal F_Z\). Their explicit current formulas are
\hyperref[eq:online-current-explicit-linearized-residual]{the residual
differential}, including the
\hyperref[eq:online-current-tracefree-variation]{variation of tracefree
projection}. The actual flat inverse is \[
\mathsf G_0^{\mathrm{src}}
=\operatorname{diag}\left(-\frac12\mathcal N^{\mathrm{src}},\mathcal V^{\mathrm{src}}\right)
=A_0^{-1},\qquad A_0=(-2\Delta,\Delta_L),
\] with its finite negative-response realization in
\hyperref[thm:online-negative-response-flat-green]{the source Green
theorem}. The current inverse is the bounded Neumann limit of
\hyperref[prop:online-current-inverse-horizontal-flatness]{Proposition
A.13}. The individual cochain words are evaluated on weighted smooth
fields; the bounded composite gives their base-space extension.

For later use, the differential of the coefficient encodings is
explicit. For \(H_1=\dot h\), \(L_1=\dot K\), put

\phantomsection\label{eq:online-source-positive-frame-variation} \[
\begin{gathered}
F\dot F+\dot F F=H_1, \\
\dot e=-e\dot F e, \\
\dot\Gamma^a{}_{ij}
=\frac12h^{ab}(\nabla_i(H_1)_{jb}+\nabla_j(H_1)_{ib}-\nabla_b(H_1)_{ij}).
\end{gathered}
\] The Sylvester equation has a unique solution because every eigenvalue
of \(F\) is positive. Direct differentiation gives

\phantomsection\label{eq:online-source-full-connection-variation} \[
\dot\Omega_i=
\begin{pmatrix}
0&(L_1)_ie+K_i\dot e\\
\dot eK_i^{\mathsf T}+e(L_1)_i^{\mathsf T}
&\dot F\Gamma_i e+F\dot\Gamma_i e+F\Gamma_i\dot e
 -(\partial_i\dot F)e-(\partial_iF)\dot e
\end{pmatrix},
\] \phantomsection\label{eq:online-source-full-gradient-variation} \[
\dot{\mathbb D}_\pm
=\sum_{a,i}E_a\otimes
\left[\dot e_a^{\,i}(\partial_i+\varrho_\pm(\Omega_i))
+e_a^{\,i}\varrho_\pm(\dot\Omega_i)\right].
\] The complete contact variation is

\phantomsection\label{eq:online-source-full-contact-variation} \[
\begin{aligned}
\dot\beta&=\frac{H_1(\ell,\ell)}{2\Delta_c}, \\
\dot\rho_n&=-\frac{\rho_n}{2}H_1(u,u), \\
u&=h^{-1}n/\rho_n, \\
\dot\eta_c
&=\frac12\left(\dot\rho_nK(\ell,\ell)
+\rho_nL_1(\ell,\ell)-n_a\dot\Gamma^a{}_{ij}\ell^i\ell^j\right) \\
&\quad-\dot\beta\left(\rho_nK(\ell,\bar z)
-n_a\Gamma^a{}_{ij}\ell^i\bar z^j\right), \\
\dot\omega&=\frac12\operatorname{tr}_hH_1.
\end{aligned}
\] For the first identity differentiate \(h(\ell,\ell)=0\) and use
\(h(\ell,\bar z)=B-\beta\bar A=\Delta_c\). The remaining formulas
differentiate \(\rho_n\), \(\ell=z-\beta\bar z\) and the defining shear.
Thus the variation of the null direction itself is retained.

\pandocstatementheading{\leavevmode\phantomsection\label{prop:online-source-normalized-vertical-response}Proposition 4.17 (Normalized source response to the curvature defect).}

At any weighted smooth state \(Z\in D\), and for \(r\in\mathcal Z_m\),
set

\phantomsection\label{eq:online-source-vertical-response} \[
\begin{gathered}
v=\mathsf A_Z^{-1}r=(v_u,v_W), \\
\mathscr R_r^{\mathrm{gr}}
=D\operatorname{Enc}_{\mathrm{gr},Z}[0,v], \\
\mathscr R_r^{\mathrm{CR}}
=D\operatorname{Enc}_{\mathrm{CR},Z}[0,v].
\end{gathered}
\] They are the unique canonical kinematic tangents with

\phantomsection\label{eq:online-source-vertical-two-readouts} \[
D\pi\,\mathscr R_r=0,\qquad
D\mathfrak d\,\mathscr R_r=r.
\] Their actual coefficients are the preceding variation formulas with

\phantomsection\label{eq:online-source-vertical-tensor-response} \[
\begin{aligned}
H_1&=v_u\delta,\qquad \dot\kappa=L_0v_W,
\qquad \vartheta_\kappa=\operatorname{tr}_h\kappa,\\
L_1&=\dot\kappa
-\frac13\left(\operatorname{tr}_h\dot\kappa
-\operatorname{tr}(h^{-1}H_1h^{-1}\kappa)\right)h
-\frac{\vartheta_\kappa}{3}H_1.
\end{aligned}
\] On retained bounded sets,
\(\|\mathscr R_r^{\mathrm{gr}}\|_{\mathrm{op},m}\asymp\|r\|_{\mathcal Z_m}\),
and the corresponding estimate holds in \(\mathcal B_{m,l}\) for every
fixed finite \(l\). Constants include the higher bounds required by the
current inverse. The operations preserve smooth finite-dimensional
families, with the same finite-order interpretation.

For a seed direction \(a\), let
\(\mathscr S_a=D\operatorname{Enc}_Z[a,0]\) be the raw seed tangent. The
horizontal source operation is

\phantomsection\label{eq:online-source-current-horizontal-operation} \[
\mathscr V_a
=\mathscr S_a-\mathscr R_{\mathsf B_Za}
=D\operatorname{Enc}_Z[a,-\mathsf A_Z^{-1}\mathsf B_Za].
\] It satisfies \(D\pi\,\mathscr V_a=a\) and
\(D\mathfrak d\,\mathscr V_a=0\). Every tangent \(T\) to either
canonical carrier has the unique decomposition

\phantomsection\label{eq:online-source-two-readout-tangent-splitting} \[
T=\mathscr V_{D\pi(T)}+\mathscr R_{D\mathfrak d(T)}.
\] In particular the vertical response, the raw seed tangent and the
corrected horizontal seed tangent have different readouts; they are not
interchangeable operations.

\emph{Proof.} The exact kinematic encodings and their inverses identify
carrier tangents with \((a,v)\). The formulas for \(H_1,L_1\) are the
variation of the actual ambient coefficients at fixed seed. The two
readouts of \(\mathscr R_r\) are consequently \(0\) and
\(\mathsf A_Zv=r\). Invertibility of \(\mathsf A_Z\) proves uniqueness.

The derivative of the defect can also be read directly from the source
commutators. In coordinates write \(P_i=\partial_i+A_i\),
\(\Pi_i=\dot P_i\) and \(\mathcal R_{ij}=[P_i,P_j]\). Then \[
\begin{aligned}
\dot{\mathcal R}_{ij}&=[P_i,\Pi_j]-[P_j,\Pi_i],\\
\dot{\mathcal R}_{ab}
&=(\dot e_a^{\,i}e_b^{\,j}+e_a^{\,i}\dot e_b^{\,j})\mathcal R_{ij}
+e_a^{\,i}e_b^{\,j}\dot{\mathcal R}_{ij},\\
\dot M_i&=\dot F^a{}_iM_a+F^a{}_i\epsilon_a{}^{bc}\dot B_{bc}.
\end{aligned}
\] The last equation uses \(M_i=F^a{}_iM_a\) in the positive frame. Its
first term is needed off the admitted set. Differentiating the Clifford
contractions with these moving-frame terms gives precisely the covariant
residual derivative displayed in the current differential lemma. Thus
the operator computation agrees with the tensor differential throughout
the kinematic carrier, not only after the defect vanishes.

The upper estimate follows from the higher current inverse and the
balanced coefficient estimates. Conversely \(D\mathfrak d\) is a bounded
differential operation of order two on \(H_1\) and order one on \(L_1\),
so the identity \(D\mathfrak d\,\mathscr R_r=r\) gives the lower
estimate. Differentiating the inverse identity and the finite
coefficient formulas proves the family statements.

Finally \(D\mathfrak d\,\mathscr S_a=\mathsf B_Za\). Subtracting its
normalized vertical response gives the displayed horizontal operation
and both its readouts. For any tangent \(T\), subtract the two terms in
the asserted decomposition. The remainder has both readouts zero; the
exact decoder gives zero seed component and a correction component in
the kernel of \(\mathsf A_Z\), hence zero. This proves the decomposition
and its uniqueness. \hfill\(\square\)

\subsection{4.8. Initial execution from the ordinary
pair}\label{initial-execution-from-the-ordinary-pair}\label{sec:online-ordinary-source-execution}

The variable input is the ordinary pair
\(c\in\mathscr C^{\mathrm{ord}}_{\gamma,\mathcal B}\). Its actual whole
seed is \(s=\mathscr P_\gamma\rho(c)\). Fix the reference state
\((0,0)\) and its flat physical and spin markings. Choose the transport
tube with the preceding graph-neighborhood reduction included, and
choose \(\mathcal B_\gamma\) inside its input ball \(S\) in
\hyperref[thm:online-current-input-transport]{the current transport
theorem}, with its one fixed base smallness margin. The initial source
operations are the ordinary normalization and root transform, the finite
coefficient formulas above, the actual negative-response word for
\(\mathsf G_0^{\mathrm{src}}\), its bounded current inverse limit, and
initialized parameter integration. The negative-response word may use
complex auxiliary calls; the coupled input condition and the actual
resulting coefficients retain Lorentz reality.

\pandocstatementheading{\leavevmode\phantomsection\label{thm:online-ordinary-negative-response-initial-source}Theorem 4.18 (Actual initial source execution from an ordinary pair).}

For each admitted ordinary pair in the selected base domain, these
operations construct a weighted smooth gradient source \(\mathbb S(c)\)
and its complete initial contact presentation \(Q(c)\) on the whole
fixed initial manifold. They satisfy the independent kinematic and
vacuum tests and the exact input identities

\phantomsection\label{eq:online-ordinary-source-exact-unit-admission} \[
\begin{aligned}
\pi\mathbb S(c)&=\pi Q(c)=\mathscr P_\gamma\rho(c),\\
\sigma_\gamma\mathscr P_\gamma^{-1}\pi\mathbb S(c)&=c,
\qquad\mathfrak d\mathbb S(c)=\mathfrak d Q(c)=0.
\end{aligned}
\] Conversely, every independently tested kinematic source in either
canonical presentation, with zero defect and decoded state in the
retained product \(\mathcal B_\gamma\times B(0,r_y)\) of the transport
theorem, is recovered exactly by executing the ordinary pair
\(\sigma_\gamma\mathscr P_\gamma^{-1}\pi\). Both assertions concern the
marked canonical presentations; the actual spin marking is retained.

For any piecewise \(C^1\) seed path in \(S\), transport on the
constructed flat branch is the actual source evolution

\phantomsection\label{eq:online-ordinary-source-actual-evolution} \[
\frac{\mathrm d}{\mathrm dt}\mathbb S(t)=\mathscr V_{\dot s(t)}(\mathbb S(t)),
\qquad
\frac{\mathrm d}{\mathrm dt}Q(t)=\mathscr V^{\mathrm{CR}}_{\dot s(t)}(Q(t)).
\] Its endpoint depends only on the endpoint seed and it obeys the
identity and composition laws. Along the radial path from zero it can be
executed by the controlled finite Neumann--Picard stages of
\hyperref[cor:online-current-finite-neumann-picard]{Corollary A.15}.
Every finite stage is an actual kinematic source with the exact
prescribed endpoint seed. Its remaining defect tends to zero with the
base error bound in
\hyperref[eq:online-current-finite-execution-error]{the finite-execution
estimate}. The limit has every finite weighted order on the same base
domain; this last assertion does not require convergence of the raw
approximants in all higher norms.

On retained sets with the required finite higher bounds, for
\(m\ge k_0\) and each fixed finite angular order \(l\),

\phantomsection\label{eq:online-ordinary-source-balanced-two-sided-bounds}
\[
\begin{aligned}
\|\mathbb S(c)-\mathbb S(\widetilde c)\|_{\mathrm{op},m}
&\asymp\|s(c)-s(\widetilde c)\|_{\mathcal X_m},\\
\|Q(c)-Q(\widetilde c)\|_{\mathcal B_{m,l}}
&\asymp\|s(c)-s(\widetilde c)\|_{\mathcal X_m}.
\end{aligned}
\] The maps preserve smooth finite-dimensional families. For compact
parameter sets \(\mathcal K\Subset\operatorname{int}\mathcal K'\), every
specified finite spatial and parameter estimate on \(\mathcal K\) uses
finitely many input seminorms on \(\mathcal K'\) and the fixed base
margins. The seed ball and the spatial domain are fixed before these
higher orders. Effective ordinary cochain comparisons give the same
canonical source, with the retained actual spin identification.

After this execution and the independent tests have been proved, the
tensor readout of either source equals the earlier datum
\(\mathcal T(\mathscr P_\gamma\rho(c))\) of Theorem 3.13. That datum is
not an additional generating input.

\emph{Proof.} Normalize the ordinary pair using Theorems 3.9--3.10 and
Corollary 3.12, and execute its root transform to obtain \(s\). Apply
the initialized transport of Theorem A.14 with \((s_*,y_*)=(0,0)\). It
constructs one actual base endpoint \(Y(s)\), by the specified
differential and negative-response operations, and puts it in every
higher weighted space for smooth input. Form \[
\mathbb S(c)=\operatorname{Enc}_{\mathrm{gr}}(s,Y(s)),\qquad
Q(c)=\operatorname{Enc}_{\mathrm{CR}}(s,Y(s)).
\] The exact coefficient results above prove that these are actual
fields on the entire initial manifold, with both chiral operators, all
angular directions, the density and the actual spin marking. They meet
every independent kinematic test before the vacuum equation is used. By
the chain rule their evolution is the displayed horizontal source
equation. Its defect derivative is zero by the normalized response
proposition. The flat initial defect is zero, so the resulting fields
satisfy the independent Weyl-image test. This proves admission by a
source differential identity.

The exact decoder returns \((s,Y(s))\). Thus its seed readout is
\(s=\mathscr P_\gamma\rho(c)\), and the identity
\(\sigma_\gamma\rho(c)=c\) on the selected ordinary domain gives the
ordinary-class unit. In particular the unit concerns an actual ordinary
cohomology class, not only its polynomial quotient or its punctured root
values.

Conversely, independently test a supplied source and use its exact
kinematic decoder to obtain \((s,y)\) in the stated product. The
commutator proposition gives \(\mathcal F(s,y)=0\). The current
transport theorem identifies the entire zero-residual set in that
product with \((s,Y(s))\). Its exact coefficient inverse therefore
returns the supplied source, including its complete symbol and
directional coefficients, or its complete contact graph and density.
Executing \(\sigma_\gamma\mathscr P_\gamma^{-1}s\) gives that same
endpoint. This proves both object-level inverse assertions on the
independently specified initial domain.

Transport and its composition laws follow by conjugating the actual
coefficient transport with the exact encodings. The construction along
an arbitrary path in \(S\) stays on the same retained tube and uses no
path-length smallness assumption. For the finite realization, let
\(Y_{J,N}(s)\) be the endpoint in Corollary A.15. The stage
\(\operatorname{Enc}(s,Y_{J,N}(s))\) is an actual weighted smooth
kinematic source and has exact seed \(s\), even when its defect is
nonzero. The balanced coefficient estimate transfers the base endpoint
error to the operator and contact norms. The identity
\(\mathfrak d\operatorname{Enc}=\mathcal F\) transfers its residual
bound. This proves the claimed convergence and its precise finite-stage
meaning. All higher regularity concerns the same limit obtained at the
base order, by Theorem A.14.

The upper difference estimates follow from the higher endpoint
difference estimate and the finite coefficient encodings. The lower
estimates follow from their exact inverse readout and the bounded
ambient seed decoder. Repeated parameter differentiation uses the same
current inverse and finite coefficient formulas, so the family
quantifiers and finite seminorm dependence are precisely those of the
current transport theorem and the ordinary normalization. Two
effectively equivalent admitted cochains have the same root and ordinary
class, hence the same whole seed and the same constructed canonical
field. Their identification retains the chosen actual spin marking; no
quotient by a central spin sign has been introduced.

Finally \(Y(s)\) solves \(\mathcal F(s,Y(s))=0\). Uniqueness in the
independently defined tensor chart identifies it with the earlier
correction \(\mathcal C(s)\), and its tensor readout is consequently
\(\mathcal T(s)\). At this point the source operations have already
produced a solution, verified the vacuum test and recovered the input
pair. The argument applies at every curvature rank in the smooth
category. \hfill\(\square\)

The construction so far takes an admitted ordinary pair to its marked
initial source. Section 5 starts with a physical germ. We choose a
maximal disk and construct a pair whose development presents that germ.
Sections 6 and 7 compare the completed sources when their initial
hypersurfaces differ.

\section{5. Local displays of arbitrary vacuum
germs}\label{local-displays-of-arbitrary-vacuum-germs}\label{sec:submission-local-coverage}

A small weighted construction can represent an arbitrary local vacuum
geometry once its physical scale and initial embedding are retained. We
first choose a local maximal disk with its point and tangent plane
fixed. We then extend its normalized data exactly across a protected
region and apply the whole-data decoder. Finally, local Cauchy
uniqueness identifies the generated development with the original
physical germ.

The exterior extension supplies the whole-space input to the
constructor. Only its agreement with the disk on the protected region
enters the identification of the physical germ.

\subsection{5.1. Conventions and analytic
inputs}\label{conventions-and-analytic-inputs}\label{sec:submission-coverage-conventions}

A spacetime germ is four-dimensional, pointed, oriented and time
oriented, and its actual local spin structure is retained. We use
signature \((+---)\). If \(j\) is a spacelike embedding and \(n\) its
future unit normal, its induced data are

\phantomsection\label{eq:submission-coverage-induced-data} \[
h=-j^*g,
\qquad
K(X,Y)=-g(\nabla_Xn,Y).
\] Thus a Gaussian presentation has \(g=\mathrm dt^2-h_t\) and
\(K=\tfrac12\partial_t h_t\). A maximal vacuum datum satisfies

\phantomsection\label{eq:submission-coverage-maximal-constraints} \[
\operatorname{tr}_hK=0,
\qquad
\operatorname{div}_hK=0,
\qquad
\operatorname{Scal}(h)=|K|_h^2.
\] Write \(B_a\subset\mathbb R^3\) for the Euclidean ball of radius
\(a\) and \(\langle x\rangle=(1+|x|^2)^{1/2}\). For a tensor in fixed
Cartesian coordinates, the weighted Sobolev norm is

\phantomsection\label{eq:submission-coverage-weighted-norm} \[
\|u\|_{H^s_\gamma}^2
=
\sum_{|\beta|\le s}
\int_{\mathbb R^3}
\langle x\rangle^{-2\gamma-3+2|\beta|}
|\partial^\beta u|^2\,\mathrm dx.
\] The notation \(\mathscr S^\mu\) means that
\(\sup_x\langle x\rangle^{-\mu+|\beta|}|\partial^\beta u(x)|\) is finite
for every multi-index \(\beta\). Weighted Hölder norms are those of
\hyperref[def:online-weighted-holder-chart-spaces]{Definition A.1}, as
used in the whole-data spatial-gauge theorem. We fix one weight
\(\gamma=-1/2-\varepsilon_\gamma\), with \(0<\varepsilon_\gamma<1/2\),
in the interval furnished by the exterior inverse theorem. No choice of
an inverse at every weight in \((-1,-1/2)\) is needed.

The following analytic inputs are used in their stated forms.

\textbf{Relative extension.} Fix the exterior normalization described in
\hyperref[app:online-relative-extension]{Appendix B}. Theorem B.1
selects an operation \(\mathfrak E\) on smooth whole prefields
\((\widehat h,\widehat K)\) satisfying
\(\operatorname{tr}_{\widehat h}\widehat K=0\) on \(\mathbb R^3\), the
maximal constraints on \(B_1\), finiteness of all weighted norms and one
finite base-smallness condition. Its output \((H,L)\) is maximal vacuum
data on \(\mathbb R^3\) and agrees exactly with the prefields on
\(B_1\). For the same selected output,

\phantomsection\label{eq:submission-relative-extension-input} \[
H-\delta\in H^M_\gamma,\qquad
L\in H^{M-1}_{\gamma-1}
\quad\text{for every integer }M\ge S_0.
\] Each norm and each specified finite parameter derivative is bounded
by finitely many input seminorms and fixed positive margins. The base
domain is chosen before those higher orders. The original extension
result is \hyperref[ref-Czimek2018]{Czimek} (2018, Theorem 3.1); the
protected relative normalization, nearby stronger weight and same-branch
assertion used here are the additional inputs specified in Appendix B.

\textbf{Ordinary classes and their whole seeds.} Let
\(\mathscr W_\gamma\) be the paired root-class space of
\hyperref[def:online-weighted-root-domain]{Definition 3.5}, with its
coupled residue reality and weighted Fourier-amplitude tests. The
isomorphism of \hyperref[thm:online-root-seed-isomorphism]{Theorem 3.6}
is

\phantomsection\label{eq:submission-root-seed-transform} \[
\mathscr P_\gamma:\mathscr W_\gamma
\longrightarrow\mathcal X^\infty_\gamma.
\] Here \(\mathcal X^\infty_\gamma\) consists of the smooth Euclidean
transverse-traceless metric and momentum seeds, of symbol orders
\(\gamma\) and \(\gamma-1\), in the indicated weighted chart spaces. All
their weighted derivatives are finite; smallness is imposed at the
chosen finite base order. The root-class complex used for this transform
is linear. Nonlinear source operations use the actual coefficient
complexes \(\mathscr E_2^\bullet\) and \(\mathscr E_{-6}^\bullet\) of
\hyperref[prop:online-actual-coefficient-exact-sequence]{Proposition
3.7}.

\hyperref[thm:online-canonical-actual-pair-lift]{Theorem 3.9} gives a
continuous real-linear section

\phantomsection\label{eq:submission-actual-ordinary-lift} \[
\sigma_\gamma:\mathscr W_\gamma
\longrightarrow
H^1(\mathscr E_2^\bullet)\oplus
H^1(\mathscr E_{-6}^\bullet),
\qquad
\rho\sigma_\gamma=\operatorname{id}.
\] The first cohomology here is formed using actual closed cochains and
actual coboundaries. The map \(\rho\) is the root readout. The section
has selected closed representatives, is compatible with smooth passive
parameters, and preserves the calibrated root reality. Its positive
representative is anchored at the origin and has symbol order
\(\gamma+1\); changing the anchor gives an actual coboundary. The
negative representative has the decaying orders stated in that theorem.

Fix the flat source as reference and choose one rotation-invariant base
neighborhood \(\mathcal B_\gamma\) of zero in
\(\mathcal X^\infty_\gamma\) on which the initialized transport of
Theorem A.14 and the source execution of
\hyperref[thm:online-ordinary-negative-response-initial-source]{Theorem
4.18} are defined. These neighborhoods contain a ball for the finitely
many base norms; a rotation-invariant ball can be taken by using their
suprema over the compact rotation group. All finite chart and inverse
thresholds are included in this choice. Define

\phantomsection\label{eq:submission-selected-ordinary-domain} \[
\mathscr C^{\mathrm{ord}}_{\gamma,\mathcal B}
=
\sigma_\gamma\bigl(\mathscr P_\gamma^{-1}
(\mathcal B_\gamma)\bigr).
\] Throughout this section, \(c=(c_+,c_-)\) denotes an actual ordinary
pair in this domain; its whole seed is

\phantomsection\label{eq:submission-ordinary-seed-readout} \[
s(c)=\mathscr P_\gamma\rho(c).
\] On the selected neighborhood, the canonical representatives supplied
by \(\sigma_\gamma\) satisfy
\hyperref[def:online-effective-ordinary-cochain-domain]{Definition
3.11}. For other supplied representatives test actual closedness, root
membership in \(\mathscr P_\gamma^{-1}(\mathcal B_\gamma)\), and
vanishing of the polynomial obstructions of
\hyperref[thm:online-actual-cochain-obstruction-test]{Theorem 3.10}. The
positive obstruction is tested in its specified quotient. Effective
arrows are actual primitives modulo closed zero-cochains that induce the
same cochain change.

\textbf{Whole-data gauge and decoding.} For smooth whole maximal vacuum
data \((H,L)\), the hypotheses of
\hyperref[cor:online-gauge-whole-data-decoding]{Corollary A.11} are

\phantomsection\label{eq:submission-whole-gauge-domain} \[
H-\delta\in\mathscr S^\gamma,\qquad
L\in\mathscr S^{\gamma-1},\qquad
\|H-\delta\|_{C^{k+2,\alpha}_\gamma}
+\|L\|_{C^{k+1,\alpha}_{\gamma-1}}<\epsilon_{\mathrm g}.
\] The based gauge of
\hyperref[thm:online-gauge-based-normalization]{Theorem A.9} has the
form

\phantomsection\label{eq:submission-based-gauge} \[
\phi_H=\operatorname{id}+X_H,\qquad
X_H(0)=0,\qquad
X_H\in\mathscr S^{\gamma+1},\qquad
\|DX_H\|_\infty<\frac12.
\] Thus its displacement is sublinear. Its isotopy from the identity has
an initialized actual spin lift. Apply the exact seed decoder of
\hyperref[prop:online-chart-exact-ambient-decoder]{Proposition A.4} to
\(\phi_H^*(H,L)\) and then apply \(\mathscr P_\gamma^{-1}\) and
\(\sigma_\gamma\). After reducing \(\epsilon_{\mathrm g}\) once so that
the seed lies in the fixed \(\mathcal B_\gamma\), the resulting ordinary
pair satisfies

\phantomsection\label{eq:submission-coverage-decoder-input} \[
\phi_H^*(H,L)=\mathcal T(s(c))
=:(h(c),K(c)).
\] The map \(\mathcal T\) is the nonlinear constraint chart of
\hyperref[thm:online-chart-base-constraint-bijection]{Theorem A.6}. The
ordinary-pair constructor of
\hyperref[thm:online-ordinary-pair-data-bijection]{Theorem 3.13} is
\(c\mapsto\mathcal T(s(c))\). The decoder, actual class lift and based
gauge have finite smooth-family estimates on the same selected base
domain. Proper Euclidean rotations preserve the chart and gauge. These
operations act on the whole extended fields.

\textbf{Source execution, geometric completion and wave completion.} For
a selected actual representative of \(c\),
\hyperref[thm:online-ordinary-negative-response-initial-source]{Theorem
4.18} constructs a canonical admitted initial gradient source and its
complete initial contact presentation, with seed \(s(c)\) and exact
input recovery. Effective ordinary arrows give the same canonical
output. The fixed reference and the supplied pair determine the source
responses and hence the constructor. The spacetime metric is recovered
when the resulting vacuum data are developed.

The initialized Cartan completion of
\hyperref[cor:online-ordinary-pair-cartan-collar]{Corollary 2.12} is the
subsequent geometric development of the readout \((h(c),K(c))\). Its
independent admission, subsidiary identities and geometric recognition
are proved in \hyperref[thm:online-cartan-collar-recognition]{Theorem
2.9}, using the exact smooth Einstein Cauchy input of
\hyperref[thm:online-smooth-local-cauchy-input]{Theorem 2.6}.
\hyperref[thm:online-marked-collar-full-hom]{Theorem 2.10} carries every
specified initial physical and actual spin comparison.
\hyperref[thm:online-common-collar-finite-families]{Theorem 2.14}
supplies a common local parent and all finite family estimates.
\hyperref[cor:online-fixed-slice-geometric-choice-independence]{Corollary
2.15} compares exterior, gauge and scale choices for the initial and
finite Cartan parts when the physical initial germ is fixed.

Finally, \hyperref[sec:online-covariant-wave-enrichment]{Section 6}
constructs the wave completion from this source collar. Work on a
smaller relatively compact causally convex globally hyperbolic collar,
retaining every nonzero null direction over its point core by a finite
angular cover of the full skies. Mixed carrier compositions use
compactly contained ray domains whose rays meet each retained spacelike
carrier once, with properly supported representatives and a positive
buffer from the patch boundaries. Its scalar kernel class
\(\Lambda_g^+\) is characterized by the independent bisolution,
commutator and positive-frequency wavefront tests of
\hyperref[def:online-independent-wave-kernel]{Definition 6.9} and
\hyperref[thm:online-canonical-wave-kernel]{Theorem 6.10}, modulo
kernels jointly smooth in the points and supplied parameters. Let
\(C_{ab}\) denote the mixed trace of
\hyperref[def:online-wave-carrier-displays]{Definition 6.12} with its
two slice half-density factors, and let \(C_{ap}\) be the one-slice
point column of \hyperref[def:online-wave-evaluation-section]{Definition
6.15}, with only the carrier half-density factor. Set
\(W_a=C_{aa}^{-1/2}\) with the specified positive principal root, and
put

\phantomsection\label{eq:submission-wave-completion-interface} \[
T_{ab}=W_aC_{ab}W_b,\qquad
e_a(p)=W_aC_{ap}.
\] \hyperref[thm:online-wave-carrier-cocycle]{Theorem 6.14} and
\hyperref[thm:online-ambient-point-connection]{Theorem 6.19} give the
complete carrier comparison and the ambient point connection on all
module sections.
\hyperref[prop:online-wave-full-evaluation-comparison]{Proposition 6.16}
proves the full evaluation-section identity, including every lower
coefficient. These are derived structures. Chosen exact smoothing
representatives are not additional generating inputs. The quantitative
assertions used below concern finite complete coefficients and joint
derivatives on the retained parents, as proved in
\hyperref[prop:online-wave-kernel-finite-families]{Proposition 6.11},
\hyperref[prop:online-wave-carrier-families-scale]{Proposition 6.17} and
\hyperref[prop:online-point-connection-finite-families]{Proposition
6.22}. \hyperref[thm:online-wave-enrichment-recognition]{Theorem 6.24}
supplies independent admission and recovery for the normalized scalar
enrichment over the constructed finite Cartan collar.

Choose \(S_0\) large enough for these inputs. If the spatial gauge is
measured in the fixed weighted \(C^{k+2,\alpha}\) and \(C^{k+1,\alpha}\)
norms, require

\phantomsection\label{eq:submission-coverage-base-order} \[
S_0>k+2+\alpha+\frac32,
\] as well as the finite lower requirements of the extension and source
construction. Take one base threshold small enough for extension, gauge,
ordinary decoding and the selected actual source-input ball. The radius
below is chosen only after this order and threshold have been fixed.

\subsection{5.2. Normal coordinates and maximal
disks}\label{normal-coordinates-and-maximal-disks}\label{sec:submission-normalization}

\pandocstatementheading{\leavevmode\phantomsection\label{lem:submission-normal-coordinate-family}Lemma 5.1 (Normalization of a smooth metric family).}

Let \((g_\lambda,p_\lambda,P_\lambda)\) be a smooth finite-dimensional
family of Lorentzian metric germs with a specified spacelike tangent
three-plane \(P_\lambda\) at \(p_\lambda\). Locally in the parameter,
choose a smooth oriented orthonormal frame
\(e_\lambda:\mathbb R^{1,3}\to T_{p_\lambda}M_\lambda\) whose time
vector is future pointing and whose spatial image is \(P_\lambda\). On
common representatives and after restriction of the parameter
neighborhood, there is \(r_0>0\) such that

\phantomsection\label{eq:submission-normalized-metric} \[
\phi_{\lambda,r}(z)
=
\exp^{g_\lambda}_{p_\lambda}(r e_\lambda z),
\qquad
g_{\lambda,r}
=
r^{-2}\phi_{\lambda,r}^*g_\lambda
\] are defined on \(Q=(-2,2)\times B_3\) for \(0<r\le r_0\). The
coefficients extend smoothly to \(r=0\) and have the form

\phantomsection\label{eq:submission-normalized-metric-remainder} \[
g_{\lambda,r}
=
\eta+r^2G(\lambda,r,z),
\qquad
\eta=\mathrm dt^2-\sum_{i=1}^3(\mathrm dx^i)^2.
\] The same \(r_0\) works at every finite derivative order. On retained
compact parameter subsets, each finite mixed derivative of \(G\) is
bounded by finitely many derivatives of the marked family, frame and
trivialization and by the retained coordinate margins. Ricci-flatness is
preserved for \(r>0\).

\emph{Proof.} Trivialize the family near its marked points. The geodesic
equation depends smoothly on its metric, initial point and initial
vector, so its normal-coordinate coefficients
\(\widetilde g_{\lambda,\mu\nu}(z)\) are smooth on one common
neighborhood. The orthonormal frame and normal-coordinate identities
give \[
\widetilde g_\lambda(0)=\eta,
\qquad
\partial_a\widetilde g_\lambda(0)=0.
\] The normalized coefficients are
\(\widetilde g_{\lambda,\mu\nu}(rz)\). Taylor's integral formula
therefore gives

\phantomsection\label{eq:submission-normal-coordinate-integral} \[
G_{\mu\nu}(\lambda,r,z)
=
z^az^b\int_0^1(1-s)
\partial_a\partial_b\widetilde g_{\lambda,\mu\nu}(srz)\,\mathrm ds.
\] This formula is smooth at \(r=0\) and gives every finite mixed
estimate on the same parent. Multiplication of a metric by a positive
constant leaves its Levi--Civita connection and covariant Ricci tensor
unchanged. Naturality under pullback consequently gives
\(\operatorname{Ric}(g_{\lambda,r})=\phi_{\lambda,r}^*\operatorname{Ric}(g_\lambda)\).
\hfill\(\square\)

\pandocstatementheading{\leavevmode\phantomsection\label{thm:submission-marked-maximal-disk}Theorem 5.2 (A maximal disk with prescribed point and tangent plane).}

Fix \(0<\alpha<1\). Each smooth metric in a sufficiently small
\(C^{4,\alpha}\) neighborhood of \(\eta\) on \(Q\) has a selected small
spacelike maximal graph

\phantomsection\label{eq:submission-marked-maximal-graph} \[
j_g(x)=(u_g(x),x),
\qquad
x\in B_2,
\qquad
u_g(0)=0,\quad Du_g(0)=0.
\] Its boundary value is \(a_g+b_g\cdot x\) on \(\partial B_2\). It is
unique among small graphs with affine boundary value and the displayed
point conditions, and is smooth near \(\overline B_{7/4}\).

For a jointly smooth finite-dimensional family in this base
neighborhood, the selection is jointly smooth on that retained interior
ball. Each finite spatial and parameter derivative has a bound in
finitely many derivatives of the metric family and the fixed base
margins. The base neighborhood and retained ball do not depend on the
requested higher order.

\emph{Proof.} Let \(E_i=\partial_i+u_i\partial_t\) be the graph tangent
vectors, let \(N\) be its future unit normal, and put
\(\nu=g(N,\partial_t)\). Near the flat zero graph, \(\nu\) is positive
and the induced metric \(h\) is positive definite. Differentiation of
\(g(N,E_j)=0\) yields

\phantomsection\label{eq:submission-maximal-graph-equation} \[
K_{ij}
=
\nu u_{ij}+g\bigl(N,\Gamma(E_i,E_j)\bigr),
\qquad
H_g(u)
=
\nu h^{ij}u_{ij}
+h^{ij}g\bigl(N,\Gamma(E_i,E_j)\bigr).
\] The coefficient \(\nu h^{ij}\) is uniformly positive on one small
graph neighborhood. At the flat metric,

\phantomsection\label{eq:submission-maximal-graph-linearization} \[
H_\eta(u)
=
\operatorname{div}\left(
\frac{Du}{\sqrt{1-|Du|^2}}
\right),
\qquad
D_uH_\eta(0)=\Delta.
\] Write \(u=v+a+b\cdot x\), with \(v\in C_D^{2,\alpha}(\overline B_2)\)
of zero Dirichlet trace. This condition imposes no boundary normal
derivative. Consider

\phantomsection\label{eq:submission-augmented-graph-map} \[
\mathcal F(g,v,a,b)
=
\bigl(
H_g(v+a+b\cdot x),\,
v(0)+a,\,
Dv(0)+b
\bigr).
\] It is continuously differentiable into
\(C^{0,\alpha}\times\mathbb R^4\) on the indicated metric and graph
spaces; the extra metric derivatives give the composition margin. Its
derivative in the unknowns at the flat zero graph is

\phantomsection\label{eq:submission-augmented-graph-derivative} \[
(v,a,b)
\longmapsto
\bigl(\Delta v,\ v(0)+a,\ Dv(0)+b\bigr).
\] For a prescribed \((f,c,d)\), solve the zero-Dirichlet problem
\(\Delta v=f\) and set \(a=c-v(0)\) and \(b=d-Dv(0)\). The Dirichlet
estimate and bounded point evaluations give a bounded inverse. The
implicit-function theorem selects the graph and its affine boundary
modes. Restrict the base neighborhood once so that the graph is
spacelike, stays in \(Q\), and the augmented derivative remains
invertible. These four boundary modes impose the point and tangent
conditions without adding conditions to a fixed zero-Dirichlet problem.

For smooth input, interior elliptic estimates improve this existing
solution. At each order the highest derivative has the same elliptic
principal part; the remaining terms contain finitely many previously
controlled derivatives. For a requested order, use a finite chain of
nested balls between \(B_{7/4}\) and \(B_2\). Every chain ends on the
same retained ball. Its constants may depend on the order, but no higher
norm must be small.

On a smooth finite-dimensional family, the augmented map restricted to
that family is smooth to each required finite order because the family
supplies the extra metric derivatives needed for composition. Its base
implicit solution is smooth into \(C_D^{2,\alpha}\times\mathbb R^4\).
Parameter differentiation gives the same invertible augmented
linearization with lower-order forcing. Interior elliptic estimates
bound spatial derivatives successively in parameter order. Difference
quotients in the base space identify those solutions with derivatives of
the already selected graph. Higher spatial estimates and interpolation
give continuity in each higher spatial space; the fundamental theorem of
calculus gives differentiability there. Induction proves joint
smoothness and the finite dependence asserted above. This is a family
statement on smooth inputs, not a same-regularity smooth action on the
full finite Hölder space. \hfill\(\square\)

\pandocstatementheading{\leavevmode\phantomsection\label{cor:submission-normalized-maximal-data}Corollary 5.3 (Normalized maximal data).}

Apply Theorem 5.2 to Lemma 5.1, including \(r=0\). On \(B_{7/4}\) the
selected graph and induced data have the form

\phantomsection\label{eq:submission-normalized-maximal-data} \[
u_{\lambda,r}=r^2U_{\lambda,r},
\qquad
\bar h_{\lambda,r}=\delta+r^2A_{\lambda,r},
\qquad
\bar K_{\lambda,r}=r^2B_{\lambda,r},
\] with jointly smooth \(U,A,B\). For Ricci-flat input the data satisfy
the maximal vacuum constraints. With
\(\Psi_{\lambda,r}=\phi_{\lambda,r}\circ j_{\lambda,r}\), their physical
normalization is

\phantomsection\label{eq:submission-physical-graph-data} \[
-\Psi_{\lambda,r}^*g_\lambda=r^2\bar h_{\lambda,r},
\qquad
K_{\Psi_{\lambda,r}}=r\bar K_{\lambda,r}.
\] \emph{Proof.} At \(r=0\) the selected graph vanishes, and the first
\(r\) derivative of the normalized metric is zero. Differentiate the
augmented equation and use its invertible derivative to obtain
\(\partial_ru_{\lambda,0}=0\). Taylor's integral formula gives
\(u=r^2U\). The induced metric and
\hyperref[eq:submission-maximal-graph-equation]{the graph equation} give
the remaining factorizations. Ricci-flatness and maximality give the
vacuum constraints.

Under \(g\mapsto r^2g\), the future unit normal is divided by \(r\), the
induced metric is multiplied by \(r^2\), and the covariant second
fundamental form is multiplied by \(r\). Pullback gives the stated
physical identities. \hfill\(\square\)

Each prescribed spatial or passive-parameter derivative of the three
vanishing quantities is \(O(r^2)\). With \(\ell\) additional \(r\)
derivatives the corresponding bound is \(O(r^{\max(2-\ell,0)})\). Thus
one positive radius can meet all finite base requirements while every
higher requested norm remains finite at that radius.

\newpage
\subsection{5.3. Protected extension and ordinary
decoding}\label{protected-extension-and-ordinary-decoding}\label{sec:submission-protected-extension}

\pandocstatementheading{\leavevmode\phantomsection\label{lem:submission-buffered-maximal-extension}Lemma 5.4 (Extension of buffered maximal data).}

Smooth maximal vacuum data \((\bar h,\bar K)\) near
\(\overline B_{3/2}\) that are sufficiently small relative to
\((\delta,0)\) in \(H^{S_0}\times H^{S_0-1}\) have a selected smooth
maximal vacuum extension \((H,L)\) on \(\mathbb R^3\) with

\phantomsection\label{eq:submission-buffered-extension} \[
(H,L)|_{B_1}=(\bar h,\bar K),
\qquad
H-\delta\in H^M_\gamma,
\quad
L\in H^{M-1}_{\gamma-1}
\quad\text{for every finite }M.
\] The base norm is bounded by a fixed constant times the local input
norm. Each higher norm is bounded using finitely many higher local norms
and the fixed base margins. Supplied smooth finite-dimensional families
have one such smooth selection, with every specified mixed estimate on
the same base domain.

\emph{Proof.} Choose \(0\le\chi\le1\), supported in \(B_{3/2}\) and
equal to one near \(\overline B_1\). Extend the cutoff products by zero
and put

\phantomsection\label{eq:submission-cutoff-prefields} \[
\widehat h=\delta+\chi(\bar h-\delta),
\qquad
\widehat K=\operatorname{TF}_{\widehat h}(\chi\bar K).
\] The metric is positive, and \(\widehat K\) is globally tracefree for
that metric. On \(B_1\) these prefields agree exactly with the supplied
constraint data. Their perturbations are compactly supported, so they
have every required weighted norm. Multiplication, metric inversion and
tracefree projection give the base estimate

\phantomsection\label{eq:submission-cutoff-base-estimate} \[
\|\widehat h-\delta\|_{H^{S_0}_\gamma}
+\|\widehat K\|_{H^{S_0-1}_{\gamma-1}}
\le
C\left(
\|\bar h-\delta\|_{H^{S_0}(B_{3/2})}
+\|\bar K\|_{H^{S_0-1}(B_{3/2})}
\right).
\] Theorem B.1 therefore applies, with its normalization fixed as in
§5.1. It gives the protected equality, the base bound and all higher
bounds for the same initialized output. The cutoff and projection
preserve smooth parameter dependence and finite mixed bounds, so the
same input also gives the family assertion. Compact support already
supplies any stronger decay required of the prefield; no additional
strict-weight assumption on the local data is needed. \hfill\(\square\)

By the fixed Sobolev-to-Hölder margin, the base output lies in the gauge
domain after one reduction of the local threshold. Its higher weighted
estimates imply the symbol bounds. Apply the whole-data gauge and
decoder to obtain

\phantomsection\label{eq:submission-extended-data-decoding} \[
(h(c),K(c))
=
\phi_H^*(H,L)
=
\mathcal T(s(c)).
\] All inverse operations act on the whole extended field. The completed
source is subsequently restricted to a local patch.

\pandocstatementheading{\leavevmode\phantomsection\label{thm:submission-marked-initial-coverage}Theorem 5.5 (Ordinary pairs for marked maximal initial germs).}

Every smooth pointed maximal vacuum initial germ \((\Sigma,h,K,o)\) with
its actual spin marking has an ordinary-pair presentation in the fixed
small weighted domain: there are \(c\), \(r>0\) and an
orientation-preserving embedding \(\iota:V\to\Sigma\), with
\(\iota(0)=o\), such that

\phantomsection\label{eq:submission-initial-coverage} \[
\iota^*h=r^2h(c)|_V,
\qquad
\iota^*K=rK(c)|_V.
\] The embedding has the specified actual spin comparison. The pair
generates its canonical admitted initial source and complete initial
contact presentation. The construction can be selected smoothly near
every parameter value of a smooth finite-dimensional family, using one
positive radius, one protected ball and one base domain.

\emph{Proof.} Choose oriented \(h\)-orthonormal frames at the marked
points and their actual spin lifts, smoothly near the chosen parameter
value. Set

\phantomsection\label{eq:submission-initial-normalization} \[
\psi_r(y)=\exp^h_o(r e y),
\qquad
\bar h_r=r^{-2}\psi_r^*h,
\qquad
\bar K_r=r^{-1}\psi_r^*K.
\] For the unscaled normal-coordinate coefficients,

\phantomsection\label{eq:submission-initial-smallness} \[
(\bar h_r)_{ij}(y)=\widetilde h_{ij}(ry)
=\delta_{ij}+r^2A_{ij}(r,y),
\qquad
(\bar K_r)_{ij}(y)=r\,\widetilde K_{ij}(ry).
\] All prescribed finite derivatives are bounded on one buffered ball.
In this general initial-data normalization the momentum is \(O(r)\);
Corollary 5.3 obtains \(O(r^2)\) for its particular spacetime-selected
disk.

The constraints are preserved. For constant \(a>0\), under \(h'=a^2h\)
and \(K'=aK\),

\phantomsection\label{eq:submission-constraint-scaling} \[
\begin{aligned}
\operatorname{Scal}(h')&=a^{-2}\operatorname{Scal}(h),
&
|K'|_{h'}^2&=a^{-2}|K|_h^2,\\
\operatorname{tr}_{h'}K'&=a^{-1}\operatorname{tr}_hK,
&
\operatorname{div}_{h'}K'&=a^{-1}\operatorname{div}_hK.
\end{aligned}
\] The spatial connection is unchanged. Apply these identities with
\(a=r^{-1}\) after pullback.

With the finite base order and all constructor thresholds already fixed,
choose \(r\) so that Lemma 5.4 and the whole gauge and decoder apply.
Let \((H_r,L_r)\) be the extension. If a particular orthonormal point
frame is retained, define

\phantomsection\label{eq:submission-polar-frame-correction} \[
R_r
=
D\phi_{H_r}(0)
\bigl((\phi_{H_r}^*H_r)(0)\bigr)^{-1/2}\in SO(3),
\qquad
O_r=R_r^{-1}.
\] Here \(H_r(0)=\delta\). Decode the whole datum
\(O_r^*\phi_{H_r}^*(H_r,L_r)\), which remains in the spatial gauge, and
call its ordinary pair \(c_r\). If no particular orthonormal frame is
marked, omit this rotation. On one smaller \(V\) such that
\(\phi_{H_r}O_r(V)\subset B_1\), set

\phantomsection\label{eq:submission-physical-initial-embedding} \[
\iota_r
=
\psi_r\circ\phi_{H_r}\circ O_r|_V.
\] Protected equality gives the two tensor identities in the statement.
The polar formula carries the canonical orthonormal frame of the rotated
datum to the selected normalized frame.

Lift the gauge isotopy and the small rotation with their initialized
spin signs, and use the frame map \(e_a\mapsto r^{-1}e_a\) for positive
metric scaling. This produces the actual spin comparison with the
prescribed central sign. On a connected local germ the lift is
determined by its value at the marked frame.

The actual-class decoder supplies canonical closed representatives and
the prescribed primitive comparisons.
\hyperref[thm:online-ordinary-negative-response-initial-source]{Theorem
4.18} constructs the admitted initial source from this pair and proves
exact input recovery. The original metric is used to obtain a
representing input in this proof; it is not a separate argument of that
source constructor.

For a family, the initial smallness estimates permit a single positive
\(r\). Relative extension, based gauge, polar normalization, decoding
and actual class transformation are smooth on the same selected domain.
Their finite estimates, followed by those of source execution, give the
claimed selection. Raising the requested order changes the required
finite norms and their bounds, not the extension or radius.
\hfill\(\square\)

\subsection{5.4. Physical scale on the generated
source}\label{physical-scale-on-the-generated-source}\label{sec:submission-coverage-physical-scale}

\pandocstatementheading{\leavevmode\phantomsection\label{prop:submission-complete-source-physical-scale}Proposition 5.6 (Physical scale and the wave completion).}

Let an admitted source collar have metric \(g\) and the wave completion
specified in §5.1. For \(r>0\) constant on spacetime, the scale
restoration \(g^{[r]}=r^2g\) has

\phantomsection\label{eq:submission-gradient-physical-scale} \[
\mathbb D_\pm^{[r]}=r^{-1}\mathbb D_\pm,\qquad
h^{[r]}=r^2h,\qquad K^{[r]}=rK.
\] The internal Lorentz and spin pairings and the future unit section
are kept fixed. The recovered Lorentz and spin connections are
unchanged, while the soldering is multiplied by \(r\). The actual spin
identification follows the positive scaling path and retains the
supplied central sign.

For the scalar wave structures in their fixed-coordinate realizations,

\phantomsection\label{eq:submission-wave-physical-scale} \[
\begin{aligned}
\Lambda_{r^2g}^+&=r^{-2}\Lambda_g^+,
&
C_{ab}^{[r]}&=rC_{ab},\\
W_a^{[r]}&=r^{-1/2}W_a,
&
T_{ab}^{[r]}&=T_{ab},
&
e_a^{[r]}&=r^{-1}e_a.
\end{aligned}
\] The ambient point connection is unchanged. These identities are
identities in the specified complete coefficient and joint-smoothing
calculus. The numerical scale remains part of the physical display.

\emph{Proof.} The scalar commutator of \(r^{-1}\mathbb D_\pm\)
multiplies the spatial symbol by \(r^{-1}\). Its metric readout is
therefore \(r^2h\) and its soldering is \(r\Theta\). Contracting the
whole directional operator against this soldering leaves the directional
connection unchanged. Its normal readout is consequently \(rK\). This is
the full-gradient calculation of
\hyperref[prop:online-cartan-physical-scale]{Proposition 2.11}; it uses
the uncontracted operators, retaining the full second fundamental form.
The proposition also gives the actual finite Cartan transport with the
physical clock \(t^{\prime}=rt\).

The positive metric rescaling identifies orthonormal frames by
\(v_a\mapsto r^{-1}v_a\). Initialization at \(r=1\) gives the spin
identification along positive scales. A supplied opposite central lift
remains distinct. Constant metric scaling leaves the unparametrized null
rays, skies and contact distribution unchanged and multiplies the
physical point density by \(r^4\).
\hyperref[prop:online-canonical-optical-completion]{Proposition 7.4} and
\hyperref[prop:online-canonical-finite-rank-models]{Proposition 7.5}
reconstruct their selected optical and canonical finite-rank fields from
the rescaled soldering and the same connection. This restores their
numerical physical scale; it does not declare an isometry between the
two unscaled metrics.

For the wave part, \(\Box_{r^2g}=r^{-2}\Box_g\) and the four-dimensional
volume form scales by \(r^4\). The scalar causal kernel therefore has
weight \(r^{-2}\). The independent kernel characterization gives that
same weight to \(\Lambda_g^+\). Each slice half-density has weight
\(r^{3/2}\), so the two-slice and one-slice kernels have weights \[
C_{ab}^{[r]}=rC_{ab},
\qquad
C_{ap}^{[r]}=r^{-1/2}C_{ap}.
\] The positive-principal inverse square root gives the stated law for
\(W_a\). Substituting into the formulas for \(T_{ab}\) and \(e_a\)
proves their weights. The conormal point families and their actual point
derivatives are unchanged in the fixed-coordinate identification, giving
the connection assertion. These are the complete-class scale laws of
\hyperref[prop:online-wave-carrier-families-scale]{Proposition 6.17} and
\hyperref[prop:online-point-connection-finite-families]{Proposition
6.22}, and also give \(\Box_{r^2g}^{\nabla^{\mathrm{amb}}}e^{[r]}
=r^{-3}\Box_g^{\nabla^{\mathrm{amb}}}e\). All comparisons of physical
germs are made after this numerical scale has been restored.
\hfill\(\square\)

\subsection{5.5. Coverage of spacetime
germs}\label{coverage-of-spacetime-germs}\label{sec:submission-spacetime-coverage}

\pandocstatementheading{\leavevmode\phantomsection\label{thm:submission-vacuum-germ-coverage}Theorem 5.7 (Ordinary-pair coverage of smooth Ricci-flat germs).}

Every smooth pointed, oriented and time-oriented Lorentzian Ricci-flat
germ, with a chosen spacelike tangent plane and actual local spin
marking, has an ordinary-pair presentation of a local maximal Cauchy
disk through that point and tangent plane. The pair generates an
admitted initial source, its geometric collar and the wave completion
specified in §5.1. The physically rescaled development is isometric to
the original germ, with the retained future normal and actual spin
comparison. Its generated complete structures have the physical scale
laws of Proposition 5.6.

The construction is locally available for smooth finite-dimensional
families on common representatives, with one positive normalization
radius and one retained parent. Local geometric Cauchy uniqueness
identifies the normalized development with the normalized germ on a
common smaller domain of dependence. Restoring the initial data to
\((r^2h,rK)\) returns the development to the physical metric scale of
the given germ.

\emph{Proof.} Lemma 5.1, Theorem 5.2 and Corollary 5.3 give maximal
vacuum data on a buffered ball. Choose one radius satisfying the finite
base conditions. Apply Lemma 5.4, the whole gauge and decoder, and the
polar frame correction when required. Compose their initial embedding
with \(\Psi_{\lambda,r}\). Its physical initial tensors are exactly
\(r^2h(c)\) and \(rK(c)\).

Execute the pair by
\hyperref[thm:online-ordinary-negative-response-initial-source]{Theorem
4.18} to obtain the admitted initial source.
\hyperref[cor:online-ordinary-pair-cartan-collar]{Corollary 2.12}
supplies its initialized Cartan completion from these derived data.
\hyperref[lem:online-incidence-identification]{Lemma 7.2} constructs its
actual full null incidence on a smaller retained core;
\hyperref[prop:online-canonical-optical-completion]{Proposition 7.4} and
\hyperref[prop:online-canonical-finite-rank-models]{Proposition 7.5}
construct its optical and canonical finite-rank fields. Shrink to a
relatively compact causally convex globally hyperbolic collar and apply
\hyperref[thm:online-wave-enrichment-recognition]{Theorem 6.24} to
construct its independently admitted scalar enrichment, using the wave
formulas of §5.1. The retained-carrier construction in
\hyperref[def:online-wave-carrier-displays]{Definition 6.12} uses
spacelike transversality and compactness of the projective sky to give
one smaller point neighborhood with full angular coverage, single
carrier intersections and boundary buffers. No additional free initial
field is supplied.

Before restoring scale, compare that development with \(g_{\lambda,r}\).
The initial tensors, future normal and actual spin identification agree.
\hyperref[thm:online-marked-collar-full-hom]{Theorem 2.10}, using local
geometric Cauchy uniqueness, gives the isometry and actual spin
comparison on a common smaller domain of dependence. The physical graph
identities of Corollary 5.3 then restore \(g_\lambda\), including its
metric scale. Proposition 5.6 restores the scale of the full generated
source and its wave completion.

For families, every initial operation is made on one smooth base
selection. \hyperref[thm:online-common-collar-finite-families]{Theorem
2.14} gives the smooth Cauchy construction on a common buffered time
interval near the chosen parameter value; actual characteristic flow and
initialized transport have common buffered domains after the same local
restriction. The next proposition states the finite-order
interpretation. \hfill\(\square\)

\subsection{5.6. Estimates on one selected
family}\label{estimates-on-one-selected-family}\label{sec:submission-selected-family-estimates}

\pandocstatementheading{\leavevmode\phantomsection\label{prop:submission-coverage-family-estimates}Proposition 5.8 (Finite estimates for one selected family).}

Let a smooth finite-dimensional family of pointed Ricci-flat germs, or
of pointed maximal initial germs, be supplied on common local
representatives with smooth markings. The selections in Theorems 5.5 and
5.7, with the scale restoration of Proposition 5.6, can be made on one
parameter neighborhood with fixed positive radius, extension branch,
initial patch and buffered spacetime and angular parents.

Fix compact parameter sets
\(\mathcal K\subset\operatorname{int}\mathcal K'\) in the selected
parameter neighborhood, finitely many source charts, a finite
complete-coefficient depth \(J\), a point derivative order \(m\) and a
parameter derivative order \(q\). The selected ordinary pair, executed
source, physical embedding and spin comparison satisfy a bound of the
form

\phantomsection\label{eq:submission-coverage-family-bound} \[
\|\mathcal Y\|_{m,q,J}
+
\|s(c)\|_{m,q}
+
\|\iota\|_{m,q}
\le
F_{m,q,J}
\left(
\mathcal N_N(\text{input and markings}),
r^{-1},
\mathfrak m^{-1}
\right).
\] Here \(\mathcal Y\) is the stipulated finite list of generated source
and wave coefficients and selected carrier and spin-identification maps.
Output seminorms are taken over \(\mathcal K\); \(\mathcal N_N\)
includes the required finite spatial and parameter seminorms of the
input and markings over \(\mathcal K'\). The index \(N\) is finite,
\(F_{m,q,J}\) is finite and locally bounded, and \(\mathfrak m\) records
the positive coordinate, spacelike, elliptic, optical, inverse and
domain-buffer margins. The seed seminorm is a finite weighted seminorm
of the whole decoded seed. The selected family gives a full source with
its physical embedding and spin identification. We choose the parent
before specifying \(m,q,J\) and estimate the indicated finite list on
that parent.

\emph{Proof.} The normal-coordinate and augmented Dirichlet
constructions give the graph and initial data on fixed buffered balls,
or one uses the direct initial normalization of Theorem 5.5. Choose the
finite base order before the radius. The cutoff prefields then lie in
one relative-extension domain.

The extension theorem concerns the same base output at every spatial
order. Its base smooth parameter dependence and uniform higher spatial
bounds imply the stipulated mixed estimates by
\hyperref[lem:online-parameter-promotion]{Lemma B.13}. The gauge, polar
rotation, decoder and ordinary-class transform have finite estimates on
that same domain. Include the finite norms of point, tangent-plane and
frame markings in \(\mathcal N_N\).

The source realization has the finite coefficient and parameter
estimates of
\hyperref[thm:online-ordinary-negative-response-initial-source]{Theorem
4.18}. Its Cauchy completion and initialized actual spin transport obey
\hyperref[thm:online-common-collar-finite-families]{Theorem 2.14}. The
time interval is fixed by lower Sobolev norms and positive hyperbolic
and spacelike margins.
\hyperref[lem:online-reduced-cauchy-finite-family-estimates]{Lemma 2.13}
proves the higher and actual parameter-derivative estimates for that
same solution, allowing the stated finite spatial derivative slack.

Null flow, transverse hit maps, Jacobi equations and finite-rank
transport are finite smooth systems on compact buffered domains. At a
fixed complete-symbol or module depth, the transport equations involve
finitely many earlier coefficients. The wave kernel, mixed carrier
transports and all-direction point connection have the
finite-coefficient estimates of
\hyperref[prop:online-wave-kernel-finite-families]{Proposition 6.11},
\hyperref[prop:online-wave-carrier-families-scale]{Proposition 6.17} and
\hyperref[prop:online-point-connection-finite-families]{Proposition
6.22}. Finite angular refinement and those initialized transport
estimates therefore give each stipulated coefficient bound on the same
parents. Compactness of the real sky supplies one finite angular
refinement; increasing the depth does not require changing the parent.

Composition of these estimates yields the displayed bound. Its order and
required input seminorms may depend on \(m,q,J\). A smoothly varying
positive numerical scale is also allowed when its finite derivatives and
positive lower bound are included. \hfill\(\square\)

\section{6. The covariant wave
enrichment}\label{the-covariant-wave-enrichment}\label{sec:online-covariant-wave-enrichment}

The Cartan collar determines the scalar wave kernel class and hence a
complete distribution module on every retained spacelike carrier. We
construct these objects from the recognized source metric, compare their
full sections across carriers and define differentiation in every point
direction. The construction applies to the collars generated from
ordinary pairs in Section 5, without supplying a further Einstein
solution.

To see the effect of changing carriers, take two spacelike hyperplanes
through the origin in Minkowski space whose future unit normals have
relative rapidity \(\theta\ne0\). Compare the two descriptions of the
same null covector at the origin under the identity map of spacetime. If
\(\rho_i=|\xi_i|\) is its spatial frequency magnitude on hyperplane
\(i\), then in boost-adapted coordinates
\(J=\rho_1/\rho_0=\cosh\theta-\sinh\theta\,(\xi_0)_1/\rho_0\). Fourier
half-density transport multiplies the amplitude by \(J^{-1/2}\). It
carries the degree-zero metric amplitude to
\(J^{-1/2}\rho_0^{-3/2}=J\rho_1^{-3/2}\). The factor \(J\) varies with
the direction at this single point, so a point scalar cannot remove it.
For the degree-one wave evaluation the same transport gives
\(J^{-1/2}\rho_0^{-1/2}=\rho_1^{-1/2}\), with the common normalization
constant suppressed. Thus the wave evaluation compares the two slices
over the physical identity.
\hyperref[thm:online-canonical-wave-kernel]{Theorem 6.10} establishes
uniqueness of the kernel class modulo joint smoothing from its own
admission conditions. The relative line in §9 retains its degree-zero
calibration.

We first specify the complete-symbol calculus, its internal adjoint and
the shifted module conventions. A normal Riccati equation then
constructs the positive wave branch; its Klein--Gordon normalization
gives a kernel characterized by independent wavefront, bisolution and
commutator tests. Mixed traces of that single kernel yield the carrier
cocycle and the full degree-one evaluation section. Actual
differentiation of the moving distribution families supplies the point
connection.

All complete classes have actual classical representatives. The quotient
removes kernels jointly smooth in the spatial, point and supplied
parameter variables. Each finite estimate controls finitely many
homogeneous coefficients and their specified derivatives on the same
selected geometric parent. The final recognition theorem concerns this
normalized scalar enrichment over an admitted finite Cartan collar. Its
comparisons under supplied physical maps leave the source's actual spin
lifts distinct.

\subsection{6.1. Actual realizations and complete coefficient
classes}\label{actual-realizations-and-complete-coefficient-classes}\label{sec:online-complete-wave-calculus}

The wave enrichment is constructed on the actual source collar of
Section 2. Its metric and physical scale are therefore source readouts.
Spatial operators act on complexified real half-densities. The
underlying representatives are smooth classical pseudodifferential
operators and Fourier integral operators, and the resulting source
objects retain their complete classes modulo smoothing. We specify this
quotient, its parameter dependence and its algebra and module
conventions before using it.

\pandocstatementheading{\leavevmode\phantomsection\label{def:online-complete-symbol-calculus}Definition 6.1 (The retained complete calculus).}

On an oriented spatial coordinate patch \(V\subset\mathbb R^3\), a
classical left symbol of upper order \(m\) has an expansion

\phantomsection\label{eq:online-classical-symbol-convention} \[
a(x,\xi;\lambda)\sim\sum_{j\le m}a_j(x,\xi;\lambda),
\qquad a_j(x,r\xi;\lambda)=r^j a_j(x,\xi;\lambda),
\qquad r>0.
\] The source algebra uses integer orders. Auxiliary positive square
roots may have half-integer orders, with the positive physical principal
branch specified each time. A fixed cutoff removes zero frequency.
Symbols, their asymptotic remainders and all prescribed spatial, angular
and finite-dimensional parameter derivatives have locally uniform symbol
bounds on larger retained compact parents. Operators are properly
supported on the retained domain, with coordinate kernel convention

\phantomsection\label{eq:online-physical-operator-kernel} \[
K_A(x,y)=(2\pi)^{-3}
\int e^{i(x-y)\cdot\xi}a(x,\xi)\,\mathrm d\xi
\,|\mathrm dx\,\mathrm dy|^{1/2}.
\] Two such operator families have the same complete class if their
difference is smoothing jointly in the spatial and supplied parameter
variables. A complete symbol is the sequence of all homogeneous
coefficients in a local presentation. Its finite seminorms concern
finitely many of those coefficients and their derivatives on a unit
angular section; they do not include every seminorm of a chosen
smoothing remainder.

For a conic Lagrangian, use actual local oscillatory distributions with
real nondegenerate homogeneous phases and classical amplitudes. Phases
and amplitudes depend smoothly on the supplied parameters on fixed
buffered charts. Equivalent phase descriptions represent the same
distribution class modulo a jointly smooth family. General Lagrangian
distributions are allowed even where their projection to the spatial
base is singular. A hypersurface-conormal graph description is used only
in a chart where that description is valid.

We use the classical conormal and pseudodifferential calculus in
\hyperref[ref-Melrose2006]{Melrose} (2006, Chapters 2--4 and 7), and the
Fourier integral calculus in
\hyperref[ref-Duistermaat1972FIO]{Duistermaat} (1972, §1, Theorem 1.1).
We use the coordinate invariance of oscillatory classes, restriction
under the wavefront transversality condition, and the symbol formulas
for composition and adjoints. For properly supported transverse
canonical graphs, composition adds the orders and multiplies the
principal symbols. The identity relation gives the pseudodifferential
calculus. An elliptic graph operator has an inverse in this local
quotient, and acts on the Lagrangian module carried by its graph. The
parameter and finite-coefficient consequences used below are explained
in the finite-calculus lemma.

All microlocal identities are asserted on open retained cones with a
larger conic buffer. Finite cones cover the whole positive-ray sphere
when a full-sky assertion is made. The full kernel characterization will
impose its wavefront and smoothing tests on the entire retained point
neighborhood squared. A statement on one angular cap alone does not
imply that full assertion.

\pandocstatementheading{\leavevmode\phantomsection\label{prop:online-formal-symbol-algebra}Proposition 6.2 (The complete spatial algebra and its internal adjoint).}

The integer-order complete operators form a filtered sheaf of algebras
on \(P_+T^*\Sigma\), where \(P_+\) means quotient by positive real
scaling. In a cotangent chart, rephase each physical coefficient by

\phantomsection\label{eq:online-symbol-degree-rephasing} \[
p_j=i^{-j}a_j.
\] The resulting composition and complex-linear internal adjoint are

\phantomsection\label{eq:online-formal-product-and-adjoint} \[
P\circ Q=
\sum_\alpha\frac1{\alpha!}
(\partial_\xi^\alpha P)(\partial_x^\alpha Q),
\qquad
P^*=
\sum_{j,\alpha}\frac{(-1)^j}{\alpha!}
\partial_\xi^\alpha\partial_x^\alpha p_j.
\] Each output degree is a finite sum. The internal adjoint is an
involutive anti-automorphism covering the same positive ray, and the
actual half-density coordinate comparisons preserve it and satisfy
strict cocycles.

We call an integer-order operator formally real when every \(p_j\) is
real. On that real subalgebra the physical Hermitian adjoint \(\dagger\)
agrees with \(*\). Physical bilinear transpose, denoted \(T\), is a
different operation: it reverses the cotangent covector before the
lower-order corrections. These statements do not identify formal reality
with reality of the physical Schwartz kernel.

\emph{Proof.} The physical left-symbol product is \[
a\mathbin{\#}b=
\sum_\alpha\frac{i^{-|\alpha|}}{\alpha!}
(\partial_\xi^\alpha a)(\partial_x^\alpha b).
\] A term from degrees \(j,k\) and \(|\alpha|\) derivatives has output
degree \(j+k-|\alpha|\). Substituting \(a_j=i^jp_j\) and \(b_k=i^kq_k\)
cancels the degree rephasing exactly and gives the stated product. In
particular, the actual differential operator \(\partial_{x^i}\) has
formal symbol \(\xi_i\).

The physical Hermitian-adjoint symbol is
\(\sum_\alpha i^{-|\alpha|}\partial_\xi^\alpha\partial_x^\alpha\overline a/\alpha!\).
Since \(\overline{i^jp_j}=(-1)^ji^j\overline{p_j}\), its rephased form
is the displayed expression for \(P^*\) with \(p_j\) replaced by
\(\overline{p_j}\). This proves the stated equality on formally real
operators, while retaining the different linearities of the two
operations.

To verify the internal involution and coordinate law without a reality
assumption, work on a cone with one nonzero signed covector component,
say \(\xi_3\). Every finite jet of a homogeneous coefficient of integer
degree \(j\) at a specified nonzero covector can be matched by \[
\xi_3^j p(x,\xi_1/\xi_3,\xi_2/\xi_3),
\] with \(p\) polynomial in the angular variables and smooth coefficient
functions in \(x\). Such homogeneous Laurent symbols are finite
expressions in multiplication, \(\partial_{x^i}\) and the microlocal
inverse of \(\partial_{x^3}\). On these expressions the internal adjoint
is product reversal with \(f^*=f\),
\((\partial_{x^i})^*=-\partial_{x^i}\), extended to inverses. It is
therefore an involutive anti-automorphism. The bilinear differential
transpose on half-densities commutes with real coordinate changes, and
so does its extension to these microlocal inverses.

At any specified output degree, every composition, coordinate
transformation and adjoint coefficient uses finitely many jets at the
corresponding cotangent point. The preceding Laurent test proves the
involution, anti-multiplication and covariance identities for arbitrary
smooth homogeneous coefficients at that degree. It applies at all
degrees, establishing the complete identities. This argument uses
physical transpose only on the differential/Laurent test algebra; it
does not replace the antipodal argument of an arbitrary angular symbol
by degree parity.

Finally, a coordinate diffeomorphism \(f\) acts by actual half-density
pushforward \(U_f\), and the operator comparison is
\(A\mapsto U_fAU_f^{-1}\). Actual pushforward satisfies
\(U_{f\circ g}=U_fU_g\), so these comparisons have strict cocycles
before and after the quotient. Local asymptotic summation realizes each
complete sequence and changes none of these quotient identities.
\hfill\(\square\)

\subsection{6.2. The whole-sky module
convention}\label{the-whole-sky-module-convention}\label{sec:online-shifted-lagrangian-modules}

\pandocstatementheading{\leavevmode\phantomsection\label{prop:online-shifted-sky-modules}Proposition 6.3 (Shifted sky modules and their local normalization).}

At \(b\in\Sigma\), use conormal half-densities with local unitary
Fourier expression

\phantomsection\label{eq:online-shifted-point-conormal} \[
(2\pi)^{-3/2}\int
 e^{i(x-b)\cdot\xi}
 \sum_{j\le m}i^j a_j(b,\xi)\,\mathrm d\xi
 \,|\mathrm dx|^{1/2},
\qquad
a_j(b,r\xi)=r^{j-3/2}a_j(b,\xi).
\] The module filtration index is \(j\). Complete classes, with smooth
parameter dependence and joint smoothing as above, form modules over the
complete algebra on the entire sky \(P_+T_b^*\Sigma\). Actual coordinate
comparisons and specified invertible Fourier integral graph comparisons
carry their full classes and satisfy the same composition laws as those
operators.

In a normalized angular contact chart with sky \(z=s(q)\), the physical
Fourier amplitude in the positive radial variable \(\tau\) has the form

\phantomsection\label{eq:online-shifted-contact-amplitude} \[
e^{-i\tau s(q)}\sum_{j\le m}i^j\tau^{j-1/2}a_j(q).
\] Writing its formal class as \(\sum_{j\le m}\tau^j a_j(q)u_s\), the
exact generator actions are

\phantomsection\label{eq:online-shifted-module-generator-actions} \[
D_z=\tau,\qquad
D_{q^a}=\partial_{q^a}-s_{q^a}\tau,
\qquad
z=-\partial_\tau+\frac1{2\tau}+s(q).
\] The generator \(u_s\) has degree zero. Any section with a
nowhere-zero highest coefficient is cyclic microlocally: every complete
module section is obtained by a complete operator acting on it. This
assertion holds for smooth families with smooth local choices of the
expressing operators. It applies in particular to a degree-one elliptic
evaluation section. Specializing the central smooth parameter functions
at a fixed parameter gives the corresponding fixed-parameter module.

\emph{Proof.} The actual pseudodifferential action preserves point
conormality. Taylor expansion at \(b\) followed by integration by parts
gives its rephased amplitude action

\phantomsection\label{eq:online-point-conormal-module-action} \[
P\cdot a=
\sum_{j,l,\alpha}\frac{(-1)^{|\alpha|}}{\alpha!}
\partial_\xi^\alpha
\bigl((\partial_x^\alpha p_j)(b,\xi)a_l(b,\xi)\bigr).
\] Every output degree uses finitely many coefficients. In this formula
\(\partial_{x^i}\) acts as \(\xi_i\) and \(x^i-b^i\) as
\(-\partial_{\xi_i}\). The principal amplitude times the fiber
half-density has homogeneity \(j\), so the convention is coordinate
invariant. Explicitly, under \(y=Ax\) the leading scalar amplitude
becomes \(|\det A|^{1/2}a(A^T\eta)\), the required fiber half-density
transformation. The general coordinate statement follows from the
conormal calculus with its successively lower integer terms.

Here is the angular normalization, including the shift. On \(\xi_3>0\)
set

\phantomsection\label{eq:online-angular-legendre-coordinates} \[
q^a=\xi_a/\xi_3,\qquad
z=x^3+q^ax^a,\qquad
\vartheta_a=-x^a,\qquad
\tau=\xi_3,
\qquad a=1,2.
\] The cotangent one-form is
\(\tau(\mathrm dz+\vartheta_a\mathrm dq^a)\). Define the local Radon
graph operator

\phantomsection\label{eq:online-normalized-angular-radon} \[
Rf(q,z)=\int_{\mathbb R^2}f(x',z-q\cdot x')\,\mathrm dx',
\qquad
\mathcal U=(2\pi)^{-1}(-i\partial_z)R.
\] With unitary Fourier transforms,
\(\widehat{\mathcal U f}(q,\tau)=\tau\widehat f(\tau q,\tau)\). The
Jacobian is \(\mathrm d\xi=\tau^2\mathrm dq\mathrm d\tau\), so this is
an invertible order-zero canonical-graph operator on the retained
positive cone. Its inverse is the inverse Fourier-coordinate change.
Applying it to the point-conormal expression gives the asserted power
\(\tau^{j-1/2}\) and phase with \(s(q)=b^3+q\cdot b'\). The same
normalized conormal form for a general graph \(z=s(q)\) follows by the
actual graph translation.

The physical actions are differentiation in \(z,q\) and multiplication
by \(z\). On the displayed Fourier expression, differentiation of the
phase gives the \(s_{q^a}\) terms. Physical multiplication by \(z\) is
\(i\partial_\tau\); its action on \(\tau^{j-1/2}\) and the degree
rephasing give exactly \(-\partial_\tau+1/(2\tau)+s\). Thus the one-half
term is fixed by the half-density convention.

For completeness, this angular comparison also respects the internal
adjoint. Its exact rephased images on coordinate generators are

\phantomsection\label{eq:online-angular-exact-generator-map} \[
\begin{aligned}
\partial_{x^a}&\longmapsto q^aD_z,&
\partial_{x^3}&\longmapsto D_z,\\
x^a&\longmapsto-D_z^{-1}D_{q^a},&
x^3&\longmapsto z+D_z^{-1}(E_q+1),
\qquad E_q=\sum_aq^aD_{q^a}.
\end{aligned}
\] They follow by integration by parts in \(R\). The derivative images
are internally skew and the multiplication images are internally
self-adjoint. For the last image use \((E_q+1)^*=-E_q-1\) and
\((D_z^{-1})^*=-D_z^{-1}\). The finite-jet Laurent argument of the
preceding proposition extends these identities to all complete smooth
angular symbols. No arbitrary contact quantization is asserted by this
calculation.

Finitely many signed component cones cover the whole sphere. Actual
coordinate maps and the specified Radon operators define the comparisons
on overlaps. Their compositions cancel to give exact cocycles in the
quotient. A supplied elliptic graph operator carries the same
distribution module to the image Lagrangian; an inverse carries it back.
This defines the corresponding shifted modules even at a singular
projection of a Lagrangian. Local conormal graph presentations are
normalized by those actual comparisons, with their half-density and
phase factors retained.

To prove cyclicity, let \(e\) have upper degree \(k\) and nonzero
leading amplitude on a retained chart. For a proposed section \(m\) of
degree at most \(l\), choose an operator of degree \(l-k\) whose leading
symbol on the Lagrangian is the ratio of their leading amplitudes.
Extend that symbol smoothly to a neighboring cotangent cone. Subtract
its action on \(e\), and repeat one degree lower. Each step divides only
by the same nonzero leading amplitude. Asymptotic summation gives a
complete operator \(P\) with \(m=Pe\). The extensions can be chosen in
one fixed local normal chart, so all coefficients depend smoothly on any
supplied parameters. This proves existence of the expressing operator;
it need not be unique away from the module. The same argument gives the
assertion after an invertible graph comparison.

For central specialization, use local parameter coordinates \(\lambda\)
near \(\lambda_0\) and a normalized phase family. Restriction evaluates
each complete amplitude coefficient at \(\lambda_0\) and is surjective
by extending that coefficient smoothly in the parameter. If every
restricted coefficient vanishes, Hadamard's integral formula writes each
one as

\phantomsection\label{eq:online-module-central-specialization} \[
a_j(\lambda,q)=\sum_\mu(\lambda^\mu-\lambda_0^\mu)
\int_0^1\partial_{\lambda^\mu}a_j
\bigl(\lambda_0+t(\lambda-\lambda_0),q\bigr)\,\mathrm dt.
\] These coefficient sequences define complete smooth module families of
the same upper degree. Thus the kernel is exactly the ideal of smooth
parameter functions vanishing at \(\lambda_0\) times the family module.
The coefficient equality is in the joint quotient and is unchanged by
phase comparisons. This proves the specialization assertion.
\hfill\(\square\)

\subsection{6.3. Finite coefficient estimates and joint
smoothing}\label{finite-coefficient-estimates-and-joint-smoothing}\label{sec:online-finite-calculus-estimates}

\pandocstatementheading{\leavevmode\phantomsection\label{lem:online-finite-symbol-family-calculus}Lemma 6.4 (Finite calculus on one retained parent).}

On fixed buffered coordinate, phase and angular parents, complete
composition, elliptic inversion with a retained principal inverse
margin, chosen simple scalar roots, canonical-graph composition and the
preceding module operations preserve smooth finite-dimensional
coefficient families. Every specified finite output depth and derivative
order is controlled by finitely many coefficient and phase derivatives
and the inverses of the retained nondegeneracy margins.

All coefficient sequences have actual smooth classical realizations on
those same parents. Different realizations of the same complete family
differ by a jointly smoothing operator or a jointly smooth distribution.
The estimates control any prescribed finite list of homogeneous
coefficients and their indicated derivatives on the retained parent. A
choice of asymptotic sum also fixes a smooth remainder whose seminorms
need not be controlled by a finite jet.

\emph{Proof.} The composition formula uses finitely many terms at a
specified output degree. To invert an elliptic symbol, choose its
inverse principal symbol, cancel the next error by multiplication with
that inverse, and proceed by decreasing degree. A simple chosen scalar
root is obtained similarly by dividing at each stage by the nonzero
derivative of its defining polynomial at the chosen principal root. Thus
a positive square root uses the denominator twice its chosen positive
principal coefficient. On retained compact angular sets all these
denominators have fixed lower bounds. Differentiating each finite
recursion proves the claimed finite estimates, with larger but finite
input orders.

For transverse canonical-graph composition, the stationary-phase
expansion at a given degree uses finitely many amplitude and phase
derivatives and inverse powers of the nondegenerate transverse Hessian.
The proper support and graph hypotheses are precisely those of the cited
composition theorem. The retained Hessian margin controls those
inverses. An elliptic operator on an inverse graph can first be chosen
with inverse principal symbol; its product is an elliptic
pseudodifferential operator with principal symbol one. The preceding
complete inverse corrects it to an actual quotient inverse. Left and
right inverses agree by associativity. This also proves the local
algebra-conjugation and module-action assertions used above, with the
same finite parameter estimates.

To realize a sequence, take homogeneous extensions of its coefficients
and multiply its successive terms by cutoffs \(\chi(\xi/R_j)\), with
\(R_j\to\infty\). Choose \(R_j\) successively so that the \(j\)th term
contributes less than \(2^{-j}\) to the first \(j\) relevant symbol
seminorms after lowering the tested order by a fixed fraction of \(j\).
Include the first \(j\) parameter derivatives on the first \(j\) compact
sets in an exhaustion. For every fixed derivative and remainder order,
all sufficiently late terms then converge in that seminorm. The
resulting sum is smooth in the parameters and has the prescribed
expansion. This is the parameter version of classical asymptotic
summation (\hyperref[ref-Melrose2006]{Melrose 2006}, Chapter 2,
Proposition 5); the cutoff choices do not change the spatial or conic
parent.

If all homogeneous coefficients vanish, the symbol and every parameter
derivative are rapidly decreasing. Its localized Fourier integral is
therefore jointly smooth, by differentiating the integral at arbitrary
finite orders. Conversely, the local complete coefficient of a smooth
distribution is zero. The same statements hold in an equivalent
nondegenerate phase by the invariant symbol calculus. They prove that
changing the asymptotic sum or its phase description leaves exactly the
stated joint quotient. They also show that a parameter derivative of a
joint smoothing remainder is again joint smoothing.

The phase of a flow graph is fixed by its actual Hamiltonian flow on the
retained tube. Its finite derivatives satisfy finite variational ODEs;
amplitude transport at each degree is a linear ODE with forcing from
finitely many earlier degrees. With a uniform flow and phase
nondegeneracy margin, the chain rule and Gronwall estimates give the
corresponding finite bounds on the same tube. No later coefficient
changes the flow domain. A finite angular refinement handles all
retained directions. These observations establish the last needed
parameter calculus statements, without a uniform estimate as the
requested depth tends to infinity. \hfill\(\square\)

In particular, a pointwise choice of smooth remainder at every parameter
value is not the equivalence relation used here. Kernel admission and
comparison will use the joint calculus just defined. The actual finite
Cartan and spin fields remain separate geometric structures; quotienting
scalar wave representatives by joint smoothing does not quotient their
physical scale or actual spin comparisons.

\subsection{6.4. The wave operator and its positive normal
branch}\label{the-wave-operator-and-its-positive-normal-branch}\label{sec:online-wave-factorization}

The metric in this construction is the metric recognized from the actual
Cartan source in \hyperref[thm:online-cartan-collar-recognition]{Theorem
2.9}. For an ordinary-pair display it has already been constructed from
the pair. The wave construction introduces no additional Einstein datum.
Its arguments apply to any smooth Lorentzian collar of the stated
signature.

Retain a Gaussian coordinate neighborhood \(g=\mathrm dt^2-h_t\), a
smaller point neighborhood and buffered initial coordinate disks. The
neighborhoods are small enough that the retained null trajectories stay
in the larger parent. A finite angular cover retains all nonzero null
directions over the point neighborhood. When invoking a scalar Cauchy or
Green operator, restrict to a smaller causally convex globally
hyperbolic neighborhood inside this parent. An arbitrary product
coordinate cylinder is not itself assumed to have that property.

We use the actual classical operator and joint-smoothing calculus of
\hyperref[def:online-complete-symbol-calculus]{the symbol conventions}.
All equalities of complete symbols and parametrices below are in that
quotient. Physical Hermitian adjoint is written \(\dagger\), physical
bilinear transpose is written \(T\), and the complex-linear internal
adjoint of the integer-order formal algebra is written \(*\). These
operations have different meanings. In particular, the frequency branch
is fixed by the flat solution \(\exp(-it|D|)\); its temporal covector is
\(-H\) in the signature \((+---)\).

\pandocstatementheading{\leavevmode\phantomsection\label{prop:online-wave-half-density-reduction}Proposition 6.5 (Half-density wave reduction and the instantaneous
defect).}

Put

\phantomsection\label{eq:online-wave-density-and-spatial-operator} \[
\begin{gathered}
\rho_t=(\det h_t)^{1/4}, \\
K_t=\tfrac12\dot h_t, \\
\tau_t=\operatorname{tr}_{h_t}K_t, \\
D_{h_t}=-\rho_t\Delta_{h_t}\rho_t^{-1}, \\
H_t(x,\xi)=|\xi|_{h_t}.
\end{gathered}
\] There is a unique complete classical square root \(A_{h_t}\) of
\(D_{h_t}\) with physical principal symbol \(H_t\). It is physically
Hermitian. The operator \(B_{h_t}=iA_{h_t}\) has real integer-degree
formal coefficients and is internally skew-adjoint. In physical left
symbols,

\phantomsection\label{eq:online-wave-root-subprincipal} \[
\sigma_0(A_{h_t}-\operatorname{Op}(H_t))
=\frac1{2i}\sum_j\partial_{x^j}\partial_{\xi_j}H_t.
\] Writing \(V_t=\dot\tau_t/2+\tau_t^2/4\), one has

\phantomsection\label{eq:online-wave-half-density-reduction} \[
\rho_t\Box_g\rho_t^{-1}
=\partial_t^2+D_{h_t}-V_t.
\] An instantaneous root parametrix satisfying
\((\partial_t+iA_{h_t})G_{t,s}=0\) consequently has wave residual

\phantomsection\label{eq:online-wave-instantaneous-defect} \[
(\rho\Box_g\rho^{-1})G_{t,s}
=(-i\dot A_{h_t}-V_t)G_{t,s}.
\] Its order-one coefficient is \(-i\dot H_t=iK_t^{ij}\xi_i\xi_j/H_t\).

\emph{Proof.} The positive principal root determines every lower
coefficient of \(A_h\) by division by \(2H\). At the step determining
degree \(1-j\), this coefficient occurs in the square only through
\(2H a_{1-j}\) at degree \(2-j\); differentiated products involve
already determined coefficients. The complete recursion is therefore
unique on the nonzero cone. The parameter-dependent asymptotic
realization gives an actual classical representative. Since
\(D_h^\dagger=D_h\), taking the physical adjoint gives another square
root with the same principal symbol. Uniqueness proves
\(A_h^\dagger=A_h\).

For \(B_h\), the square equation is \(B_h^2=-D_h\). Its formal principal
symbol is \(H\), and the differential operator \(-D_h\) has real formal
coefficients. The same recursion in the rephased product therefore gives
real coefficients at every integer degree. Physical Hermitian adjoint
and internal adjoint agree on such operators. The equality
\((iA_h)^\dagger=-iA_h\) proves \(B_h^*=-B_h\).

In coordinate half-densities the degree-one physical symbol of \(D_h\)
is \((1/i)(\partial_{x^i}h^{ij})\xi_j\). It equals
\((2i)^{-1}\sum_j\partial_{x^j}\partial_{\xi_j}(H^2)\). The degree-one
coefficient of the square of \(H+a_0+\cdots\) is \[
2Ha_0+\frac1i\sum_j
(\partial_{\xi_j}H)(\partial_{x^j}H).
\] Comparison gives \hyperref[eq:online-wave-root-subprincipal]{the
stated coefficient}. Equivalently, the physical subprincipal symbol of
\(A_h\) on half-densities is zero. The root is taken in the local
complete-symbol calculus.

The scalar divergence formula gives
\(\Box_g=\partial_t^2+\tau\partial_t-\Delta_h\), and
\(\dot\rho/\rho=\tau/2\). Direct conjugation gives the displayed
reduction. Twice differentiating the instantaneous evolution gives its
residual. Finally \(\partial_t h^{ij}=-2K^{ij}\) implies
\(\dot H=-K^{ij}\xi_i\xi_j/H\), proving the last statement.
\hfill\(\square\)

For a scalar wave solution let \(v=\rho u\) and \(\pi=\rho\partial_tu\).
The two field variables are half-density coefficients on the same
coordinate disk. Their equations are

\phantomsection\label{eq:online-wave-canonical-cauchy-system} \[
\dot v=\frac{\tau}{2}v+\pi,
\qquad
\dot\pi=-D_hv-\frac{\tau}{2}\pi.
\] The normal branch below corrects the instantaneous defect while
retaining the same real null characteristics.

\pandocstatementheading{\leavevmode\phantomsection\label{thm:online-wave-normal-riccati}Theorem 6.6 (The complete positive normal symbol).}

There is a unique complete classical symbol \(Q_t\) of order one with
physical principal symbol \(iH_t\) whose graph \(\pi=-Q_tv\) is
invariant under the scalar wave equation modulo joint smoothing. It
obeys

\phantomsection\label{eq:online-wave-normal-riccati} \[
Q_t^2-\dot Q_t-Q_t\frac{\tau_t}{2}
-\frac{\tau_t}{2}Q_t+D_{h_t}=0.
\] Its integer-degree formal coefficients are real, and \(Q_t^T=Q_t\).
Put \(B_t=Q_t-\tau_t/2\). There are initialized order-zero Fourier
integral parametrices \(E_{t,s}\) with

\phantomsection\label{eq:online-wave-positive-propagator} \[
(\partial_t+B_t)E_{t,s}=0,\qquad E_{s,s}=I,
\qquad E_{t,r}E_{r,s}=E_{t,s}.
\] They solve the wave equation modulo joint smoothing and have the null
canonical graph selected by \(\exp(-itH)\). In Gaussian coordinates,

\phantomsection\label{eq:online-wave-corrected-generator} \[
B_t=B_{h_t}
+\operatorname{Op}\!\left(\frac{\dot H_t}{2H_t}\right)
\pmod{\Psi^{-1}}.
\] \emph{Proof.} Substitute \(\pi=-Qv\) into
\hyperref[eq:online-wave-canonical-cauchy-system]{the canonical system}.
The equation for \(\dot\pi\) gives
\(Q^2-\dot Q-Q\tau/2-\tau Q/2+D_h=0\), with the indicated order of the
noncommuting factors.

In the rephased calculus the highest equation is \(q_1^2-H^2=0\). The
prescribed branch is \(q_1=H\). Suppose the coefficients above degree
\(1-j\) have been determined. At degree \(2-j\) the new coefficient
occurs only as \(2Hq_{1-j}\). Time differentiation does not raise symbol
order, the terms containing \(\tau\) have order one, and every
differentiated symbol product lowers total order. All their
contributions at the current degree are therefore already known.
Division by \(2H\) determines \(q_{1-j}\). The argument also removes
successively every degree of the difference of two possible solutions.
The real coefficients of the equation in the formal product show that
every \(q_{1-j}\) is real.

Physical transpose reverses products, commutes with time differentiation
and interchanges the two \(\tau\) terms. Also \(D_h^T=D_h\). Its action
on the principal coefficient gives \(iH(x,-\xi)=iH(x,\xi)\).
Consequently \(Q^T\) satisfies the same branch equation. On the retained
full angular cover, uniqueness proves \(Q^T=Q\).

The complete recursion is realized by a classical family using
\hyperref[lem:online-finite-symbol-family-calculus]{the finite-symbol
family calculus}. For the actual parametrix, in a noncaustic chart solve

\phantomsection\label{eq:online-wave-positive-phase} \[
\partial_tS+H_t(x,\partial_xS)=0,
\qquad S(s,x,\eta)=x\cdot\eta,
\] by the actual Hamiltonian flow. The initialized mixed projection is
invertible for a common sufficiently short interval on each compact
unit-direction chart. A finite cover supplies one interval for all
retained directions. Seek a kernel

\phantomsection\label{eq:online-wave-positive-fio-kernel} \[
E_{t,s}(x,y)
=(2\pi)^{-3}\int
e^{i(S(t,x,\eta)-y\cdot\eta)}
w(t,s,x,\eta)\,\mathrm d\eta\,
|\mathrm dx\,\mathrm dy|^{1/2}.
\] Composition with \(\partial_t+B_t\) first gives the Hamilton--Jacobi
equation. At the next order it gives a linear transport equation for the
leading amplitude along the same characteristics, with initial value
one. Every successive homogeneous amplitude satisfies that same
transport operator with a forcing involving finitely many derivatives of
earlier amplitudes and already determined coefficients of \(B_t\); its
initial value is zero. Integration along the retained flow determines
each coefficient as an actual smooth function. The leading solution is
nonzero because it is an initialized exponential transport. Uniform
asymptotic summation produces a classical amplitude, with a jointly
smoothing residual and identity initial value modulo joint smoothing.

These are the local geometric-optics parametrices in the shared Fourier
integral calculus. The wave factorization is also obtained in
\hyperref[ref-gerard-wrochna2014]{Gérard and Wrochna} (2014, Theorem
6.3, Remark 6.5 and Appendix A.3); the negative time phase fixes which
of its two branches is used here. The required localization is explicit.
Extend \(h_t\) from a larger buffered coordinate patch to a positive
metric on \(\mathbb R^3\) by convex interpolation with a fixed Euclidean
metric, and extend it smoothly in time outside the retained interval.
The coefficients of \(D_h-V\) then have bounded spatial derivatives and
a uniformly positive principal spatial matrix on each retained time
interval. Thus they satisfy the model conditions in that paper's §5.1.
No vacuum equation is required of this auxiliary extension. Proposition
4.6 and Corollary 6.6 there give the corresponding characteristic
propagation; in the local construction it is the canonical relation of
the displayed phase. Every retained trajectory stays inside the original
metric patch.

The amplitude recursion proves uniqueness of the initialized parametrix
class. Composition on common buffered canonical graphs gives another
parametrix with the same initialization, so it proves the composition
law and the inverse \(E_{s,t}\). Substitution of its equation and the
Riccati identity into the canonical system proves the wave equation. The
smoothing ideal is preserved under these properly supported
compositions.

For the first correction write \(B=ia\). The conjugated wave equation
gives \(a^2+i\dot a=A_h^2-V\). At degree one,
\(2H\sigma_0(a-A_h)+i\dot H=0\). Multiplication by \(i\) gives
\hyperref[eq:online-wave-corrected-generator]{the stated correction}.
Its principal Hamiltonian is unchanged. \hfill\(\square\)

\pandocstatementheading{\leavevmode\phantomsection\label{prop:online-wave-charge-normalization}Proposition 6.7 (Klein--Gordon normalization of the positive branch).}

Define

\phantomsection\label{eq:online-wave-charge-normalization} \[
R_t=i(Q_t^\dagger-Q_t),\qquad
W_t=R_t^{1/2},\qquad
U_{t,s}=W_tE_{t,s}W_s^{-1},
\] where the complete square root has physical principal symbol
\(\sqrt{2H_t}\). The operator \(R_t\) is physically Hermitian and has
real, transpose-symmetric kernel modulo joint smoothing. Its principal
symbol is \(2H_t\). The normalized propagator is physically unitary and
composes in the quotient. Its integer-order generator satisfies

\phantomsection\label{eq:online-wave-normalized-generator} \[
\begin{gathered}
B_t^{\mathrm{cov}}
=W_tB_tW_t^{-1}-\dot W_tW_t^{-1}, \\
(B_t^{\mathrm{cov}})^*=-B_t^{\mathrm{cov}}, \\
B_t^{\mathrm{cov}}=B_{h_t}\pmod{\Psi^{-1}}.
\end{gathered}
\] \emph{Proof.} From \(Q^T=Q\) one obtains \(Q^\dagger=\overline Q\),
where the bar is physical kernel conjugation. Hence
\(R=i(\overline Q-Q)\) is real and transpose symmetric, as well as
Hermitian, with principal symbol \(2H\). Inverse and square-root
recursion give \(W\) and \(W^{-1}\) in the complete calculus. Adjointing
\(W^2=R\) preserves the prescribed principal root, so uniqueness gives
\(W^\dagger=W\). A globally positive exact spectral representative is
not required.

The Klein--Gordon form in canonical half-density variables is

\phantomsection\label{eq:online-wave-klein-gordon-form} \[
i\int_\Sigma
\bigl(\overline v\,\pi-\overline\pi\,v\bigr).
\] On the positive graph it equals \(\langle v,Rv\rangle\). Direct
subtraction of the Riccati equation and its Hermitian adjoint gives

\phantomsection\label{eq:online-wave-charge-conservation} \[
\dot R=B^\dagger R+RB,
\qquad
E_{t,s}^\dagger R_tE_{t,s}=R_s.
\] The second equality follows by differentiating its left-hand side and
using its initial value. It proves \(U_{t,s}^\dagger U_{t,s}=I\) and
physical anti-Hermiticity of \(B^{\mathrm{cov}}\). Composition follows
from the corresponding law for \(E\).

For the internal assertion, formal reality of \(Q\) gives
\(Q^\dagger=Q^*\). Put \(F=Q-Q^*\) in the integer-order formal algebra.
Its complete coefficients are real and its formal principal coefficient
is \(2H\). The formal half-order root \(F^{1/2}\) has real coefficients
by its recursion. With \(i^{1/2}=e^{i\pi/4}\) in the physical degree
conversion, the positive physical root has formal expression

\phantomsection\label{eq:online-wave-charge-half-order-phase} \[
W=e^{-i\pi/4}F^{1/2},\qquad
W^2=-iF=R.
\] The scalar phase cancels from both \(WBW^{-1}\) and
\(\dot W W^{-1}\). Their difference has real coefficients at every
integer degree. Physical Hermitian adjoint and internal adjoint
therefore agree on this generator, giving the asserted internal
skew-adjointness. This reasoning does not extend the internal involution
to arbitrary half-integer degrees.

At degree zero, the correction \(\dot H/(2H)\) in \(B\) cancels the
principal coefficient of \(\dot W W^{-1}\). The possible degree-zero
term in \(WB_hW^{-1}-B_h\) is proportional to the Poisson bracket
\(\{\sqrt{2H},H\}=0\). All other terms have degree at most \(-1\). This
proves the final formula. \hfill\(\square\)

\subsection{6.5. The independent scalar kernel
class}\label{the-independent-scalar-kernel-class}\label{sec:online-wave-kernel-admission}

We first fix the causal-kernel normalization used in the source test.
The point slots of the kernel are scalar; its integral action uses the
physical volume measure. Half-density factors are inserted only when
taking traces on spacelike disks.

\pandocstatementheading{\leavevmode\phantomsection\label{lem:online-local-wave-green}Lemma 6.8 (Local causal kernel and its double Cauchy traces).}

On a sufficiently small protected Cauchy neighborhood there are local
retarded and advanced inverses for \(\Box_g\) on compactly supported
smooth sources. Their difference
\(E_g=E_g^{\mathrm{ret}}-E_g^{\mathrm{adv}}\) is a real antisymmetric
scalar bisolution kernel. With field and future-normal half-density
traces on a common initial disk, its four double traces, in field/normal
order, are

\phantomsection\label{eq:online-causal-kernel-four-traces} \[
\bigl((\gamma_i\otimes\gamma_j)E_g\bigr)_{i,j=0,1}
=\begin{pmatrix}0&-I\\ I&0\end{pmatrix}.
\] Here \(\gamma_0u=\rho u|_\Sigma\) and
\(\gamma_1u=\rho n(u)|_\Sigma\). The row index \(i\) refers to the first
kernel variable and the column index \(j\) to the second. For a kernel
admitting these noncharacteristic Cauchy traces, smooth errors in both
wave equations and all four double traces determine a jointly smooth
error on a smaller product neighborhood. The required restrictions exist
automatically for the characteristic wavefront classes used below. The
same conclusions hold jointly with supplied smooth finite-dimensional
metric parameters.

\emph{Proof.} Use the scalar linear specialization of the hyperbolic
input specified in
\hyperref[thm:online-smooth-local-cauchy-input]{Theorem 2.6}, with the
same protected coordinate extension. Its estimates are the linear wave
energy estimates written in
\hyperref[lem:online-reduced-cauchy-finite-family-estimates]{Lemma
2.13}. One may obtain the smooth scalar evolution directly by spatial
Friedrichs regularization of the first-order wave system: the
differentiated energy inequality is uniform in the regularization,
compactness gives a solution, and the zero-data energy inequality gives
uniqueness. Applying this at every Sobolev order gives the canonical
Cauchy evolution for smooth compactly supported pairs \((v,\pi)\) on one
common short interval. Duality, followed by localization, extends the
solution operators to compactly supported distributions. Their kernels
are distributions by the Schwartz kernel theorem.

For a compactly supported forcing, choose an initial time before its
support and solve the inhomogeneous equation with zero data by Duhamel's
formula. This gives the retarded inverse on the retained domain. Solving
backwards from a time after the support gives the advanced inverse.
Finite propagation makes both independent of the auxiliary endpoint and
extension when restricted to the smaller protected neighborhood. The
local energy estimate gives uniqueness with these support conditions.
Integration by parts for the scalar wave operator against the physical
volume gives \((E^{\mathrm{ret}})^T=E^{\mathrm{adv}}\). The coefficients
are real, so both kernels are real. Their difference is therefore real,
antisymmetric and a bisolution.

For the trace normalization, let \(\mathcal S(t,s)\) be the canonical
evolution matrix for
\hyperref[eq:online-wave-canonical-cauchy-system]{the system},
initialized by the identity. The field part of its response to initial
normal data is its upper-right block. Duhamel's formula identifies the
half-density field-field kernel of \(E_g(t,s)\) with this block. It
vanishes at \(t=s\), and its first normal trace there is the identity.
Antisymmetry gives the opposite field-normal trace. At equal times the
canonical generator has blocks \[
\begin{pmatrix}\tau/2&I\\-D_h&-\tau/2\end{pmatrix}.
\] The upper-right block of its square is zero. Equivalently,
differentiate the initialized field-normal trace, including the
\(\tau/2\) corrections that convert \(\partial_t v\) to \(\pi\). This
gives the zero normal-normal trace. These are the four stated
identities, with the sign of \(E^{\mathrm{ret}}-E^{\mathrm{adv}}\) fixed
by the forward normal response.

For the smooth-error assertion, first restrict a proposed kernel error
and its normal derivative to the first initial disk. Each is a
distribution in the second spacetime variable, depending on the first
spatial variable and any external parameters. Its second-variable wave
equation and both initial traces are jointly smooth. The linear Cauchy
estimate at all orders, with these variables retained as parameters,
makes both restricted errors jointly smooth. The first-variable Cauchy
equation, its now smooth initial data and its smooth forcing then make
the entire error jointly smooth. Derivatives in the passive parameters
satisfy linear inhomogeneous equations with finitely many earlier
derivatives. The same energy induction proves the joint parameter
assertion on the same smaller parent. \hfill\(\square\)

Let \(\mathcal C_g^+\) denote the primed null relation whose first
temporal covector in a Gaussian chart is \(-H\) and whose two covectors
lie on the same null bicharacteristic. This notation fixes a phase
convention; it is independent of the auxiliary Gaussian chart under
time-oriented physical comparisons. Its transposed unprimed relation
corresponds, after priming, to the opposite null component. For a
quasifree Klein--Gordon field on a fixed globally hyperbolic spacetime,
\hyperref[ref-Radzikowski1996]{Radzikowski (1996)} characterized the
Hadamard condition by the wavefront set of the two-point distribution.
Definition 6.9 specifies the positive-frequency condition in our phase
convention and its extension to joint parameter variables.

\pandocstatementheading{\leavevmode\phantomsection\label{def:online-independent-wave-kernel}Definition 6.9 (Independent admission of the wave kernel).}

Over an actual source collar, supply a scalar distribution kernel class
\([\Lambda]\) modulo jointly smooth kernels. Require actual classical
oscillatory representatives on the retained null relation, in the
uniform family calculus, and test

\phantomsection\label{eq:online-independent-wave-kernel-tests} \[
\Box_{g,p}\Lambda=0,\qquad
\Box_{g,p'}\Lambda=0,\qquad
\operatorname{WF}'(\Lambda)\subset\mathcal C_g^+,
\qquad
\Lambda-\Lambda^T=-iE_g,
\] where the three kernel equations are modulo joint smoothing. All
nonzero null directions over the retained point neighborhood are
included.

For a parameter domain \(\mathcal P\), every representative and residual
uses the joint variables \((\lambda,p,p')\). In particular, its joint
wavefront set is required to lie in

\phantomsection\label{eq:online-wave-joint-positive-relation} \[
\left\{
(\lambda,p,p';\sigma,\xi,-\xi'):
(p,\xi;p',\xi')\in\mathcal C_{g_\lambda}^+
\right\}.
\] The spacetime covectors in this set are nonzero. Pure parameter
singularities are excluded. The parameter covector \(\sigma\) is the one
furnished by the uniform oscillatory family; it is not set to zero when
the metric or phase varies. These conditions concern the supplied kernel
and recognized source fields. Membership in a constructed kernel image,
a choice of quantum state and a chosen exact smoothing tail are not
admission conditions.

\pandocstatementheading{\leavevmode\phantomsection\label{thm:online-canonical-wave-kernel}Theorem 6.10 (Canonical admitted wave kernel).}

The source metric determines exactly one independently admitted scalar
kernel class \(\Lambda_g^+\) on a sufficiently small retained
neighborhood squared. It is physically Hermitian modulo joint smoothing,
natural under restriction and time-oriented physical isometry, and
independent of the auxiliary Gaussian disk, quantization, cutoff and
asymptotic realization.

On a Gaussian collar with initial time \(s\), the class is represented
by

\phantomsection\label{eq:online-wave-kernel-factorization} \[
\widetilde\Lambda^+(t,t')
=E_{t,s}R_s^{-1}E_{t',s}^\dagger,
\qquad
\Lambda^+(t,x;t',y)
=\rho_t(x)^{-1}\widetilde\Lambda^+(t,x;t',y)\rho_{t'}(y)^{-1}.
\] The uniqueness, equations and naturality hold in the joint class for
smooth finite-dimensional families.

\emph{Proof.} The normal-branch and charge results give the displayed
expression. The two first-order evolutions and the Riccati identity make
it a bisolution modulo joint smoothing. Its real phase is the retained
null canonical relation, with the fixed frequency component. Products
and proper support are taken on compactly contained ray domains; the
construction covers a smaller neighborhood in all directions by the
finite angular cover. Changes of low-frequency cutoff contribute a
jointly smooth kernel. The construction and its Hermitian transpose
agree because \(R_s^{-1}\) is Hermitian.

Put \(C=R_s^{-1}\). It is real and transpose symmetric. We verify all
four double Cauchy traces of the antisymmetric difference, using
\(\pi=-Qv\) and \(Q^T=Q\). They are

\phantomsection\label{eq:online-wave-commutator-four-traces} \[
\begin{aligned}
\Lambda_{00}-\Lambda_{00}^T&=C-C^T=0,\\
(\Lambda-\Lambda^T)_{10}
&=(-Q+Q^\dagger)C=-iI,\\
(\Lambda-\Lambda^T)_{01}
&=C(Q-Q^\dagger)=iI,\\
(\Lambda-\Lambda^T)_{11}
&=QCQ^\dagger-Q^\dagger CQ=0.
\end{aligned}
\] The last identity keeps the noncommuting order of every factor: since
\(Q-Q^\dagger=iR\) and \(C=R^{-1}\), its difference is
\(iRCQ^\dagger-iQ^\dagger CR=0\). The preceding Green lemma identifies
these traces with those of \(-iE_g\). Its successive Cauchy argument in
the two variables makes the remaining bisolution error jointly smooth.
Thus the commutator equation holds as a complete kernel identity, not
just at principal order.

For uniqueness, let \(D\) be the difference of two admitted kernels.
Their commutator tests give \(D-D^T\) jointly smooth. Their joint
wavefront conditions put \(\operatorname{WF}(D)\) in the positive
relation of \hyperref[eq:online-wave-joint-positive-relation]{the
definition}. Transposition sends its nonzero spacetime covectors to the
opposite component. Since equality modulo joint smoothing identifies the
wavefront sets of \(D\) and \(D^T\), that wavefront set lies in the
intersection of the two components. The intersection is empty. The
exclusion of pure parameter covectors makes this a statement about the
full joint wavefront set. Hence \(D\) is jointly smooth.

All auxiliary constructions satisfy the same tests, so uniqueness
removes their choices in the stated quotient. Restriction to a smaller
causally convex neighborhood preserves the causal kernel by finite
propagation and preserves the other tests. A physical isometry carries
the volume-normalized wave equation, causal kernel and chosen null
component to their counterparts. Pullback therefore satisfies the same
tests and gives the unique class. For a smooth isometry family its
parameter derivative changes the parameter covectors of the phase, but
it preserves the nonzero spacetime component and the joint smooth ideal.
This proves joint naturality as well. \hfill\(\square\)

\pandocstatementheading{\leavevmode\phantomsection\label{prop:online-wave-kernel-finite-families}Proposition 6.11 (Finite coefficient estimates on one wave parent).}

Let a smooth finite-dimensional family of actual source collars be given
on the common parent of
\hyperref[thm:online-common-collar-finite-families]{Theorem 2.14}. After
one further restriction near the chosen parameter value, the normal
symbols, charge roots, propagator classes and admitted kernels above are
constructed on one point parent and a finite angular phase cover, chosen
before any prescribed coefficient depth or derivative order.

For each finite coefficient depth \(J\), spacetime/angular derivative
order \(m\) and parameter derivative order \(q\), the corresponding
finite list \(\mathcal Z_J\) of symbols, phase functions and amplitudes
satisfies a bound

\phantomsection\label{eq:online-wave-kernel-finite-family-bound} \[
\|\mathcal Z_J\|_{m,q}
\le F_{J,m,q}\bigl(
\mathcal N_N(g,\text{markings};\mathcal K'\times\text{larger parent}),
\mathfrak m^{-1}\bigr),
\] for a finite \(N\) and a finite locally bounded function
\(F_{J,m,q}\). Here \(\mathcal K\Subset\operatorname{int}\mathcal K'\)
are the retained parameter sets, angular norms use compact unit-covector
sections, and \(\mathfrak m\) includes spacelike, inverse-phase,
positive-root and domain-buffer margins. Every finite coefficient of
\(Q,R,W\) is a local expression in a finite jet of the metric. The wave
kernel has its unique joint class; the estimate selects no exact
smoothing remainder.

\emph{Proof.} On the compact normalized cones, positivity of \(h\) gives
a uniform lower bound for \(H\). The recursive divisions defining
\(A_h,Q,R^{-1}\) and \(W\) have denominators consisting of positive
powers of \(H\) or \(\sqrt{2H}\). At each fixed depth only finitely many
coefficients and derivatives occur. The product, inverse and derivative
estimates of the shared symbol calculus therefore give the indicated
finite bounds for this part.

The Hamiltonian flow and its mixed generating map satisfy finite ODE
systems. Their low-order bounds select the retained time and phase
parents, including a positive inverse margin, before higher derivatives
are requested. Smooth dependence and their differentiated ODEs give
every finite mixed estimate on those same parents. Compactness of the
unit-direction sphere gives a finite cover with a common positive
buffer. Increasing coefficient depth requires higher input norms but
does not change that cover or the chosen trajectories.

Each amplitude coefficient solves a linear transport equation along this
already selected flow. Its forcing is a finite differential expression
in earlier coefficients, the phase and the operator symbol. The
initialized leading coefficient is a nonzero exponential transport.
Induction over coefficient depth and differentiation in the finite
parameters give the required bounds by the corresponding integral
equations. Composition and stationary phase at fixed depth use finitely
many further derivatives and the retained inverse-phase margins.
Multiplication by the scalar density factors has the same finite
dependence.

Lemma 6.4 realizes the amplitude families by asymptotic summation on the
selected phase charts. To verify the joint wavefront condition, observe
that the nonzero spacetime covectors control the homogeneous phase
variable on each chart, while the parameter derivatives of the phase are
bounded on its compact unit section. These bounds exclude a pure
parameter wavefront direction. The kernel theorem then identifies every
such realization with the same joint class. This proves the stated
finite-order family conclusion. \hfill\(\square\)

\subsection{6.6. Spacelike carriers and the complete comparison
operators}\label{spacelike-carriers-and-the-complete-comparison-operators}\label{sec:online-wave-carriers}

The scalar kernel is determined on the actual source collar. A spacelike
disk gives a carrier on which its traces become operators on
half-densities. We compare different carriers by the same kernel,
including carriers that intersect. All operator identities in this
subsection are identities in the complete-symbol and joint-smoothing
quotient of \hyperref[def:online-complete-symbol-calculus]{the calculus
convention}.

\pandocstatementheading{\leavevmode\phantomsection\label{def:online-wave-carrier-displays}Definition 6.12 (Retained wave carriers).}

Let \(B\) be a relatively compact point neighborhood in a source collar
and let \(F_a:\Sigma_a\to M\) be smooth spacelike embeddings with future
unit normals. Retain open homogeneous ray domains on which every
relevant null ray meets each carrier once, transversely and away from
its boundary. Compositions use a smaller common ray domain with compact
angular closure and properly supported representatives. To obtain a
full-sky assertion over \(B\), these domains must contain every nonzero
null direction over \(B\), using a finite angular cover.

For each comparison use a single open wave neighborhood \(W\) in the
developed collar. It is globally hyperbolic and causally convex in the
collar and contains the point core, the retained carrier hits and the
ray segments joining them. The kernel tests on \(W\times W\) use the
retarded-minus-advanced propagator for its physical volume. Only the
retained portions of each carrier enter the comparison.

For a smooth family the preceding conditions hold for
\((W_\lambda,g_\lambda)\) at each parameter. Use a smooth spacelike
Cauchy hypersurface \(S_\lambda\) of each \(W_\lambda\) as an auxiliary
initial surface. Near a reference parameter \(\lambda_0\), require
smooth identifications \(\Phi_\lambda:N\to W_\lambda\) onto their images
on neighborhoods \(N\subset W_{\lambda_0}\) of the compact sets used,
with \(\Phi_{\lambda_0}\) the inclusion and
\(\Phi_\lambda^{-1}(S_\lambda)=N\cap S_{\lambda_0}\). Retain positive
spacelike, transversality and interior margins there. At a single
metric, the existence of a smooth spacelike Cauchy surface follows from
\hyperref[ref-BernalSanchez2005]{Bernal and Sánchez (2005), Theorem
1.1}.

For the local displays of §5 one can choose these neighborhoods about
any marked point. Inside the causally convex development take a small
coordinate cylinder around the point with a spacelike coordinate slice
through it. On a smaller cylinder a constant Lorentz cone strictly
contains all the metric causal cones, including those for nearby
parameters. Use the development time as coordinate time. On a compact
subcylinder the spatial speeds of causal curves in this time are
uniformly bounded. Choose the time span of a central wider-cone diamond
so small that this bound times its span is less than the spatial buffer
to the subcylinder wall. Time is monotone on the development. A causal
curve between points of the diamond must therefore stay in the
subcylinder, by the speed bound up to a first exit. Cone inclusion then
keeps it in the diamond. The central slice of the diamond is Cauchy for
each metric: toward that slice the sections of the diamond expand faster
than any causal curve moves in the spatial coordinates, so such a curve
continues to the slice. Strict time monotonicity gives a unique
crossing. A smaller point core and an interior disk of the slice retain
all directions, by compactness of the unit sky and smooth transverse
dependence of the hits. The comparison with any further carrier uses the
common rays for which the specified margins hold.

Write \(h_a=-F_a^*g\) and, in local coordinates,
\(\rho_a=(\det h_a)^{1/4}\). With the coordinate half-density factors
understood, define

\phantomsection\label{eq:online-wave-carrier-covariances} \[
\begin{aligned}
C_{ab}(x,y)
&=\rho_a(x)\rho_b(y)
  \Lambda_g^+(F_a(x),F_b(y)),\\
C_a&=C_{aa},\qquad W_a=C_a^{-1/2}.
\end{aligned}
\] Thus \(C_{ab}\) acts from half-densities on \(\Sigma_b\) to
half-densities on \(\Sigma_a\). On the retained nonzero cone \(W_a\) is
the complete elliptic inverse square root of \(C_a\) with positive
physical principal symbol.

The restrictions are defined. A nonzero null covector cannot annihilate
the tangent space of a spacelike hypersurface, so the kernel wavefront
relation misses the conormal obstruction to both pullbacks. By
constructing the canonical kernel in a Gaussian collar based at
\(\Sigma_a\), its diagonal trace is

\phantomsection\label{eq:online-wave-carrier-diagonal-symbol} \[
C_a=R_a^{-1},\qquad
\sigma_{-1}^{\mathrm{phys}}(C_a)=(2H_a)^{-1},
\qquad H_a=|\xi|_{h_a}.
\] Here \(R_a=i(Q_a^\dagger-Q_a)\) is the normal charge operator of
\hyperref[prop:online-wave-charge-normalization]{the charge
normalization}. Uniqueness of
\hyperref[thm:online-canonical-wave-kernel]{the canonical kernel} makes
this the same \(C_a\) for every auxiliary Gaussian construction. In
particular \(C_a^\dagger=C_a\), and the square-root recursion gives
\(W_a^\dagger=W_a\) and \(W_aC_aW_a=I\).

The common ray domains identify \(T^*\Sigma_b\setminus0\) with
\(T^*\Sigma_a\setminus0\) by the null characteristic flow. This is a
homogeneous canonical graph: restriction of the characteristic
symplectic form to either spacelike section gives its cotangent
symplectic form. The field trace of the chosen wave branch is an
elliptic order-zero Fourier integral operator on this graph. In a short
phase chart this follows from the real Hamilton--Jacobi phase and its
nonvanishing initialized transport amplitude. Finite composition of such
charts gives the retained graph operator and its inverse in the
quotient. The graph composition and inverse statements use
\hyperref[lem:online-finite-symbol-family-calculus]{the stated
finite-symbol calculus}.

\textbf{Mixed traces on a common wave neighborhood.} Write \((W,S)\) for
one of these wave neighborhoods and its Cauchy surface. Construct the
kernel on \(W^2\) from its four scalar Cauchy traces. In Gaussian
patches of \(S\) let \(J_s=(I,-Q_s)^{\mathsf t}\) be the positive Cauchy
column in the field and physical normal half-density traces \((v,\pi)\).
By Theorem 6.10 the four initial traces form the matrix \[
A_s=J_s R_s^{-1}J_s^\dagger .
\] Here the row is \((I,-Q_s^\dagger)\), with the physical adjoint. On
overlapping initial patches the four kernels agree near the diagonal
modulo joint smoothing. Indeed, restrict the two short spacetime kernels
to a common causally convex neighborhood of each point of the overlap
and use uniqueness in Theorem 6.10 before taking traces. On \(S^2\)
every possible singular pair lies on the diagonal since \(S\) is
acausal. Glue the four properly supported classes there by a locally
finite partition in the first variable and right cutoffs equal to one
near each partition support, extending each term within its patch. The
resulting matrix \(A\) has the four specified complete symbols. Changing
cutoffs or partitions changes it by a smooth matrix kernel, jointly in
parameters.

Let \(\mathcal U\) send the two initial half-density traces to the
spacetime field by ordinary scalar Cauchy evolution. For a fixed metric,
\(\Box_g\) is normally hyperbolic, \(W\) is globally hyperbolic and
\(S\) is a smooth spacelike Cauchy hypersurface.
\hyperref[ref-BarGinouxPfaffle2007]{Bär, Ginoux and Pfäffle (2007),
Theorems 3.2.11--3.2.12} give smooth well-posedness, continuous
dependence on the data and propagation inside their causal future and
past. Testing against compactly supported smooth sources and solving the
adjoint inhomogeneous equation gives the distributional evolution by
duality. Evolve the matrix \(A\) in its two variables to obtain
\(\mathcal U A\mathcal U^\dagger\) with the volume and surface pairings.
The second evolution solves the conjugate equation. The evolutions
commute. Smooth changes in all four traces have smooth double evolutions
by Lemma 6.8. Applying this to the Hermitian difference and to the four
antisymmetric traces of Theorem 6.10 gives physical Hermiticity and the
commutator \(-iE_g\) in the quotient.

\textbf{Localization of the initial data.} Take compact neighborhoods
\(K_a,K_b\) of the endpoint patches in \(W\), with room to restrict the
patches, and set \(K=K_a\cup K_b\). At a reference parameter the set
swept out by causal segments from \(K\) toward \(S\) up to their first
hit is compact. Indeed it is the union of \(J^-(K)\cap J^+(S)\) and
\(J^+(K)\cap J^-(S)\), both compact by
\hyperref[ref-BarGinouxPfaffle2007]{Bär, Ginoux and Pfäffle, Corollary
A.5.4}. Its intersection with \(S\) is the causal shadow of \(K\). A
spatial cutoff equal to one near this shadow leaves the solution on the
endpoint patches unchanged by finite propagation. Apply it in each
initial variable.

Here the same compact localization works near the reference parameter.
Choose a relatively compact open neighborhood \(O\) of the swept set,
with closure in the reference wave neighborhood. Use the local
identifications on a neighborhood of its closure, straightening the
varying surface there to \(S\). At the reference metric a temporal
function \(t\) with \(S=t^{-1}(0)\) exists by
\hyperref[ref-BernalSanchez2006]{Bernal and Sánchez (2006), Theorem
1.2}. On the closure of \(O\), \(\mathrm dt\) remains timelike for
nearby metrics; hence future causal vectors normalized by
\(\mathrm dt(v)=1\) have uniformly bounded auxiliary Riemannian length.
If segments from the endpoint compacts toward \(S\) could first exit
\(O\) before hitting it for parameters tending to the reference value,
parameterization by \(t\) would yield a uniformly Lipschitz subsequence
up to first exit. Its limit joins \(K\) to the boundary on the same side
of \(t=0\), allowing a boundary endpoint at \(t=0\). The limit is
causal: on short coordinate subarcs the tangent vectors lie in
arbitrarily small convex enlargements of the reference causal cones, and
the increment bound passes to the limit. Extend the limiting segment
toward the Cauchy surface if necessary. Its boundary endpoint belongs to
the reference swept set, which lies in the interior of \(O\). This
contradiction gives the common localization. One can take the initial
cutoff equal to one on \(S\cap\overline O\) with slightly larger compact
support. The argument uses the metric only on compact neighborhoods and
works in either time direction.

\textbf{Positive branch on the localized data.} Cover the entire unit
cosphere over that compact support by finitely many coordinate cones. On
smaller cones use order-zero symbols forming a partition of unity at
high frequency, quantized to sum to the identity modulo smoothing there.
In each larger cone the initial matrix is \(J_sR_s^{-1}J_s^\dagger\) to
every order. Insert the partition \emph{inside} this factorization, on
the left of \(R_s^{-1}\), with a right cutoff equal to one on a
neighborhood of the microsupport of the left cutoff. Thus each term has
initial columns \(J_s\) on both sides, with the second column physically
adjointed. The nested cutoffs preserve the complete matrix on summation
by pseudolocality; the low-frequency term is jointly smooth. The
transition regions of the outer spatial cutoffs lie beyond the causal
shadow, so commutators supported there have no effect on the endpoint
patches.

On a short Gaussian neighborhood the ordinary evolution of \(J_s\)
agrees with the positive parametrix of Theorem 6.6 by its field and
normal traces. At fixed metric
\hyperref[ref-DuistermaatHormander1972]{Duistermaat and Hörmander
(1972), Theorem 6.1.1} propagates the wavefront of a solution with
smooth forcing along the null bicharacteristics: the real principal
symbol \(g^{-1}(\xi,\xi)\) has nonzero fiber derivative on its nonzero
characteristic set. In passing to another spacelike slice the two short
branch columns split the full Cauchy data, since the difference of their
normal roots is elliptic. An opposite-branch coefficient there would
propagate elliptically into the opposite cone of the overlap, where the
incoming solution is smooth. Successive short evolutions therefore carry
the positive column to the endpoint, including across changes of phase
coordinates.

For the family argument the phase representations can be taken on a
finite cover \emph{before} selecting an endpoint cone. Subdivide the
bounded \(t\) interval in \(\overline O\) into short slabs. With
spacelike \(t\)-sections, the two local phase constructions apply in
flow boxes by using Gaussian coordinates and taking traces on the
sections. Cover the compact part of each intermediate cosphere by
smaller boxes inside larger ones. The nonzero normal roots stay
separated and the short flow projections remain invertible for a common
sufficiently short step, by compactness and smooth dependence. Compose
the two-trace Cauchy parametrices across these sections with a partition
on \textbf{all} intermediate directions; equivalently the paths needed
at an endpoint are obtained by tracing its two null lifts back toward
\(S\). Finite propagation permits spatial localization at the successive
sections in neighborhoods of the swept set. With cutoffs equal to one
near each required domain of dependence, the short parametrices have
smooth equation and initial errors there to every order. The
differentiated Cauchy estimates of Lemma 6.8 carry these errors through
the finite number of steps, also for smooth parameter derivatives. In
each normalized short graph chart the opposite-branch symbol has every
homogeneous coefficient zero at each parameter by the fixed-metric
argument above, so all their parameter derivatives vanish and the
uniform classical remainder estimates to arbitrary depth give joint
smoothing after every finite mixed derivative. The properly supported
graph compositions are transverse on the spacelike sections, as required
in \hyperref[ref-Duistermaat1972FIO]{Duistermaat (1972), §1, Theorem
1.1}. \hyperref[lem:online-finite-symbol-family-calculus]{Lemma 6.4}
gives the corresponding parameter-dependent symbol compositions and
remainder estimates on these buffered charts, including the terms
subsequently discarded at the endpoints.

\textbf{Mixed field traces.} Fix smaller conic cutoffs at the two
endpoints. A stationary point of a localized product must lift each
endpoint to its positive null covector, follow the two Hamiltonian
trajectories to \(S\), and give the same point and the same tangential
covector there in the primed convention. The last condition is the
diagonal canonical relation of \(R_s^{-1}\). On a patch about a matched
pair, choose a disk \(D\subset S\) and an initial cone containing both
returned covectors, with a cutoff equal to one on a smaller neighborhood
of their possible values. Shrink the endpoint cutoffs inside that
neighborhood under the return maps. In the terms of the finite partition
outside this initial cutoff the matching equations have no solution.
Overlapping terms inside it sum to the identity microlocally.

To see the regularity of the discarded terms, localize to dyadic
frequencies of size \(L\). For comparable frequencies the joined phases,
rescaled to unit frequency, have gradients bounded away from zero in the
integration variables on the compact supports separated from the
matching equations. For noncomparable adjacent frequencies,
differentiation in the common spatial variable separates their covectors
since the short graph maps preserve frequency up to bounded positive
factors. One can equivalently subdivide into dyadic shells for each
phase variable, integrating first at the largest unmatched scale.
Repeated integration by parts in these variables gives arbitrary inverse
powers of the largest frequency; bounded frequencies give smooth
kernels. The separation persists on smaller parameter neighborhoods.
Each specified endpoint or parameter derivative costs at most a fixed
power of that frequency by the phase and amplitude estimates of
\hyperref[prop:online-wave-kernel-finite-families]{Proposition 6.11}.
Choose the number of integrations above this loss, the symbol orders and
the integration dimensions. The shells then sum absolutely with every
prescribed derivative. This proves joint smoothing off the matched
graph. On it, both nonzero endpoint covectors control the phase
frequency, and parameter covectors are bounded by that frequency on the
compact unit sections. These are the joint wavefront conditions of
Definition 6.9.

Write \(E_{as}\) for the field trace on \(a\) of the ordinary evolution
of \(J_s\) from \(D\) on this cone, including the carrier half-density
factor \(\rho_a\). The localized product is
\(C_{ab}=E_{as}R_s^{-1}E_{bs}^\dagger\) in the joint quotient. Taking
the spacetime field in the second slot instead gives the kernel family
used in \(C_{ap}\), with the point retained as a variable. The graph
operators for the carrier traces are elliptic by short-time transport
and transverse restriction. On intersections of reference disks and ray
patches all the products are traces of the same double evolution.
Applying the construction on compact neighborhoods in \(W^2\) gives the
positive wavefront test everywhere; the transposition argument of
Theorem 6.10 identifies this class with every independently admitted
kernel on \(W^2\).

To compare two choices \(W,W'\) in the same collar, restrict to a
connected globally hyperbolic open set \(U\subset W\cap W'\) causally
convex in both. One may use the intersection component containing the
retained ray segment: causal convexity in the collar keeps each causal
diamond between its points inside that component, and the diamond is
compact by global hyperbolicity of \(W\). On \(U\) the scalar Green
operators from either neighborhood restrict to the unique Green
operators of \(U\) by zero extension of test sources and causal
compatibility (\hyperref[ref-BarGinouxPfaffle2007]{Bär, Ginoux and
Pfäffle}, Proposition 3.5.1 and Corollary 3.4.3). The two kernel classes
thus have the same commutator on \(U^2\) and the same positive wavefront
constraint there. The transposition argument of
\hyperref[thm:online-canonical-wave-kernel]{Theorem 6.10} identifies
them on all of \(U^2\) modulo joint smoothing. For displaced carriers
retain a common ray patch whose two hits and joining segments lie in
\(U\), with the margins of Definition 6.12; on this patch the mixed
traces, diagonal charge roots and hence the normalized transports agree.
For a family these comparisons use the same local identifications and
buffers on the retained compact sets.

\pandocstatementheading{\leavevmode\phantomsection\label{lem:online-wave-spacelike-transport}Lemma 6.13 (Normalized transport through a spacelike isotopy).}

Let \(F_s\) be a smooth local isotopy of spacelike embeddings of a fixed
disk. Decompose its velocity as \(N_sn_s+Z_s\), with \(n_s\) the future
unit normal and \(Z_s\) tangential. The lapse \(N_s\) may vanish or
change sign. For a scalar wave solution define its half-density traces
\(v_s=\rho_sF_s^*u\) and \(\pi_s=\rho_sF_s^*(n_su)\). Their equations
are

\phantomsection\label{eq:online-wave-spacelike-hamiltonian-system} \[
\begin{aligned}
\dot v_s&=L_sv_s+N_s\pi_s,\\
\dot\pi_s&=-D_sv_s-L_s^\dagger\pi_s,\\
L_s&=\mathcal L_{Z_s}^{(1/2)}+\frac{N_s\tau_s}{2},\\
D_s&=-\rho_s\operatorname{div}_{h_s}
       \bigl(N_s\operatorname{grad}_{h_s}\bigr)\rho_s^{-1},
\qquad \tau_s=\operatorname{tr}_{h_s}K_s.
\end{aligned}
\] The normal positive-branch symbol \(Q_s\) gives \(\pi_s=-Q_sv_s\) and
hence

\phantomsection\label{eq:online-wave-spacelike-trace-generator} \[
\dot v_s=-B_sv_s,\qquad B_s=N_sQ_s-L_s.
\] Charge normalization by \(W_s=R_s^{1/2}\) gives an integer-order
generator

\phantomsection\label{eq:online-wave-spacelike-normalized-generator} \[
B_s^{\mathrm{cov}}
=W_sB_sW_s^{-1}-\dot W_sW_s^{-1}
\] whose complete formal symbol is real and which satisfies
\((B_s^{\mathrm{cov}})^*=-B_s^{\mathrm{cov}}\). Two sufficiently small
spacelike disks can be joined by such an isotopy on smaller common
domains, including when they intersect.

\emph{Proof.} Differentiation of the induced metric gives
\(\dot h_s=2N_sK_s+\mathcal L_{Z_s}h_s\). Consequently
\(\dot\rho_s/\rho_s=N_s\tau_s/2+\operatorname{div}_{h_s}Z_s/2\).
Differentiating the scalar trace therefore gives the first equation,
with the usual half-density Lie derivative. The wave divergence formula,
or differentiation of the normal momentum density, gives \[
\frac{\mathrm d}{\mathrm ds}
  \bigl(\rho_s^2F_s^*(n_su)\bigr)
=\rho_s^2\operatorname{div}_{h_s}
       \bigl(N_s\operatorname{grad}_{h_s}F_s^*u\bigr)
  +\mathcal L_{Z_s}^{(1)}
       \bigl(\rho_s^2F_s^*(n_su)\bigr).
\] Dividing by \(\rho_s\) and using its displayed derivative gives the
second equation. In particular \(D_s^T=D_s\), and the half-density Lie
derivative of a real vector field is skew for both physical adjoints.
These formulas involve multiplication by \(N_s\) and its derivatives;
they do not divide by \(N_s\).

For each \(s\), construct \(Q_s\) in a separate normal Gaussian collar
by \hyperref[thm:online-wave-normal-riccati]{the normal Riccati
theorem}. Its denominators are \(2H_s\). The kernel built in that collar
agrees with the canonical kernel by its independent characterization.
Thus the same positive wave branch has graph \(\pi_s=-Q_sv_s\) on every
embedding. Substitution in the Hamiltonian system gives

\phantomsection\label{eq:online-wave-spacelike-riccati} \[
Q_sN_sQ_s-\dot Q_s-Q_sL_s-L_s^\dagger Q_s+D_s=0.
\] This is a compatibility identity for the already constructed branch.
It is not a recursion divided by the deformation lapse. The construction
therefore applies also where \(N_s=0\).

The Hamiltonian system preserves the Klein--Gordon form. Equivalently,
subtract the displayed equation from its physical adjoint and use
\(R_s=i(Q_s^\dagger-Q_s)\) to obtain

\phantomsection\label{eq:online-wave-spacelike-charge-identity} \[
\dot R_s=B_s^\dagger R_s+R_sB_s.
\] If \(E_{s,s_0}\) is the field-trace evolution, this gives
\(E_{s,s_0}^\dagger R_sE_{s,s_0}=R_{s_0}\). Thus
\(W_sE_{s,s_0}W_{s_0}^{-1}\) is physically unitary and its generator
\(B_s^{\mathrm{cov}}\) is physically anti-Hermitian.

For the internal adjoint, \(N_s\), \(Q_s\) and \(L_s\) have real
complete formal coefficients, so the same is true of \(B_s\). The
half-order calculation of
\hyperref[prop:online-wave-charge-normalization]{the charge
normalization} gives, with \(\mathcal F_s=Q_s-Q_s^*\),
\(W_s=e^{-i\pi/4}\mathcal F_s^{1/2}\) in rephased symbols. The recursive
coefficients of \(\mathcal F_s^{1/2}\) are real; the constant phase
cancels in both terms of \(B_s^{\mathrm{cov}}\). Its integer-degree
formal coefficients are therefore real. On those coefficients the
physical Hermitian adjoint agrees with the internal adjoint, proving the
asserted skew identity. No internal involution on half-integer orders is
required.

Finally put the two disks into graph form over a common spacelike
coordinate disk. Choose this local metric neighborhood sufficiently
small relative to the retained endpoint spacelike margins that both
graph conormals belong to the future timelike cone of one fixed
reference Lorentz form. Their convex interpolations then have a positive
timelike margin on the compact graph and interpolation parameters.
Shrinking the common neighborhood further preserves that margin for the
varying metric along the interpolated graphs. Interpolating the graph
functions therefore gives the required spacelike isotopy. Its lapse can
vanish where the graphs meet; the preceding proof already permits this.
\hfill\(\square\)

\pandocstatementheading{\leavevmode\phantomsection\label{thm:online-wave-carrier-cocycle}Theorem 6.14 (Complete carrier comparison and cocycle).}

On common retained ray domains define

\phantomsection\label{eq:online-wave-normalized-carrier-map} \[
T_{ab}=W_aC_{ab}W_b.
\] Then \(T_{ab}\) is an elliptic order-zero Fourier integral operator
on the null canonical graph, and

\phantomsection\label{eq:online-wave-carrier-complete-cocycle} \[
\begin{aligned}
C_{ab}C_b^{-1}C_{bc}&=C_{ac},\\
T_{aa}&=I,\qquad T_{ab}^\dagger=T_{ba},\qquad
T_{ab}T_{bc}=T_{ac}.
\end{aligned}
\] Conjugation \(\mathcal A_{ab}(P)=T_{ab}PT_{ba}\) is a filtered
isomorphism of the complete carrier pseudodifferential algebras and
preserves their internal adjoints. The map
\(\mathcal M_{ab}(m)=T_{ab}m\) is an invertible semilinear map between
the complete modules of the same retained sky:

\phantomsection\label{eq:online-wave-carrier-algebra-module-identities}
\[
\mathcal M_{ab}(Pm)=\mathcal A_{ab}(P)\mathcal M_{ab}(m),
\qquad
\mathcal A_{ab}(P^*)=\mathcal A_{ab}(P)^*.
\] These maps commute with restriction, have exact identity, inverse and
triple-composition laws in the quotient, and are natural under supplied
physical isometries and their induced carrier maps.

\emph{Proof.} On each of the ray patches constructed before Lemma 6.13,
take a common reference disk in the Cauchy surface for all carriers in
this composition. The double Cauchy evolution proved there gives

\phantomsection\label{eq:online-wave-carrier-common-factorization} \[
C_{ab}=E_{as}R_s^{-1}E_{bs}^\dagger.
\] The trace maps are invertible graph operators in the quotient.
Inserting the complete inverse \[
C_b^{-1}=(E_{bs}^\dagger)^{-1}R_sE_{bs}^{-1}
\] between \(C_{ab}\) and \(C_{bc}\) cancels all four adjacent inverse
factors and proves the first identity. In particular, if
\(E_{ab}=E_{as}E_{bs}^{-1}\), then

\phantomsection\label{eq:online-wave-carrier-normalized-trace} \[
T_{ab}=W_aE_{ab}W_b^{-1}.
\] This is order zero and elliptic. Hermitian symmetry of the
factorization gives \(C_{ab}^\dagger=C_{ba}\). Since the \(W_a\) are
selfadjoint, it gives \(T_{ab}^\dagger=T_{ba}\). The identities
\(W_b^2=C_b^{-1}\) and \(W_aC_aW_a=I\), together with the mixed identity
already proved, give the identity and cocycle formulas. Consequently
\(T_{ba}\) is the exact quotient inverse. No undetermined scalar occurs
in these cancellations.

The canonical-graph composition theorem makes conjugation by \(T_{ab}\)
a filtered pseudodifferential algebra isomorphism. First consider a pair
of carrier patches joined by a spacelike isotopy as in Lemma 6.13, with
interior transverse hits throughout. The normalized trace formula above
identifies the transport from its initial patch to each moving patch
with the normalized evolution of that lemma. In fixed coordinates on its
parameter disk, the normalized evolution has the generator
\(B_s^{\mathrm{cov}}\). Conjugation obeys \[
\partial_s\mathcal A_s(P)
=[-B_s^{\mathrm{cov}},\mathcal A_s(P)].
\] Since \((B_s^{\mathrm{cov}})^*=-B_s^{\mathrm{cov}}\), applying the
internal anti-involution to this equation gives the same equation for
\(\mathcal A_s(P)^*\) as for \(\mathcal A_s(P^*)\). Their initial values
agree. The complete transport recursion, degree by degree along the
retained flow, has a unique initialized solution, so the two expressions
agree at every degree.

For separated hits, fix a ray of Definition 6.12 and follow its compact
segment from the hit on \(b\) to the hit on \(a\) inside the retained
parent. Subdivide the segment into finitely many pieces, each in a
Gaussian coordinate cylinder, with each subdivision point interior to
both adjacent cylinders. Within each piece use the Gaussian time slices
through its two ends. In the space coordinates on these slices take
disks centered along the segment, transporting the disk coordinates by
translation of the center. The time slices are spacelike and Gaussian
time is strictly monotone along the ray. By compact containment one can
choose a positive disk radius and a buffer throughout that piece. These
moving disks form a spacelike isotopy with the given ray meeting their
centers. At each subdivision point the incoming and outgoing disks pass
through the same point. Join their restrictions by the local graph
interpolation of Lemma 6.13 in the overlap. Use the same construction at
the two endpoint hits to join the given carriers to the first and last
Gaussian disks. The interpolation at a shared hit keeps that point
fixed. Intersecting endpoint carriers with coincident hits require only
this last local comparison.

On the unit section about the chosen ray, shrink once for this finite
list of isotopies. Spacelikeness and transverse intersection give a
positive normal-pairing margin, and the central hits have positive
distance from the disk boundaries. Compactness of each closed parameter
interval and continuous dependence of the hits preserve these margins on
the smaller ray patch. Label the successive patches
\(D_0=b,\ldots,D_m=a\), with restrictions understood, and use
\(T_{D_jD_{j-1}}\) from \(D_{j-1}\) to \(D_j\). The kernel cocycle
proved above gives \[
T_{ab}=T_{D_mD_{m-1}}\cdots T_{D_1D_0}.
\] Conjugate first by the rightmost factor and then by the successive
factors. Each conjugation preserves the internal adjoint by the
initialized transport equation, hence so does their composition. On
intersections of ray patches each product equals the same mixed-kernel
operator. A finite patch cover of the retained compact unit section
therefore gives the identity throughout the common domain.

The graph operator maps the Lagrangian distribution class of a sky on
\(\Sigma_b\) to that same sky on \(\Sigma_a\). The action includes its
complete amplitude, as prescribed by
\hyperref[prop:online-shifted-sky-modules]{the shifted module
convention}. Properly supported graph composition sends jointly smooth
remainders to jointly smooth remainders. The semilinear identity is now
ordinary operator composition. Invertibility and the operator cocycle
prove the module inverse and cocycle laws as well.

A supplied physical isometry preserves the scalar kernel by its
independent characterization. Its induced carrier half-density
pushforwards therefore carry \(C_{ab}\) to the corresponding mixed
kernel. Uniqueness of the positive-principal complete square root
carries \(W_a\), and the defining product carries \(T_{ab}\). This
proves naturality. All constructions and cancellations are valid after
restriction to common buffered ray domains, proving the remaining laws.
\hfill\(\square\)

The scalar kernel and these scalar module maps do not detect the central
spin sign. The actual spin maps of the geometric source remain
separately retained comparisons.

\subsection{6.7. Evaluation on the whole moving
sky}\label{evaluation-on-the-whole-moving-sky}\label{sec:online-complete-wave-evaluation}

\pandocstatementheading{\leavevmode\phantomsection\label{def:online-wave-evaluation-section}Definition 6.15 (The full wave evaluation section).}

The point slot remains scalar. For \(p\in B\) set

\phantomsection\label{eq:online-wave-complete-evaluation} \[
C_{ap}(x)=\rho_a(x)\Lambda_g^+(F_a(x),p),
\qquad e_a(p)=W_aC_{ap}.
\] Only the carrier variable has a half-density factor. These are actual
distribution families modulo kernels jointly smooth in the point and
carrier variables, and in any supplied passive parameters. Let
\(L_{a,p}\) be the homogeneous Lagrangian obtained by following the
chosen-branch null covectors at \(p\) to their carrier hits and
restricting the transported covectors to \(T\Sigma_a\). The section
\(e_a(p)\) belongs to its complete Lagrangian module on
\(T^*\Sigma_a\setminus0\).

This definition is valid when \(p\) lies on the carrier. The pullback
with fixed \(p\) has no wavefront obstruction because the first, carrier
covector cannot vanish. The same fact gives smooth dependence on \(p\)
as a distribution-valued family. A finite phase cover describes the
moving Lagrangians uniformly. If \(p=F_a(b)\), then
\(L_{a,p}=T_b^*\Sigma_a\setminus0\); a spatial light-cone cut need not
remain a smooth hypersurface there. In a short Gaussian phase chart the
relevant phase is \(S(t,x,\xi)-b\cdot\xi\), with
\(S(0,x,\xi)=x\cdot\xi\). At \(t=0\) its critical equations are
\(x-b=0\) and have full differential rank in \(x\). Thus it remains a
nondegenerate three-frequency phase through the carrier. Propagating and
composing such phase charts gives the stated moving Lagrangian family.

A contact graph \(z=s(p,q)\) can be obtained near any fixed point
\(p_0\) by first transporting to a spacelike disk through \(p_0\) with
the already constructed \(T_{ab}\), and then using the normalized Radon
charts of \hyperref[prop:online-shifted-sky-modules]{the shifted module
convention} on its point fiber. The graph projection remains invertible
for nearby points after restriction. This supplies the graph form
through actual carrier transport and the specified Radon normalization,
with their established algebra and internal-adjoint laws.

\pandocstatementheading{\leavevmode\phantomsection\label{prop:online-wave-full-evaluation-comparison}Proposition 6.16 (Full evaluation and carrier covariance).}

For every retained point,

\phantomsection\label{eq:online-wave-full-evaluation-intertwining} \[
T_{ab}e_b(p)=e_a(p).
\] The comparing operator is independent of \(p\). If \(p=(t,b)\) in a
Gaussian collar, and \(U_{s,t}=W_sE_{s,t}W_t^{-1}\) denotes normalized
trace transport, then

\phantomsection\label{eq:online-wave-evaluation-gaussian-normal-form} \[
e_s(t,b)
=U_{s,t}W_t^{-1}
  \bigl(\rho_t(b)^{-1}\delta_b\bigr).
\] Here \(\delta_b\) denotes the coordinate delta half-density. At
\(p=F_a(b)\), the physical Fourier amplitude of \(e_a(p)\), using the
unitary Fourier normalization of the shifted module, has principal
coefficient

\phantomsection\label{eq:online-wave-evaluation-principal-amplitude} \[
(2\pi)^{-3/2}(2H_a(b,\xi))^{-1/2}\rho_a(b)^{-1}.
\] Its formal degree-one coefficient is \(i^{-1}\) times this
expression. The full section is elliptic and locally cyclic, of module
filtration degree one. These are equalities and properties of complete
classes; all lower coefficients are determined by the canonical kernel
prescription.

\emph{Proof.} The common factorization also holds with its last variable
fixed at \(p\). The same inverse cancellations used for the carrier Gram
identity give \(C_{ab}C_b^{-1}C_{bp}=C_{ap}\). Since \(W_b^2=C_b^{-1}\),
multiplication by \(W_a\) proves the evaluation identity at every point.
The definition of \(T_{ab}\) involves only the source metric and the two
carrier embeddings, so it is independent of the evaluated point.

On \(\Sigma_t\), the point column is \(C_t\delta_b/\rho_t(b)\):
multiplying a scalar kernel by two slice half-densities and then fixing
one variable must remove exactly the half-density factor in that fixed
variable. Hence on \(\Sigma_s\) it is \(E_{s,t}C_t\delta_b/\rho_t(b)\).
Multiplication by \(W_s\), followed by \(C_t=W_t^{-2}\), gives the
displayed normal form.

The principal symbol of \(W_t^{-1}\) is \((2H_t)^{-1/2}\). Its operator
kernel uses the prefactor \((2\pi)^{-3}\), whereas the conormal section
uses \((2\pi)^{-3/2}\). Thus its scalar conormal amplitude has the extra
fixed factor \((2\pi)^{-3/2}\) stated above. The scalar amplitude order
is \(-1/2\), which is \(j-3/2\) for \(j=1\) in the shifted filtration.
Dividing by the degree rephasing \(i^1\) gives the formal coefficient.
This amplitude is nonzero on every retained direction.

The lower coefficients arise from the full operator action, with its
finite Taylor and integration-by-parts terms; no replacement by the
principal amplitude has been made. An elliptic conormal amplitude is
locally a cyclic generator by division and successive correction in the
complete symbols. Transport by the invertible graph operators preserves
that property and the filtration degree. This proves ellipticity and
cyclicity on the entire retained moving sky, using its finite angular
cover. The exact kernel and Gram identities prove the full-class
assertions. \hfill\(\square\)

\subsection{6.8. Finite families and physical
scale}\label{finite-families-and-physical-scale}\label{sec:online-wave-carrier-estimates-scale}

\pandocstatementheading{\leavevmode\phantomsection\label{prop:online-wave-carrier-families-scale}Proposition 6.17 (Carrier families, finite estimates and scale).}

Let source collars and retained carrier embeddings vary in a smooth
finite-dimensional family, with the common parents and positive margins
of \hyperref[thm:online-common-collar-finite-families]{the geometric
family theorem}. Include uniform spacelike, ray-intersection,
phase-inverse and boundary-buffer margins. The kernels \(C_{ab}\), roots
\(W_a\), comparisons \(T_{ab}\) and complete evaluations \(e_a\) have
smooth complete-coefficient families on those parents.

For every finite coefficient depth \(J\), base and angular derivative
order \(m\), and passive-parameter derivative order \(q\), there are a
finite input order \(N\) and a finite locally bounded function
\(F_{J,m,q}\) such that

\phantomsection\label{eq:online-wave-carrier-finite-family-bound} \[
\|\mathcal Y\|_{J,m,q}
\le
F_{J,m,q}\left(
\mathcal N_N(g,F_a;\mathcal K'\times\text{larger parents}),
\mathfrak m^{-1}\right).
\] Here \(\mathcal Y\) is any specified finite list of these
coefficients and transport amplitudes, the output is measured over
\(\mathcal K\Subset\operatorname{int}\mathcal K'\), and \(\mathfrak m\)
denotes the retained positive margins. The parent is selected before
\(J,m,q\). The same statement holds after natural comparison by a
supplied smooth family of physical maps, with its finite map and inverse
bounds included.

In four spacetime dimensions, constant physical scaling has the exact
complete-class laws

\phantomsection\label{eq:online-wave-carrier-physical-scale} \[
\begin{aligned}
\Lambda_{r^2g}^+&=r^{-2}\Lambda_g^+,&
C_{ab}^{(r)}&=rC_{ab},&
W_a^{(r)}&=r^{-1/2}W_a,\\
C_{ap}^{(r)}&=r^{-1/2}C_{ap},&
T_{ab}^{(r)}&=T_{ab},&
e_a^{(r)}&=r^{-1}e_a.
\end{aligned}
\] The numerical scale is retained. It may vary smoothly with passive
parameters, with its finite derivatives and positive lower bound
included in the estimates, but is constant on each spacetime member.

\emph{Proof.} The kernel family's finite-coefficient construction is
provided by \hyperref[prop:online-wave-kernel-finite-families]{the
kernel family estimate}. Restriction to a spacelike carrier uses the
retained transverse phase charts and the finite derivatives of its
embedding. On a unit angular section, the metric bounds give the
required upper and lower bounds for \(H_a\). The elliptic inverse and
positive square-root recursions divide only by their nonzero leading
coefficients. Thus each fixed coefficient of \(W_a\) and each of its
prescribed derivatives uses finitely many previously determined
coefficients and derivatives.

The ray maps solve the actual null Hamilton equations, and the
transverse hit times follow from the implicit function theorem with the
stated intersection margin. Their finite variations satisfy finite ODE
systems and inverse-function identities. Phase coefficients consequently
have finite bounds from the source metric, embeddings and retained
margins. The amplitudes satisfy linear transport recursions on this same
ray tube; differentiating finitely many of their equations gives finite
Gronwall bounds. Stationary phase in a retained nondegenerate
composition uses finitely many derivatives at every output degree. These
facts, together with the operator and evaluation products above, prove
the displayed estimate. No successive reduction of the parent is needed
when the coefficient depth increases.

Uniform asymptotic summation on compact parameter sets supplies actual
classical representatives of the coefficient families. The cutoffs for a
full sum may depend on the entire sequence of symbol bounds. Lemma 6.4
gives the quantitative estimates for the prescribed finite list of
carrier coefficients. Differences of two such sums are uniformly rapidly
decreasing with every retained derivative. Integration by parts, or
direct integration after choosing a sufficiently large decay order,
shows that their kernels are jointly smooth in all variables and passive
parameters. Properly supported composition and restriction with the
stated transverse margins preserve this joint-smoothing quotient.
Supplied physical comparison maps act by the same finite pullback,
inverse and density operations, proving the natural family assertion.

For scale, \(\Box_{r^2g}=r^{-2}\Box_g\) and the physical four-volume is
multiplied by \(r^4\). The scalar causal kernel normalized against that
measure is therefore multiplied by \(r^{-2}\). The independent
canonical-kernel characterization gives the same law for \(\Lambda^+\).
A carrier half-density has weight \(r^{3/2}\), so two carrier factors
give \(C_{ab}^{(r)}=rC_{ab}\), while one carrier factor and the scalar
point slot give \(C_{ap}^{(r)}=r^{-1/2}C_{ap}\). Uniqueness of the
positive-principal inverse square root gives the \(W_a\) law.
Substitution into \(T_{ab}\) and \(e_a\) proves their two weights. The
chain rule gives the finite passive-parameter estimates when \(r\)
varies in such a family. \hfill\(\square\)

The wave evaluation is a full section with its canonical kernel
normalization. Its point derivatives and covariant wave equation are
defined next on the actual moving distribution families. The source's
finite geometric and actual spin structures retain the separate
admission and comparison laws already proved.

\subsection{6.9. Moving modules and actual point
differentiation}\label{moving-modules-and-actual-point-differentiation}\label{sec:online-moving-wave-modules}

Fix a retained carrier disk \(\Sigma_a\) and a buffered spacetime point
neighborhood \(B\) as in \hyperref[def:online-wave-carrier-displays]{the
carrier construction}. The carrier is fixed while the evaluated point
\(p\) varies. Let \(L_{a,p}\subset T^*\Sigma_a\setminus0\) be the conic
Lagrangian determined by the full sky of \(p\). The point slot is
scalar; the distribution on \(\Sigma_a\) carries its spatial
half-density. All geometric coefficients are read from the admitted
finite Cartan collar.

The family of Lagrangians is smooth even when \(p=F_a(b)\) lies on the
carrier. At such a point it is \(T_b^*\Sigma_a\setminus0\). The
nondegenerate phase-family construction in
\hyperref[def:online-wave-evaluation-section]{the evaluation
construction} covers this case. A graph phase \(z=s(p,q)\) below refers
to a fixed angular contact presentation in which the Legendrian is a
graph; it does not require a smooth hypersurface cut in the physical
carrier coordinates. In these presentations \(\tau>0\) denotes the
radial covariable dual to \(z\).

\pandocstatementheading{\leavevmode\phantomsection\label{def:online-moving-wave-module}Definition 6.18 (The moving complete module).}

Let \(\mathscr M_a\) be the sheaf on \(B\) of smooth families of
classical Lagrangian distributions valued in spatial half-densities and
associated with \(L_{a,p}\), in the shifted filtration of
\hyperref[prop:online-shifted-sky-modules]{the sky-module construction},
modulo jointly smooth families. Locally these families have actual
oscillatory representatives

\phantomsection\label{eq:online-moving-module-oscillatory-family} \[
m(p,x)=\int e^{i\phi(p,x,\vartheta)}
          a(p,x,\vartheta)\,\mathrm d\vartheta
          |\mathrm dx|^{1/2},
\] where \(\phi\) is a nondegenerate homogeneous phase smoothly
parametrizing \(L_{a,p}\), and \(a\) is classical with a locally bounded
upper order and symbol estimates uniform with every prescribed finite
point derivative. Different phase presentations are identified by the
actual distribution calculus. The quotient is by
\(C^\infty(B\times\Sigma_a;\Omega_{\Sigma_a}^{1/2})\) on the retained
microlocal patches. With an additional finite-dimensional parameter
\(\lambda\in D\), the discarded remainder must be jointly smooth on
\(D\times B\times\Sigma_a\).

Write \(F_m\mathscr M_a\) for the filtration. The complete carrier
algebra \(\widehat{\mathscr E}_a\) acts on this sheaf. A smooth operator
family is \emph{basic in the point variable} if its actual
representative is independent of \(p\). The point derivative of any
smooth operator family is taken on this same fixed carrier. The actual
distribution realization and its normalization are part of the
definition.

The uniform symbol-family hypotheses provide actual representatives by
\hyperref[lem:online-finite-symbol-family-calculus]{the finite
symbol-family calculus}. They also ensure that differences with zero
complete expansion have jointly smooth kernels after all prescribed
finite parameter derivatives. Choosing unrelated fiberwise smooth
remainders would not define this quotient.

\pandocstatementheading{\leavevmode\phantomsection\label{thm:online-ambient-point-connection}Theorem 6.19 (The full ambient point connection).}

Actual differentiation in \(p\) induces a connection
\(\nabla^{\mathrm{amb}}\) on the union of the filtered sheaves
\(\mathscr M_a\). For a smooth point vector field \(X\), a smooth point
scalar \(f\), a smooth complete carrier-operator family \(P\) and a
module family \(m\), it satisfies

\phantomsection\label{eq:online-ambient-point-connection-laws} \[
\begin{aligned}
\nabla_X^{\mathrm{amb}}(fm)
 &=X(f)m+f\nabla_X^{\mathrm{amb}}m,\\
\nabla_X^{\mathrm{amb}}(Pm)
 &=(XP)m+P\nabla_X^{\mathrm{amb}}m,\\
\nabla_X^{\mathrm{amb}}F_\ell\mathscr M_a
 &\subset F_{\ell+1}\mathscr M_a,\\
[\nabla_X^{\mathrm{amb}},\nabla_Y^{\mathrm{amb}}]
 &=\nabla_{[X,Y]}^{\mathrm{amb}}.
\end{aligned}
\] The corresponding algebra connection is
\(\nabla_X^{\mathrm{alg}}P=XP\). In a fixed angular contact chart with
sky graph \(z=s(p,q)\), write a complete section as \(a u_s\) in the
rephased convention. Then

\phantomsection\label{eq:online-ambient-point-graph-formula} \[
\nabla_X^{\mathrm{amb}}(a u_s)
=\bigl(Xa-\tau(Xs)a\bigr)u_s.
\] These are identities on every complete module section. The connection
is determined by the actual realization. Equivalently, its algebra
Leibniz rule together with the actual derivative of the full elliptic
evaluation family \(e_a\) determines it uniquely.

\emph{Proof.} Differentiate an actual oscillatory representative. Its
derivative has the same phase and amplitude \(Xa+i(X\phi)a\).
Homogeneity of \(\phi\) raises the amplitude order by at most one.
Nondegeneracy and the parameter-uniform symbol conditions remain valid
on the same retained phase patch. A derivative of a jointly smooth
remainder is jointly smooth. Actual differentiation is independent of
the phase chosen to represent the distribution, so it induces a
well-defined sheaf operator with the asserted filtration loss. Its
linearity in \(X\) and scalar Leibniz identity follow before taking the
quotient.

Properly supported carrier operators act on these distribution families
by the actual parameter-dependent calculus. Differentiating their action
gives \(X(Pm)=(XP)m+P(Xm)\). Every term respects the joint smoothing
quotient. The equality therefore descends to the complete algebras and
modules and to their angular presentations. The algebra derivative obeys
its product rule and commutes with the internal adjoint, because that
adjoint is defined by finite coefficient operations independent of
\(p\).

In the graph presentation the physical Fourier amplitude is

\phantomsection\label{eq:online-ambient-physical-graph-amplitude} \[
e^{-i\tau s(p,q)}
\sum_{j\le m}i^j\tau^{j-1/2}a_j(p,q).
\] Differentiating its phase produces \(-i^{j+1}\tau^{j+1-1/2}(Xs)a_j\).
After rephasing degree \(j+1\), this is the second term in the graph
formula. In particular the complete coefficient at degree \(j\) becomes

\phantomsection\label{eq:online-ambient-point-coefficient-derivative} \[
a_j\longmapsto Xa_j-(Xs)a_{j-1},
\] with coefficients above the upper order set to zero.

The exact module actions in this chart are

\phantomsection\label{eq:online-ambient-graph-module-generators} \[
D_z=\tau,\qquad
D_{q^j}=\partial_{q^j}-s_{q^j}\tau,
\qquad
z=-\partial_\tau+\frac1{2\tau}+s.
\] They give a direct check of the operator identity. For \(D_{q^j}\),
the point derivative of \(-s_{q^j}\tau\) is canceled by the commutator
of \(-\tau Xs\) with \(\partial_{q^j}\). For \(z\), the point derivative
\(Xs\) is canceled by \([-\tau Xs,-\partial_\tau]=-Xs\). The
\(1/(2\tau)\) term is retained throughout, and \(D_z\) commutes with the
derivative. These checks use full module formulas, rather than their
principal symbols.

Ordinary derivatives of actual distribution families satisfy
\([X,Y]m=X(Ym)-Y(Xm)\). Passing to the joint quotient proves flatness.
In graph coefficients the possible curvature multiplier is \[
-\tau\bigl(X(Ys)-Y(Xs)-[X,Y]s\bigr)=0.
\] For the uniqueness statement let \(\nabla'\) obey the same algebra
Leibniz rule and have the same derivative on the full evaluation family.
The difference \(L_X=\nabla'_X-\nabla_X^{\mathrm{amb}}\) is linear over
point functions and commutes with every smooth carrier-operator family.
The evaluation amplitude has a nonzero leading coefficient of degree one
by \hyperref[prop:online-wave-full-evaluation-comparison]{the full
evaluation comparison}. In a graph chart write \(e_a=a_eu_s\) and
\(m=a_mu_s\). Complete Laurent division gives a smooth coefficient
series \(b=a_m/a_e\). Each coefficient is obtained by finitely many
divisions by the nonvanishing leading coefficient and finite operations
on the preceding coefficients. The normal-ordered operator
\(P=b(q,D_z;p)\) acts on \(a_eu_s\) by multiplication by
\(b(q,\tau;p)\), so \(m=Pe_a\) in that chart. A classical representative
of \(P\) exists by asymptotic summation. Therefore
\(L_Xm=L_X(Pe_a)=P L_Xe_a=0\). The local conclusion agrees on overlaps,
proving uniqueness. \hfill\(\square\)

Point differentiation raises the order by one, and flatness is the
commutator identity on the resulting filtered union. Parallel transport
of arbitrary module elements between distinct skies would require an
additional construction. The finite Cartan and spin connections retain
their respective curvatures. Realization by distributions fixes the
ambiguity left by algebra compatibility and flatness, namely addition of
a closed point one-form times the identity.

\subsection{6.10. Carrier and stored-chart
comparisons}\label{carrier-and-stored-chart-comparisons}\label{sec:online-point-connection-covariance}

\pandocstatementheading{\leavevmode\phantomsection\label{prop:online-ambient-carrier-covariance}Proposition 6.20 (Covariance in all point directions).}

The normalized carrier comparison of
\hyperref[thm:online-wave-carrier-cocycle]{the carrier cocycle theorem}
obeys

\phantomsection\label{eq:online-ambient-basic-carrier-comparison} \[
\nabla_X^{\mathrm{amb},a}(T_{ab}m)
=T_{ab}\nabla_X^{\mathrm{amb},b}m
\] on every smooth complete module family. It also intertwines the
algebra point connections under \(P\mapsto T_{ab}PT_{ba}\).

If a specified physical isometry \(f:B_0\to B_1\) compares two
recognized collars, compose half-density pushforward to the image
carrier with a normalized target-carrier comparison. Its induced map
\(\widehat F\) satisfies

\phantomsection\label{eq:online-ambient-physical-point-comparison} \[
\nabla_{f_*X}^{\mathrm{amb},1}\widehat F
=\widehat F\nabla_X^{\mathrm{amb},0}.
\] These constructions have exact restriction, identity, inverse and
composition laws on common retained parents. In point-dependent stored
coordinates \(m'=F_pm\), a connection written as \(X+A_X\) transforms by

\phantomsection\label{eq:online-ambient-stored-connection-law} \[
A'_X=F_pA_XF_p^{-1}-(XF_p)F_p^{-1}.
\] \emph{Proof.} The kernel defining \(T_{ab}\) uses the two fixed
carrier embeddings and the source collar. It is independent of the
evaluated point \(p\). Differentiation of its properly supported
oscillatory realization therefore commutes with its action on a moving
module family. Jointly smooth errors remain jointly smooth, which proves
the first identity. Differentiating \(T_{ab}PT_{ba}\) proves the algebra
identity for the same reason.

Under \(f\), naturality of the canonical scalar kernel and the actual
carrier half-density pushforward give the full carrier comparison
already established in
\hyperref[prop:online-wave-full-evaluation-comparison]{the evaluation
comparison}. The fixed operator depends on \(f\) and the carrier
embeddings, not on a separately variable evaluated point. Applying the
ordinary chain rule in that point proves the displayed physical
comparison. In particular, when \(f_*X\) has both normal and tangential
components, all those components occur in
\(\nabla_{f_*X}^{\mathrm{amb}}\). The connection has been defined in
every direction.

The carrier cocycle gives the identity, inverse and composition laws;
the actual constructions commute with restriction. Finally, if
\(m'=F_pm\), write the transported derivative as
\(F_p(X+A_X)(F_p^{-1}m')\). The derivative of \(F_pF_p^{-1}=I\) gives
\(F_p(XF_p^{-1})=-(XF_p)F_p^{-1}\), and expansion yields the stated
stored-chart formula. This also gives the full derivative term in
compositions of varying stored presentations. \hfill\(\square\)

The physical isometry in this proposition is supplied. Recognizing every
physical map from the complete source-side arrow equations is a separate
full-Hom assertion. Here the proposition constructs and verifies the
point-connection part of the induced comparison.

\subsection{6.11. The differential law and finite
estimates}\label{the-differential-law-and-finite-estimates}\label{sec:online-evaluation-wave-differential-law}

\pandocstatementheading{\leavevmode\phantomsection\label{prop:online-evaluation-wave-equation}Proposition 6.21 (The intrinsic wave equation of the full evaluation
family).}

The normalized evaluation family \(e_a=W_aC_{ap}\) satisfies

\phantomsection\label{eq:online-evaluation-ambient-wave-equation} \[
\Box_g^{\nabla^{\mathrm{amb}}}e_a=0,
\qquad
\Box_g^{\nabla^{\mathrm{amb}}}m
=|\det g|^{-1/2}
\nabla^{\mathrm{amb}}_{\partial_\mu}
\left(|\det g|^{1/2}g^{\mu\nu}
\nabla^{\mathrm{amb}}_{\partial_\nu}m\right)
\] in the joint complete-module quotient. The equation is coordinate
independent and is preserved by the physical comparisons above.

For a positive constant physical scale \(g^{[r]}=r^2g\), the
fixed-coordinate ambient connection is unchanged, and

\phantomsection\label{eq:online-ambient-wave-physical-scale} \[
e_a^{[r]}=r^{-1}e_a,
\qquad
\Box_{r^2g}^{\nabla^{\mathrm{amb}}}e_a^{[r]}
=r^{-3}\Box_g^{\nabla^{\mathrm{amb}}}e_a.
\] \emph{Proof.} The carrier restriction of the canonical kernel has no
wavefront covector with zero carrier component: a nonzero null covector
cannot annihilate the tangent space of a spacelike carrier. Its local
phase constructions therefore give the smooth distribution family
\(C_{ap}\) used in the definition. The bisolution test gives
\(\Box_{g,p}C_{ap}=0\) modulo a jointly smooth family. Here only the
carrier factor \(\rho_a(x)\) has been inserted; there is no density in
the scalar point slot.

The normalizing operator \(W_a\) acts on \(x\in\Sigma_a\) and is
independent of \(p\). It commutes both with point differentiation and
with the coefficients \(g^{\mu\nu}(p)\) and \(|\det g(p)|^{1/2}\) in the
displayed expression. Proper support makes its action preserve joint
smoothness. Thus \[
\Box_{g,p}(W_aC_{ap})
=W_a\Box_{g,p}C_{ap}=0.
\] The ambient-connection theorem identifies this actual derivative
calculation with the stated module equation. The scalar divergence
expression is invariant under point coordinate changes; the proposition
on all-direction covariance and preservation of the physical metric then
give its covariance.

For scaling, \hyperref[prop:online-wave-carrier-families-scale]{the
carrier scale law} gives \(T_{ab}^{[r]}=T_{ab}\) and
\(e_a^{[r]}=r^{-1}e_a\). Constant scaling preserves the conic null
relation and the underlying moving distribution realizations, so their
actual point derivative is unchanged. Since \(\Box_{r^2g}=r^{-2}\Box_g\)
and \(r\) is constant in the point variable, the second identity
follows. The physical-clock pullback of
\hyperref[prop:online-cartan-physical-scale]{Proposition 2.11} may
subsequently be applied by the already proved coordinate covariance.
\hfill\(\square\)

The wave equation is a differential law of the fully normalized
evaluation family. Its complete initial value is supplied by the
canonical kernel and carrier-root prescription. It is not a
horizontality condition. Its prolonged identities retain derivatives of
the metric coefficients. For example, write
\(b^\nu=|\det g|^{-1/2}\partial_\mu(|\det g|^{1/2}g^{\mu\nu})\).
Flatness in a point coordinate chart and differentiation of the wave
identity give

\phantomsection\label{eq:online-evaluation-prolonged-wave-equation} \[
\Box_g^{\nabla^{\mathrm{amb}}}
   \nabla^{\mathrm{amb}}_{\partial_\kappa}e_a
=-(\partial_\kappa g^{\mu\nu})
    \nabla^{\mathrm{amb}}_{\partial_\mu}
    \nabla^{\mathrm{amb}}_{\partial_\nu}e_a
  -(\partial_\kappa b^\nu)
    \nabla^{\mathrm{amb}}_{\partial_\nu}e_a.
\] The full comparison carries these identities and every finite point
jet because it intertwines the connection on all sections.

\pandocstatementheading{\leavevmode\phantomsection\label{prop:online-point-connection-finite-families}Proposition 6.22 (Finite estimates on one point-family parent).}

Take a smooth finite-dimensional family of admitted collars and carrier
embeddings on the common parents and with the positive margins of
\hyperref[prop:online-wave-carrier-families-scale]{the carrier-family
result}. Fix a finite angular cover, a homogeneous depth \(J\), a number
\(k\) of iterated coordinate point derivatives, a coefficient derivative
order \(m\) and an external parameter derivative order \(q\). On
retained compact subsets, the corresponding evaluation and comparison
coefficients satisfy

\phantomsection\label{eq:online-point-connection-finite-bound} \[
\begin{multlined}
\|\nabla^{\mathrm{amb},k}e_a\|_{m,q,J}
\\
\le F_{m,q,J,k}\left(
\mathcal N_N(\text{source coefficients and carrier embeddings}),
\mathfrak m^{-1}\right)
\end{multlined}
\] for some finite \(N\) and finite locally bounded function
\(F_{m,q,J,k}\). The filtration origin on the left is the upper degree
\(1+k\). The same kind of estimate holds for any stipulated finite list
of the connection formulas and their normalized comparisons, including
derivatives of supplied stored-coordinate maps with their inverse
margins. The common parent is chosen before \(m,q,J,k\).

\emph{Proof.} In a graph chart one derivative acts on coefficients by
the exact rule in
\hyperref[eq:online-ambient-point-coefficient-derivative]{the
coefficient formula}. Iteration uses only finitely many coefficients of
\(e_a\), finitely many derivatives of the same phase graph and the
coefficients of the selected vector fields. The top possible degree
increases by one at each step. Finite product estimates therefore bound
every requested coefficient and mixed derivative on the retained compact
chart.

The carrier-family result bounds the phase, evaluation and transition
coefficients by finite source and embedding seminorms with fixed
ellipticity, spacelike and domain-buffer margins. Finite-depth
stationary phase for an angular change uses finitely many earlier
coefficients and finite derivatives. The finite symbol-family calculus
consequently gives the same estimate on the finite angular cover. This
argument applies to the complete conormal phase families through points
on the carrier as well; a physical spatial cut is not needed.

Stored-coordinate formulas use composition, inversion on the supplied
nondegenerate domains and the explicit derivative term in the preceding
proposition. Their finite-depth estimates therefore use only the stated
finite extra seminorms and inverse margins. Jointly smooth quotient
equalities may be differentiated with all retained parameters, so
comparison identities do not require a new representative branch or a
smaller parent at every order.

A positive scale may depend on an external parameter. Its finite
derivatives and a positive lower bound are included in the input
seminorms and margins, and the chain rule gives the same conclusion.
Here we differentiate the supplied family with respect to its external
parameters. A canonical connection is defined only in the point
variables in this construction. For the point connection, Lemma 6.4
gives the claimed bounds on the coefficients at each prescribed
filtration depth. \hfill\(\square\)

\subsection{6.12. Independent scalar wave enrichment of a Cartan
collar}\label{independent-scalar-wave-enrichment-of-a-cartan-collar}\label{sec:online-independent-wave-enrichment}

The following admission isolates the scalar wave part over a geometric
carrier whose independent recognition was proved in Section 2. It is
used before the final category with all geometric and algebroid
comparison data is assembled.

\pandocstatementheading{\leavevmode\phantomsection\label{def:online-independent-wave-enrichment}Definition 6.23 (Independent wave enrichment on a retained carrier).}

Supply an independently admitted finite Cartan collar of
\hyperref[def:online-finite-cartan-collar]{Definition 2.7}, retained
spacelike carrier disks and buffered ray and point domains in a common
wave neighborhood \(W\) as in Definition 6.12. Form
\(g=\eta(\theta,\theta)\) from its supplied soldering. A normalized
scalar wave enrichment consists of the following supplied classes,
realizations and tests.

Supply a scalar kernel class \(\Lambda^+\) satisfying
\hyperref[def:online-independent-wave-kernel]{the independent kernel
tests}: in the joint smoothing quotient,

\phantomsection\label{eq:online-independent-enrichment-kernel-tests} \[
\Box_{g,p}\Lambda^+=\Box_{g,q}\Lambda^+=0,
\qquad
\Lambda^+-(\Lambda^+)^T=-iE_g,
\qquad
\operatorname{WF}'(\Lambda^+)\subset\mathcal C_g^+.
\] Here \(E_g\) is the retarded-minus-advanced scalar kernel on \(W\) of
the differential operator computed from \(\theta\), normalized by its
physical volume. The chosen positive null relation and its joint-family
condition are the ones in that definition. They are read from the source
null flow. No additional metric or wave state is supplied.

Supply complete carrier algebra and whole-sky module presentations with
specified invertible normalizations to the half-density algebra and
moving distribution modules above. Test these maps and their inverses on
the complete products, module actions, filtrations, internal adjoints
and homogeneous characteristic maps; test their overlap and
central-specialization identities in the stated quotient. These are
specified normalizations to the known cotangent and distribution models.
They are not an existential membership condition for an ordinary-pair
constructor.

Use the actual carrier half-densities to form \(C_{ab}\) and \(C_{ap}\)
from \(\Lambda^+\). Require the supplied \(C_a^{-1}\) and \(W_a\) to
satisfy their complete inverse and square-root identities, with the
positive physical principal symbol for \(W_a=C_a^{-1/2}\). Require the
complete normalized carrier maps to be \(T_{ab}=W_aC_{ab}W_b\). In each
supplied normalization the full point section and connections must obey

\phantomsection\label{eq:online-independent-enrichment-section-connection-tests}
\[
e_a=W_aC_{ap},\qquad
\nabla_X^{\mathrm{alg}}P=XP,\qquad
\nabla_X^{\mathrm{amb}}(a u_s)
=\bigl(Xa-\tau(Xs)a\bigr)u_s
\] for every point vector field and every complete module coefficient
family. The graph test is imposed in fixed angular actual realizations;
equivalent general phase realizations use actual point differentiation.
Point-dependent stored normalizations carry the full derivative term of
the stored-chart formula.

All statements with external parameters use the uniform classical family
and joint smoothing conditions already specified. Smoothing
representatives are not retained as additional data. The actual spin
bundle, its initial sign and the Cartan fields remain the supplied
geometric part of the object. The scalar normalization does not identify
their distinct spin lifts.

\pandocstatementheading{\leavevmode\phantomsection\label{thm:online-wave-enrichment-recognition}Theorem 6.24 (Construction and recovery of the wave enrichment).}

Every admitted finite Cartan collar has the wave enrichment of the
preceding definition on sufficiently small retained parents. The kernel,
complete carrier maps, full point section and all-direction point
connection are determined by the supplied geometric collar and the
specified actual normalizations.

Conversely, every independently admitted enrichment agrees, in its
supplied normalizations, with this construction from the same collar.
Constructing and then reading these normalized fields, or reading and
then constructing them, gives the identity on their complete classes.
The identifications commute with restriction, normalized carrier changes
and supplied physical isometries, and preserve smooth finite-dimensional
families with the finite estimates above.

In particular an ordinary-pair display followed by
\hyperref[cor:online-ordinary-pair-cartan-collar]{Corollary 2.12} has
this derived enrichment. Its wave fields require no further generating
class or variable Einstein input.

\emph{Proof.} The finite Cartan recognition theorem gives an actual
Lorentzian Ricci-flat metric from the supplied fields, with the retained
physical scale and actual spin structure. Choose buffered local domains
and carriers as in the carrier construction.
\hyperref[thm:online-canonical-wave-kernel]{The canonical-kernel
theorem} constructs a class satisfying the independent bisolution,
commutator and characteristic tests. For displaced hits the double
Cauchy construction preceding Lemma 6.13 extends this class to the
common wave neighborhood. Its actual representative families have the
required joint regularity.

The complete inverse and positive-principal square-root recursions give
\(C_a^{-1}\) and \(W_a\) uniquely. The carrier cocycle theorem and the
full evaluation comparison give \(T_{ab}\) and \(e_a\), including their
complete coefficients and all retained angular directions. The
ambient-connection theorem gives the required operator on every smooth
module family. Thus every admission condition is satisfied, with the
actual distribution normalizations supplied by the construction.

Now start with a supplied admitted enrichment. Kernel uniqueness
identifies its \(\Lambda^+\) with the constructed class, including joint
parameters. The trace-density formulas then identify every \(C_{ab}\)
and \(C_{ap}\). Uniqueness of the complete inverse and
positive-principal root identifies \(C_a^{-1}\) and \(W_a\), and hence
the supplied \(T_{ab}\) and full \(e_a\). The connection test on all
module families identifies the supplied connection with actual point
differentiation; the algebra test identifies the algebra connection. The
supplied normalizations and their inverse equations carry these
equalities back to each stored presentation. These arguments prove both
normalized identity laws without assuming representability in the
admission definition.

Restriction commutes with every construction. Carrier changes obey the
exact cocycles and all-direction connection identities already proved. A
supplied physical isometry carries the kernel by its independent
uniqueness and carries the remaining structures by the exact density and
mixed-kernel formulas, followed by the point chain rule. The
finite-family assertions follow by composition of
\hyperref[thm:online-common-collar-finite-families]{the common
geometric-collar estimates}, the kernel and carrier finite-coefficient
estimates, and the preceding point-connection estimates.

Finally, the ordinary construction supplies its finite Cartan collar
from the executed initial source. The metric used in this proof is its
readout \(g=\eta(\theta,\theta)\). The kernel and every subsequent field
are determined from that readout and its carrier markings. This proves
the last assertion. \hfill\(\square\)

This theorem concerns the normalized scalar enrichment over an admitted
geometric collar. Recognizing all abstract source arrows, carrying the
complete algebroid associators, and obtaining effective descent for the
final unmarked atlas require their full comparison results. The analytic
coupled Real operation also retains its separate structured comparison;
it is not inferred from the scalar wave construction.

\section{7. Geometric admission and full source
comparisons}\label{geometric-admission-and-full-source-comparisons}\label{sec:online-complete-source-comparisons}

The scalar enrichment of Section 6 is covariant under a supplied
physical isometry. For independently supplied sources, the labelled
skies and the physical density recover the underlying physical map from
a source comparison. We then use the full evaluation family over an open
neighborhood of the point to show that the physical map and its
specified spin lift determine the scalar and finite-rank comparison.

All geometric representatives are smooth, Hausdorff and second
countable. Point cores are connected and have buffered parents; null
rays mean their retained connected plaques. When the point core is
restricted, these plaques and their local quotients are taken again.
Scalar equalities use the complete calculus and joint smoothing quotient
of Section 6. Actual spin maps and geometric connections are not taken
modulo that quotient.

\subsection{7.1. Independent real incidence over the Cartan
collar}\label{independent-real-incidence-over-the-cartan-collar}\label{sec:online-actual-incidence-admission}

\hyperref[ref-Low2001]{Low} studied the geometry of the space of null
geodesics, and \hyperref[ref-BautistaIbortLafuente2014]{Bautista, Ibort
and Lafuente} gave a reconstruction theorem using light rays and skies.
Here the point manifold and the comparison map \(f\) on points are
supplied. On the retained ray plaques the null cones determine the
conformal class. If \(f^*g_1=cg_0\) with \(c>0\), the point density
scales by \(c^2\) in four dimensions. Its preservation gives \(c^2=1\)
and hence \(c=1\).

\pandocstatementheading{\leavevmode\phantomsection\label{def:online-complete-geometric-incidence}Definition 7.1 (Geometric incidence tests).}

Let an admitted finite Cartan collar have point manifold \(B\),
soldering \(\theta\), connection \(A\) and metric
\(g=\eta(\theta,\theta)\). Supply a cooriented contact five-manifold
\((Z,\mathcal H)\) and a double fibration

\phantomsection\label{eq:online-geometric-double-fibration} \[
B\xleftarrow{\ p\ }F\xrightarrow{\ q\ }Z.
\] Here \(F\) has dimension six, \(p\) is a proper surjective submersion
with connected fibers \(S^2\), and \(q\) is a surjective submersion with
connected one-dimensional fibers. Each \(q|_{F_b}\) is an embedded
Legendrian sphere. Supply a positive point density \(\mu\) and require
\(\mu=\operatorname{vol}_g\).

For every positive local contact form \(\alpha\), require the semibasic
covector \(\ell\) defined by

\phantomsection\label{eq:online-incidence-normal-covector} \[
q^*\alpha=\ell\circ\mathrm dp
\] to be nonzero, past directed and null for \(g\) on every point fiber.
Changing \(\alpha\) by a positive ray function changes \(\ell\) by that
function. These tests use the supplied incidence and the readout of its
supplied Cartan fields. They require neither an ordinary-class
presentation nor a weighted inverse.

A geometric comparison supplies an orientation-preserving point
diffeomorphism \(f:B_0\to B_1\), a coorientation-preserving contact
diffeomorphism \(\Psi:Z_0\to Z_1\), and their incidence comparison, with

\phantomsection\label{eq:online-geometric-arrow-tests} \[
f^*\mu_1=\mu_0,\qquad
\Psi\bigl(q_0(F_{0,b})\bigr)=q_1(F_{1,f(b)}).
\] An actual principal spin comparison is separately retained, with its
compatibility with the supplied soldering and frame reductions. No
metric-isometry equation is included as a substitute for the incidence
and density tests.

\pandocstatementheading{\leavevmode\phantomsection\label{lem:online-incidence-identification}Lemma 7.2 (Full null incidence).}

The preceding tests identify \(F\) with the positive projectivization of
the past null covectors of \(g\). They identify \(\ker\mathrm dq\) with
the null-geodesic line distribution and \(Z\) with its quotient by the
retained connected leaves. In particular each supplied sky contains
every real null direction exactly once. Conversely every admitted Cartan
collar supplies this incidence after restriction to a sufficiently small
causally convex point core with a buffered full-sky ray parent.

\emph{Proof.} Put \(\lambda=q^*\alpha\). It is nonzero because \(q\) is
a submersion. It vanishes on \(\ker\mathrm dp\) by the Legendrian
hypothesis, so the covector \(\ell\) exists smoothly and is unique. The
radical of \(\mathrm d\lambda\) on \(\ker\lambda\) is
\(\ker\mathrm dq\): the differential of \(q\) maps that hyperplane onto
\(\mathcal H\), where \(\mathrm d\alpha\) is symplectic.

The map \(u\mapsto[\ell_u]_+\) is an immersion on a point fiber. Indeed,
if \(v\in\ker\mathrm dp\) differentiates \(\ell\) only by
\(v(\ell)=c\ell\), then for every \(w\in\ker\lambda\),
\(\mathrm d\lambda(v,w)=c\lambda(w)=0\). Thus \(v\in\ker\mathrm dq\).
The sky embedding has injective vertical derivative, so \(v=0\). Its
image lies in the past projective null sphere of \(g\). The fiber map is
consequently a local diffeomorphism between two-dimensional manifolds.
Compactness makes its image closed as well as open. It covers the whole
sphere, and is a one-sheeted covering because that sphere is simply
connected and the source fiber is connected. These fiberwise
diffeomorphisms give a diffeomorphism over \(B\).

On \(T^*B\setminus0\), let \(\vartheta=\xi_a\,\mathrm db^a\) and
\(G=\tfrac12g^{ab}\xi_a\xi_b\), with
\(\iota_{X_G}\mathrm d\vartheta=-\mathrm dG\). On \(G=0\), the radical
of \(\mathrm d\vartheta\) on the common kernel of \(\vartheta\) and
\(\mathrm dG\) is spanned by the radial Euler field and \(X_G\).
Positive projectivization removes the Euler field. The remaining
characteristic line pulls back to the radical of \(\mathrm d\lambda\)
just computed. Hamilton's equations for \(G\) are the affinely
parametrized null-geodesic equations. This proves the line-distribution
assertion.

The connected fibers of \(q\) are precisely the maximal connected leaves
within the retained domain. A surjective submersion is open, and its
local sections identify its quotient topology and smooth structure with
those of \(Z\). Thus this is the asserted local null-ray quotient, with
its contact distribution.

For existence, take the actual null Hamiltonian flow of the recognized
metric on a buffered small Gaussian collar. A spacelike disk transverse
to that flow supplies its ray carrier as an open subset of
\(P_+T^*\Sigma\). Choose a smaller causally convex core so that every
direction over it has one transverse hit inside the buffered disk;
compactness of the sky and smooth transverse dependence permit this
choice. Put \(F=P_+\mathcal N_g^-|_B\) and let \(q\) be the hit map.
Projection to \(B\) is a proper sphere bundle. The hit map is a
submersion, and its fibers are the connected null plaques; causal
convexity and the short flow parent exclude separated returns. Variation
of a ray through a fixed point gives an embedded sky: injectivity
follows from flow uniqueness, and the vertical derivative is injective
by the same uniqueness in phase space. The tautological form vanishes on
variations at that point, so the sky is Legendrian. The ray contact
structure is the quotient of the tautological hyperplane field;
Hamiltonian transport preserves it. Its positive contact frames give
precisely the past null covectors required in the definition. Taking
\(\mu=\operatorname{vol}_g\) completes all tests. \hfill\(\square\)

\pandocstatementheading{\leavevmode\phantomsection\label{thm:online-incidence-physical-recognition}Theorem 7.3 (Physical-map recognition from the source).}

A geometric comparison in Definition 7.1 is a
time-orientation-preserving local isometry of the recognized metrics.
Its ray map is uniquely determined by its point map. Conversely, every
orientation- and time-orientation-preserving local isometry, with
specified actual spin lift, induces the unique comparison of these
incidence and density fields on common retained domains.

\emph{Proof.} The map \((p,q):F\to B\times Z\) is an embedding. Its
derivative is injective because its kernel is vertical for \(p\), where
the sky map is an immersion. It is injective because each sky is
embedded, and is proper because \(p\) is proper. The sky equation
therefore induces a smooth incidence diffeomorphism \(\Phi\) covering
\(f\) and \(\Psi\).

Write \(\Psi^*\alpha_1=a\alpha_0\), with \(a>0\). Pulling back the
contact-normal equation gives

\phantomsection\label{eq:online-incidence-arrow-covector-law} \[
\ell_{1,\Phi(u)}\circ\mathrm df_b
=a(q_0(u))\ell_{0,u}.
\] Lemma 7.2 shows that these are the complete past null cones. Their
quadratic annihilator is one-dimensional. To see this directly in a
Lorentz frame, a quadratic form vanishing on \((1,n)\) for every
\(|n|=1\) has the form \(a_0+2c\cdot n+n^{\mathsf T}Dn\). Comparison at
\(n\) and \(-n\) gives \(c=0\); evaluation at the coordinate vectors and
their pairwise normalized sums gives \(D=-a_0I\). Thus it is a multiple
of the Lorentz form. Consequently \(f^*g_1=c_0g_0\) with \(c_0>0\). In
dimension four its density is \(c_0^2\mu_0\). Preservation of \(\mu\)
forces \(c_0=1\). Positivity in the covector law preserves the selected
past component, hence time orientation.

The covector lift is now necessarily
\((b,[\xi]_+)\mapsto(f(b),[(\mathrm df_b^{-1})^*\xi]_+)\). It determines
the incidence map and, by surjectivity of \(q_0\), the ray map.
Conversely this lift of an isometry preserves the tautological
hyperplane and the null-geodesic line. It descends to the retained
connected-leaf quotients. Local sections of \(q_0,q_1\) make the
descended maps and their inverses smooth. They are contact and obey the
sky and density equations. The specified principal spin map is included
unchanged. Composition and restriction are the corresponding actual
covector and principal-bundle operations. \hfill\(\square\)

The use of the entire sky is essential. Finitely many directions can
recover a candidate quadratic form, but do not test its equality on a
supplied arbitrary sphere family. The density also matters: the contact
incidence alone determines the conformal class, rather than its
numerical metric scale.

\subsection{7.2. Optical scale and the canonical finite-rank
fields}\label{optical-scale-and-the-canonical-finite-rank-fields}\label{sec:online-optical-and-finite-rank-fields}

\pandocstatementheading{\leavevmode\phantomsection\label{prop:online-canonical-optical-completion}Proposition 7.4 (Actual optical and affine completion).}

The incidence of Definition 7.1 has a canonical physical optical scale
with value \(y=1\). For a contact form \(\alpha\), there is a unique
null generator \(D\) with \(\mathrm dp(D)=g^{-1}\ell\). It is affine for
\(g\). If the tangent-sky curve along a ray is written in a constant
symplectic graph chart as a symmetric matrix \(S(t)\), then \(S'\) is
invertible and

\phantomsection\label{eq:online-canonical-optical-potential} \[
\kappa_\alpha
=\frac12\operatorname{tr}\left(
\frac12(S')^{-1}S'''-
\frac34\bigl((S')^{-1}S''\bigr)^2\right)
=\frac12\operatorname{Ric}_g(k,k)=0,
\qquad k=g^{-1}\ell.
\] The selected value is constant on the whole point sky and solves
\(D^2y+\kappa_\alpha y=0\). Its value, slope and positive branch
determine the selected optical states and roots. Initialized fundamental
matrices and companion solutions give the full affine transition maps.
They have exact inverse and cocycle laws and are natural for the
comparisons of Theorem 7.3.

\emph{Proof.} Lift the incidence to \(T^*B\setminus0\) by
\(u\mapsto(p(u),\ell_u)\). Its Liouville pullback is \(q^*\alpha\). The
lift of \(D\) and \(X_G\) have the same base projection. Their
difference is vertical and pairs to zero under \(\mathrm d\vartheta\)
with every lifted incidence tangent. Surjectivity of \(\mathrm dp\)
makes that vertical difference zero. Thus \(D\) is the affine
Hamiltonian generator, including its covector normalization.

The contact plane along a ray is the four-dimensional space of screen
Jacobi solutions; the tangent sky at a point is the kernel of their
evaluation there. In a parallel orthonormal screen frame, let \(Y,Z\)
solve \(j''+\mathcal Rj=0\) with identity and zero initial position and
velocity blocks. Their symplectic Wronskian identities give, where \(Y\)
is invertible, \[
S=-Y^{-1}Z,\qquad
S'=-Y^{-1}Y^{-\mathsf T},\qquad
A_0:=(S')^{-1}S''=-2(Y^{-1}Y')^{\mathsf T},
\] \phantomsection\label{eq:online-optical-jacobi-calculation} \[
\frac12A_0'-\frac14A_0^2
=(Y^{-1}\mathcal RY)^{\mathsf T}.
\] The screen tidal operator is \(\mathcal R[J]=[R(J,k)k]\). Its trace
is \(\operatorname{Ric}_g(k,k)\): the two extra null slots in the full
curvature trace vanish. Taking the trace proves the formula. Any
constant symplectic change gives another fundamental pair with the same
Wronskians and Jacobi equation, so the scalar is independent of the
graph chart. Such charts cover the curve. Ricci-flatness was already
recognized from the actual Cartan source in Section 2, and gives its
last equality.

For \(\alpha_i=c_{ij}\alpha_j\), the factor \(c_{ij}\) is pulled back
from the ray. Hence \(D_i=c_{ij}D_j\) and \(D_jc_{ij}=0\), so
\(\kappa_i=c_{ij}^2\kappa_j\). The scalar optical equations agree on
overlaps. Let \(F_j(t)\) be the identity-initialized fundamental matrix
for the state \((y,D_jy)\). A change of hit and contact frame has matrix
\(B_{ij}=\operatorname{diag}(1,c_{ij})F_j(\tau_{ij})\), with the hit
time and coefficients evaluated in their moving arguments. Uniqueness of
that finite ODE proves its inverse and triple identities.

More generally write the selected nonzero solution in an auxiliary scale
as \(\sigma\), and put

\phantomsection\label{eq:online-optical-affine-states} \[
G_i=\begin{pmatrix}\sigma_i&0\\D_i\sigma_i&\sigma_i^{-1}\end{pmatrix},
\qquad
(G_i\circ\Phi_{ij})^{-1}B_{ij}G_j
=\begin{pmatrix}1&b_{ij}\\0&c_{ij}\end{pmatrix}.
\] The first column follows from preservation of the selected solution;
the second diagonal entry follows from the determinant. The
companion-to-selected ratio \(\tau_i\) satisfies
\(D_i\tau_i=\sigma_i^{-2}\). Its affine change has slope \(c_{ij}^{-1}\)
and translation \(-b_{ij}/c_{ij}\). These give all affine entries, not
only their linear parts. In the physical splitting \(\sigma=1\) and the
selected state is \((1,0)\).

For a physical incidence comparison with contact multiplier \(d\), the
physical affine generators obey \(\Phi_*D_0=d^{-1}D_1\). Therefore
\(\tau_1\circ\Phi-d^{-1}\tau_0\) is constant on the connected ray
plaque. Evaluate it at a retained hit to obtain the affine translation.
It is independent of the hit, and these translations compose exactly.
The positive root is the specified physical branch, with its initialized
sign. The same ODE uniqueness carries every optical state and its
inverse.

If a stored chart uses \(\widehat g=\Omega^2g\), its selected value is
\(\widehat y=\Omega\). With \(u=D\log\Omega\), direct differentiation
gives \(\widehat D=\Omega^{-2}D\) and
\(\widehat\kappa=\Omega^{-4}(\kappa-Du+u^2)\). Thus
\((\widehat D^2+\widehat\kappa)(\Omega y)=\Omega^{-3}(D^2+\kappa)y\).
This transports the same physical optical scale, rather than supplying
an independent direction-dependent solution. Every construction above
uses actual flow, differentiation and finite linear ODEs on the retained
tubes, which proves its naturality. \hfill\(\square\)

\pandocstatementheading{\leavevmode\phantomsection\label{prop:online-canonical-finite-rank-models}Proposition 7.5 (Canonical tractor, local-twistor and screen data).}

The supplied Cartan soldering and actual spin connection determine the
following finite-rank models in the physical Ricci-flat splitting. On
\(\mathcal T=\mathbb R\oplus\mathbb E\oplus\mathbb R\) set

\phantomsection\label{eq:online-canonical-tractor-model} \[
\begin{aligned}
H((\sigma,v,\rho),(\sigma',v',\rho'))
&=\eta(v,v')+\sigma\rho'+\sigma'\rho,\\
\nabla_X^{\mathcal T}(\sigma,v,\rho)
&=(X\sigma-\eta(\theta_X,v),\ D_Xv+\rho\theta_X,\ X\rho),\\
\mathsf I&=(1,0,0),\qquad \mathsf X=(0,0,1).
\end{aligned}
\] It preserves \(H\), makes \(\mathsf I\) parallel, and has
\(H(\mathsf I,\mathsf X)=1\). Here \(D\) is the Lorentz connection
induced by the actual spin connection; it is not the affine ray
generator of the preceding proposition.

These are the explicit Ricci-flat normalization formulas for the
finite-rank diagram (\hyperref[ref-IsthmusLocal2026]{Isthmus 2026a},
Proposition I.3). In the ordered rank-four spinor sum
\((\omega^A,\pi_{A'})\), put
\(\vartheta(v)=2^{-1/2}(v^0I+v^a\sigma_a)\), where \(\sigma_a\) are the
Pauli matrices, and
\(J=\left(\begin{smallmatrix}0&1\\-1&0\end{smallmatrix}\right)\). With
this soldering matrix convention, set

\phantomsection\label{eq:online-canonical-local-twistor-model} \[
\nabla_X^{\mathrm{LT}}
=P_X+\begin{pmatrix}0&i\vartheta(\theta_X)\\0&0\end{pmatrix},
\qquad
\mathcal F(\sigma,v,\rho)
=\begin{pmatrix}
\rho J&-iJ\vartheta(v)\\
-i\vartheta(v)^{\mathsf T}J&\sigma J
\end{pmatrix}.
\] The map \(\mathcal F\) intertwines the tractor connection with the
induced connection on covariant alternating two-forms. Its normalization
is \(2\operatorname{Pf}\mathcal F=2\sigma\rho+\eta(v,v)\). The rank-two
infinity pair is

\phantomsection\label{eq:online-canonical-infinity-pair} \[
I_{\mathrm{cov}}=\begin{pmatrix}0&0\\0&J\end{pmatrix},\qquad
I_{\mathrm{contra}}=\begin{pmatrix}J&0\\0&0\end{pmatrix},\qquad
I_{\mathrm{cov}}I_{\mathrm{contra}}=0.
\] Both are parallel, as are the spin volume and the Hermitian pairing
with matrix
\(\left(\begin{smallmatrix}0&I\\I&0\end{smallmatrix}\right)\). For a
nonzero null \(k\), the optical plane is
\(\operatorname{span}\{\mathsf X,(0,k,0)\}\), with affine slice
\(\mathsf X+\mathbb R(0,k,0)\); its complexification has the
corresponding complex affine slice and spinor pencil. Along an affine
null ray these data obey \(\nabla_k\mathsf X=(0,k,0)\),
\(\nabla_k(0,k,0)=0\).

The screen is \(k^\perp/\mathbb Rk\) with positive metric \(-\eta\), its
induced orientation and quarter-turn, and its induced connection. Its
complete Jacobi state satisfies \(D_kY=Q\), \(D_kQ=-\mathcal RY\). These
constructions and their actual bundle maps are natural under physical
comparisons with specified spin lifts.

\emph{Proof.} Differentiating \(H\) cancels the two soldering terms
against the two \(\rho\) terms; the remaining equality is \(D\eta=0\).
Direct substitution gives the assertions about \(\mathsf I,\mathsf X\).
The torsion identity \(D\theta=0\) gives curvature
\(R^{\mathcal T}(X,Y)(\sigma,v,\rho)=(0,R^D(X,Y)v,0)\). This is the
normal tractor formula since the Schouten tensor is zero in the
recognized Ricci-flat metric.

The matrix convention gives \(2\det\vartheta(v)=\eta(v,v)\) and
\(\operatorname{tr}(\vartheta(\theta_X)\operatorname{adj}\vartheta(v))=\eta(\theta_X,v)\).
For the upper triangular soldering term \(\Omega_X\), multiplication of
the displayed matrices gives

\phantomsection\label{eq:online-local-twistor-intertwining-check} \[
-\Omega_X^{\mathsf T}\mathcal F(\sigma,v,\rho)
-\mathcal F(\sigma,v,\rho)\Omega_X
=\mathcal F(-\eta(\theta_X,v),\rho\theta_X,0).
\] The spin part is equivariant because \(A^{\mathsf T}JA=J\) for
determinant-one spin frames. This proves the full connection
intertwining. Expanding the four-by-four Pfaffian gives
\(\sigma\rho+\det\vartheta(v)\). The selected infinity form is
\(\mathcal F(\mathsf I)\), and its volume dual is the displayed
contravariant form. Their ranks, product and parallelism follow from
their matrices and volume preservation. For real \(X\),
\(\vartheta(\theta_X)\) is Hermitian, so
\(\Omega_X^\dagger H_{\mathrm{LT}}+H_{\mathrm{LT}}\Omega_X=0\); the spin
part preserves that pairing as well.

The two generators of the optical plane are independent and have all
three mutual pairings zero. Evaluation by \(\mathsf I\) gives its affine
slice. Every nonzero alternating form in its image under \(\mathcal F\)
has rank two. In a spin frame where \(\vartheta(k)\) has one nonzero
diagonal entry, its kernels contain one common line and lie in one
common three-plane; solving the remaining linear equation gives the
stated projective pencil. The tractor derivative formulas and the Jacobi
system follow directly from affine geodesic transport and the induced
screen curvature.

Every formula uses the supplied soldering, spin connection and fixed
determinant normalization. A physical map with its actual spin lift
preserves precisely these data, hence intertwines all the formulas. The
central spin element acts by \(-1\) on the rank-four bundle and by
\(+1\) on its alternating square. Retaining the rank-four map keeps
those two lifts distinct. The rank-two optical plane and the
four-dimensional screen Jacobi state have different roles and are not
identified. \hfill\(\square\)

\pandocstatementheading{\leavevmode\phantomsection\label{lem:online-initialized-finite-rank-test}Lemma 7.6 (Independent test for supplied finite-rank fields).}

Let a retained null flow tube have generator \(K_0\), transverse section
\(S\) and finite-rank connections \(D^g,D^s\) on two supplied bundles.
Fix an actual initial isomorphism \(J_0\). Define \(J\) by

\phantomsection\label{eq:online-finite-rank-initialized-test} \[
D^s_{K_0}J-JD^g_{K_0}=0,\qquad J|_S=J_0.
\] It intertwines the full connections if and only if their tangential
initial connections agree under \(J_0\) and \(R^s(K_0,W)J=JR^g(K_0,W)\)
throughout the actual tube for every transverse direction \(W\).
Parallel tensors agree if they agree initially. Selected tractor
positions and screen Jacobi states agree under the additional actual
equations and initial values of Proposition 7.5, including equality of
the transported tidal endomorphisms. These are independent tests for
supplied copies of precisely those canonical models.

\emph{Proof.} If \(P^g_{t,0}\) and \(P^s_{t,0}\) are parallel transports
on the flow lines, the unique candidate is
\(J=P^s_{t,0}J_0(P^g_{t,0})^{-1}\). For a flow-transported transverse
frame \([K_0,W]=0\), its transverse defect \(E(W)=D^s_WJ-JD^g_W\)
satisfies

\phantomsection\label{eq:online-finite-rank-transverse-defect} \[
\nabla_{K_0}^{\operatorname{Hom}}E(W)
=R^s(K_0,W)J-JR^g(K_0,W).
\] The two tests make this a finite homogeneous ODE with zero initial
value. Thus the defect vanishes in every direction. Necessity follows by
taking the curvature of a connection intertwiner. If a different frame
is used, the curvature formula contains its full bracket term
\(-D_{[K_0,W]}\); omitting that term is not an equivalent test.

Tensor products, contractions, duality and the stipulated real
structures commute with parallel transport. Their initially zero
differences therefore remain zero. For positions, the difference system
is \(D e_X=e_V\), \(D e_V=0\), \(D e_I=0\). For Jacobi states it is
\(D e_Y=e_Q\), \(D e_Q=-\mathcal R^s e_Y\) after the tidal test. Both
are finite initialized systems. Principal spin transport gives the same
assertion for an actual initial principal spin map and preserves its
central sign.

These arguments also prove restriction, inverse and composition laws. At
a fixed derivative order their differentiated finite ODEs give finite
estimates on the same buffered tube, with constants depending on that
order, the finite coefficient norms and the retained flow margins. The
curvature tests hold throughout the tube in every retained angular
direction. \hfill\(\square\)

\subsection{7.3. The complete independently admitted
source}\label{the-complete-independently-admitted-source}\label{sec:online-independent-full-wave-source}

\pandocstatementheading{\leavevmode\phantomsection\label{def:online-complete-independent-wave-source}Definition 7.7 (Complete local wave source and its arrows).}

An object consists of the following supplied fields and normalizations,
subject to the indicated equations.

Its initial gradients satisfy Definition 2.1 with maximal readout, and
its initialized finite Cartan collar satisfies Definition 2.7 with
exactly those initial values. Retain its actual principal spin bundle,
physical scale and future orientation. The Cartan fields determine \(g\)
by Theorem 2.9. Supply the incidence of Definition 7.1, the physical
optical states and affine transitions of Proposition 7.4, and the
canonical finite-rank diagrams of Proposition 7.5. Independently
supplied realizations of those diagrams have specified initial
comparison maps and satisfy Lemma 7.6, including the selected tensor,
position and tidal tests. No additional unnormalized finite-rank object
is included.

The object supplies the normalized scalar enrichment of Definition 6.23.
The normalization maps for its fixed-carrier algebra and whole-sky
modules are supplied with their inverses, full homogeneous
characteristic maps, native half-density roots and
central-specialization maps. They satisfy the inverse equations. The
kernel, full evaluation section and point connections satisfy the
independent tests of that definition.

The parent collar contains the marked maximal initial patch and carries
the initialized evolution. The pointed spacetime core is a connected
open restriction of this development. Its marked point may lie away from
the initial surface. Germ comparisons restrict the core at that point.
The parent retains the portions of evolution and the transverse carrier
hits needed to define the compared fields. The choice of maximal surface
and Gaussian clock belongs to the presentation and may change under an
arrow. An auxiliary scalar carrier may be displaced from the point core,
provided its retained hits and segments lie in the common wave
neighborhood required by Definition 6.12.

Stored presentations retain their complete algebra comparisons
\(\varphi_{ij}\), normalized associators \(C_{ijk}\) and semilinear
module comparisons \(\tau_{ij}\). The supplied normalizations, including
their overlap comparison cells, identify them with the corresponding
cotangent and actual distribution presentations of Section 6. In these
normalizations the supplied maps and overlap cells satisfy the complete
identities for products, inverses, filtration, internal adjoints, roots,
characteristic maps, sections, point connections and central
specializations. Their descent equations include

\phantomsection\label{eq:online-stored-source-descent-equations} \[
\begin{aligned}
\varphi_{ij}\varphi_{jk}&=\operatorname{Ad}(C_{ijk})\varphi_{ik},\\
C_{ijk}C_{ikl}&=\varphi_{ij}(C_{jkl})C_{ijl},\\
\tau_{ij}\tau_{jk}&=C_{ijk}\tau_{ik}.
\end{aligned}
\] The normalization data include the native density-line squares and
their chosen branches. The degree-one evaluation and its connection in
every point direction are those of Section 6. They have no independent
scalar connection or degree-zero full metric vector as additional data.
The object retains the complete classes modulo joint smoothing.

An arrow supplies the actual geometric comparison of Definition 7.1 and
its principal spin map, together with filtered
internal-adjoint-preserving algebra maps and semilinear whole-module
maps over its homogeneous ray comparison. On every graded piece the
algebra map is the full homogeneous cotangent lift of that ray
comparison. A freely adjustable radial lift of a projective map is not
allowed. The maps preserve the canonical optical and finite-rank data,
native density roots, full evaluation sections and both algebra and
module point connections on every section. Thus, with pullback by the
point map \(f\) understood,

\phantomsection\label{eq:online-complete-source-arrow-point-tests} \[
\Xi e_0=e_1\circ f,\qquad
\nabla^{\mathrm{amb},1}_{f_*X}\Xi=\Xi\nabla^{\mathrm{amb},0}_X,\qquad
\nabla^{\mathrm{alg},1}_{f_*X}k=k\nabla^{\mathrm{alg},0}_X.
\] In stored charts its vertices \(k_i,\Xi_i\) and normalized overlap
cells \(T_{ij}\) obey

\phantomsection\label{eq:online-complete-source-two-associator-tests} \[
\begin{aligned}
k_i\varphi^0_{ij}
 &=\operatorname{Ad}(T_{ij})\varphi^1_{ij}k_j,\\
T_{ij}\varphi^1_{ij}(T_{jk})C^1_{ijk}
 &=k_i(C^0_{ijk})T_{ik},\\
\Xi_i\tau^0_{ij}
 &=T_{ij}\tau^1_{ij}\Xi_j.
\end{aligned}
\] Both source and target associators occur here. The stored sections
and connections are carried by the same vertices; for a point-dependent
module normalization \(F\), its connection matrix is
\(A'_X=FA_XF^{-1}-(XF)F^{-1}\). A representative sends the marked
spacetime point to the marked spacetime point. Two representatives
define the same arrow germ if on a common smaller pointed core their
point and principal spin maps agree and, on a common refinement of the
tested ray charts there, their remaining data agree up to the normalized
presentation cells. The transverse transports are performed on the
corresponding buffered parents. Write
\(\mathsf{WaveSrc}^{\mathrm{loc}}_0\) for this groupoid.

More explicitly, compare vertices on the source core \(B\) after pulling
target coefficient sheaves back by \(f\) and reindexing passive
parameters by the supplied base map (the identity for families over the
same base). Algebra coefficients basic in the point remain basic. The
module maps are over the point map: multiplication by a smooth point or
passive function on the target acts on the source by its pullback. The
tested conic domain on the source carrier is
\(\Omega_B=\bigcup_{p\in B}L_p\); its image is the tested domain on the
target. Restrict all charts and overlap equations to these open sets.
With the point and actual spin maps fixed, pass to one common
refinement. A normalized change of target presentation by units
\(U_i\in 1+F_{-1}\), \(U_i^*U_i=1\), acts on the vertices and overlap
cells by \[
\begin{aligned}
k'_i&=\operatorname{Ad}(U_i) k_i,\qquad \Xi'_i=U_i\Xi_i,\\
T'_{ij}&=U_iT_{ij}\varphi^1_{ij}(U_j^{-1}).
\end{aligned}
\] The sections and connections are expressed in the changed
presentations, including the derivative term above. This formula follows
by substitution in the three overlap equations; a change of source
presentation is pulled through the invertible vertices. Refinements and
these compatible changes identify representatives. Once the
presentations are fixed, Lemmas 7.8--7.10 test the remaining ambiguity
on \(\Omega_B\). Equality also requires equality of the actual principal
spin maps, including their central sign.

The complete algebra and module maps preserve their displayed
filtrations and act on the coefficientwise completions. Algebra
comparisons are morphisms of the smooth coefficient sheaves. Smooth
dependence on additional parameters means that their maps on every fixed
finite filtration quotient act smoothly on smooth coefficient families,
as do the supplied inverse maps and module normalizations. Admission
uses only equations on supplied source fields, actual maps and complete
coefficient families. It includes no ordinary pair, weighted smallness,
global inverse or membership in a constructor image. The original
analytic calibrated Hom-line and its coupled Real operation are
addressed by the later analytic comparison; they are not identified with
the new evaluation by this definition.

\subsection{7.4. Normalized algebra and module
ambiguities}\label{normalized-algebra-and-module-ambiguities}\label{sec:online-normalized-comparison-rigidity}

\pandocstatementheading{\leavevmode\phantomsection\label{lem:online-normalized-inner-rigidity}Lemma 7.8 (Normalized inner comparison).}

Let \(\Phi\) be an automorphism of the complete algebra of Proposition
6.2, covering the identity, fixing passive functions, preserving the
internal adjoint and acting identically on every graded piece. There is
a unique \(U\in1+F_{-1}\) with \(\Phi=\operatorname{Ad}(U)\). It
satisfies \(U^*U=1\). For smooth comparison families, its every finite
coefficient is smooth and determined by finitely many coefficients and
derivatives of \(\Phi\). The construction commutes with restriction.

\emph{Proof.} Suppose the first nonzero difference of \(\Phi\) lowers
order by \(k\). Its leading action \(\delta\) is a derivation of the
commutative homogeneous algebra, by multiplicativity. Preservation of
commutators makes it a Poisson derivation. The leading internal adjoint
on degree \(m\) is \((-1)^m\). Commutation with that adjoint gives
\((-1)^m\delta a=(-1)^{m-k}\delta a\). Thus the difference vanishes at
every odd \(k\), and the first possible \(k\) is even and at least two.

The local derivation is a smooth homogeneous vector field \(V\) on the
symplectic cone. Indeed, the Hadamard formula for a smooth function
minus its affine first-order part makes a derivation depend only on the
coordinate derivatives; apply this in degree-zero angular coordinates
and an invertible radial coordinate. The Poisson equation says that
\(V\) is symplectic, and \([E,V]=-kV\) for the Euler field. If
\(\lambda=\xi\,\mathrm dx\) and \(\omega=\mathrm d\lambda\), then
\(\iota_V\omega\) is closed and has Euler weight \(1-k\). Cartan's
formula gives the actual primitive

\phantomsection\label{eq:online-negative-weight-inner-primitive} \[
b=\frac{\lambda(V)}{1-k},\qquad
\mathrm db=-\iota_V\omega,\qquad
\delta a=\{b,a\}.
\] The nonzero weight removes any topological primitive ambiguity. The
degree of \(b\) is the negative odd integer \(1-k\). Take a symbol with
this leading coefficient and skew-symmetrize it to obtain \(B^*=-B\)
with \(\sigma_{1-k}B=b\). Conjugating \(\Phi\) by \(\exp(-B)\) removes
its leading difference. Repeat at each next nonzero even lowering
degree. The skew generators have orders tending to minus infinity, so
their ordered exponential product converges in every finite filtration
quotient to the required \(U\). Its finite products are internally
unitary, proving \(U^*U=1\).

For uniqueness, the normal product shows that commutation with all
coordinate multiplications gives \(\partial_\xi P=0\), and commutation
with all coordinate derivatives gives \(\partial_xP=0\). The local
center consists of degree-zero constants, or passive functions. A
central unit of \(1+F_{-1}\) is therefore one. The ratio of two
implementing units has this property. Each finite recursive coefficient
uses finite differentiation and division by a nonzero integer;
uniqueness makes the local constructions independent of the coordinate
presentation and makes them glue. This proves all family and restriction
assertions. \hfill\(\square\)

\pandocstatementheading{\leavevmode\phantomsection\label{lem:online-module-comparison-scalar}Lemma 7.9 (The remaining module scalar).}

On a connected sky, a module automorphism over the identity algebra map
is multiplication by a nowhere-zero passive-parameter scalar.
Consequently, if two algebra comparisons differ by
\(\operatorname{Ad}(U)\) as in Lemma 7.8 and \(F_0\) is a module
comparison over the first, every module comparison over the second is
\(cUF_0\) for such a scalar \(c\).

\emph{Proof.} Flatten the sky in the normalized angular model of
Proposition 6.3. Write a module endomorphism on its generator as
\(u\mapsto a(q,\tau)u\). Commutation with multiplication by the \(q\)
coordinates and with \(D_z=\tau\), followed by coefficientwise
completeness, makes its action on every complete coefficient
multiplication by \(a\). Commutation with \(D_{q^a}=\partial_{q^a}\)
gives \(\partial_{q^a}a=0\). The generator is annihilated by
\(zD_z+\tfrac12=-\tau\partial_\tau\), so \(\tau\partial_\tau a=0\). Its
complete homogeneous expansion therefore contains only a degree-zero
coefficient independent of \(q\). It may depend only on the passive
parameters. Connected angular overlaps make it the same scalar on the
whole sky. Invertibility is equivalent to its nonvanishing. Undoing the
graph and normalized Fourier integral comparisons proves the general
statement. Multiplication by \(U\) changes semilinearity to the second
algebra map; the ratio with any other comparison is the automorphism
just classified. \hfill\(\square\)

\pandocstatementheading{\leavevmode\phantomsection\label{lem:online-full-point-family-faithfulness}Lemma 7.10 (Faithfulness on the tested ray domain).}

Let \(B\) be a connected open point core, let \(\Sigma\) be a retained
scalar carrier, and let \(L_p\subset T^*\Sigma\setminus0\) be the conic
sky of \(p\). The union

\phantomsection\label{eq:online-tested-open-ray-domain} \[
\Omega_B=\bigcup_{p\in B}L_p
\] is open and conic. If a complete carrier operator \(P\), basic in the
point variable, satisfies \(P e_\Sigma(p)=0\) for every \(p\in B\) as a
joint complete-family identity, then \(P=0\) on \(\Omega_B\). This does
not require \(B\) to meet \(\Sigma\).

\emph{Proof.} Follow a carrier covector to a point \(p\in B\) at its
retained flow time. At that same local flow time, every sufficiently
nearby carrier covector reaches a point still in the open set \(B\).
Those covectors therefore belong to the union. This proves openness.
Positive radial dilation preserves the unparametrized ray, proving
conicity. The argument is on connected retained plaques, rather than an
unrestricted long-ray quotient.

The evaluation has nonzero principal module symbol on every \(L_p\) by
Proposition 6.16 and the ellipticity of carrier transport. If \(P\) has
upper order \(m\), the leading symbol of its annihilation identity is
\(\sigma_m(P)|_{L_p}\,\sigma_1(e_\Sigma(p))=0\). Hence the same actual
homogeneous function \(\sigma_m(P)\) vanishes on every \(L_p\), and thus
on their open union. Basicness is what makes it the same function for
all \(p\). The order of \(P\) drops by one. Repeating gives
\(P\in F_{m-j}\) for every finite \(j\), so the complete class is zero.
Every degree is tested on the same domain. Joint smoothing makes these
coefficient identities compatible in the point and passive parameters. A
single sky tests the symbol only along its own conic Lagrangian. As the
point varies, these tests cover \(\Omega_B\), which is the domain of the
conclusion. \hfill\(\square\)

\subsection{7.5. Complete comparisons across surfaces and
carriers}\label{complete-comparisons-across-surfaces-and-carriers}\label{sec:online-full-unmarked-hom}

\pandocstatementheading{\leavevmode\phantomsection\label{thm:online-full-unmarked-source-hom}Theorem 7.11 (Full source Hom).}

Every geometric comparison of two independently admitted sources of
Definition 7.7, with its specified actual spin lift, has a unique full
source lift modulo normalized presentation cells. Conversely every full
source arrow has the physical point map recognized in Theorem 7.3 and is
that lift. The lift preserves the full evaluation, all point directions
of both connections, the internal adjoint, native density roots, source
and target associators, and canonical optical and finite-rank diagrams.
It has exact identity, inverse, composition and restriction laws on
common retained domains. Its two initial hypersurfaces need not agree or
be disjoint, and its scalar carriers need not meet the point cores.

\emph{Proof.} Theorem 7.3 first recognizes the point map \(f\) as an
actual isometry. Propositions 7.4--7.5 and Lemma 7.6 give its canonical
optical and finite-rank comparisons, with the stipulated actual spin
map. This step uses source incidence and density tests and the
independently admitted Cartan fields.

Work first in subneighborhoods on which the restriction of \(f\)
identifies both causal propagators. To obtain them choose a causally
convex target diamond in the range of \(f\), then a causally convex
source diamond in its preimage, and finally a smaller target diamond
inside the image of the source diamond. Its preimage is causally convex
in the source wave neighborhood: a causal segment between its points
stays in the source diamond and its image stays in the smaller target
diamond by causal convexity there. The two isometric final domains are
globally hyperbolic. Choose a small spacelike disk \(D\) through the
marked source point in this common isometry domain. After one common
restriction of the pointed cores, the null segment for each direction at
a source point \(p\) meets \(D\) transversely, and its image segment at
\(f(p)\) meets \(f(D)\) transversely in the target. Compactness of the
full sky gives uniform transversality on unit directions and a positive
margin to both disk boundaries. Choose the restriction so that each
segment from \(p\) to its hit on \(D\) lies inside the domain of \(f\).
Let \(C_f\) be actual half-density pushforward from \(D\) to \(f(D)\).
In each source separately, use the carrier maps of Theorem 6.14 to put

\phantomsection\label{eq:online-full-source-displaced-carrier-lift} \[
F_f=T^1_{\Sigma_1,f(D)}\,C_f\,T^0_{D,\Sigma_0},
\qquad k_f(P)=F_fPF_f^{-1}.
\] The outside factors require only the individual source neighborhoods.
No value of \(f\) on the possibly displaced \(\Sigma_0\) is required.

Kernel uniqueness in Theorem 6.10 identifies the kernels on the small
isometric subparents. Together with Proposition 6.16 and the carrier
cocycle, it gives \(F_fe_0(p)=e_1(f(p))\) in the complete joint
quotient. Each factor is independent of the evaluation point. Thus
actual point differentiation and its full chain rule give both
point-connection identities on every section. The internal-adjoint and
semilinearity identities follow from Theorem 6.14 and half-density
pushforward. Its root statement is the native one: the initialized
leading amplitude is the half-density Jacobian root. In the normalized
spacelike evolution of Lemma 6.13, internal skewness gives the
left-symbol relation \(2b_0=\partial_x\partial_\xi b_1\), so the
physical subprincipal symbol is zero and no additional subprincipal
phase changes that root. The kernel formula fixes its branch and its
composition. This is not a comparison with the separately calibrated
analytic root.

For a second disk \(D'\) through the marked point, work on a small
causally convex subparent in the domain of \(f\) and its isometric
image. Restrict the two disks and the point core so that the intervening
ray segments stay there. Write \(C_f^D\) for pushforward on \(D\).
Naturality of the mixed kernel and its diagonal square roots gives \[
C_f^{D'}T^0_{D',D}=T^1_{f(D'),f(D)}C_f^D .
\] Precompose with \(T^0_{D,\Sigma_0}\) and postcompose with
\(T^1_{\Sigma_1,f(D')}\). The two carrier cocycles give exactly the lift
above, computed on \(D'\) on the left and on \(D\) on the right. For a
second isometry \(g\), take \(D\) small enough that \(f(D)\) is an
auxiliary disk for \(g\). In \(F_gF_f\) the consecutive transports from
\(f(D)\) to \(\Sigma_1\) and back cancel. Half-density pushforwards
compose as \(C_g^{f(D)}C_f^D=C_{gf}^D\), giving \(F_gF_f=F_{gf}\) on the
common domain. The principal spin maps compose in the same order.

Conjugate by the supplied stored normalizations to obtain the stored
algebra and module vertices, carrying their section and connection by
the same maps. Substitution in the supplied normalization overlap
equations gives the first and third equations in the two-associator
test. The two ways to compose three algebra vertices implement the same
comparison; Lemma 7.8 makes their normalized units identical. This gives
its middle equation with both supplied associators. The module ratio in
Lemma 7.9 is fixed by the transported full nonzero evaluation, so its
overlap equation holds with the same units. To see the composition in
these presentations, use one refinement and the same intermediate charts
for \(f\) and \(g\), suppressing point pullbacks. The composite has
vertices \(k_i^g k_i^f\) and \(\Xi_i^g\Xi_i^f\). Its overlap unit is \[
T^{gf}_{ij}=k_i^g(T^f_{ij})T^g_{ij}.
\] Indeed, moving \(k_i^g\) past \(\varphi^1_{ij}\) in the algebra
equation gives conjugation by \(T^g_{ij}\), and moving \(\Xi_i^g\) past
\(T^f_{ij}\) in the module equation gives multiplication by
\(k_i^g(T^f_{ij})\). For the triple overlap, apply the middle equation
of Definition 7.7 first to \(g\) and then to \(f\): \[
\begin{aligned}
T^{gf}_{ij}\varphi^2_{ij}(T^{gf}_{jk})C^2_{ijk}
 &= k_i^g\!\left(T^f_{ij}\varphi^1_{ij}(T^f_{jk})C^1_{ijk}\right)T^g_{ik}\\
 &= (k_i^g k_i^f)(C^0_{ijk})T^{gf}_{ik}.
\end{aligned}
\] Here the intermediate associator \(C^1_{ijk}\) is the one carried by
the intermediate normalization. More explicitly, a normalized change
\(U_i\) of a target vertex sends \(k_i\) to
\(\operatorname{Ad}(U_i)k_i\) and \(\Xi_i\) to \(U_i\Xi_i\). The unit in
the new charts is \(U_iT_{ij}\varphi^1_{ij}(U_j^{-1})\), by the first
and third equations of Definition 7.7. Substitution in the middle
equation cancels the two \(U_j\) factors and moves the last \(U_k^{-1}\)
past \(C^1_{ijk}\) by the target descent law. The source descent law
gives the inverse change upon precomposition. For two different
presentations of the intermediate object, insert their supplied
transition in one arrow and its inverse in the other on a common
refinement. Their vertices cancel there, and the displayed unit law
gives the same composite up to the normalized endpoint cells. These
formulas carry the section and the point connection by the same change
of vertex.

For uniqueness, compose another lift with the inverse of this one. Its
actual geometric and spin labels are identities. In fixed-carrier
normalizations the algebra point-connection equation makes its action on
every basic operator independent of \(p\). Coefficientwise asymptotic
summation then supplies a point-independent representative of each
resulting complete class, so basicness has the actual meaning of
Definition 6.18. Its full graded map is the identity. Lemma 7.8 gives a
basic normalized unit \(U\) implementing it. The module comparison has
the form \(c(p)U\) by Lemma 7.9. Intertwining the point connection on
every module section gives \((Xc)m=0\) for every \(X,m\), because \(U\)
is basic. Hence \(c\) is constant on the connected point core. The
evaluation equation is now

\phantomsection\label{eq:online-full-source-uniqueness-family-test} \[
(cU-I)e_1(p)=0\qquad(p\in B_1).
\] Lemma 7.10 makes \(cU-I=0\) on the tested open ray domain, even when
the carrier is displaced. Its degree-zero symbol gives \(c=1\), then
\(U=I\) in every degree. Thus the normalized scalar maps agree. The
supplied canonical finite-rank tests fix their remaining bundle maps,
and the actual spin label has been retained throughout. This proves
uniqueness including the stored normalized cells.

The disk and composition formulas hold modulo jointly smooth kernels on
the common tested ray charts over the restricted core. For the identity
isometry the outside transports cancel, and the lift of \(f^{-1}\)
composes on either side to the identity there. Restriction of the
parents restricts the same kernel formulas. The uniqueness just proved
identifies their stored representatives on \(\Omega_{B_1}\) up to the
normalized cells. The two central spin lifts are not identified: their
scalar maps can coincide, but their actual principal and rank-four maps
remain distinct. \hfill\(\square\)

\pandocstatementheading{\leavevmode\phantomsection\label{prop:online-full-comparison-finite-families}Proposition 7.12 (Finite families of supplied comparisons).}

For smooth finite-dimensional families of admitted sources and supplied
smooth geometric comparisons with actual spin lifts, Theorem 7.11 gives
smooth full comparisons locally on one retained parent. Every specified
finite coefficient depth, point derivative order and passive-parameter
order is bounded by finitely many source, comparison and normalization
seminorms on larger buffered parents and their positive inverse margins.
The parent is selected before these orders.

\emph{Proof.} Near a fixed parameter value choose \(D\) and \(f(D)\)
with positive spacelike margins on one compact buffered neighborhood.
The supplied maps and inverse maps remain defined there after one
parameter restriction. Compactness of the full sky gives one finite
transverse-hit cover and one common smaller output core. No coefficient
recursion changes those geometric flow domains.

The geometric comparisons use finite differentiation, initialized
geodesic and bundle transport, matrix inverses and the optical state
formulas. Their derivatives at any fixed order obey differentiated
finite ODEs on the same tubes. The middle half-density map in the
full-lift formula uses finitely many derivatives of \(f\) and a square
root of a positive determinant. The outer carrier maps have the
estimates of Proposition 6.17; the point connections have those of
Proposition 6.22. Finite composition has the estimates of Lemma 6.4. The
stored normalization maps and their inverses contribute their stated
finite norms. Finally every fixed coefficient of a normalized comparison
cell uses only finitely many recursion steps in Lemma 7.8. Combining
these bounds proves the assertion, with order-dependent constants and
finite input orders. These estimates start with a supplied smooth family
of physical isometries and spin lifts. \hfill\(\square\)

The full comparison identifies arrows even when the chosen surfaces
differ. We now use it to construct inverses on objects and arrows.

\section{8. Local presentations, inverse functors and
descent}\label{local-presentations-inverse-functors-and-descent}\label{sec:local-source-equivalence}

Different whole inputs can present the same physical germ. The inverse
therefore returns the comparison class of its local displays. At a
point, the displays are compared by full source arrows over physical
isometries with specified spin lifts. On an open cover, these arrows
give the identifications on overlaps.

\subsection{8.1. Execution and the two
inverses}\label{execution-and-the-two-inverses}\label{sec:local-display-inverses}

\pandocstatementheading{\leavevmode\phantomsection\label{def:local-display-groupoids}Definition 8.1 (The three local groupoids).}

Write \(\mathcal V=\mathsf{Vac}^{\mathrm{loc}}_0\) for the groupoid of
pointed, oriented and time-oriented smooth Lorentzian four-dimensional
Ricci-flat germs with their numerical metric and actual local spin
structures. An arrow is the germ of a pointed orientation- and
time-orientation-preserving isometry together with a specified lift to
the principal spin bundles, commuting with their orthonormal frame
projections. Restriction to a smaller neighborhood defines the same
germ. Write \(\mathcal S=\mathsf{WaveSrc}^{\mathrm{loc}}_0\) for the
independent sources and comparisons of
\hyperref[def:online-complete-independent-wave-source]{Definition 7.7},
on the corresponding pointed cores.

An ordinary-pair display \(P=(c,r,V,\mathfrak m)\) consists of a whole
ordinary pair \(c=(c_+,c_-)\) in the small class domain of §§3--4, a
positive number \(r\), a pointed initial patch \(V\) and markings
\(\mathfrak m\). These include the spin marking and a spacetime point
\(x_*\) in the retained core of the developed collar. The spacetime base
point of \(EP\) is \(x_*\); the marking of \(V\) specifies an
initial-data point. For a disk chosen through \(x_*\) one may take its
initial marking at \(x_*\). Execute the whole pair, develop the
resulting initial source, restore its scale and retain the local fields
on a buffered parent of \(V\). The wave marking retains a neighborhood
and carriers satisfying Definition 6.12. Such a neighborhood exists for
every physical germ in Theorem 5.7: take a small causally convex diamond
in its developed collar, a spacelike disk through the marked point, and
a smaller point core whose whole sky has interior hits. Retain the full
wave enrichment there. The domain of a comparison can be reduced about
the spacetime point. For these reductions one keeps representatives on
buffered parents long enough to perform all transverse carrier
transports. Denote the resulting source by \(E P\). An arrow \(P\to P'\)
is a full local source comparison \(E P\to E P'\) in \(\mathcal S\).
This defines \(\mathcal P=\mathsf{WavePair}^{\mathrm{loc}}_0\) with
composition inherited from \(\mathcal S\). The objects retain their
whole inputs, and the arrows compare their outputs on a local core.

\pandocstatementheading{\leavevmode\phantomsection\label{prop:local-display-execution-recognition}Proposition 8.2 (Execution, recognition and actual local embeddings).}

Execution defines a fully faithful functor
\(E:\mathcal P\to\mathcal S\). Recognition defines a fully faithful
functor \(R:\mathcal S\to\mathcal V\). Every independent source has a
display and a full comparison from it on the pointed core. One may use a
display of the source's own marked initial patch when the comparison
point belongs to the common Cauchy development of the patch on which the
initial data agree. For every \(M\in\mathcal V\) there is a display
\(P_M\) and an actual physical and spin identification
\(b_M:R E P_M\to M\).

\emph{Proof.}
\hyperref[thm:online-ordinary-negative-response-initial-source]{Theorem
4.18} constructs the two gradient operators and reads back the ordinary
inputs on their whole domain. Corollary 2.12 constructs the collar from
the derived data. Lemma 7.2 and Propositions 7.4--7.5 construct the
geometric incidence, optical data and finite-rank fields. Section 6
constructs the kernel, the carrier module, its full evaluation and its
point connection. Proposition 5.6 restores their numerical scale, with
\(h\mapsto r^2h\), \(K\mapsto rK\) and \(e\mapsto r^{-1}e\). These
operations give an object of Definition 7.7. On arrows \(E\) uses
precisely the specified full comparison, which proves the first Hom
assertion.

For a supplied source \(S\), first recover the initial data by Theorem
2.3. The initial maximality test and the initialized Cartan equations
apply on its own collar. Theorem 2.9 recognizes the metric
\(g_S=\eta(\theta,\theta)\) there as Ricci-flat. Lemma 7.2 identifies
its actual null rays and Theorem 6.24 identifies the supplied scalar
class and associated fields on their normalized domains. This procedure
defines \(R S\) on the entire independent domain. On an arrow use the
point map recognized by
\hyperref[thm:online-incidence-physical-recognition]{Theorem 7.3} and
the supplied principal spin map. Conversely, an arrow of \(\mathcal V\)
between \(R S\) and \(R S'\) supplies that geometric comparison, and
\hyperref[thm:online-full-unmarked-source-hom]{Theorem 7.11} gives its
unique full lift. Write this lift as

\phantomsection\label{eq:local-unique-physical-lift} \[
L_{S,S'}(f):S\longrightarrow S',\qquad R L_{S,S'}(f)=f .
\] Here and below \(f\) includes its spin lift. The restriction,
identity, inverse and composition assertions of that theorem make \(R\)
a functor and give its Hom bijection.

First retain the marked initial patch of \(S\).
\hyperref[thm:submission-marked-initial-coverage]{Theorem 5.5} supplies
a whole input whose data agree on a smaller patch, including the future
normal. The initialized spin identification and the Cauchy isometry of
Theorem 2.10 identify a common development of that patch. When it
contains the comparison point, restrict the isometry there and apply
\(L\) to compare the executed display to \(S\).

For an arbitrary pointed physical germ \(M\), choose a spacelike tangent
plane and a spin frame at its point.
\hyperref[thm:submission-vacuum-germ-coverage]{Theorem 5.7} supplies a
new maximal disk through that point, a whole input and a physical
identification \(b_M:R E P_M\to M\) with its continued spin lift. In
particular, take \(M=RS\) at any marked point of an independent source.
The full comparison \(L_{E P_M,S}(b_M)\) of
\hyperref[thm:online-full-unmarked-source-hom]{Theorem 7.11} then
compares the display with \(S\) on a smaller pointed core, including any
displaced scalar carriers. This gives essential surjectivity at every
point and supplies the local choices used in Theorems 8.3 and 8.5.
\hfill\(\square\)

\pandocstatementheading{\leavevmode\phantomsection\label{thm:local-two-inverses}Theorem 8.3 (The inverses and their natural comparisons).}

The functors \(E\) and \(R\) are equivalences. A local selection of the
embeddings \(b_M\) gives inverse functors for both, with natural units
and counits satisfying both triangle identities. All comparisons lift
the actual physical arrows, including comparisons between displays with
different maximal surfaces, clocks and scalar carriers.

\emph{Proof.} Put \(Q=RE\). Proposition 8.2 makes \(Q\) fully faithful.
Choose \((P_M,b_M)\) locally as in that proposition. Define
\(K:\mathcal V\to\mathcal P\) by \(K M=P_M\) and by the unique lift

\phantomsection\label{eq:local-selected-inverse-on-arrows} \[
Q(K f)= b_N^{-1}\,f\, b_M
\qquad (f:M\longrightarrow N).
\] The right side is a composition of germs on a common restriction and
of their actual spin maps. Identity and composition follow by
cancellation and faithfulness of \(Q\). Thus
\(b:QK\Rightarrow 1_{\mathcal V}\) is natural. For \(P\in\mathcal P\)
let \(\eta_P:P\to KQP\) be the unique arrow over \(b_{QP}^{-1}\). For
\(h:P\to P'\), applying \(Q\) to \((KQh)\eta_P\) gives
\(b_{QP'}^{-1} Q h\), the image of \(\eta_{P'}h\). Faithfulness gives
naturality of \(\eta\). The triangles are

\phantomsection\label{eq:local-selected-inverse-triangles} \[
b_{QP}\,Q\eta_P=1_{QP},\qquad K(b_M)\,\eta_{KM}=1_{KM}.
\] The first is the defining inverse identity. In the second,
\(QK(b_M)=b_M^{-1}b_M b_{QKM}=b_{QKM}\), so the image of the left side
is the identity, including its spin action. This proves the two laws for
the composite.

For the equivalence \(R\) take the inverse \(E K\). Its counit is \(b\)
and its unit at \(S\) is \(u_S=L_{S,EKRS}(b_{RS}^{-1})\). Naturality
follows from the arrow formula and the Hom bijection for \(R\). The
first triangle is \(b_{RS} R u_S=1\). For the second triangle the image
of \((EK)(b_M)u_{EKM}\) under \(R\) is \(b_{QKM}b_{QKM}^{-1}=1\). For
\(E\) take the inverse \(K R\). Its unit is \(\eta\) and its counit at
\(S\) is \(v_S=u_S^{-1}:EKRS\to S\). At \(EP\) the arrow
\(v_{EP} E\eta_P\) lies over \(b_{QP}b_{QP}^{-1}=1\). At \(KRS\) the
arrow \((KR)(v_S)\eta_{KRS}\) has image \(b_{QKRS}b_{QKRS}^{-1}=1\)
under \(Q\). Faithfulness gives the remaining two triangles. Thus each
unit and counit in a source category is an arrow of the complete source
comparison, with all of Definition 7.7's structure.

For a second selection \((P'_M,b'_M)\), lift \((b'_M)^{-1}b_M\) through
\(Q\). Functoriality gives the natural comparison \(P_M\to P'_M\) and
its uniqueness over the specified physical identification. A third
selection gives the composite of these lifts. In particular exterior
choices, gauges and distinct initial disks of a fixed physical germ
compare through their embeddings. The intermediate disk used in Theorem
7.11 lies inside the domain of the physical arrow, even when neither
retained carrier lies there. The inverse returns a local display up to
this comparison. An automorphism with the nontrivial central spin lift
remains distinct from the identity. \hfill\(\square\)

\subsection{8.2. Families and descent over open
covers}\label{families-and-descent-over-open-covers}\label{sec:families-and-descent}

\pandocstatementheading{\leavevmode\phantomsection\label{prop:local-inverse-families}Proposition 8.4 (A smooth selection on one parent).}

For a smooth finite-dimensional family on common representatives with
smooth markings, selections and the inverse identifications above exist
locally in parameter space. On compact subsets of a sufficiently small
selected parameter neighborhood, one positive radius, extension branch
and set of buffered parents works for all prescribed finite orders. The
same conclusion holds for the lifts of supplied smooth comparison and
spin families.

\emph{Proof.} The plane and lifted frame used to select a maximal disk
can be chosen in smooth local bundle charts. Fix the base norm and its
threshold in Proposition 5.8 before choosing the radius and extension.
That proposition controls the whole input, embedding, continued spin
frame and the derived wave fields. Positivity of the spacelike and
inverse margins and compactness of the full sky give one finite set of
transverse-hit patches for the selected representatives. For two
selections of the same family the physical map is
\((b'_\lambda)^{-1}b_\lambda\). Its finite derivative bounds follow from
the inverse function formula and the two embedding estimates, on a
compact subset with positive inverse and image margins. Their actual
spin maps compose smoothly, with the chosen initial sign. For an
independently supplied comparison \(f_\lambda\) the middle factor in the
same formula is \(f_\lambda\).

\hyperref[prop:online-full-comparison-finite-families]{Proposition 7.12}
estimates the complete lifts of these maps, including conjugation by the
stored normalizations and every required inverse. Consequently each
stipulated point order \(m\), parameter order \(q\) and homogeneous
depth \(J\) has an estimate of the form in
\hyperref[prop:submission-coverage-family-estimates]{Proposition 5.8},
with finite norms of any supplied maps and inverse normalizations on
those parents included. Increasing \((m,q,J)\) uses stronger finite
input norms on the selected parents. The equations for the triangles
concern those same maps at every degree. For a smoothly varying scale
one includes its derivatives and a positive lower bound. The
construction requires a smooth comparison as input only when a
comparison between two different physical families is requested.
\hfill\(\square\)

\pandocstatementheading{\leavevmode\phantomsection\label{thm:local-display-atlas-descent}Theorem 8.5 (Effective descent of local atlases).}

On open covers of Hausdorff spacetime representatives, with effective
descent for the actual spin bundles, local ordinary-pair atlases with
their full comparison maps have effective descent. Arrows descend
uniquely. The assertion holds locally in the smooth families of
Proposition 8.4.

\emph{Proof.} Use charts centered at variable points of an open
representative, marking a distinguished point on every chart that
contains it. Restriction refines the spacetime cores, the connected null
plaques and their buffered parents. An atlas assigns a display locally
and full arrows on overlaps, over the specified base identifications,
satisfying the triple identity in the groupoid of Definition 7.7.
Presentation cells obey the two associator equations there. Morphisms
are compatible families after common refinement. Passing to germs at a
point recovers Definition 8.1.

Given descent data for such atlases, apply \(R E\) to their charts and
overlap arrows. On the given base, the metric, orientations and point
markings agree under pullback on intersections. They therefore define
smooth global fields on that representative. Ricci flatness holds
locally. Glue the actual principal spin bundles by the specified lifted
transitions. The cocycle identity holds for principal maps and therefore
retains the central sign on each overlap. This gives the physical object
and identifications with the input charts. The same argument applies to
an atlas presented over an already effective Hausdorff geometric gluing.
For arrows the compatible local point and spin maps glue over the
prescribed base arrow.

About each point of the descended representative select a display by
Proposition 8.2. On small overlap cores lift the induced physical and
spin identifications by Theorem 7.11. The lift has the required
evaluation, connections, finite-rank fields and stored cells, and agrees
on further restrictions by uniqueness. On triple intersections the two
composites lie over the same actual physical and spin map and hence
agree by full faithfulness. Equations in Definition 7.7 give the
corresponding local representatives with both associators. This
constructs a descended atlas. Comparison of its restrictions with the
originally given charts uses the same lift of their geometric
identification. Applying the faithful recognition functor to the overlap
diagram verifies compatibility with the original descent data. Similarly
the lift of a glued physical arrow coincides on every member of the
cover with its prescribed local arrow, which proves uniqueness and
effectiveness for arrows.

On each small causally convex parent, Theorems 6.10 and 6.24 determine
the kernel class and its normalized fields. The overlap maps identify
these local data even when the union is not causally convex. The whole
exterior of each generating pair belongs to its own chart. For a smooth
family the same argument uses Proposition 8.4 on retained compact
parameter subsets and finite refinements. On a noncompact representative
each chart has its own positive width. \hfill\(\square\)

\section{9. Comparison on the analytic
locus}\label{comparison-on-the-analytic-locus}\label{sec:analytic-wave-comparison}

The analytic correspondence uses a holomorphic operator algebra and a
framed family of simple modules on complex skies. Its framing fixes the
scalar of a module comparison. The wave source has a full evaluation
section of degree one and a connection in every point direction. We
first identify the two underlying module presentations. The analytic
framing then fixes the scalar of the comparison map, which in turn
determines the transport of the lower wave coefficients and the point
connection.

The analytic input throughout is the original projective correspondence
(\hyperref[ref-IsthmusLocal2026]{Isthmus 2026a}, Definition H.1 and
Theorem 9.10), on its Lorentz locus at \(\Lambda=0\). Its independently
admitted object has a separated contact germ along a compact ordered
quadric, the first-negative-order unit torsor, calibrated affine and
optical data, an analytic filtered half-density algebroid with
complex-linear adjoint, and a locally standard simple module. The
original projective correspondence recovers the complete support family
from the compact characteristic. The analytic object retains the actual
relative root square and a frame of the rank-one local-system factor
over the point and passive-parameter family. In particular a
frame-carrying module comparison can have a nontrivial scalar between
different framed objects. Normalized \emph{algebra} units have principal
symbol one.

For the comparison we use v1.1, Theorem F.4 and Lemma F.8 for analytic
quantization and normalized inner comparisons, Lemma G.2 and Theorem G.5
for module transport and the spin and Real lifts, and Lemma H.2,
Proposition H.3 and Theorem H.4 for supplied comparisons, calibrated
supports and descent. The actual counit of the original Theorem 9.10
identifies this independent domain with the analytic source data and
carries these structures. Its use here requires the stated hypotheses on
parents, finite refinements, calibrated root maps and full spin labels.

\subsection{9.1. Real sectors and their module
presentations}\label{real-sectors-and-their-module-presentations}\label{sec:analytic-sector-restriction}

\pandocstatementheading{\leavevmode\phantomsection\label{lem:analytic-physical-maps-and-carrier}Lemma 9.1 (Physical comparisons and analytic carriers).}

A smooth isometry between analytic pseudo-Riemannian germs, together
with either specified spin lift, is analytic on a smaller
representative. Through a point of an analytic Lorentzian germ one can
choose an analytic spacelike maximal disk. A smooth scalar carrier can
subsequently be used for the wave realization.

\emph{Proof.} If \(f\) is an isometry, preservation of the Levi-Civita
connection gives \(f=\exp_{f(p)}\,d f_p\,\exp_p^{-1}\) on a normal
neighborhood. Analytic geodesic ODEs and the analytic inverse function
theorem make the expression analytic. In analytic orthonormal frame
charts the lift through the spin covering is analytic on its specified
sheet.

For the disk start with the smooth maximal graph from Theorem 5.2 in
analytic ambient normal coordinates. Shrink and dilate so that the graph
has small first and second derivatives and the ambient metric is close
to the constant metric on a buffered ball. This is possible in
coordinates adapted to its tangent plane, by the scaling in §5. The
linearization in the graph variable with fixed boundary data is then an
invertible small perturbation of the flat Dirichlet Laplacian from
\(C^{2,\alpha}_0\) (zero trace here) to \(C^{0,\alpha}\). For an
interior cutoff \(\chi\) put \(\psi_z(x)=x+\chi(x)z\). Regard the mean
curvature of \(x\mapsto(u(x),\psi_z(x))\) as an equation with the actual
fixed boundary value of the graph \(u_0\). Its coefficients use only the
holomorphic ambient metric on a smaller complex neighborhood, \(u\),
\(Du\), \(D^2u\) and the explicit first two derivatives of \(\psi_z\).
The positive spacelike margin preserves the initialized square-root
branch. It is an analytic Banach-space equation in \((z,u)\) also for
complex \(z\) near zero. The implicit function theorem gives \(u_z\).
For real \(z\), the reparametrized graph has \(u_z=u_0\circ\psi_z\) by
the fixed-boundary uniqueness. Evaluating where \(\chi=1\) gives a
holomorphic extension in the translation variable, which proves interior
real analyticity of the graph. For a displaced smooth carrier Theorem
6.14 subsequently transports its wave fields through the real geometry.
\hfill\(\square\)

In real-adapted coordinates on a conic sector, restrict the analytic
coefficient functions and their homogeneous variable to the chosen real
cone. Its \emph{complete formal algebra} allows smooth complex
homogeneous coefficients at every successively lower integer degree,
with the rephased product and internal adjoint of Proposition 6.2. Use
the same coefficient convention for the shifted simple module in a
flattened sky chart (Proposition 6.3), with its local half-form
generator. For mixed coordinates make the analytic change to an adapted
chart first. The overlap maps of these local models are the finite-order
extensions of the analytic transitions constructed below.

\pandocstatementheading{\leavevmode\phantomsection\label{prop:analytic-real-sector-interface}Proposition 9.2 (Real sector presentations).}

For a Lorentz analytic object of the original H.1 with \(\Lambda=0\),
choose a local analytic Cauchy disk \(D\) as in Lemma 9.1. On a common
finite refinement about the whole real sky, after shrinking the point
neighborhood, there are filtered adjoint-preserving isomorphisms \(k_i\)
from the algebra of each analytic object chart to the algebra of the
corresponding adapted chart on \(D\), over the full homogeneous ray
maps. Over each \(k_i\) there is a same-order semilinear isomorphism
\(F_i\) from the simple module of the object chart to the shifted
conormal module of the adapted chart. Their analytic sector maps and
inverses extend to the complete formal models above and the
corresponding modules on \(D\). At each degree the output depends on
finite jets of the input coefficients at the real preimage of the
endpoint, smoothly in the point and passive parameters, and the maps
restrict to the given analytic actions on analytic inputs. A supplied
analytic quantization may be used when, on each strict conic compact of
the common analytic coordinate and operator domains, the actual
coordinate images under the comparison and its inverse have total-symbol
terms of degree \(m-j\) bounded by \(M A^j j!\) for \(j\ge0\), with
\(m\) the order of the image and \(M,A\) uniform over the retained
family (v1.1, §§F.1 and H.1--H.2).

On overlaps the extended maps satisfy the two-associator and module
comparison equations of
\hyperref[def:online-complete-independent-wave-source]{Definition 7.7}
with unique normalized algebra units. For fixed algebra maps and their
overlap cells, the full module comparisons form a torsor for
nowhere-zero smooth complex functions of the point and passive
parameters: any two differ by exactly one such function before the
framing of Proposition 9.5 is imposed. Composition with \(T_{aD}\)
carries the comparison to a smooth carrier \(a\) on the common ray
domain of \(D\) and \(a\) in
\hyperref[def:online-wave-carrier-displays]{Definition 6.12}.

\emph{Proof.} Use the geometric comparison of the original counit and
the disk of Lemma 9.1. On a real Cauchy carrier the past covector over
\(\xi\) is \((-H(\xi),\xi)\), \(H=(h^{ij}\xi_i\xi_j)^{1/2}\). Its
positive root, the null flow and the transverse hit map extend
holomorphically on sectors about each retained direction. The past
extensions of \(\xi\) and \(-\xi\) describe different real rays. The
sector choice retains both when the complex spatial projectivization
identifies their spatial lines. Compactness gives a common finite
refinement of the whole real sky, also over a small closed point core.

Use an analytic disk \(D\) through the center of the point core. The
transverse hit map identifies the two real ray germs on this disk,
including their homogeneous lifts. Choose real angular coordinates there
and use their holomorphic extensions on the analytic side. In a mixed
analytic chart first make this coordinate change within the analytic
algebra. The real formal coefficients will be taken in the resulting
adapted coordinates. In these coordinates the rephased product and
complex-linear adjoint of Proposition 6.2 have the finite coefficient
formulas of the original Appendix F.1.

The analytic simple presentation in a flattened contact chart has cyclic
vector \(u\) with \(D_{q^a}u=0\) and \((zD_z+\tfrac12)u=0\). On its
coefficient functions the normal actions are \(D_z=\tau\),
\(D_{q^a}=\partial_{q^a}\) and \(z=-\partial_\tau+1/(2\tau)\). They are
exactly the actions of Proposition 6.3, with the translated equations
for \(z=s(q)\) on moving skies. The radial normalization can be seen on
\(\xi=\tau(q^1,q^2,1)\), \(\tau>0\). Then

\phantomsection\label{eq:analytic-sector-half-form} \[
\tau^{-3/2}|d^3\xi|^{1/2}
   =|dq^1dq^2\,d\tau/\tau|^{1/2}.
\] It follows that the degree-zero analytic generator and the shifted
conormal generator define the same local model. The normalized Radon map
of Proposition 6.3 and the graph projection in Definition 6.15 put the
skies of all nearby points in the form \(z=s(q)\) on \(D\). For the
analytic disk their graph functions are analytic by the holomorphic hit
map. The substitution \(z'=z-s(q)\), with its cotangent lift and
half-density action, flattens each graph. On the analytic side use the
locally standard presentations and moving supports of the original H.1
and H.3. In the two local standard presentations choose the standard
cyclic generators, each satisfying \(D_{q^a}u=0\) and
\((z\,D_z+\tfrac12)u=0\) in the flattened coordinates. The displayed
normal actions then agree on coefficients, so sending one generator to
the other with its coefficient fixed intertwines the actions over the
identity and preserves order. Exchanging the generators gives the
inverse.

We now lift the given algebra comparison in the flattened charts. Write
the contact form as \(\alpha=dz+p^t dq\) and the common reference
Legendrian as \(\Lambda_*=\{z=0,\ p=0\}\). Here \(p\) is an angular
momentum coordinate; the physical point and the passive parameters are
held as parameters in this construction. The flattened contact map \(h\)
carries \(\Lambda_*\) onto itself. Write \(h_0\) for \(h\) followed by
the translation in \(q\) that removes \(h(0)\) along \(\Lambda_*\). This
translation varies analytically with the family. For \(h_0\) use
\(\delta_t(z,q,p)=(t^2z,tq,tp)\). Its weighted leading map
\(S=\lim_{t\to0}\delta_t^{-1}h_0\delta_t\) is the contact lift of a
conformal symplectic linear map with multiplier \(\lambda\ne0\), as in
the original H.2. This map preserves \(p=0\). Its \((q,p)\) matrix thus
has blocks \((A,B;0,\lambda A^{-t})\), with \(A\) invertible and
\(K=A^{-1}B\) symmetric, by the conformal symplectic equation.
Explicitly \(S\) is the shear followed by the diagonal map \[
\begin{aligned}
(z,q,p)&\longmapsto(z-\tfrac12 p^t Kp,\ q+Kp,\ p),\\
(z,q,p)&\longmapsto(\lambda z,\ Aq,\ \lambda A^{-t}p).
\end{aligned}
\] Indeed the shear pulls \(\alpha\) back to itself, the diagonal map
multiplies it by \(\lambda\), and the value of \(z\) at the origin fixes
the integration constant in the lift. Replacing \(K\) by \(tK\) gives a
shear path. For the diagonal path, at one reference parameter join the
identity to the reference \((A,\lambda)\) in
\(\mathrm{GL}(2,\mathbb C)\times\mathbb C^\times\) by matrix and scalar
exponentials. On a smaller parameter chart append the exponential of
\(t\) times the local logarithm of the remaining near-identity factor.
Translation, shear and each diagonal path preserve \(\Lambda_*\)
throughout.

The residual map \(R=S^{-1}h_0\) has weighted leading identity. Both
\(R\) and \(R^{-1}\) preserve \(\Lambda_*\). For \(t\ne0\) the same is
true of \(R_t=\delta_t^{-1}R\delta_t\) and its inverse. Beyond the
identity the terms in \((q,p)\) have weight at least two and those in
\(z\) have weight at least three. Hence substitution gives normally
convergent power series with positive powers of \(t\) for the remainders
of both maps, and \(R_0=R_0^{-1}=I\). More precisely, on a bounded
complex neighborhood of \([0,1]\) choose a smaller weighted polydisk
whose dilates stay inside the common convergence polydisks for \(R\) and
\(R^{-1}\). The scaled series give uniform holomorphic maps there. Their
inverse identities extend across \(t=0\). After a further fixed
restriction their images lie in the prescribed larger charts, and the
contact multipliers and their inverses stay bounded away from zero. For
the translations use charts about their moving centers; for the finitely
many linear paths the matrices and inverses are bounded on smaller
complex neighborhoods of their parameter intervals. These choices give
flow and inverse-flow tubes about the retained cone sectors, using the
contact multipliers to lift the radial variable. They can be made
simultaneously on compact subsets of each conditioned parameter chart. A
finite refinement covers the retained angular and parameter cores.

For any path piece let \(V_t=\dot h_t h_t^{-1}\). Its restriction to
\(\Lambda_*\) is tangent to \(\Lambda_*\), so \(\alpha(V_t)=0\) there.
The homogeneous Hamiltonian on the symplectization is
\(\tau\alpha(V_t)\) up to the pullback sign, and vanishes on the
reference cone. Use the skew order-one lift \(B(H)=(Q_0(H)-Q_0(H)^*)/2\)
of the original F.1 for this Hamiltonian \(H\), as in H.2. Its analytic
symbol bounds on the smaller parents and the tubes just constructed meet
the hypotheses of F.4 and G.2. Starting with the identity of the
standard modules, G.2 yields invertible semilinear evolution for the
algebra evolution of F.4. Compose the finitely many evolutions and undo
the flattenings. Denote the resulting algebra and module maps by
\(k_0,F_0\). Reverse evolution gives their inverses. This gives
existence with \(k=k_0\) and \(F=F_0\).

Now let \(k\) be a supplied filtered adjoint-preserving analytic
comparison with the same full graded action. For a conditioned diagram
H.1 requires uniform factorial bounds on the actual operator coordinate
images and inverse comparisons on common parents. In particular
\(k(y_i)\) satisfies the source-coordinate image bounds used in the
second part of H.2. For a single analytic germ the finitely many
coordinate images are analytic symbols; restrict their parents together
with those of the constructed inverse. For a family the bounds must hold
on the stipulated retained parameter set. Form the source automorphism
\(\psi=k_0^{-1}k\). It fixes the full graded algebra and preserves the
adjoint. On the smaller buffered parents its coordinate differences are
\[
\Delta_i=\psi(y_i)-y_i=k_0^{-1}\bigl(k(y_i)\bigr)-y_i .
\] The controlled inverse \(k_0^{-1}\) acts on the bounded
coordinate-image symbols, so the \(\Delta_i\) have the uniform factorial
estimates required by F.8. This is the left-quotient argument of H.2.
Lemma F.8 gives a unique normalized source unit \(Q\) with
\(\psi=\operatorname{Ad}(Q)\). Taking \(U=k_0(Q)\) gives
\(k=\operatorname{Ad}(U)k_0\). Thus \(F=UF_0\) is a semilinear
isomorphism over the supplied \(k\) with inverse \(F_0^{-1}U^{-1}\). The
estimates for \(Q\) and its inverse pass to \(U,U^{-1}\) through the
constructed quantization.

To pass to smooth complete coefficients, factor out the leading pullback
along the endpoint cone map and the transported half-form in G.2. In the
transported coordinates the remaining generator strictly lowers degree.
Its term lowering degree by \(j\ge1\) differentiates an input at most
\(j+1\) times. At a prescribed total drop \(n\) the iterated integral
therefore contains at most \(n\) factors, of total differential order at
most \(2n\). Time integration takes the coefficients of these
differential operators; the input jet is taken at the single preimage
under the endpoint map. The algebra evolution has the same finite-jet
property by the positive-drop expansion in F.4. The normalized unit, its
inverse and the graph substitutions also use finite jets in each degree.
In real adapted endpoint coordinates these formulas act on arbitrary
smooth homogeneous coefficients by real differentiation with complex
coefficients. Intermediate complex paths enter only the coefficients of
the formulas. Products, semilinearity and both inverse identities
extend: at each fixed degree and point they involve finite jets, and
Taylor polynomials in angular and base coordinates, multiplied by the
prescribed radial powers, realize those jets by analytic germs. The
operations depend smoothly on the physical point and passive parameters,
with finite derivative bounds on the stated parents. This constructs the
map of the complete modules on \(D\). More generally, fix any smooth
formal algebra comparison \(k'\) in these presentations with the same
full graded map, adjoint and action on passive functions. By Lemma 7.8
its quotient with the extended \(k\) is conjugation by a unique
normalized target unit \(V\). Then \(VF\) is a module map over \(k'\),
with inverse \(F^{-1}V^{-1}\). The finite-coefficient construction of
\(V\) in that lemma preserves smooth dependence in families.

On an overlap, apply the second part of H.2 and Lemma F.8 to the
supplied transition \(\varphi^0_{ij}\) with the H.1 coordinate-image
bounds. As for \(k\) above, this gives
\(\varphi^0_{ij}=\operatorname{Ad}(P_{ij})\circ K_{ij}\), where
\(K_{ij}\) is an F.4 quantization of the same graded map and \(P_{ij}\)
is a normalized unit in the target chart. The quantization and
conjugation, and their inverses, use only finite jets of the input
coefficients at each degree after the endpoint pullback is factored out.
Thus \(k_i\varphi^0_{ij}k_j^{-1}\) and its inverse have this property at
their real endpoints in the adapted ray charts and recover the analytic
transitions on analytic inputs. Let \(T_{ij}\) be the normalized unit
comparing \(k_i\varphi^0_{ij}\) with \(\varphi^1_{ij}k_j\) as in
Definition 7.7. These maps have the same homogeneous action and preserve
the adjoint, so Lemma 7.8 supplies \(T_{ij}\) uniquely. On a triple the
two ways of composing the algebra maps have equal inner action, hence \[
T_{ij}\varphi^1_{ij}(T_{jk})C^1_{ijk}
   =k_i(C^0_{ijk})T_{ik}.
\] Consequently the module transports, corrected by these units, act on
the lines of local intertwiners with an ordinary cocycle: in a triple
composition the two module associators cancel by this equality. Over the
fixed algebra comparison any two local intertwiners differ by a scalar
constant in the sky directions by Lemma 7.9. The resulting flat Hom line
on the positive cone has trivial monodromy since the cone retracts to
\(S^2\). Choose a nonzero value at one angular section over a small
point-parameter chart and propagate it on a finite good cover. The
transition scalars vary smoothly on the base, and the cocycle and
trivial monodromy give path independence. The resulting maps obey
\(F_i\tau^0_{ij}=T_{ij}\tau^1_{ij}F_j\) and retain precisely one nonzero
scalar freedom on the base.

For a displaced smooth carrier \(a\), compose \(F_D\) with \(T_{aD}\)
and \(k_D\) with \(\operatorname{Ad}(T_{aD})\). With both hits and the
intervening segments retained in the wave neighborhood,
\hyperref[thm:online-wave-carrier-cocycle]{Theorem 6.14} gives the
inverse and carrier cocycle in the joint smoothing quotient.
\hyperref[prop:online-wave-full-evaluation-comparison]{Proposition 6.16}
gives \(T_{aD}e_D(x)=e_a(x)\), and the transport between the two fixed
carriers is independent of \(x\), so it also transports point
derivatives. Proposition 9.5 will select the remaining scalar using the
relative root and the initialized unit of Propositions 9.3--9.4. The
resulting family is the interface used to transfer the full section and
connection in Theorem 9.7. \hfill\(\square\)

\subsection{9.2. Calibrating the root and its local-system
factor}\label{calibrating-the-root-and-its-local-system-factor}\label{sec:analytic-framing}

\pandocstatementheading{\leavevmode\phantomsection\label{prop:relative-residue-calibration}Proposition 9.3 (The relative root and the wave frame).}

Let \(S\to B\) be the analytic point family over its passive base and
\(p:F\to S\) the compact complex quadric family. Write \(\mathcal L\)
for the contact line pulled to \(F\). The physical metric gives
\(R_g=\mathcal L^{-1}\) with an actual square map
\(\kappa_g:R_g^2\to K_{F/S}\). The analytic root \((R,\kappa_R)\)
factors as

\phantomsection\label{eq:relative-analytic-quadratic-line} \[
R=R_g\otimes p^*A,\qquad q_A:A^2\longrightarrow\mathcal O_S,
\] where the product square is \(\kappa_R\). Relative to \(R_g\), the
principal evaluation supplies a frame \(n_{\mathrm W}\) of the native
wave local-system factor \(N_{\mathrm W}\).

\emph{Proof.} The first-jet normal evaluation in the proof of the
original Proposition H.3, with the null incidence identification of D.7,
identifies the symplectization \((\mathcal L^{-1})^\times\) with the
nonzero null cotangent cone in \(E=T^*_{S/B}\). In particular this
identification has the actual homogeneous radial coordinate. Put
\(q_g(k)=g^{-1}(k,k)\) and let \(\mu_g^\vee\) be the oriented metric
volume on the cotangent fibers, continued analytically from the chosen
Lorentz phase. With \(\mathscr E\) the fiber Euler vector field put

\phantomsection\label{eq:relative-analytic-null-residue} \[
\eta_g=-\operatorname{Res}_{q_g=0}(\mu_g^\vee/q_g),\qquad
\omega_g=\iota_{\mathscr E}\eta_g .
\] At a nonzero null covector \(dq_g\ne0\). Division by the derivative
in an eliminated coordinate gives the residue independently of that
coordinate by its wedge identity. In Gaussian covectors
\(q_g=k_0^2-H^2\), \(\mu_g^\vee=(\det h)^{-1/2}dk_0\wedge d^3\xi\). On
\(k_0=-H\) this gives \((\det h)^{-1/2}d^3\xi/(2H)\) with the indicated
sign. The residue has weight two. In coordinates \((\tau,z^1,z^2)\) it
is \(\tau d\tau\wedge a_g\,dz^1\wedge dz^2\) with \(a_g\) a unit, so
\(\omega_g\) is a nowhere-zero section of
\(K_{F/S}\otimes\mathcal L^2\). A change of radial coordinate adds a
projective differential to \(d\log\tau\) which wedges to zero against
the projective top form. Multiplying by \(\omega_g\) defines
\(\kappa_g\). The same residue, before choosing the metric volume,
identifies \(p^*\det E^*\) with \(K_{F/S}\otimes\mathcal L^2\) and so
retains the point-base canonical factor.

On each complete \(\mathbb P^1\times\mathbb P^1\) the line
\(\mathcal L\otimes R\) is of type \((0,0)\). Its fiberwise constant
sections form \(A=p_*(\mathcal L\otimes R)\) with evaluation isomorphism
\(p^*A\cong\mathcal L\otimes R\). Indeed, \(\omega_g\) and \(\kappa_R\)
trivialize the square on the ruled charts, and the two local roots
differ by signs with trivial monodromy on each simply connected quadric.
Local base sections and their transitions give the stated evaluation. In
a base frame \(b\), the quotient \[
q_A(b^2)=(1\otimes\kappa_R)(b^2)/\omega_g
\] is holomorphic and constant on every compact fiber. It is a nonzero
function on the base. For example, if \(\ell\) is the contact frame
evaluated by \(\tau\) and \(\kappa_R(\rho^2)=s\,dz^1\wedge dz^2\), the
two factors in the square of \(\ell^{-1}\otimes(\ell\otimes\rho)\) give
\(a_g\) and \(s/a_g\). Their product is the supplied \(s\). This retains
each base transition and any two-torsion factor.

Let \(w_g=1\) in \(\mathcal L\otimes R_g\) and let
\(d_0=(2\pi)^{-3/2}\). On the chosen real cone the half-form
representing \(w_g\) is the physically normalized square root of the
residue density. Indeed multiplying the degree-zero half-form
\(\sqrt{a_g\,dz^1dz^2\,d\tau/\tau}\) by \(\tau\) squares to \(\eta_g\).
Proposition 6.16 gives through the point the principal physical conormal
half-density \(d_0(\det h)^{-1/4}(2H)^{-1/2}|d^3\xi|^{1/2}\). The
rephasing of degree one therefore defines the frame by

\phantomsection\label{eq:relative-wave-native-frame} \[
\sigma^{\mathrm{for}}_1(e)=-i d_0\,w_g\otimes n_{\mathrm W}.
\] In the standard module the quotient on the right has constant
local-system coefficient on every angular chart. Transport by the
complete carrier comparison gives the identical section on every other
carrier. Both the residue and \(e\) are natural under physical arrows,
hence so is \(n_{\mathrm W}\). Relative to \(R\), its local-system
factor is \(A^{-1}\otimes N_{\mathrm W}\). This is a tensor
factorization, retaining the square and the original root comparison
even when \(q_A(b^2)\) is negative in a Real frame. \hfill\(\square\)

\textbf{Functoriality of the calibrated root maps.} In Proposition 9.3
the fiber of \(A_X\) consists of the fiberwise constant sections of
\(\operatorname{Hom}(R_{g_X},R_X)\) in its quadratic trivializations.
Use the forward direction for an arrow \(f:X\to Y\), identifying its sky
fibers by the induced map. Write \(I_R(f)\) for the calibrated
square-preserving root comparison of the original H.1 and H.3, and
\(I_g(f)\) for the native root comparison given by the contact line and
metric residue. On these sections put \[
I_A(f)b=I_R(f)\circ b\circ I_g(f)^{-1}.
\] This is a map from \(A_X\) to \(A_Y\) over the point map and respects
\(q_A\) by the two square identities. The composite of two calibrated
root maps carries the same square and calibration as the root map of the
composite, so they agree by calibrated uniqueness in H.3. Equivalently,
a normalized comparison cell acts as one on the leading symbol and the
separately framed local-system factor is carried by both maps.
Cancelling that factor identifies the root maps. The native residue maps
compose by pullback of the residue, hence \(I_A(gf)=I_A(g)I_A(f)\);
identities and restrictions follow from the same formula. For a coupled
conjugate arrow use the conjugate source lines. This gives a map of
quadratic lines before choosing any section of prescribed square.

\pandocstatementheading{\leavevmode\phantomsection\label{prop:initialized-analytic-unit}Proposition 9.4 (An initialized unit on the independent analytic
domain).}

With the flat Lorentz calibration of the original source there is,
locally on the retained point and parameter charts of every admitted
analytic object \(X\), a section \(a_X\) of \(A_X\) with
\(q_{A_X}(a_X^2)=\varepsilon_0\). Here \(\varepsilon_0=\pm1\) is the
actual sign of the normalized flat square. Physical comparisons and the
original coupled operation carry these units. A simultaneous change of
physical scale to \(r^2g\), with the original root square held in fixed
cotangent coordinates, gives \(\kappa_{r^2g}=r^{-2}\kappa_g\),
\(q_A^{[r]}=r^2q_A\), \(a^{[r]}=r^{-1}a\).

\emph{Proof.} In a frame \(b\) with \(q_A(b^2)=\lambda\), write
\(a=sb\). The equation \(s^2\lambda=\varepsilon_0\) defines an
unramified double cover, with differential
\(ds=-\tfrac12s\,d\log\lambda\). The change \(b\mapsto ub\) gives
\(\lambda\mapsto u^2\lambda\), so this equation is intrinsic. On the
product of an amplitude disk and a point-parameter parent, a choice on
the entire zero-amplitude family lifts uniquely to the product by its
retraction. Local square roots show that the lift is holomorphic.
Coupled conjugation transports a compatible initialized value to the
compatible extension by the same uniqueness.

We specify the zero-amplitude choice, since an automorphism of an
uninitialized quadratic line may act by \(-1\). Write a flat nonzero
null covector as \(\tau(1,z)\otimes(1,w)\) in the two ordered spin
factors. If determinant-one spin transformations contribute multipliers
\(k,l\) to the two spinors, then \(dz'=k^{-2}dz\), \(dw'=l^{-2}dw\),
\(\tau'=kl\,\tau\). In the product with \(dz'\wedge dw'\) the extra term
in \(d\log\tau'\) vanishes. The actual spin lift on the degree-zero cone
half-form \(\rho_0\) thus has multiplier \((kl)^{-1}\), and the contact
frame has multiplier \(kl\). Their product is invariant under the two
spin transformations and under translations. Choose the compatible
half-form phase and initial calibrated sign of the original Theorem G.5.
The product gives a Real flat section \(b_0\) with constant
\(\lambda_0=q_{A_0}(b_0^2)\ne0\). Set \(a_0=b_0/\sqrt{|\lambda_0|}\) and
\(\varepsilon_0=\operatorname{sgn}\lambda_0\). These are statements
about the supplied phase and square.

For a constructed analytic source \(U\), continue \(a_0\) along its
constructor amplitude and evaluate at one. To compare two physical
presentations use precisely the original G.5 family
\(U_\epsilon=(\epsilon U;0)\), \(V_\epsilon=R_{d,H}(U_\epsilon)\), where
\((d,H)\) denotes the displacement and full spin transformation.
Construct the unit for \(V_\epsilon\) along its own internal family
\(tV_\epsilon\), keeping \(\epsilon\) passive, then evaluate \(t=1\).
The conditioned two-parameter construction in G.5 allows these two paths
on common parents. The square-preserving physical root comparison acts
on the two units by a holomorphic sign. At \(\epsilon=0\) it is the
translation and determinant-one spin comparison just computed. The sign
is therefore \(+1\) on the connected amplitude domain. In particular
this argument uses the actual transformed \(V_\epsilon\), which
generally differs from \(\epsilon V_1\). The original coupled root map
has the compatible initial Real sign and the same continuation proof.
The normalized square cell acts by its principal symbol one on these
factor maps.

For arbitrary independently admitted \(X\) let
\(C_X:X^{\mathrm{can}}\to X\) be the actual source counit of the
original Theorem 9.10, where \(X^{\mathrm{can}}\) is obtained by its
constructor after its decoder. Define
\(a_X=I_A(C_X)a_{X^{\mathrm{can}}}\) using its actual base-line
comparison. For \(f:X\to Y\) the original counit naturality cell
compares \(f C_X\) and \(C_Y f^{\mathrm{can}}\) with the same spin
label. The normalized unit acts by one on degree zero, and the
local-system factor frame is carried separately (original Definition H.1
and Proposition H.3). It follows that the root maps agree. A sign
discrepancy would require changing the separately framed factor to
obtain equal total module maps. Dividing by the common native
residue-root comparison gives equality of the two \(I_A\) maps. Applying
them to the initialized unit proves \(I_A(f)a_X=a_Y\). The same counit
argument for coupled conjugation gives
\(I_A(J_X)\overline{a_X}=a_{\mathfrak c X}\), where \(\mathfrak c X\) is
the coupled conjugate. It also proves independence from the
presentations used for the counit and compatibility on restrictions and
finite diagrams. Thus the deformation used for a canonical presentation
supplies the unit for an independent source by transport through its
counit.

Finally \(\mu_{r^2g}^\vee=r^{-4}\mu_g^\vee\) and \(q_{r^2g}=r^{-2}q_g\).
Their residue ratio gives the asserted square laws. The positive factor
\(r^{-1}\) preserves the initialized branch. The two line comparisons
multiply by \(r\) on \(R_g\) and by \(r^{-1}\) on \(A\), with identity
product on the fixed original root. \hfill\(\square\)

\pandocstatementheading{\leavevmode\phantomsection\label{prop:analytic-scalar-one}Proposition 9.5 (The calibrated comparison and its scalar).}

Let \(\nu_X\) be the original nonzero frame of the local-system factor.
Take the constructed interface of Proposition 9.2 or its comparison over
a supplied quantization on the conditioned parents specified there. Over
the resulting formal algebra map \(k_X\) there is a unique module
interface \(F_X\) satisfying the relative framing condition

\phantomsection\label{eq:analytic-factor-calibration} \[
\nu_X\longmapsto a_X^{-1}\otimes n_{\mathrm W,X}.
\] Here \(a_X^{-1}\) is the dual vector paired to one with \(a_X\). If
\((\alpha_f,\Xi_f)\) and \((\mathcal A_f,T_f)\) are the analytic and
wave representatives of a physical arrow \(f:X\to Y\) in the same
direction, there is a unique normalized unit
\(U_f\in1+F_{-1}\mathscr E_Y\) with \(U_f^*U_f=1\) and

\phantomsection\label{eq:analytic-physical-scalar-one} \[
k_Y\alpha_f=\operatorname{Ad}(U_f)\mathcal A_f k_X,\qquad
F_Y\Xi_f=U_f T_f F_X .
\] \emph{Proof.} The analytic physical comparison of the original G.5
acts on the real ray cone by the lift of the physical isometry.
Proposition 9.2 therefore extends it in adapted coordinates; for a
supplied representative on a bounded diagram we use H.1--H.2 as in that
proposition. Both sides of the calibration belong to local-system
factors relative to the actual original root. The proposed image is
nonzero, horizontally constant in the cone directions and defined for
the full family. The Hom line in Proposition 9.2 therefore supplies the
unique required scalar. The numerical phase of the native symbol in
Proposition 9.3 and the original \(\kappa_R\) both remain in this rule.
Reselecting an auxiliary radial section uses parallel transport of the
original frame along the cone. Its value as a flat factor section is
independent of the new radial coordinate. A true change
\(\nu'_X=u\nu_X\) of the framed object scales its interface to
\(u^{-1}F_X\). Its induced connection in the reframed presentation is
\(u\nabla^F u^{-1}\), which adds the scalar one-form
\(-u^{-1}\mathrm du\) to its coefficient.

The two algebra maps have the same full homogeneous map on the retained
covectors. Lemma 7.8 gives \(U_f\). After this correction both module
maps have the same semilinearity, so their ratio is a point-parameter
scalar by Lemma 7.9. The original map carries \(\nu_X\) to \(\nu_Y\),
Proposition 9.4 carries the dual of \(a_X\) to the dual of \(a_Y\), and
Proposition 9.3 carries \(n_{\mathrm W,X}\) to \(n_{\mathrm W,Y}\). The
unit has principal symbol one. Both module maps thus have the same value
on the prescribed nonzero factor frame, proving that their ratio is
exactly one. \hfill\(\square\)

\subsection{9.3. Real and physical
comparisons}\label{real-and-physical-comparisons}\label{sec:analytic-real-and-physical-maps}

\pandocstatementheading{\leavevmode\phantomsection\label{prop:analytic-native-real-parity}Proposition 9.6 (Native conjugation of the complete evaluation).}

On a compatible real presentation, coefficient conjugation defines maps
\(\kappa_X:\overline{\mathscr E_X}\to\mathscr E_{\mathfrak cX}\) and
\(j_X^0:\overline{\mathscr M_X}\to\mathscr M_{\mathfrak cX}\) with
strict square. They commute with the point connections and carrier
comparisons. The evaluation has the law
\(j_X^0(\overline e_X)=-e_{\mathfrak cX}\) at every homogeneous degree.

\emph{Proof.} In Proposition 6.2 the rephased product and internal
adjoint have real coefficients. Proposition 6.3 gives real formulas for
the module action and radial Jacobian. In a real graph chart
\(\nabla_Y^{\mathrm{amb}}(a u_s)=(Ya-\tau(Ys)a)u_s\) for real \(Y\).
These formulas prove the algebra, square and differential assertions.
Real coordinate pushforwards obey these formulas by the real
half-density substitution and its Taylor expansion. Along a spacelike
isotopy the normalized generator of Lemma 6.13 is formally real. After
flattening the real moving graph its recursive coefficient transports
have real coefficients and real identity initial values. Uniqueness at
each depth, followed by the carrier cocycle, proves commutation with
every retained \(T_{ab}\). Complex stored charts carry the conjugates of
these maps through their normalizations.

On a carrier through the evaluated point \(b\), the exact column is
\(e(b)=C^{1/2}(\rho(b)^{-1}\delta_b)\). With \(\mathcal F=Q-Q^*\), §6
gives \(C=i\mathcal F^{-1}\) and \(C^{1/2}=e^{i\pi/4}\mathcal F^{-1/2}\)
in the formal convention. The coefficients of the latter inverse root
before its constant phase are real by recursion. Its physical left
symbol has degree \(-1/2-\ell\) term \(i^{-\ell}b_\ell\) with \(b_\ell\)
real, by the half-integer rephasing. Taylor expand at \(b\) and move
each \((x-b)^\beta\) to the frequency coefficient by integration by
parts. At degree \(-1/2-k\) the physical column has phase \(i^{-k}\)
times the real sum \[
\sum_{\ell+|\beta|=k}\frac{(-1)^{|\beta|}}{\beta!}
\partial_\xi^\beta\partial_x^\beta b_\ell(b,\xi)
\] times the real density and Fourier factors. Indeed integration by
parts gives \(i^{|\beta|-\ell}=i^{-k}(-1)^{|\beta|}\). This is module
degree \(1-k\), so its formal coefficient acquires the factor
\(i^{-(1-k)}i^{-k}=-i\). Every coefficient transforms with the odd sign.
Transport from a disk through each point gives the same identity for the
full joint family on a fixed carrier. The
\hyperref[eq:relative-wave-native-frame]{principal-symbol formula} makes
\(n_{\mathrm W}\) Real in its compatible root presentation. Formal
conjugation acts on the coefficients over the same ray, whereas
conjugation of the physical distribution reverses the frequency.
\hfill\(\square\)

\pandocstatementheading{\leavevmode\phantomsection\label{thm:analytic-structured-wave-comparison}Theorem 9.7 (Structured comparison on the analytic locus).}

The independently admitted analytic Lorentz operator sources at
\(\Lambda=0\) and the independently admitted wave sources with analytic
physical germ have a natural comparison after sectorwise real
restriction and complete formal enrichment. It has both inverse
comparisons and preserves the two chiral operations, the actual root and
its square and base factors, both associators, physical and optical
scale, the finite-rank diagrams and the actual spin label. It compares
the full evaluation and all point derivatives on every complete module
section, carrying the original coupled operation with its own square
cell. The calibrated interfaces over a fixed physical identification,
with their simultaneous presentation changes, form a contractible
groupoid locally on point-parameter charts. Every stipulated finite
coefficient and derivative has a finite bound on common conditioned
parents.

\emph{Proof.} Work in the direction of the actual geometric arrow, using
the inverse of an original covariant source instruction where necessary
as in original Theorem G.5. For a physical comparison let
\(P_f=k_Y^{-1}(U_f)\) in the real smooth formal algebra of \(Y\) from
Proposition 9.2, and put

\phantomsection\label{eq:analytic-formal-representatives} \[
\begin{gathered}
\alpha_f^{\mathrm W}=\operatorname{Ad}(P_f^{-1})\alpha_f, \\
\Xi_f^{\mathrm W}=P_f^{-1}\Xi_f, \\
e_X^F=F_X^{-1}e_X, \\
\nabla_X^F=F_X^{-1}\nabla_X^{\mathrm{amb}}F_X.
\end{gathered}
\] Proposition 9.5 gives \(k_Y\alpha_f^{\mathrm W}=\mathcal A_f k_X\)
and \(F_Y\Xi_f^{\mathrm W}=T_fF_X\). Applying the latter first to
\(e_X^F\) and then to an arbitrary section and its covariant derivative
proves both full intertwining equations. The corresponding algebra
derivative is transported through \(k_X\). An operator
\(D_Y=\alpha_f(D_X)\) in either original chirality now has
representative \(P_f^{-1}D_Y P_f\). Under a varying gauge \(G\), a
connection coefficient \(B_Z\) goes to \(GB_ZG^{-1}-(ZG)G^{-1}\). These
equations transport every original coefficient. The first-unit affine
torsor is preserved because commutators of a unit in \(1+F_{-1}\) with
the first negative filtration piece vanish at first negative order.
Proposition 7.5 identifies the canonical finite-rank comparison over the
same spin and metric. The actual optical comparison is the one
determined by the same physically selected scale.

For composable original representatives let \(D_{g,f}\) denote their
normalized composition cell, so \(\Xi_g\Xi_f=D_{g,f}\Xi_{gf}\) and the
algebra composite uses \(\operatorname{Ad}(D_{g,f})\). Expanding the two
squares with the exact wave composite gives

\phantomsection\label{eq:analytic-physical-cells} \[
U_g\mathcal A_g(U_f)=k_Z(D_{g,f})\,U_{gf},\qquad
\alpha_g(P_f)P_g=D_{g,f}P_{gf}
\] for target \(Z\). Equality of implementers follows from
normalized-unit uniqueness. Substitution gives exact composition for the
corrected maps, including inverses and identities with their normalized
cells. In stored charts put
\(T'_{ij}=P_i^{-1}T_{ij}\varphi^Y_{ij}(P_j)\), where \(T_{ij}\) is the
original overlap unit. Moving \(P_j\) across \(\varphi_{ij}^Y\) and
using the original overlap equations gives

\phantomsection\label{eq:analytic-wave-both-associators} \[
\begin{aligned}
\alpha'_i\varphi^X_{ij}
 &=\operatorname{Ad}(T'_{ij})\,\varphi^Y_{ij}\alpha'_j,\\
T'_{ij}\varphi^Y_{ij}(T'_{jk})C^Y_{ijk}
 &=\alpha'_i(C^X_{ijk})T'_{ik},\\
\Xi'_i\tau^X_{ij}&=T'_{ij}\tau^Y_{ij}\Xi'_j .
\end{aligned}
\] For example, the middle equality uses
\(\varphi^Y_{ij}\varphi^Y_{jk}(P_k)C^Y_{ijk}=C^Y_{ijk}\varphi^Y_{ik}(P_k)\)
before applying the original comparison identity. This retains both
original associators on every common refinement. If
\(k'_X=\operatorname{Ad}(V)k_X\), the same calibrated Hom computation
gives \(F'_X=VF_X\). The unique change on the analytic presentation is
\((F'_X)^{-1}F_X\). When the carriers differ, the comparison inserts the
canonical wave transport between the two interfaces. The multiplication
and cancellation of these comparisons give identity, inverses and all
finite coherences, proving the asserted contractibility over the fixed
physical identification.

For the coupled operation write \((J_X,j_X)\) for its original
representatives with normalized square \(Z_X\). Transport them to the
wave presentation as
\(\widehat J_X=k_{\mathfrak cX}J_X\overline{k_X}^{-1}\) and
\(\widehat j_X=F_{\mathfrak cX}j_X\overline{F_X}^{-1}\). The original
contact-valued normal evaluation embeds \(\mathcal L^{-1}\) in
\(T^*_{S/B}\) and commutes with the physical cotangent comparison.
Spin-factor exchange and conjugation fix the Hermitian real null
covector, so there is no additional radial action on its injective
image. The two real algebra comparisons have the same full homogeneous
map. By Lemma 7.8 a unique normalized unit \(V_X\) gives
\(\widehat J_X=\operatorname{Ad}(V_X)\kappa_X\). The original root
operation carries the initialized base unit by Proposition 9.4, and the
factor \(n_{\mathrm W}\) is Real by Proposition 9.6. The original map
also carries the separately specified frame. Evaluating on
\hyperref[eq:analytic-factor-calibration]{the calibrated factor} as in
Proposition 9.5 fixes the residual module scalar, hence
\(\widehat j_X=V_X j_X^0\). The two squares then yield

\phantomsection\label{eq:analytic-wave-coupled-square} \[
V_{\mathfrak c X}\,\kappa_{\mathfrak cX}(\overline{V_X})=k_X(Z_X),
\qquad
\widehat j_X(\overline e_X)=-V_X e_{\mathfrak c X}.
\] The first equation follows from normalized uniqueness on the algebra,
and on the module the same unit acts by left multiplication. In
fixed-carrier coordinates the connection defect is
\((YV_X)V_X^{-1}\widehat j_X\) for a real point direction \(Y\), by
differentiating \(V_Xj_X^0\) using Proposition 9.6. Equivalently one
transports the target connection by \(V_X\) along with the target
vector. In other charts the derivative is the transported covariant
derivative of the unit. The odd marking on the conjugate native package
is \(-\overline e\), or equivalently \(ie\) provides a Real full vector.
Two applications give the original square cell on the original
representatives. All these comparisons retain both chiralities of the
coupled operation.

For the two inverses, start with the source counit \(C_X\) of original
Theorem 9.10. The actual root and framed-module part of this counit is
what defined \(a_X\) in Proposition 9.4. Apply
\hyperref[eq:analytic-formal-representatives]{the corrected
representatives} to that very arrow, obtaining an arrow that carries the
added full fields and the original operators in their simultaneous
presentations. Conversely take any independent wave source whose
recognized metric admits analytic coordinates. Apply the original local
correspondence to this analytic geometry. Its original calibrated source
gives the interfaces of Propositions 9.2--9.5. Theorem 7.11 then
compares the output over the chosen physical identification to the given
wave source. These arrows use the full Hom on both independent domains.
A smooth physical map in this restricted target and its specified lift
are analytic by Lemma 9.1, so neither inverse discards physical arrows.
Over a physical identity the composites in the unit and counit triangles
agree with the original triangles after inserting the cells in
\hyperref[eq:analytic-physical-cells]{the physical composition law}, and
the wave composites coincide by Theorem 8.3. The uniqueness of the
normalized algebra implementers and of the framed module comparisons
gives the representative cells for those equalities. The same argument
with a physical arrow between two objects proves naturality. The whole
exterior ordinary inputs used to present either geometry may be smooth.

For estimates retain the compact point, angular and parameter cores of
the stipulated finite analytic diagram, with the root, inverse, phase
and image margins in the original H.1. Its finite buffered refinement
covers those cores. The original quantizers, tangent transports, counit
maps and inverses have their stated analytic estimates on those parents.
Finite real coefficient restrictions therefore have finite bounds.
Derivatives of \(s^2\lambda=\varepsilon_0\) involve finitely many
derivatives of \(\lambda\) and divisions controlled by the initialized
root margins. The inverse-frame margin controls normalization of the Hom
scalar. The wave transports on a smooth displaced carrier have the
bounds of Proposition 7.12. At each negative degree the Euler recursion
of Lemma 7.8, composition, inversion and the differentiated gauge
equations use only finitely many earlier coefficients and finite jets.
This proves every specified bound on the retained parents. The new
comparisons and full evaluations live in the complete formal category
after real restriction. The original analytic inputs and their analytic
comparisons carry the factorial estimates. \hfill\(\square\)

\subsection{9.4. The complete local
correspondence}\label{the-complete-local-correspondence}\label{sec:analytic-correspondence-assembly}

\pandocstatementheading{\leavevmode\phantomsection\label{thm:combined-local-atlas}Theorem 9.8 (The two-class atlas with its analytic comparison).}

For smooth pointed oriented and time-oriented Lorentzian Ricci-flat
four-germs with actual spin, the execution and recognition functors of
Definition 8.1 give the two equivalences in
\hyperref[thm:local-main-correspondence]{Theorem 1.1}. Every independent
source on a wave neighborhood of Definition 6.12 and every physical germ
has a display from the two ordinary classes on the selected whole-input
domain, with numerical scale restored. Physical isometries and their
specified spin lifts are in bijection with complete local comparisons,
including changes of initial surface and scalar carrier. The inverse
returns the class of local displays with its natural identifications.
The common-parent finite-family bounds and effective descent hold under
the hypotheses of Proposition 8.4 and Theorem 8.5. On the analytic
geometric locus, using the v1.1 operator correspondence
(\hyperref[ref-IsthmusLocal2026]{Isthmus 2026a}, Theorem 9.10) and the
analytic inputs of Section 9, these identifications carry the calibrated
analytic operator data, after real restriction and complete formal
enrichment, through the two inverse comparisons of Theorem 9.7,
including its root, associator and coupled Real maps.

\emph{Proof.} First consider whole inputs and their physical
restrictions. \hyperref[thm:online-ordinary-pair-data-bijection]{Theorem
3.13} identifies the ordinary classes with maximal data in a fixed
spatial gauge.
\hyperref[thm:online-ordinary-negative-response-initial-source]{Theorem
4.18} constructs the actual gradients from the classes by the curvature
response and recovers the same whole input. It uses the negative and
positive class operations before comparing their output with that tensor
chart. \hyperref[thm:online-cartan-collar-recognition]{Theorem 2.9}
develops and recognizes its finite Cartan fields.
\hyperref[thm:submission-vacuum-germ-coverage]{Theorem 5.7} constructs a
display of a given physical germ by choosing a maximal disk through its
point, normalizing at one finite order, applying the protected vacuum
extension of Appendix B and decoding the resulting whole data. Physical
Cauchy uniqueness gives an actual local embedding with the restored
metric scale and spin lift.

On \emph{local presentations}, a different inverse uses these
embeddings. The independent wave tests determine the scalar kernel and
the derived carrier fields by
\hyperref[thm:online-wave-enrichment-recognition]{Theorem 6.24}. The
geometry of the skies and the density recognize the isometry by
\hyperref[thm:online-incidence-physical-recognition]{Theorem 7.3}. The
full lift in \hyperref[thm:online-full-unmarked-source-hom]{Theorem
7.11} uses an intermediate disk inside its domain and tests the whole
open family of moving skies. \hyperref[thm:local-two-inverses]{Theorem
8.3} uses that Hom bijection to prove the two natural inverse
identifications and their triangles. Hence the inverse of a physical
germ recovers its local comparison class, even when its representatives
have different whole exterior inputs.
\hyperref[prop:local-inverse-families]{Proposition 8.4} gives the family
quantifiers and \hyperref[thm:local-display-atlas-descent]{Theorem 8.5}
proves effective descent.

Finally, on the analytic locus
\hyperref[prop:relative-residue-calibration]{Propositions 9.3} and
\hyperref[prop:initialized-analytic-unit]{9.4} identify the actual
relative quadratic line and its initialized unit through the original
counit. They fix the scalar of the complete module comparison.
\hyperref[thm:analytic-structured-wave-comparison]{Theorem 9.7}
constructs the corrected physical and Real representatives, proves their
two inverse comparisons and gives the stated finite estimates.
\hfill\(\square\)

\section{Conclusion}\label{conclusion}\label{sec:conclusion}

In the chosen gauge the two ordinary classes determine a whole
initial-data representative. Different extensions of the same disk can
produce different pairs for one physical germ, so the inverse on germs
takes values in a local comparison class. The full lift of a specified
spin isometry makes this class independent of the initial disk and of
the carrier used for the scalar wave kernel.

The smooth correspondence retains both chiralities at the physical
metric scale. Its comparison with analytic twistor data uses real
restriction and the complete formal symbol and module. For a smooth
family, one common parent supports the estimates at every prescribed
finite order. On the analytic locus, the additional hypotheses on the
analytic inputs supply the factorial estimates.

\textbf{Computational supplement.} The \href{../computations/README.md}{supplement guide} describes the code, inputs and recorded results for the following checks. The archive \href{https://github.com/lioristhmus/smooth-twistor-germs/releases/download/v1.0/smooth-twistor-germs-computations-v1.0.zip}{\nolinkurl{smooth-twistor-germs-computations-v1.0.zip}} accompanies version 1.0 at \url{https://github.com/lioristhmus/smooth-twistor-germs/releases/tag/v1.0}.

\begin{itemize}
\tightlist
\item
  \textbf{Scalar and York inverses (§4).} At the nonaxial frequency
  \((1,2,2)\), both helicities give the same single-response multiplier.
  At that frequency the contractions after three responses in Theorem
  4.12 agree with the direct scalar inverse and the full York inverse,
  including its longitudinal and transverse components.
\item
  \textbf{Exterior momentum range (Lemma B.10).} For spherical harmonic
  degrees \(l\ge2\), the magnetic calculation derives the range identity
  of Lemma B.10 by differential elimination. Dropping the angular shift
  \(+2\) changes the residual, and a divergence-free example lies
  outside the selected range. On the electric side the script
  substitutes one transcribed Hodge equation into the other. For the
  boundary profiles it prescribes division by \(r\) to pass to the
  physical frame and checks the resulting ratios at the same cutoff
  argument for two degrees. Appendix B proves the low-mode equations,
  \(M_\delta J=I\), and the uniform estimates and higher regularity
  across all modes.
\item
  \textbf{Wave transport (§6).} In flat spacetime, three boosted rays
  distinguish the degree-zero Doppler factor from the invariant
  degree-one normalization, and a Gaussian-coordinate mode distinguishes
  the two frequency branches by their opposite charges.
\item
  \textbf{Analytic comparison (§9).} In a flat analytic example the
  analytic and wave frames transported by the same spin lift have ratio
  \(+1\). In a nonlinear chart both the leading evaluation coefficient
  and its nonzero degree-zero term change sign under formal conjugation.
\end{itemize}

\phantomsection\label{app:online-weighted-construction}

\section{Appendix A. Weighted construction and based
gauge}\label{appendix-a.-weighted-construction-and-based-gauge}

This appendix proves the analytic operations used in the ordinary-pair
initial-data constructor. It first supplies the whole decaying scalar
and York inverses, then the nonlinear maximal constraint chart with its
exact ambient decoder. We obtain the higher weighted and parameter
estimates for the solution selected at base order. A separate anchored
inverse normalizes the spatial coordinates and retains their initialized
actual spin comparison.

The residual estimates of \hyperref[ref-LeFlochNguyen2023]{LeFloch and
Nguyen} for their seed-to-solution map imply that the correction
vanishes for an already exact whole-space seed (Theorem 2.7, (2.6)). Our
chart uses transverse trace-free seeds and a based spatial gauge. Its
exact decoder recovers the seeds that determine the ordinary classes,
with finite-order dependence along the selected branch. Agreement with
data prescribed only on an interior ball is the extension problem of
Appendix B.

All small neighborhoods below are chosen using finitely many base norms.
Every subsequent higher-order assertion has finite dependence on higher
input seminorms, without a new smallness requirement at those orders.

\subsection{A.1. Weighted spaces and the decaying flat
inverses}\label{a.1.-weighted-spaces-and-the-decaying-flat-inverses}\label{sec:online-weighted-chart-spaces}

The initial fields in this appendix are defined on the whole marked
\(\mathbb R^3\) in fixed Euclidean coordinates. Write \(\delta\) for its
Euclidean metric and \(\langle x\rangle=(1+|x|^2)^{1/2}\). All tensor
norms refer to Cartesian components. A negative weight denotes decay.
The inverse used to decode the momentum is the decaying whole-space
inverse below; the anchored sublinear inverse used for a change of
spatial coordinates has a different domain. The longitudinal and
transverse trace-free parts of the momentum decomposition follow
\hyperref[ref-York1974]{York}.

\pandocstatementheading{\leavevmode\phantomsection\label{def:online-weighted-holder-chart-spaces}Definition A.1 (Weighted Hölder spaces).}

For an integer \(m\ge0\), \(0<\alpha<1\) and a real weight \(\omega\),
set

\phantomsection\label{eq:online-chart-weighted-holder-norm} \[
\begin{gathered}
\|u\|_{C^{m,\alpha}_\omega}
=\|u\|_{C^{m,\alpha}(B_2)}
+\sup_{R\ge1}R^{-\omega}
\|u(R\,\cdot)\|_{C^{m,\alpha}(\overline A)}, \\
A=\{x:1<|x|<2\}.
\end{gathered}
\] The resulting Banach space is denoted by \(C^{m,\alpha}_\omega\).
Tensor-valued spaces are defined componentwise. Taking dyadic radii or
using slightly larger fixed annuli gives equivalent norms.
Differentiation maps \(C^{m,\alpha}_\omega\) into
\(C^{m-1,\alpha}_{\omega-1}\), and

\phantomsection\label{eq:online-chart-weighted-product} \[
\|uv\|_{C^{m,\alpha}_{\omega+\nu}}
\le C_{m,\alpha}\|u\|_{C^{m,\alpha}_\omega}
\|v\|_{C^{m,\alpha}_\nu}.
\] For different differentiability orders, the product estimate holds at
the smaller order. If \(\omega\le\nu\), then
\(C^{m,\alpha}_\omega\subset C^{m,\alpha}_\nu\) continuously. These
assertions follow on the inner ball and on each rescaled annulus from
the ordinary Hölder estimates, followed by the indicated supremum.

We write \(\mathscr S^\omega\) for smooth fields with \[
\sup_x\langle x\rangle^{-\omega+|\beta|}
|\partial^\beta u(x)|<\infty
\quad\text{for every multi-index }\beta.
\] For a fixed \(\alpha\), this is the intersection of the spaces
\(C^{m,\alpha}_\omega\) over all finite \(m\), with equivalent seminorm
systems. One direction is immediate from the weighted Hölder norm. For
the other, one additional derivative bounds the Hölder seminorm on each
fixed rescaled annulus. Thus passing between the symbol notation and all
weighted Hölder orders uses finitely many extra seminorms.

Define

\phantomsection\label{eq:online-chart-flat-york-operators} \[
\begin{gathered}
(L_0W)_{ij}
=\partial_iW_j+\partial_jW_i
-\frac23\delta_{ij}\partial_aW_a, \\
\Delta_LW=\operatorname{div}_\delta L_0W
=\Delta W+\frac13\nabla\operatorname{div}W.
\end{gathered}
\] \textbf{The interior estimate.} We prove the interior Schauder
estimate needed for our second-order systems. Douglis and Nirenberg
describe the extension of interior Schauder estimates to elliptic
systems (\hyperref[ref-DouglisNirenberg1955]{Douglis and Nirenberg
1955}, p.~503). Here every row has weight zero and every unknown has
weight two. Write
\(P=a^{ij}(x)\partial_i\partial_j+b^i(x)\partial_i+c(x)\) with matrix
coefficients on \(A'\subset\mathbb R^3\). Choose the sign of each row as
in the applications below. A uniform inverse bound for the principal
matrix \(\sum a^{ij}(x)\xi_i\xi_j\) on \(|\xi|=1\) suffices. For \(N\)
components and bounded coefficients this is equivalent to
\(|\det (\sum a^{ij}(x)\xi_i\xi_j)|\ge c|\xi|^{2N}\), since the other
singular values have a uniform upper bound on the unit sphere. Strong
ellipticity implies this bound. On fixed nested bounded domains
\(A\Subset A'\), for \(0<\alpha<1\), \(j\) a nonnegative integer, and
coefficients bounded in \(C^{j,\alpha}(A')\), the estimate is

\phantomsection\label{eq:online-chart-interior-schauder-input} \[
\|v\|_{C^{j+2,\alpha}(A)}
\le C\bigl(\|Pv\|_{C^{j,\alpha}(A')}
+\|v\|_{C^0(A')}\bigr).
\] For variable coefficients we start with an existing
\(C^{2,\alpha}_{\mathrm{loc}}\) solution. A locally bounded
distributional solution is sufficient for the constant-coefficient
applications, with the estimate on nested precompact sets. The constant
depends on the domains, \(\alpha\), \(j\), the inverse-symbol margin and
the indicated coefficient bounds. For the strongly elliptic variable
systems below, the interior regularity conclusion also follows from
\hyperref[ref-GiaquintaMartinazzi2012]{Giaquinta and Martinazzi} (2012,
Theorem 5.22). At each bootstrap step the known solution belongs to
\(W^{2,2}_{\mathrm{loc}}\) and the terms \(b^i\partial_i v+cv\) are
included in the forcing at the order used. For \(\Delta\) and
\(\Delta_L\) the coefficients are constant. The positive symbol of
\(-\Delta_L\) is \(|\xi|^2I+\frac13\xi\otimes\xi\), with eigenvalues
\(|\xi|^2\) and \(\frac43|\xi|^2\).

Here is a derivation that also explains which regularity is used in
passing to higher order on the same branch. Freeze the principal
coefficients at \(x_0\) and call the resulting homogeneous operator
\(P_0\). Take a smooth inhomogeneous dyadic Fourier partition, with low
block \(S_0\) and blocks \(\Delta_l\), \(l\ge0\), supported in fixed
dilates of a compact annulus. For nonintegral \(s>0\) the norm on
bounded functions on \(\mathbb R^3\) given by \[
\|S_0v\|_\infty+\sup_{l\ge0}2^{ls}\|\Delta_lv\|_\infty
\] is equivalent to \(C^{\lfloor s\rfloor,s-\lfloor s\rfloor}\). Indeed,
the annular kernel has vanishing moments. Subtracting the Taylor
polynomial of degree \(\lfloor s\rfloor\) bounds its convolution by
\(C2^{-ls}\|v\|_{C^s}\). Conversely derivatives of order \(k\) of a
block cost \(C_k2^{lk}\) by convolution with a slightly larger frequency
cutoff. Sum the blocks through \(k=\lfloor s\rfloor\). For the top
derivative split a difference at \(2^l|h|=1\), using one more derivative
below that scale and twice the supremum above it. The two geometric
series give \(C|h|^{s-k}\). This proof also constructs the indicated
representative of a distribution from its blocks.

The symbol matrix of \(P_0\) is invertible on the unit frequency
annulus. There, its inverse times an enlarged smooth cutoff has bounded
derivatives to any prescribed finite order, by matrix inversion and the
stated margins. Its inverse Fourier transform is integrable with a
uniform bound, by integration by parts. Scaling the annulus shows \[
\|\Delta_lv\|_\infty\le C2^{-2l}\|\widetilde\Delta_l P_0v\|_\infty
\] with \(\widetilde\Delta_l\) an enlarged block equal to one on the
smaller support. The low block is bounded by \(\|v\|_\infty\). At
\(s=2+\alpha\) this proves the full-space estimate for \(P_0\) on
compactly supported fields.

For the local estimate choose \(A\Subset D\Subset A'\) and put
\(d(x)=\allowbreak\min(r_0,\allowbreak\operatorname{dist}(x,\partial D))\), \(t=d(x)/8\). In
coordinates \(x+ty\), \(|y|<4\), the \(C^{0,\alpha}\) norm of
\(a(x+ty)-a(x)\) is bounded by \(C r_0^\alpha[a]_\alpha\). Choose
\(r_0\) once to make this coefficient small compared with the frozen
estimate. Apply that estimate to a cutoff of \(v(x+t\,\cdot)\) equal to
one on \(B_1\) and supported in \(B_2\). An extension of the coefficient
difference by a second cutoff inside \(B_3\) has the same small bound.
Its product with the second derivatives of the cut-off unknown is
absorbed on the left. All cutoff commutators and lower terms are bounded
by a constant times \(\|v(x+t\,\cdot)\|_{C^{1,\alpha}(B_3)}\), besides
the rescaled forcing \(t^2(Pv)(x+t\,\cdot)\). The lower coefficients
receive factors \(t,t^2\).

Let \(M\) be the supremum over \(x\in D\) of the rescaled
\(C^{2,\alpha}(B_1)\) norm. It is finite for the existing solution since
\(\overline D\Subset A'\). On \(B(x,3t)\) the radii \(d(y)/8\) are
comparable to \(t\). A bounded number of such balls, using smaller
concentric subballs for the cover, bounds the rescaled
\(C^{2,\alpha}(B_3)\) norm by \(C M\). Pairs of points farther apart
than a fixed fraction of the unit radius use the supremum bound in the
Hölder quotient. On \(B_3\) interpolation bounds the preceding
\(C^{1,\alpha}\) norm by \(\delta C M+C_\delta\|v\|_\infty\) for each
\(\delta>0\), by extension from the fixed ball and a split into low and
high frequency blocks at a sufficiently large index. Taking the supremum
and choosing \(\delta\) small absorbs \(M\) on the right. On \(A\) the
chosen radii have a positive lower bound. This proves the displayed
interior estimate for \(j=0\).

For \(j\ge1\) use nested domains and induction. If
\(Pv=f\in C^{j,\alpha}\) and \(v\) is already \(C^{j+1,\alpha}\)
locally, differentiate the equation \(j-1\) times and take a first
spatial difference quotient. The resulting unknown is
\(D_h\partial^\beta v\), \(|\beta|=j-1\), with
\(D_h w=(w(\,\cdot+h e_i)-w)/h\). For each fixed nonzero \(h\) this
unknown has the base regularity. Its principal matrix is a translate of
\(a\) or \(a\) itself, according to the product rule used. Every
commutator term of highest order has a differentiated coefficient and at
most \(j+1\) derivatives of \(v\). Difference quotients of
differentiated coefficients and of \(f\) have uniform \(C^{0,\alpha}\)
bounds from their next derivatives. Applying the base estimate bounds
the unknown uniformly in \(C^{2,\alpha}\) on a smaller domain, including
its bounded supremum. Compactness in \(C^2\) on smaller sets and passage
of the Hölder inequality to the limit give the bound for
\(\partial_i\partial^\beta v\) in \(C^{2,\alpha}\). The preceding
induction bounds its lower-order terms by the right side of the
displayed estimate on a larger domain. For each finite \(j\) choose the
intermediate domains between the same two endpoints. This gives the
displayed bound on the prescribed inner domain. For a bounded
distributional solution of a constant-coefficient equation, local
convolution on a slightly smaller domain commutes with \(P_0\). Applying
the estimates to the mollifications and passing to their distributional
limit gives the same conclusion. This covers the flat potentials in
Lemmas A.2 and A.8. The nonlinear applications in Theorems A.7, A.9 and
A.14 and the intervening linearized or averaged systems all start above
the base regularity. Their rescaled principal symbols and coefficient
bounds are checked where those equations are formed.

\pandocstatementheading{\leavevmode\phantomsection\label{lem:online-chart-decaying-flat-inverses}Lemma A.2 (The decaying scalar and York inverses).}

For \(m\ge2\), \(0<\alpha<1\) and \(-1<\omega<0\), the maps \[
\Delta:C^{m,\alpha}_\omega
\longrightarrow C^{m-2,\alpha}_{\omega-2},
\qquad
\Delta_L:C^{m,\alpha}_\omega(T\mathbb R^3)
\longrightarrow C^{m-2,\alpha}_{\omega-2}(T\mathbb R^3)
\] are bounded isomorphisms. Their inverses are convolution with

\phantomsection\label{eq:online-chart-decaying-green-kernels} \[
\Phi(x)=-\frac1{4\pi|x|},
\qquad
G_{ij}(x)=-\frac1{32\pi}
\left(\frac{7\delta_{ij}}{|x|}
+\frac{x_ix_j}{|x|^3}\right),
\] respectively. They are the same operators at every finite order and
satisfy

\phantomsection\label{eq:online-chart-decaying-inverse-bound} \[
\|\Delta^{-1}f\|_{C^{m,\alpha}_\omega}
+\|\Delta_L^{-1}F\|_{C^{m,\alpha}_\omega}
\le C_{m,\alpha,\omega}
\bigl(\|f\|_{C^{m-2,\alpha}_{\omega-2}}
+\|F\|_{C^{m-2,\alpha}_{\omega-2}}\bigr).
\] \emph{Proof.} The differential maps are bounded by the weighted
differentiation rule. Distributionally, \(\Delta\Phi=\delta_0\). Put
\(\Psi(x)=-|x|/(8\pi)\), so that \(\Delta\Psi=\Phi\). Then \[
G_{ij}=\delta_{ij}\Phi-\frac14\partial_i\partial_j\Psi,
\qquad
\partial_iG_{ij}=\frac34\partial_j\Phi.
\] Applying \(\Delta_L\) cancels the two terms involving
\(\partial_i\partial_j\Phi\) and gives
\((\Delta_LG)_{ij}=\delta_{ij}\delta_0\). Differentiating \(|x|\) gives
the displayed matrix kernel. Both kernels are bounded by \(C|x|^{-1}\).

Suppose \(|f(y)|\le B\langle y\rangle^{\omega-2}\). The convolution is
locally integrable at \(y=x\). At infinity its absolute integrand,
including the volume element, is bounded by a multiple of
\(r^{\omega-1}\,\mathrm dr\), which is integrable because \(\omega<0\).
This proves absolute convergence, also componentwise for \(G*F\). The
distributional kernel identity gives the equation; testing against a
compactly supported smooth function and using the same absolute bounds
justifies the interchange of integrals.

For the weighted zeroth-order estimate let \(R=|x|\ge2\). On
\(|y|<R/2\), \[
\frac1R\int_{|y|<R/2}
\langle y\rangle^{\omega-2}\,\mathrm dy
\le C_\omega R^\omega,
\] where \(\omega>-1\) is used. On \(|x-y|<R/2\), one has
\(\langle y\rangle\asymp R\), and \[
R^{\omega-2}\int_{|z|<R/2}|z|^{-1}\,\mathrm dz
\le C R^\omega.
\] On the remaining region, both \(|y|\) and \(|x-y|\) are at least
\(R/2\). The substitution \(y=Rz\) gives \(R^\omega\) times an integral
bounded uniformly for \(x/R\in S^2\); its tail is integrable because
\(\omega<0\). The bounded region \(|x|\le2\) follows from local
integrability and the same tail estimate. Thus \[
\|\Phi*f\|_{C^0_\omega}
+\|G*F\|_{C^0_\omega}
\le C_\omega
\bigl(\|f\|_{C^0_{\omega-2}}
+\|F\|_{C^0_{\omega-2}}\bigr).
\] Splitting off a small neighborhood of the kernel singularity proves
continuity of the convolutions. The
\hyperref[eq:online-chart-interior-schauder-input]{interior estimate}
gives their local higher regularity. On nested fixed annuli use
\(v_R(z)=v(Rz)\) and \(Pv_R=R^2f(Rz)\). The zeroth-order bound and the
source norm give \[
\|v_R\|_{C^{m,\alpha}(\overline A)}
\le C R^\omega
\|f\|_{C^{m-2,\alpha}_{\omega-2}}.
\] An interior-ball estimate controls the remaining part of the weighted
norm. The same calculation applies to the vector system, proving
existence and the claimed bounds.

A harmonic scalar in \(C^{m,\alpha}_\omega\) tends to zero at infinity,
hence is zero by the maximum principle. If \(\Delta_LW=0\), then taking
the divergence gives \(\frac43\Delta\operatorname{div}W=0\)
distributionally. Local harmonic regularity and its decay imply
\(\operatorname{div}W=0\). Hence \(\Delta W=0\), and each component
tends to zero. Thus \(W=0\). Uniqueness proves that the inverse formula
and all its finite-order realizations agree. \hfill\(\square\)

No integral or moment condition is imposed on the residual in this
lemma. For \(-1<\omega<0\) a \(1/|x|\) tail has the allowed decay at
infinity; the solution itself is smooth on the whole space. Constants
and other nondecaying homogeneous solutions are excluded. The
restriction \(\gamma<-1/2\) used below is compatible with these tails.

\subsection{A.2. The two seed coordinates and the independent
target}\label{a.2.-the-two-seed-coordinates-and-the-independent-target}\label{sec:online-chart-ambient-coordinates}

Fix once and for all

\phantomsection\label{eq:online-chart-base-indices} \[
k_0\ge4,\qquad 0<\alpha<1,\qquad -1<\gamma<-\frac12.
\] The base index can be increased before choosing the construction
domain when a later finite-order operation requires it.

\pandocstatementheading{\leavevmode\phantomsection\label{def:online-chart-seeds-and-target}Definition A.3 (Seed spaces and the maximal target in spatial gauge).}

For \(m\ge k_0\), let

\phantomsection\label{eq:online-chart-seed-correction-residual-spaces}
\[
\begin{aligned}
\mathcal Q_m
&=\{q\in C^{m,\alpha}_\gamma(\operatorname{Sym}^2T^*\mathbb R^3):
\operatorname{tr}_\delta q=0,\ 
\operatorname{div}_\delta q=0\},\\
\mathcal P_{m-1}
&=\{p\in C^{m-1,\alpha}_{\gamma-1}(\operatorname{Sym}^2T^*\mathbb R^3):
\operatorname{tr}_\delta p=0,\
\operatorname{div}_\delta p=0\},\\
\mathcal X_m&=\mathcal Q_m\oplus\mathcal P_{m-1},\\
\mathcal Y_m&=C^{m,\alpha}_\gamma
\oplus C^{m,\alpha}_\gamma(T\mathbb R^3),\\
\mathcal Z_m&=C^{m-2,\alpha}_{\gamma-2}
\oplus C^{m-2,\alpha}_{\gamma-2}(T^*\mathbb R^3).
\end{aligned}
\] The variables are \(s=(q,p)\in\mathcal X_m\) and
\(y=(u,W)\in\mathcal Y_m\). Give finite direct sums their sum norms.

The ambient data domain consists of positive whole metrics and covariant
symmetric tensors satisfying

\phantomsection\label{eq:online-chart-ambient-gauge-domain} \[
\begin{gathered}
h-\delta\in C^{m,\alpha}_\gamma,\qquad
K\in C^{m-1,\alpha}_{\gamma-1},\qquad
\frac12\delta<h<\frac32\delta,\\
\operatorname{div}_\delta(\operatorname{TF}_\delta h)=0,
\qquad \operatorname{tr}_hK=0,
\end{gathered}
\] where \(\operatorname{TF}_h T=T-\frac13(\operatorname{tr}_hT)h\). Its
maximal vacuum subset is defined independently by

\phantomsection\label{eq:online-chart-independent-constraints} \[
\mathcal H(h,K)=R(h)-|K|_h^2=0,\qquad
\mathcal M(h,K)=\operatorname{div}_hK=0.
\] These conventions agree with \(g=\mathrm dt^2-h_t\) and
\(K=\frac12\partial_t h_t\). The target conditions contain no
requirement that a datum already be generated by a seed or a cohomology
class.

\pandocstatementheading{\leavevmode\phantomsection\label{prop:online-chart-exact-ambient-decoder}Proposition A.4 (Ambient coordinates and exact decoder).}

On a neighborhood of zero the map

\phantomsection\label{eq:online-chart-ambient-parametrization} \[
\mathcal A(q,p,u,W)=(h,K),\qquad
h=\delta+q+u\delta,\qquad
\kappa=p+L_0W,\qquad
K=\operatorname{TF}_h\kappa
\] gives smooth local coordinates on the ambient gauge domain. Its
inverse is

\phantomsection\label{eq:online-chart-exact-decoder} \[
\begin{aligned}
q&=\operatorname{TF}_\delta h,
&u&=\frac13\operatorname{tr}_\delta h-1,\\
\kappa&=K-\frac{\operatorname{tr}_\delta K}
{\operatorname{tr}_\delta h}h,
&W&=\Delta_L^{-1}\operatorname{div}_\delta\kappa,\\
p&=\kappa-L_0W.
\end{aligned}
\] This inverse is defined before imposing the vacuum constraints. The
second seed is the flat transverse-traceless part of the displayed
\(\kappa\).

\emph{Proof.} The tensors \(p\) and \(L_0W\) are Euclidean tracefree, so
\(\operatorname{tr}_\delta\kappa=0\). The metric formula gives
\(\operatorname{TF}_\delta h=q\) and its required divergence condition.
The definition of \(K\) gives its \(h\)-tracefreeness.

Conversely, the first two decoder formulas return a Euclidean
transverse-traceless \(q\) of weight \(\gamma\) and a scalar \(u\) of
that weight. Put
\(c=(\operatorname{tr}_\delta K)/(\operatorname{tr}_\delta h)\). The
denominator is uniformly positive on the metric neighborhood. Then
\(\kappa=K-ch\) is Euclidean tracefree and, using
\(\operatorname{tr}_hK=0\), has \(\operatorname{tr}_h\kappa=-3c\).
Consequently \(\operatorname{TF}_h\kappa=K\).

Weighted multiplication gives \(\kappa\in C^{m-1,\alpha}_{\gamma-1}\),
so its Euclidean divergence belongs to \(C^{m-2,\alpha}_{\gamma-2}\).
Lemma A.2 gives \(W\in C^{m,\alpha}_\gamma\). Thus \(p=\kappa-L_0W\) has
the required weight and regularity, and \[
\operatorname{tr}_\delta p=0,\qquad
\operatorname{div}_\delta p
=\operatorname{div}_\delta\kappa-\Delta_LW=0.
\] Starting from Euclidean tracefree \(\kappa\) and forming
\(K=\operatorname{TF}_h\kappa\), one obtains \[
\frac{\operatorname{tr}_\delta K}
{\operatorname{tr}_\delta h}
=-\frac13\operatorname{tr}_h\kappa.
\] The decoder therefore returns the same \(\kappa\). It recovers \(W\)
by the uniqueness of the decaying York inverse, then \(p\). It plainly
recovers \(q\) and \(u\). The opposite composition is the identity by
the preceding construction.

Matrix inversion near \(\delta\), the scalar division above and
tracefree projection are smooth weighted algebraic operations. Lemma A.2
is bounded at the specified orders. Both coordinate maps are therefore
smooth. The assertions hold at every finite order with the same
formulas. \hfill\(\square\)

\subsection{A.3. The selected nonlinear constraint
chart}\label{a.3.-the-selected-nonlinear-constraint-chart}\label{sec:online-chart-constraint-solution}

\pandocstatementheading{\leavevmode\phantomsection\label{lem:online-chart-constraint-remainder}Lemma A.5 (Constraint map and quadratic remainder).}

The map \[
\mathcal F(s,y)
=\bigl(R(h)-|K|_h^2,\operatorname{div}_hK\bigr),
\qquad (h,K)=\mathcal A(s,y),
\] is smooth from a neighborhood of zero in
\(\mathcal X_m\oplus\mathcal Y_m\) to \(\mathcal Z_m\). With

\phantomsection\label{eq:online-chart-flat-transverse-operator} \[
A_0(u,W)=(-2\Delta u,\Delta_LW),
\] it has the decomposition

\phantomsection\label{eq:online-chart-constraint-split} \[
\mathcal F(s,y)=A_0y+\mathcal N(s,y).
\] The two components of \(\mathcal N\) take values in
\(C^{m-2,\alpha}_{2\gamma-2}\), vanish to second order at the origin,
and obey a quadratic norm bound on a smaller neighborhood. In
particular, \[
D_y\mathcal F(0,0)=A_0,\qquad
D_s\mathcal F(0,0)=0.
\] \emph{Proof.} For a covariant perturbation \(r\), \[
DR_\delta(r)
=\partial_i\partial_jr_{ij}
-\Delta\operatorname{tr}_\delta r.
\] The transverse-traceless conditions give \(DR_\delta(q)=0\), while
\(DR_\delta(u\delta)=-2\Delta u\). Writing \(r=q+u\delta\), the
nonlinear scalar-curvature remainder is a finite sum of smooth
contractions of the forms \[
O(r)\,\partial^2r,\qquad
O(1)\,\partial r\,\partial r.
\] Here the coefficients are smooth matrix functions with uniformly
bounded derivatives on a retained positive metric set. Both terms have
weight \(2\gamma-2\) and regularity \(C^{m-2,\alpha}\).

Since \(\operatorname{tr}_\delta\kappa=0\), \[
K-\kappa
=-\frac13\bigl((h^{-1}-\delta^{-1}):\kappa\bigr)h.
\] It is quadratic and belongs to \(C^{m-1,\alpha}_{2\gamma-1}\). The
term \(|K|_h^2\) belongs to \(C^{m-2,\alpha}_{2\gamma-2}\). The linear
momentum term is \(\operatorname{div}_\delta\kappa=\Delta_LW\). The
other momentum terms have the forms \(O(r)\partial\kappa\) and
\(O(1)\partial r\,\kappa\), hence the same quadratic order, weight and
required regularity.

The weighted product estimate and smooth matrix inversion prove
smoothness into the stated faster-decay spaces and give \[
\|\mathcal N(s,y)\|_{C^{m-2,\alpha}_{2\gamma-2}
\oplus C^{m-2,\alpha}_{2\gamma-2}}
\le C_m(\|s\|_{\mathcal X_m}+\|y\|_{\mathcal Y_m})^2.
\] Since \(\gamma<0\), the faster target embeds into \(\mathcal Z_m\).
The derivative formulas follow from the displayed linear terms.
\hfill\(\square\)

\pandocstatementheading{\leavevmode\phantomsection\label{thm:online-chart-base-constraint-bijection}Theorem A.6 (The base constraint chart).}

At the fixed base order \(k_0\), choose sufficiently small neighborhoods
\(\mathcal U\subset\mathcal X_{k_0}\) and
\(\mathcal V\subset\mathcal Y_{k_0}\) of zero. There is a unique smooth
correction

\phantomsection\label{eq:online-chart-selected-correction} \[
\mathcal C:\mathcal U\longrightarrow\mathcal V,\qquad
\mathcal F(s,\mathcal C(s))=0,\qquad
\mathcal C(0)=0,
\] such that the constraint solutions in \(\mathcal U\times\mathcal V\)
are precisely its graph. The map

\phantomsection\label{eq:online-chart-nonlinear-constructor} \[
\mathcal T(s)=\mathcal A(s,\mathcal C(s))
\] is a smooth bijection onto a relatively open neighborhood of
\((\delta,0)\) in the independent constrained target, with smooth
inverse given by the seed part of Proposition A.4. Moreover,

\phantomsection\label{eq:online-chart-quadratic-correction-bound} \[
D\mathcal C(0)=0,\qquad
\|\mathcal C(s)\|_{\mathcal Y_{k_0}}
\le C\|s\|_{\mathcal X_{k_0}}^2,
\qquad D\mathcal T(0)[q,p]=(q,p).
\] The neighborhoods can be chosen as convex base neighborhoods with a
fixed transverse inverse margin: throughout their product,

\phantomsection\label{eq:online-chart-uniform-base-inverse-margin} \[
\|D_y(A_0^{-1}\mathcal N)(s,y)\|
_{\mathcal L(\mathcal Y_{k_0})}\le\kappa<\frac12.
\] Every resulting metric is complete and asymptotically flat in the
stated coordinates.

\emph{Proof.} Lemma A.2 makes
\(A_0:\mathcal Y_{k_0}\to\mathcal Z_{k_0}\) a bounded isomorphism. Lemma
A.5 and the Banach implicit function theorem give the selected graph.
Since the derivative of \(\mathcal N\) vanishes at zero, continuity
permits a convex product neighborhood on which the last displayed
operator bound holds, while preserving the graph and metric positivity.
Equivalently, the selected correction is the zero-initialized fixed
point of \[
y\longmapsto-A_0^{-1}\mathcal N(s,y)
\] on a sufficiently small invariant ball. The product neighborhood and
ball are chosen once at the base order.

Differentiating \(\mathcal F(s,\mathcal C(s))=0\) at zero gives
\(A_0D\mathcal C(0)=0\). Hence \(D\mathcal C(0)=0\), and a bounded
second derivative on a smaller neighborhood gives the quadratic
estimate. The ambient formula then gives the derivative of
\(\mathcal T\).

Proposition A.4 supplies independent ambient coordinates with their
exact inverse. In those coordinates the constraint subset in the chosen
open product is the graph of \(\mathcal C\). Its projection to the seed
coordinates is therefore a smooth bijection onto \(\mathcal U\) with the
stated inverse. This proves both directions using the independently
imposed gauge, maximality and vacuum equations.

The fixed metric bounds make the Riemannian distance uniformly
comparable with Euclidean distance, so completeness follows. The
weighted estimates give \(h\to\delta\) and \(K\to0\) with their
specified derivatives. No condition on the higher norms has been used in
choosing this base chart. \hfill\(\square\)

\subsection{A.4. All weighted orders on the same
branch}\label{a.4.-all-weighted-orders-on-the-same-branch}\label{sec:online-chart-same-branch-smoothness}

Let \(\mathcal X^\infty_\gamma=\bigcap_{m\ge k_0}\mathcal X_m\) and
\(\mathcal Y^\infty_\gamma=\bigcap_{m\ge k_0}\mathcal Y_m\), with their
indicated seminorms. Their smooth elements are the corresponding
symbol-decay fields. The selected domain is
\(\mathcal U\cap\mathcal X^\infty_\gamma\): smallness is imposed only
through \(\mathcal U\) at the base order.

A smooth finite-dimensional family means a map from a finite-dimensional
parameter manifold into this intersection whose realization in every
\(\mathcal X_m\) is smooth. On each retained compact parameter set,
every specified finite parameter derivative therefore has finite
\(\mathcal X_m\) norm. These bounds are required separately at each
finite order; no single bound over all orders is assumed.

\pandocstatementheading{\leavevmode\phantomsection\label{thm:online-chart-one-branch-weighted-smoothness}Theorem A.7 (Weighted smoothness, continuity and finite families).}

For every \(s\in\mathcal U\cap\mathcal X^\infty_\gamma\), the same
correction \(\mathcal C(s)\) of Theorem A.6 belongs to
\(\mathcal Y^\infty_\gamma\). Its metric and momentum satisfy \[
h-\delta\in\mathscr S^\gamma,\qquad
K\in\mathscr S^{\gamma-1}.
\] At each finite output order the norms are bounded by finitely many
input seminorms and the fixed base margins. Higher input norms need only
be finite. The maps \(\mathcal C,\mathcal T\) and the ambient seed
decoder are continuous in these full seminorm topologies.

For every supplied smooth finite-dimensional family in the selected
domain, these maps are smooth into each finite weighted space. On a
compact parameter set contained in the interior of a larger retained
compact set, each specified finite parameter derivative at a specified
spatial order is bounded by finitely many input spatial and parameter
seminorms on that larger set and the fixed margins. The selected base
domain precedes the choice of these orders.

\emph{Proof.} We first establish the spatial improvement of the existing
base correction, then the estimates for its transverse linearization.

Write \(\mathcal C(s)=(u,W)\) and \(h=\delta+q+u\delta\). The
coefficient of the second derivatives of \(u\) in the Hamiltonian
equation is

\phantomsection\label{eq:online-chart-scalar-principal-coefficient} \[
a^{ij}(h)
=(h^{-1}\delta h^{-1})^{ij}
-(\operatorname{tr}_h\delta)h^{ij}.
\] Indeed these are the coefficients of
\(h^{ia}h^{jb}\partial_i\partial_j r_{ab}
-h^{ij}h^{ab}\partial_i\partial_jr_{ab}\) when \(r=u\delta\). At the
flat metric \(a^{ij}=-2\delta^{ij}\), so its negative stays uniformly
positive under the base bounds.

For the momentum equation, retain the second derivatives of \(W\) in
\(\operatorname{div}_h\operatorname{TF}_h(L_0W)\). Its coefficients
depend algebraically on \(h\) and are uniformly close to those of
\(\Delta_L\). The Hamiltonian equation contains \(W\) only through first
derivatives, and the momentum equation contains \(u\) only through first
derivatives. Thus, after changing the sign of the scalar row, the
second-order principal part in \((u,W)\) is block diagonal and uniformly
strongly elliptic. The fixed sufficiently small base neighborhood of
\(\delta\) supplies its ellipticity margin.

Suppose the existing correction has \(C^{j,\alpha}\) regularity on a
bounded region, with \(j\ge k_0\). The principal coefficients then have
the required \(C^{j-1,\alpha}\) regularity. After moving the prescribed
\(\partial^2q\) and \(\partial p\) terms to the right, all remaining
nonlinear terms use at most first derivatives of \((u,W)\). A finite
additional collection of seed derivatives therefore makes the right-hand
side \(C^{j-1,\alpha}\). The interior estimate gives \(C^{j+1,\alpha}\)
regularity on a smaller region. Difference quotients, or the regularity
part of that estimate, apply to the already existing solution; no
higher-order contraction is being asserted. Induction proves local
smoothness.

For the weighted improvement, apply the same argument on fixed nested
rescaled annuli. With \(x=Rz\), normalize the seed and correction
components as

\phantomsection\label{eq:online-chart-annular-normalized-fields} \[
q_R(z)=R^{-\gamma}q(Rz),\qquad
p_R(z)=R^{1-\gamma}p(Rz),\qquad
y_R(z)=R^{-\gamma}(u,W)(Rz).
\] The metric on the rescaled annulus is
\(\delta+R^\gamma(q_R+u_R\delta)\), and the scaled momentum expression
uses \(R^\gamma(p_R+L_0W_R)\). Multiplying the constraint equations by
\(R^{2-\gamma}\) yields the same principal elliptic system with
uniformly controlled rescaled coefficients. The base norms bound \(y_R\)
at order \(k_0\), uniformly in \(R\ge1\). At each later step the
required finite seed and already obtained correction norms are uniformly
bounded by their weighted seminorms. Factors \(R^\gamma\) do not grow,
since \(\gamma<0\). The interior estimate consequently gives the
next-order bound uniformly in \(R\). Taking the supremum over annuli and
adding the interior-ball estimate proves all weighted orders. Each
finite induction has finite input dependence and uses only the original
ellipticity margin.

Now let \[
\mathsf A_{s,y}=D_y\mathcal F(s,y).
\] At every weighted smooth point in the base product neighborhood, its
base inverse is

\phantomsection\label{eq:online-chart-base-transverse-inverse} \[
\mathsf A_{s,y}^{-1}
=\bigl(I+D_y(A_0^{-1}\mathcal N)(s,y)\bigr)^{-1}A_0^{-1},
\] with the first inverse bounded by \((1-\kappa)^{-1}\). This follows
directly from the base inverse margin.

After the scalar sign change, write this linearized operator as
\(a^{ij}\partial_i\partial_j+b^i\partial_i+c\). Its principal blocks
have the ellipticity margin just verified; the off-diagonal blocks have
order at most one. The coefficients \(a-a_\infty\), \(b\) and \(c\) have
weights \(\gamma\), \(\gamma-1\) and \(\gamma-2\), respectively. This
follows by differentiating the constraint expressions: the principal
coefficients are smooth algebraic functions of \(h\); a first-order
coefficient contains one derivative of the metric or one momentum
factor; a zeroth-order coefficient contains two metric derivatives, one
momentum derivative, or products with at least the same decay.

Under \(v_R(z)=R^{-\gamma}v(Rz)\) the coefficients become \(a(Rz)\),
\(Rb(Rz)\) and \(R^2c(Rz)\). They have uniformly bounded finite
\(C^{m-2,\alpha}\) norms on fixed nested annuli. The interior estimate
and the corresponding inner-ball estimate give

\phantomsection\label{eq:online-chart-transverse-weighted-estimate} \[
\|v\|_{\mathcal Y_m}
\le C_m\bigl(\|\mathsf A_{s,y}v\|_{\mathcal Z_m}
+\|v\|_{C^0_\gamma}\bigr).
\] The constant depends on the ellipticity margin and finitely many
higher background seminorms. The base inverse controls the last term for
a solution of \(\mathsf A_{s,y}v=f\). Local elliptic regularity first
improves this base solution, and the annular estimate makes the improved
norm global. Thus

\phantomsection\label{eq:online-chart-transverse-all-order-inverse} \[
\|\mathsf A_{s,y}^{-1}f\|_{\mathcal Y_m}
\le C'_m\|f\|_{\mathcal Z_m},
\qquad m\ge k_0.
\] Base uniqueness identifies these inverses at every order. The bound
uses a constant \(C'_m\) for each fixed \(m\).

For continuity, let \(s_n\to s\) in every \(\mathcal X_m\), and put
\(y_n=\mathcal C(s_n)\), \(y=\mathcal C(s)\). Base smoothness gives
\(y_n\to y\) in \(\mathcal Y_{k_0}\). The spatial induction just proved
gives uniform bounds for \(y_n\) at every finite higher order. Set
\(d_n=s_n-s\) and \(e_n=y_n-y\). On the convex product neighborhood use
the segment \((s_\theta,y_\theta)=(s+\theta d_n,y+\theta e_n)\), and
define \[
L_n=\int_0^1D_y\mathcal F(s_\theta,y_\theta)\,\mathrm d\theta,
\qquad
M_n=\int_0^1D_s\mathcal F(s_\theta,y_\theta)\,\mathrm d\theta.
\] The exact difference equation is

\phantomsection\label{eq:online-chart-exact-difference-equation} \[
L_ne_n=-M_nd_n.
\] The segment is an ambient segment; it need not lie in the constrained
subset. Averaging preserves the fixed ellipticity margin and the order
of the off-diagonal terms. The finite rescaled coefficient bounds are
uniform in \(n\), by the preceding spatial bounds. Therefore the same
annular estimate gives

\phantomsection\label{eq:online-chart-higher-difference-estimate} \[
\|e_n\|_{\mathcal Y_m}
\le C_m\bigl(\|M_nd_n\|_{\mathcal Z_m}
+\|e_n\|_{\mathcal Y_{k_0}}\bigr)
\le C'_m\bigl(\|d_n\|_{\mathcal X_m}
+\|e_n\|_{\mathcal Y_{k_0}}\bigr).
\] The last bound uses that \(D_s\mathcal F\) has at most two
derivatives of the metric seed and one of the momentum seed, together
with the uniform finite background norms. The right-hand side tends to
zero. This proves continuity of \(\mathcal C\) at all weighted orders.
The ambient algebraic formulas give continuity of \(\mathcal T\).

For a supplied smooth parameter \(a\), base differentiation gives

\phantomsection\label{eq:online-chart-parameter-linearization} \[
D_y\mathcal F(s,\mathcal C(s))\,\partial_a\mathcal C(s)
=-D_s\mathcal F(s,\mathcal C(s))\,\partial_a s.
\] The all-order inverse estimate applies to this actual derivative.
Repeated differentiation retains the same invertible left side; its
right side is a finite sum of input derivatives and previously obtained
correction derivatives. At a fixed parameter point the Banach implicit
function theorem can also be applied in any higher space, because the
just-proved higher inverse is an isomorphism there. Base uniqueness
identifies its local solution with the already selected \(\mathcal C\).
This use of a higher-space neighborhood around an existing solution
imposes no additional smallness on that solution's higher norms. It
identifies the derivatives and their continuity in the high space.
Induction gives smooth parameter dependence at each finite order.

On a retained compact parameter set, the higher estimates have uniform
constants when the requisite finite input seminorms and the base margins
are uniform. A slightly larger retained compact set supplies the
parameter neighborhoods and derivative bounds used in this argument. A
finite cover and the differentiated equations give the asserted finite
estimates on the smaller set.

Finally, the decoder has only algebraic operations, one divergence and
the fixed decaying inverse \(\Delta_L^{-1}\). The orders in Proposition
A.4 and Lemma A.2 show its continuity and smooth finite-family
dependence at every weighted level. Its exact identities hold before the
high-order argument and consequently hold for this same smooth branch.
\hfill\(\square\)

\subsection{A.5. Based spatial gauge}\label{a.5.-based-spatial-gauge}\label{sec:online-anchored-spatial-gauge}

The inverse used for a coordinate displacement has a different weight
from the decaying inverse used in the constraint decoder. Fix

\phantomsection\label{eq:online-gauge-fixed-weights} \[
k=k_0\ge4,\qquad 0<\alpha<1,\qquad -1<\gamma<-\frac12,
\qquad d=\gamma+1\in(0,1/2).
\] The displacement is permitted to grow with order \(d\); its
derivative decays with order \(\gamma\). Its value at the marked point
\(o=0\) fixes the translation ambiguity. We use the same base index as
the constraint chart, with the weighted Hölder spaces and Cartesian
tensor conventions specified above.

\pandocstatementheading{\leavevmode\phantomsection\label{lem:online-gauge-anchored-york-inverse}Lemma A.8 (The anchored sublinear York inverse).}

For \(0<d<1\), the operator

\phantomsection\label{eq:online-gauge-anchored-isomorphism} \[
\Delta_L:
\{X\in C^{k+1,\alpha}_d(T\mathbb R^3):X(0)=0\}
\longrightarrow C^{k-1,\alpha}_{d-2}(T\mathbb R^3)
\] is a bounded isomorphism. With \(G\) the matrix Green kernel of the
decaying York inverse, its inverse is the actual integral

\phantomsection\label{eq:online-gauge-anchored-kernel} \[
(\mathcal G_oF)_i(x)
=\int_{\mathbb R^3}
\bigl(G_{ij}(x-y)-G_{ij}(-y)\bigr)F_j(y)\,\mathrm dy,
\] and

\phantomsection\label{eq:online-gauge-anchored-bound} \[
\|\mathcal G_oF\|_{C^{k+1,\alpha}_d}
\le C_{k,\alpha,d}\|F\|_{C^{k-1,\alpha}_{d-2}}.
\] \emph{Proof.} The two possible singularities, at \(y=x\) and \(y=0\),
are locally integrable in three dimensions. For \(|y|>2|x|+2\), the
degree \(-1\) homogeneity of \(G\) and the mean-value formula imply \[
|G(x-y)-G(-y)|\le C|x|\,|y|^{-2}.
\] The absolute tail is bounded by
\(C\|F\|_{C^0_{d-2}}|x|\int_{2|x|+2}^{\infty}t^{d-2}\,\mathrm dt\),
which is finite because \(d<1\). The integral vanishes at \(x=0\).

For the weighted zeroth-order estimate, let \(R=|x|\ge2\). On
\(|y|\le2R\), the subtracted kernel term has integral bounded by \[
C\int_0^{2R}t\langle t\rangle^{d-2}\,\mathrm dt
\le C_dR^d.
\] For \(G(x-y)\) split this region into \(|y|<R/2\), \(|x-y|<R/2\) and
the remaining comparable annulus. On the first part the kernel is
bounded by \(C/R\), giving
\(CR^{-1}\int_0^{R/2}t^2\langle t\rangle^{d-2}\,\mathrm dt\le CR^d\). On
the second part the source is bounded by \(CR^{d-2}\) and the kernel has
integral \(O(R^2)\). On the remaining part rescaling by \(R\) gives the
same bound. The outer difference term is bounded by
\(CR\int_{2R}^{\infty}t^{d-2}\,\mathrm dt\le C_dR^d\). The region
\(|x|\le2\) has a uniform bound by local integrability and the tail
estimate. Splitting small neighborhoods of the integrable singularities
also proves continuity. Hence \[
\|\mathcal G_oF\|_{C^0_d}\le C_d\|F\|_{C^0_{d-2}}.
\] Applying \(\Delta_L\) in the \(x\) variable kills the subtracted term
and gives \(F\). This holds distributionally: on a compact test support
the same tail difference is integrable, so the kernel calculation can be
tested and integrated. The local elliptic regularity and interior
estimate stated above give the higher regularity. On nested annuli set
\(v_R(y)=X(Ry)\); the equation is \(\Delta_{L,y}v_R=R^2F(Ry)\).
Multiplication of the local estimate by \(R^{-d}\) gives exactly the
asserted higher weighted bound. A fixed ball handles the origin.

For uniqueness, suppose \(\Delta_LX=0\) in the stated domain. Taking
divergence gives \((4/3)\Delta\operatorname{div}X=0\). The divergence is
a smooth harmonic scalar of weight \(d-1<0\), so it is zero by the
maximum principle and decay at infinity. Thus every component of \(X\)
is harmonic. Each first derivative is harmonic and has the negative
weight \(d-1\), so every first derivative is zero. The field is constant
and its value at the anchor is zero. This proves injectivity.
\hfill\(\square\)

For a based displacement \(X\), put \(\phi_X(x)=x+X(x)\). If
\(\|DX\|_{L^\infty}<1\), then \(\phi_X\) is a global diffeomorphism. For
every \(y\) the map \(x\mapsto y-X(x)\) is a contraction of complete
Euclidean space and has one fixed point. The differential of \(\phi_X\)
is invertible, so its inverse is smooth to the available order. The path
\(x\mapsto x+tX(x)\) has invertible differential for \(0\le t\le1\),
hence the diffeomorphism preserves orientation. Since \(X(0)=0\),

\phantomsection\label{eq:online-gauge-radius-comparability} \[
(1-\|DX\|_\infty)|x|
\le|\phi_X(x)|
\le(1+\|DX\|_\infty)|x|.
\] This comparability and the chain rule on rescaled annuli give the
weighted composition estimates. In particular, if \(X\in\mathscr S^d\)
and \(d<1\), the inverse displacement is also in \(\mathscr S^d\). To
see the weight explicitly, write \(Y=-X\circ\phi_X^{-1}\). Its
zeroth-order estimate follows from comparability. Differentiating the
inverse equation first gives
\(D\phi_X^{-1}=(I+DX\circ\phi_X^{-1})^{-1}\); subsequent derivatives are
finite sums of derivatives of \(X\) and lower derivatives of the
inverse. Induction gives the claimed symbol order for \(Y\). On an
annulus the normalized displacement is \(R^{-1}X(R\cdot)=O(R^{d-1})\),
so the same estimates are uniform in its radius.

For a positive metric \(h\), the actual gauge residual is

\phantomsection\label{eq:online-gauge-residual} \[
\mathcal B_h(X)
=\operatorname{div}_\delta\operatorname{TF}_\delta(\phi_X^*h).
\] In Cartesian coordinates, with
\(A_i{}^a=\delta_i{}^a+\partial_iX^a\), its input tensor is

\phantomsection\label{eq:online-gauge-actual-pullback} \[
(\phi_X^*h)_{ij}(x)
=h_{ab}(x+X(x))A_i{}^aA_j{}^b.
\] Both Jacobian factors belong to this definition.

\pandocstatementheading{\leavevmode\phantomsection\label{thm:online-gauge-based-normalization}Theorem A.9 (Based normalization of the metric gauge).}

There are positive constants \(\epsilon\) and \(\rho\), fixed at the
finite order in \hyperref[eq:online-gauge-fixed-weights]{the weight
convention}, such that every metric satisfying

\phantomsection\label{eq:online-gauge-base-metric-domain} \[
\|h-\delta\|_{C^{k+2,\alpha}_\gamma}<\epsilon
\] has a unique based displacement \(X_h\in C^{k+1,\alpha}_d\) in the
ball \(\|X_h\|<\rho\) with \(\mathcal B_h(X_h)=0\). Its derivative has
norm less than \(1/2\). Thus \(\phi_h=x+X_h\) is a global
orientation-preserving diffeomorphism fixing the origin. Moreover,

\phantomsection\label{eq:online-gauge-residual-displacement-bound} \[
\|X_h\|_{C^{k+1,\alpha}_d}
\le C\|\operatorname{div}_\delta\operatorname{TF}_\delta h
       \|_{C^{k-1,\alpha}_{\gamma-1}}.
\] It follows that \(\phi_h\) is the identity whenever \(h\) already
satisfies the gauge. If \(h-\delta\in\mathscr S^\gamma\), then
\(X_h\in\mathscr S^d\). For a finite-dimensional family smooth in all
these weighted seminorms and remaining in the one base domain, the same
\(X_h\) is smooth in every weighted spatial order. Each specified finite
spatial and parameter estimate uses finitely many higher input seminorms
and the fixed base margins. Higher seminorms need not be small.

\emph{Proof.} At the flat zero input,

\phantomsection\label{eq:online-gauge-flat-linearization} \[
D_X\mathcal B_\delta(0)V
=\operatorname{div}_\delta\operatorname{TF}_\delta
 (\mathcal L_V\delta)
=\Delta_LV.
\] Write \(r=h-\delta\) and \[
\mathcal B_h(X)=b_h+\Delta_LX+\mathcal N_h(X),
\qquad
b_h=\operatorname{div}_\delta\operatorname{TF}_\delta h,
\qquad \mathcal N_h(0)=0.
\] On a sufficiently small based displacement ball, the actual pullback
formula gives

\phantomsection\label{eq:online-gauge-remainder-estimate} \[
\begin{aligned}
&\|\mathcal N_h(X)-\mathcal N_h(Y)\|_{C^{k-1,\alpha}_{\gamma-1}}
\\[-.3ex]&\quad\le C\bigl(\|r\|_{C^{k+2,\alpha}_\gamma}
+\|X\|_{C^{k+1,\alpha}_d}
+\|Y\|_{C^{k+1,\alpha}_d}\bigr)
\|X-Y\|_{C^{k+1,\alpha}_d}.
\end{aligned}
\] Here is why the stated orders suffice. Differentiate the pullback
once in space, expand its two Jacobian factors, and subtract the flat
terms linear in \(X\). Every remaining displacement term contains a
metric perturbation or at least two displacement factors. Subtracting
two such expressions and using the mean-value formula for composed
coefficients gives terms of the types \[
rD^2X,\qquad Dr\,DX,\qquad D^2r\,X,
\qquad DX\,D^2X,
\] with smooth contractions and composed versions of the coefficients.
Their weights are at most \(2\gamma-1\), which embeds into \(\gamma-1\)
because \(\gamma<0\). On an annulus of radius \(R\), all arguments lie
in one enlarged comparable annulus by
\hyperref[eq:online-gauge-radius-comparability]{radius comparability}.
The normalized displacement is \(R^{-1}X(R\cdot)\), of order
\(R^\gamma\). Ordinary Hölder product and composition estimates on fixed
annuli therefore have uniform radius-independent constants. Taking
\(k-1\) derivatives and their Hölder differences is covered by the
displayed \(k+2\) metric norm and \(k+1\) displacement norm. The
additional metric derivatives supply the mean-value bound in the highest
Hölder norm retained here. This accounts for the derivative cost of the
diffeomorphism action. The bounded region is treated on fixed balls.
These observations prove the remainder estimate.

Use Lemma A.8 to define

\phantomsection\label{eq:online-gauge-initialized-iteration} \[
\mathcal T_h^{\mathrm g}(X)
=-\mathcal G_o\bigl(b_h+\mathcal N_h(X)\bigr).
\] Choose \(\rho\) and then \(\epsilon\) so that this map has Lipschitz
constant at most \(1/2\) and sends the closed based ball of radius
\(\rho\) into itself. This follows from the remainder bound,
\(\mathcal N_h(0)=0\), and the small bound on \(b_h\). Its fixed point
is unique in that ball. The inverse identity of Lemma A.8 makes its
equation equivalent to \(\mathcal B_h(X_h)=0\). Applying the norm
estimate to the fixed-point equation and absorbing its Lipschitz part
proves \hyperref[eq:online-gauge-residual-displacement-bound]{the
residual bound}. If \(b_h=0\), the initialized zero displacement is a
fixed point, so uniqueness makes \(X_h=0\). Decrease the one base radius
so that the derivative bound is less than \(1/2\). The preceding global
inverse argument then applies.

We next improve this same solution. The coordinate formula writes its
equation as a quasilinear second-order system for \(X\). The
second-order coefficients depend on \(h(x+X)\) and \(DX\); all terms
involving derivatives of \(h\) use at most its first derivative before
the equation is differentiated. The positive principal symbol of
\(-D_X\mathcal B_h\) at \((h,X)=(\delta,0)\) is \[
|\xi|^2I+\frac13\xi\otimes\xi.
\] Its eigenvalues are \(|\xi|^2\) and \((4/3)|\xi|^2\). The small base
metric and first-derivative bounds make the actual symbol a uniformly
strongly elliptic perturbation. This ellipticity margin is fixed before
every higher derivative order.

Suppose the existing \(X\) has weighted order \(C^{m,\alpha}_d\), with
\(m\ge k+1\). On a fixed annulus the principal coefficients then have
\(C^{m-1,\alpha}\) regularity, and the remaining terms have that same
regularity using finitely many derivatives of the supplied smooth
metric. The stated interior system Schauder estimate, applied to the
existing solution on nested annuli, gives one additional derivative. In
rescaled coordinates the equation has the same form for \[
X_R(y)=R^{-1}X(Ry),\qquad h_R(y)=h(Ry).
\] Both \(h_R-\delta\) and \(X_R\) have size \(O(R^\gamma)\) in the
corresponding already controlled norms. Their nonlinear products have at
least this decay, and the normalized composition arguments stay in fixed
comparable annuli. The interior estimates consequently give a finite
bound for \(R^{-\gamma}\|X_R\|_{C^{m+1,\alpha}}\) using finitely many
input seminorms and the preceding bound. Taking the supremum over
\(R\ge1\) proves the weighted gain. A fixed ball handles the origin.
Induction gives every finite weighted order without decreasing the base
neighborhood.

For the linearization at the selected solution, put \(d=\gamma+1\). In
this proof write \(E_m=\{v\in C_d^{m,\alpha}:v(0)=0\}\) and
\(F_m=C_{d-2}^{m-2,\alpha}\) for spaces of vector fields, with
\(m\ge k+1\). Then
\(L=D_X\mathcal B_h(X_h)=\Delta_L+D_X\mathcal N_h(X_h)\) maps \(E_m\)
boundedly to \(F_m\) at the regularity just obtained. At \(m=k+1\),
Lemma A.8 and the contraction estimate give the inverse
\((I+\mathcal G_o D_X\mathcal N_h(X_h))^{-1}\mathcal G_o\). For
\(f\in F_m\) apply this base inverse to obtain \(v\). Interior Schauder
regularity for \(Lv=f\) on a fixed ball and on rescaled annuli gives \[
\|v\|_{E_m}\le C_m\bigl(\|f\|_{F_m}+\|v\|_{C_d^0}\bigr).
\] Indeed rescaling \(v(R\cdot)\) by \(R^{-d}\) and \(f(R\cdot)\) by
\(R^{2-d}\) gives uniform estimates on nested unit annuli by the
coefficient bounds above. The base inverse bounds the last term by
\(C\|f\|_{F_{k+1}}\le C'\|f\|_{F_m}\). Hence the same based solution
lies in \(E_m\) and satisfies \(\|v\|_{E_m}\le C'_m\|f\|_{F_m}\). This
proves surjectivity at order \(m\), and injectivity follows by inclusion
in the base space.

Fix a finite differentiation order \(j\ge1\). On the open set of
positive metrics and based displacements satisfying the strict base
bounds and the diffeomorphism condition, \((h,X)\mapsto\mathcal B_h(X)\)
is \(C^j\) from \((\delta+C_\gamma^{m+j,\alpha})\times E_m\) to \(F_m\).
Before spatial differentiation the highest metric derivative in the
\(j\)th displacement derivative is \((D^{j+1}h)\circ\phi_X\). Taking the
\(m-2\) target derivatives uses at most \(m+j-1\) metric derivatives.
The remaining derivative gives continuity of composition in the top
Hölder norm by the mean-value estimate on enlarged annuli. The same
argument for a metric variation, together with the rescaled product
estimates above for the Jacobian factors and the displacement
variations, gives the mixed Fréchet derivatives with the stated weights.
The Banach implicit-function theorem at \((h,X_h)\) now gives a \(C^j\)
solution map into \(E_m\) locally in this higher metric norm. Its values
solve the base equation in the uniqueness ball, hence coincide with
\(X_h\) there. In particular the parameter derivatives through order
\(j\) are continuous in \(E_m\). Differentiating the equation expresses
each highest parameter derivative through the bounded inverse of \(L\)
and terms of lower parameter order.

These differentiability arguments are applied at points of the already
fixed parameter domain where the base smallness holds. For any compact
subset, a finite cover and a slightly larger retained compact set bound
the finitely many input seminorms and derivatives used at a specified
order. The output is therefore smooth at every order on the original
retained domain; no new solution or higher-order smallness condition is
selected. \hfill\(\square\)

\pandocstatementheading{\leavevmode\phantomsection\label{prop:online-gauge-rotation-and-spin}Proposition A.10 (Rotation and actual spin markings).}

Choose equivalent rotation-invariant base norms and the preceding
domains invariant under proper Euclidean rotations. For
\(A\in\operatorname{SO}(3)\),

\phantomsection\label{eq:online-gauge-rotation-equivariance} \[
X_{A^*h}(x)=A^{-1}X_h(Ax),
\qquad
\phi_{A^*h}=A^{-1}\circ\phi_h\circ A.
\] The path \(\phi_{h,t}(x)=x+tX_h(x)\), \(0\le t\le1\), gives the
initialized actual spin lift of the normalization. It varies smoothly
with every stipulated finite-dimensional smooth family. If a comparison
is supplied with the opposite central lift, that lift remains a distinct
marking.

\emph{Proof.} The Euclidean trace, divergence and anchor are preserved
by proper rotations. The matrix kernel obeys \(G(Ax)=AG(x)A^{-1}\), so
the anchored inverse commutes with the corresponding tensor pullbacks.
Alternatively, direct substitution in the gauge equation shows that the
displayed conjugated displacement solves the rotated equation. It has
the same rotation-invariant norm and belongs to the same based
uniqueness ball. Theorem A.9 gives the two identities.

Every \(\phi_{h,t}\) is a based orientation-preserving diffeomorphism by
the fixed derivative bound. Transport orthonormal frames by the
differential and the positive polar identification between the
corresponding metrics. This yields a smooth path in the oriented frame
bundle starting at the stipulated identity marking. The path-lifting
property of its spin double cover gives one actual lift with that
initial value. Smooth dependence follows from local trivializations of
the covering and the same initialization. The two lifts differ by the
central deck transformation; fixing the initial lift determines which
one is transported. \hfill\(\square\)

\pandocstatementheading{\leavevmode\phantomsection\label{cor:online-gauge-whole-data-decoding}Corollary A.11 (Access to the exact maximal-data decoder).}

Let \((H,L)\) be whole smooth maximal vacuum data on the marked
\(\mathbb R^3\), with \[
H-\delta\in\mathscr S^\gamma,
\qquad L\in\mathscr S^{\gamma-1},
\] and sufficiently small in
\(C^{k+2,\alpha}_\gamma\oplus C^{k+1,\alpha}_{\gamma-1}\). After
applying the based gauge \(\phi_H\), the normalized data

\phantomsection\label{eq:online-gauge-normalized-whole-data} \[
(h,K)=\phi_H^*(H,L)
\] belong to the independently defined gauged maximal constraint
neighborhood. Their exact ambient decoder gives a unique smooth pair of
transverse-traceless seeds \((q,p)\) in its selected branch, and

\phantomsection\label{eq:online-gauge-exact-chart-recovery} \[
(h,K)=\mathcal T(q,p).
\] All indicated higher spatial and finite-family estimates hold for
this same gauge and decoded pair. The actual spin comparison is the one
in Proposition A.10.

\emph{Proof.} Pullback preserves the three maximal vacuum constraints.
Theorem A.9 imposes the metric gauge. Weighted composition with the
sublinear displacement preserves the symbol orders of both tensors; the
Jacobian factors are the identity plus a decaying tensor. The gauge
estimate and the smallness of the whole pair put the normalized fields
in the ambient decoder neighborhood after one restriction of the finite
base threshold.

The exact decoder gives

\phantomsection\label{eq:online-gauge-decoded-seeds} \[
\begin{aligned}
q&=\operatorname{TF}_\delta h,
\qquad u=\frac13\operatorname{tr}_\delta h-1,\\
\kappa&=K-\frac{\operatorname{tr}_\delta K}{\operatorname{tr}_\delta h}\,h,
\qquad W=\Delta_L^{-1}\operatorname{div}_\delta\kappa,
\qquad p=\kappa-L_0W.
\end{aligned}
\] Its algebraic identities and the decaying York inverse put \(q,p\) in
the stated transverse-traceless spaces. The independently satisfied
constraints force \((u,W)=\mathcal C(q,p)\) by uniqueness of the
nonlinear correction in its base branch. This proves the tensor
equality. The same-branch weighted estimates for the chart and gauge,
and continuity of the exact decoder at every finite order, give all
stated finite-family bounds. Proposition A.10 supplies the initialized
actual spin marking. \hfill\(\square\)

The gauge and the decoder operate on the whole fields before the
resulting data or source are restricted to a local patch. The allowed
displacement is based and sublinear. The source comparisons for
different whole extensions and different initial hypersurfaces use the
later local comparison theorem; they are not prerequisites for this
normalization.

\subsection{A.6. Current differential
operations}\label{a.6.-current-differential-operations}\label{sec:online-current-residual-derivatives}

Use the spaces \(\mathcal X_m,\mathcal Y_m,\mathcal Z_m\) of
\hyperref[def:online-chart-seeds-and-target]{Definition A.3}, with the
same \(k_0,\alpha,\gamma\). A current state is \(Z=(s,y)\), where
\(s=(q,p)\) and \(y=(u,W)\). Its coefficients are

\phantomsection\label{eq:online-current-ambient-coefficients} \[
h=\delta+q+u\delta,\qquad
\kappa=p+L_0W,\qquad
K=\kappa-\frac13\vartheta_\kappa h,\qquad
\vartheta_\kappa=\operatorname{tr}_h\kappa.
\] The residual \(\mathcal F(Z)=(R(h)-|K|_h^2,\operatorname{div}_hK)\)
is defined throughout the positive ambient product. Its current
derivatives are finite differential operations, including off the
zero-residual set. The inverse used below is the completed operation

\phantomsection\label{eq:online-current-source-flat-inverse} \[
\mathsf G_0^{\mathrm{src}}
=\operatorname{diag}\left(-\frac12\Delta^{-1},\Delta_L^{-1}\right)
=A_0^{-1}
\] constructed from the finite negative-response word in
\hyperref[thm:online-negative-response-flat-green]{the source Green
theorem}. Its bounded realization on \(\mathcal Z_m\to\mathcal Y_m\) is
the decaying inverse of Lemma A.2. Thus the following limits have a
specified Banach-space meaning as well as an actual finite-word
realization on weighted smooth fields.

\pandocstatementheading{\leavevmode\phantomsection\label{lem:online-current-residual-differential}Lemma A.12 (The current residual differential).}

For a seed variation \(a=(a_q,a_p)\) and a correction variation
\(v=(v_u,v_W)\) put

\phantomsection\label{eq:online-current-tracefree-variation} \[
\begin{gathered}
r=a_q+v_u\delta, \\
z=a_p+L_0v_W, \\
\dot\vartheta_\kappa=\operatorname{tr}_h z-r^{ij}\kappa_{ij}, \\
k=z-\frac13\dot\vartheta_\kappa h-\frac13\vartheta_\kappa r,
\end{gathered}
\] where \(r^{ij}=h^{ia}h^{jb}r_{ab}\). All covariant derivatives below
use \(h\). Define \[
C^\ell{}_{ij}(r)=\frac12h^{\ell a}
\bigl(\nabla_i r_{ja}+\nabla_jr_{ia}-\nabla_ar_{ij}\bigr).
\] Then the scalar and one-form components of \(D\mathcal F_Z[a,v]\) are

\phantomsection\label{eq:online-current-explicit-linearized-residual} \[
\begin{aligned}
\mathfrak H_Z(r,k)
={}&\nabla^i\nabla^jr_{ij}
-\Delta_h\operatorname{tr}_hr
-\operatorname{Ric}(h)^{ij}r_{ij}\\
&-2K^{ij}k_{ij}
+2(K\circ K)^{ij}r_{ij},\\
\mathfrak M_Z(r,k)_i
={}&\nabla^jk_{ji}
-r^{jk}\nabla_jK_{ki}\\
&-h^{jk}\bigl(C^\ell{}_{jk}(r)K_{\ell i}
             +C^\ell{}_{ji}(r)K_{k\ell}\bigr),
\end{aligned}
\] where \((K\circ K)^{ij}=K^i{}_\ell K^{\ell j}\). In particular,

\phantomsection\label{eq:online-current-transverse-and-seed-derivatives}
\[
\mathsf A_Z=D_y\mathcal F_Z,\qquad
\mathsf B_Z=D_s\mathcal F_Z
\] are obtained by setting \(a=0\) or \(v=0\) in these formulas. At the
flat state, \(\mathsf A_0=A_0=(-2\Delta,\Delta_L)\) and
\(\mathsf B_0=0\). The maps are continuous at each indicated weighted
order and real analytic on the positive base neighborhood.

\emph{Proof.} Differentiating \(h^{-1}h=I\) gives
\(\dot h^{ij}=-r^{ij}\). Differentiate
\(\vartheta_\kappa=h^{ij}\kappa_{ij}\) and
\(K=\kappa-\vartheta_\kappa h/3\) to obtain the first displayed
variation. In particular it retains the variation of the tracefree
projection.

Differentiating the Koszul formula gives \(C^\ell{}_{ij}\). Its Ricci
contraction gives
\(DR_h[r]=\nabla^i\nabla^jr_{ij}-\Delta_h\operatorname{tr}_hr-\operatorname{Ric}^{ij}r_{ij}\).
Differentiation of the two inverse metrics in \(|K|_h^2\) gives
\(D|K|_h^2[r,k]=2K^{ij}k_{ij}-2(K\circ K)^{ij}r_{ij}\). Finally,
differentiate \(h^{jk}\nabla_jK_{ki}\). The inverse-metric term is the
displayed \(-r^{jk}\nabla_jK_{ki}\), and differentiating each covariant
slot of \(K\) gives the two terms involving \(C\). This proves the
formulas for the covariant momentum one-form.

The flat identities follow from the TT conditions and the definition of
\(L_0\). Each formula contains at most two derivatives of \(r\) and one
of \(k\). Weighted multiplication therefore gives the asserted maps into
\(\mathcal Z_m\). Matrix inversion is real analytic in the fixed
positive neighborhood; the other operations are finite differentiations,
contractions and products. This proves the final assertion.
\hfill\(\square\)

\subsection{A.7. The current inverse and horizontal
fields}\label{a.7.-the-current-inverse-and-horizontal-fields}\label{sec:online-current-inverse-and-flatness}

At a current state set

\phantomsection\label{eq:online-current-horizontal-fields} \[
\begin{aligned}
\mathsf E_Z&=\mathsf G_0^{\mathrm{src}}(\mathsf A_Z-A_0),\\
\mathsf G_Z&=\sum_{j=0}^{\infty}(-\mathsf E_Z)^j\mathsf G_0^{\mathrm{src}},\\
\mathsf V_Za&=-\mathsf G_Z\mathsf B_Za,\qquad
\mathscr H_a(Z)=(a,\mathsf V_Za).
\end{aligned}
\] The parameter direction \(a\) changes the whole input seed. These are
vector fields on the space of source coefficients, distinct from the
physical point directions of a spacetime.

\pandocstatementheading{\leavevmode\phantomsection\label{prop:online-current-inverse-horizontal-flatness}Proposition A.13 (Current inverse and flatness).}

There is a convex base product
\(D\subset\mathcal X_{k_0}\times\mathcal Y_{k_0}\) around zero on which
the metric has a fixed positive margin and

\phantomsection\label{eq:online-current-fixed-inverse-margin} \[
\|\mathsf E_Z\|_{\mathcal L(\mathcal Y_{k_0})}
\le\varkappa<\frac12.
\] The series defining \(\mathsf G_Z\) converges in the base operator
norm and equals \(\mathsf A_Z^{-1}\). The maps
\(\mathsf G_Z,\mathsf V_Z\) are real analytic there. For constant seed
directions \(a,b\),

\phantomsection\label{eq:online-current-horizontal-conservation-and-flatness}
\[
D\mathcal F_Z\,\mathscr H_a(Z)=0,\qquad
[\mathscr H_a,\mathscr H_b]=0.
\] At a weighted smooth current state, this same inverse maps
\(\mathcal Z_m\) boundedly to \(\mathcal Y_m\) at every finite
\(m\ge k_0\). Each bound depends on finitely many higher state seminorms
and the fixed base margins. The same higher inverse estimate holds for
the average of \(\mathsf A\) along any segment in \(D\) whose endpoints
have the stated higher bounds.

\emph{Proof.} The quadratic residual estimate in Lemma A.5 gives, on a
smaller product, \[
\|\mathsf E_Z\|
+\|\mathsf G_0^{\mathrm{src}}\mathsf B_Z\|
\le C\bigl(\|s\|_{\mathcal X_{k_0}}+\|y\|_{\mathcal Y_{k_0}}\bigr).
\] Choose the product convex and small enough for the stated margin.
Since \(\mathsf A_Z=A_0(I+\mathsf E_Z)\), the Neumann series gives the
actual inverse with norm at most
\((1-\varkappa)^{-1}\|\mathsf G_0^{\mathrm{src}}\|\). The current
derivatives in Lemma A.12 are analytic. The convergent inverse, or the
inverse identity differentiated on this neighborhood, therefore gives
analytic dependence. In particular, for a state variation \(\zeta\),

\phantomsection\label{eq:online-current-horizontal-differential} \[
D\mathsf V_Z[\zeta]a
=-\mathsf G_Z\bigl(D\mathsf A_Z[\zeta]\mathsf V_Za
                   +D\mathsf B_Z[\zeta]a\bigr).
\] Residual conservation is the identity
\(\mathsf B_Za+\mathsf A_Z\mathsf V_Za=0\). Differentiate it in the
direction \(\mathscr H_b\) and subtract the corresponding equation with
\(a,b\) exchanged. The Hessian terms cancel by symmetry, giving
\(D\mathcal F_Z[\mathscr H_a,\mathscr H_b]=0\). The seed component of
the bracket vanishes because \(a,b\) are constant. Its remaining
component is killed by the invertible \(\mathsf A_Z\), proving flatness.
No zero-residual hypothesis is used here.

The proof of Theorem A.7 already gives the required transverse estimate
at every smooth point in a sufficiently small ambient product, rather
than only on its constraint graph. After changing the sign of the scalar
row, the principal blocks are uniformly strongly elliptic and the
off-diagonal terms have order at most one. Relative to the flat blocks,
their second-, first- and zeroth-order coefficients have weights
\(\gamma,\gamma-1,\gamma-2\). On nested rescaled annuli the estimate is

\phantomsection\label{eq:online-current-higher-transverse-estimate} \[
\|v\|_{\mathcal Y_m}
\le C_m\bigl(\|\mathsf A_Zv\|_{\mathcal Z_m}
                  +\|v\|_{C^0_\gamma}\bigr).
\] Apply local regularity to the existing base solution, then this
annular estimate and the interior-ball estimate. The base inverse
controls the last term. This proves the higher bound for the same
inverse, without requiring high-order convergence of the raw Neumann
partial sums.

For an average \(\overline{\mathsf A}\) along a segment, the normalized
perturbation is the average of \(\mathsf E\), so it has the same base
norm margin. Convexity also preserves the near-flat principal
ellipticity margin, and the finite coefficient estimates are bounded by
the stipulated finite endpoint bounds. The same proof applies.
\hfill\(\square\)

\subsection{A.8. Transport on one base
tube}\label{a.8.-transport-on-one-base-tube}\label{sec:online-current-flat-input-transport}

Fix a weighted smooth reference state \((s_*,y_*)\in D\). Choose outer
balls \(B(s_*,r_s)\times B(y_*,r_y)\) contained in \(D\) and constants
\(\|\mathsf V\|\le M\), \(\|D_y\mathsf V\|\le L\) there. Take \(\rho>0\)
such that

\phantomsection\label{eq:online-current-transport-base-radii} \[
\rho<r_s,\qquad M\rho<r_y/4,\qquad L\rho<1/4,
\] omitting a condition with zero left coefficient. Set
\(S=B_{\mathcal X_{k_0}}(s_*,\rho)\) and
\(J=B_{\mathcal Y_{k_0}}(y_*,r_y/4)\). The ordinary-pair constructor
uses the fixed flat choice \((s_*,y_*)=(0,0)\). The general fiber labels
below describe transport within the chosen coefficient neighborhood;
they are not additional generating inputs to that constructor.

\pandocstatementheading{\leavevmode\phantomsection\label{thm:online-current-input-transport}Theorem A.14 (Actual current transport and its endpoint).}

For \(s\in S\), \(y_0\in J\) and \(a=s-s_*\), the iteration

\phantomsection\label{eq:online-current-picard-iteration} \[
\begin{aligned}
y^{[0]}(t)&=y_0,\\
y^{[j+1]}(t)
&=y_0+\int_0^t
\mathsf V_{(s_*+\tau a,y^{[j]}(\tau))}a\,\mathrm d\tau
\end{aligned}
\] converges in \(C^0([0,1];\mathcal Y_{k_0})\). Its endpoint
\(\Psi(s,y_0)\) satisfies

\phantomsection\label{eq:online-current-endpoint-identities} \[
\begin{aligned}
\mathcal F(s,\Psi(s,y_0))&=\mathcal F(s_*,y_0),\\
D_s\Psi(s,y_0)&=\mathsf V_{(s,\Psi(s,y_0))},\\
D_{y_0}\Psi(s,y_0)&=
\mathsf G_{(s,\Psi(s,y_0))}\mathsf A_{(s_*,y_0)}.
\end{aligned}
\] The map \((s,y_0)\mapsto(s,\Psi(s,y_0))\) is a diffeomorphism onto an
open tube in \(D\). Every piecewise \(C^1\) seed path remaining in \(S\)
has a unique horizontal lift from any point of that tube, with its fiber
label fixed. Its endpoint depends only on the endpoint seed, and the
transports satisfy

\phantomsection\label{eq:online-current-transport-composition} \[
\mathscr T_{s_2,s_1}\mathscr T_{s_1,s_0}
=\mathscr T_{s_2,s_0},\qquad
\mathscr T_{s,s}=I.
\] These assertions apply to paths of arbitrary finite length in \(S\).

For weighted smooth seeds and fiber labels, the same base endpoint
belongs to every \(\mathcal Y_m\). On sets with the required finite
higher bounds,

\phantomsection\label{eq:online-current-higher-endpoint-difference} \[
\|\Psi(s,y_0)-\Psi(t,\widetilde y_0)\|_{\mathcal Y_m}
\le C_m\bigl(\|s-t\|_{\mathcal X_m}
                   +\|y_0-\widetilde y_0\|_{\mathcal Y_m}\bigr).
\] Smooth finite-dimensional families are preserved. For compact
parameter sets \(\mathcal K\Subset\operatorname{int}\mathcal K'\), every
specified finite spatial and parameter estimate on \(\mathcal K\) uses
finitely many input seminorms on \(\mathcal K'\) and the fixed base
margins. The input ball and the spatial domain are fixed before these
higher orders.

If \(\mathcal F(s_*,y_*)=0\), put \(Y(s)=\Psi(s,y_*)\). The entire
independently tested zero-residual set in \(S\times B(y_*,r_y)\)
consists precisely of \((s,Y(s))\).

\emph{Proof.} On the closed path ball
\(\sup_t\|y(t)-y_0\|_{\mathcal Y_{k_0}}\le r_y/4\), the integral map
takes values in the same ball and has contraction constant
\(q_{\mathrm P}=L\rho<1/4\). All its states stay in the outer product.
Its iterates converge with error at most
\(q_{\mathrm P}^jM\|a\|/(1-q_{\mathrm P})\). The limit is \(C^1\) in
\(t\) and solves the integral equation. The integral difference estimate
gives uniqueness while a competing solution remains in the ambient
domain.

The derivative in the path variable has norm at most \(q_{\mathrm P}\).
The differentiated fixed-point equation therefore has the bounded
inverse of \(I\) minus that derivative. Repeated differentiation proves
smooth dependence in the base Banach spaces. Along the actual limit the
chain rule gives
\(\partial_t\mathcal F=\mathsf B a+\mathsf A\mathsf V a=0\).
Differentiate this conserved identity in \(s\) and \(y_0\) and use the
current inverse to obtain all three endpoint identities.

There is an independent injectivity estimate on every convex fiber
segment in \(D\):

\phantomsection\label{eq:online-current-fiber-injectivity} \[
\begin{aligned}
\mathsf G_0^{\mathrm{src}}\bigl(\mathcal F(s,y)-\mathcal F(s,z)\bigr)
&=\left(I+\int_0^1
\mathsf E_{(s,z+\tau(y-z))}\,\mathrm d\tau\right)(y-z),\\
(1-\varkappa)\|y-z\|_{\mathcal Y_{k_0}}
&\le
\|\mathsf G_0^{\mathrm{src}}(\mathcal F(s,y)-\mathcal F(s,z))\|_{\mathcal Y_{k_0}}.
\end{aligned}
\] The derivative in the fiber label in the endpoint identities is
invertible, so the product map is a local diffeomorphism. If two labels
have the same output over \(s\), residual conservation and this
injectivity estimate at \(s_*\) show that the labels agree. Its image is
consequently an open tube and its local inverses agree to give a global
inverse there.

For an arbitrary seed path \(s(t)\), the field \(\Psi(s(t),y_0)\) stays
in this tube and, by its seed derivative, solves the horizontal
equation. Local uniqueness gives its unique lift along the whole path.
Thus the label remains fixed, proving endpoint independence and the
composition laws without a length restriction beyond the stated finite
path domain. If the reference residual is zero, conservation gives
\(\mathcal F(s,Y(s))=0\). The fiber estimate identifies every other
zero-residual state in the larger outer fiber ball with \(Y(s)\).

The constructed base endpoint solves the actual equation

\phantomsection\label{eq:online-current-endpoint-elliptic-equation} \[
\mathcal F(s,y)=\mathcal F(s_*,y_0).
\] For weighted smooth input, its right side lies in every
\(\mathcal Z_m\). Apply the same spatial induction as in Theorem A.7,
now with this prescribed right side. If the existing correction has
order \(j\), the principal coefficients have the required order \(j-1\);
the prescribed metric seed is used through order \(j+1\), the momentum
seed through order \(j\), and the residual through order \(j-1\). All
other terms contain only the already controlled correction derivatives.
Interior system regularity improves the existing solution by one order.
Rescale nested annuli exactly as in that theorem; the factors
\(R^\gamma\) stay bounded and the fixed base ellipticity margin is
unchanged. Taking the supremum and adding the inner-ball estimate gives
the next weighted norm. Each finite induction uses finitely many input
norms. This proves all spatial orders for the one base endpoint.

For two endpoints \(y,\widetilde y\), integrate the derivative of
\(\mathcal F\) along the ambient segment joining \((s,y)\) to
\((t,\widetilde y)\). The exact equation is

\phantomsection\label{eq:online-current-exact-endpoint-difference} \[
\overline{\mathsf A}(y-\widetilde y)
=-\overline{\mathsf B}(s-t)
+\mathcal F(s_*,y_0)-\mathcal F(s_*,\widetilde y_0).
\] The segment lies in \(D\). Proposition A.13 applies to its averaged
transverse operator; the preceding spatial bounds control the finite
coefficients. The residual difference at fixed \(s_*\) has
\(\mathcal Z_m\) norm controlled by the \(\mathcal Y_m\) label
difference, using finitely many reference and label bounds. The seed
term has the corresponding \(\mathcal X_m\to\mathcal Z_m\) estimate.
This proves the higher difference bound and continuity at every weighted
order.

At any already obtained smooth endpoint, the same transverse operator is
an isomorphism at each higher order. Apply the Banach implicit function
theorem in that higher space to the endpoint equation near this
endpoint; base uniqueness identifies its local solution with the
constructed one. This identifies higher-space parameter derivatives
without imposing smallness on high norms. Differentiating at each finite
parameter order leaves the same invertible transverse operator on the
left and finite expressions in input and lower output derivatives on the
right. Induction proves smooth family dependence. The finite input
dependence is uniform on \(\mathcal K\) when measured on the larger
retained set \(\mathcal K'\), by a finite cover and the fixed margins.
\hfill\(\square\)

For the flat initial label, the proof constructs \(Y(s)\) from current
differential operations and the executed inverse. Uniqueness in Theorem
A.6 then gives \(Y(s)=\mathcal C(s)\). This is an independent check of
the endpoint against the earlier constraint chart. The transport
iteration itself uses only the current operations.

\pandocstatementheading{\leavevmode\phantomsection\label{cor:online-current-finite-neumann-picard}Corollary A.15 (Controlled finite execution on the flat branch).}

Take \((s_*,y_*)=(0,0)\). Truncate the current inverse through power
\(N\) and set

\phantomsection\label{eq:online-current-finite-neumann-field} \[
\mathsf V_Z^{[N]}a
=-\sum_{j=0}^{N}(-\mathsf E_Z)^j
\mathsf G_0^{\mathrm{src}}\mathsf B_Za.
\] Starting at the zero path, perform \(J\) Picard steps along the
radial seed path \(ts\) with this field, and denote the endpoint by
\(Y_{J,N}(s)\). After one fixed reduction of the seed radius and for
\(N\) sufficiently large, all these iterates remain in the retained
product and

\phantomsection\label{eq:online-current-finite-execution-error} \[
\|Y_{J,N}(s)-Y(s)\|_{\mathcal Y_{k_0}}
\le C\|s\|_{\mathcal X_{k_0}}
\bigl(q_{\mathrm P}^{J}+\varkappa^{N+1}\bigr).
\] Their endpoint residuals obey the same bound in \(\mathcal Z_{k_0}\).
Every finite endpoint has exactly the seed \(s\). For weighted smooth
input, each finite stage is an actual weighted smooth field obtained by
finite current differential operations, finite negative-response words
and parameter integration. The higher bounds may depend on \(J,N\). The
convergence estimate is in the base norm.

\emph{Proof.} Uniformly on the fixed base product, the geometric tail
gives \[
\|(\mathsf V_Z^{[N]}-\mathsf V_Z)a\|_{\mathcal Y_{k_0}}
\le\frac{\varkappa^{N+1}}{1-\varkappa}
\|\mathsf G_0^{\mathrm{src}}\mathsf B_Za\|_{\mathcal Y_{k_0}}
\le C\varkappa^{N+1}\|a\|_{\mathcal X_{k_0}}.
\] Choose a fixed margin \(\epsilon_0>0\), reduce the seed radius once
so that \((M+\epsilon_0)\rho<r_y/4\), and choose \(N\) with the last
operator error at most \(\epsilon_0\). The approximate integral map
preserves the same path ball. Comparing an approximate iterate to the
exact path, with error \(e_J\) in the base path norm, gives \[
e_{J+1}\le q_{\mathrm P}e_J
+C\varkappa^{N+1}\|s\|_{\mathcal X_{k_0}},
\qquad e_0\le C\|s\|_{\mathcal X_{k_0}}.
\] Summation proves the error bound. The residual bound follows from the
bounded derivative of \(\mathcal F\) in the correction variable and the
zero residual of the exact endpoint.

The iteration changes only the correction coordinate over the prescribed
seed path. Its endpoint therefore has seed \(s\) independently of its
residual. Each finite inverse word on smooth residuals is actual by the
source Green theorem, including the domain of every auxiliary call. The
coefficient derivatives preserve weighted smoothness, and integration on
the finite parameter interval preserves every fixed seminorm. Induction
proves actual smoothness of every finite stage. Theorem A.14 gives
higher regularity directly for the exact endpoint. \hfill\(\square\)

\phantomsection\label{app:online-relative-extension}

\section{Appendix B. Relative maximal extension on one smooth
branch}\label{appendix-b.-relative-maximal-extension-on-one-smooth-branch}

\hyperref[ref-ChruscielCong2023]{Chruściel and Cong} extend smooth
vacuum data across a boundary, preserving the original metric and second
fundamental form (Theorem 6.4). For sufficiently small maximal vacuum
data on a ball, \hyperref[ref-Czimek2018]{Czimek} already gives a
maximal asymptotically flat extension on \(\mathbb R^3\) agreeing with
the data on that ball. \hyperref[ref-MaoOhTao2023]{Mao, Oh and Tao}
construct constraint gluings using linear solution operators with
prescribed support.

Here the correction is normalized relative to a supplied whole-space
extension of the interior fields and vanishes when that extension
already solves the constraints. We keep the interior data fixed and
allow a decaying exterior tail, without prescribing exterior fields or
ADM charges. Czimek's explicit exterior right inverses and the nonlinear
construction below give existence. Their ranges supply the estimates at
weights slightly below \(-1/2\) and the higher spatial and parameter
estimates on the same selected branch. Write
\(R(h)=\operatorname{Scal}(h)\).

\subsection{B.1. Weighted spaces and fixed exterior
ranges}\label{b.1.-weighted-spaces-and-fixed-exterior-ranges}\label{sec:online-extension-spaces}

Put \(\Omega=\mathbb R^3\setminus\overline B_1\) and write
\(\varrho=|x|\) for the exterior radial coordinate. All tensor
coefficients and derivatives in the weighted norm use the fixed
Euclidean coordinates. For an integer \(s\ge0\) set

\phantomsection\label{eq:online-extension-weighted-norm} \[
\|u\|_{H^s_\omega}^2
=\sum_{|\alpha|\le s}
\int_{\mathbb R^3}
\langle x\rangle^{-2\omega-3+2|\alpha|}
|\partial^\alpha u(x)|^2\,\mathrm d x.
\] On \(\Omega\), \(\overline H^s_\omega\) is the closure of
\(C_c^\infty(\Omega)\) in this norm. Equivalently, the zero extension
belongs to the whole-space weighted Sobolev space. The bar records the
vanishing boundary traces; the resulting field can have a noncompact
decaying tail.

The first flat operator is \[
P_1v=\operatorname{div}_\delta v,
\qquad \operatorname{tr}_\delta v=0.
\] For the second, write a metric in polar form as \[
\widehat h=a^2\,\mathrm d\varrho^2
+\widehat\gamma_{AB}
(b^A\mathrm d\varrho+\mathrm d\theta^A)
(b^B\mathrm d\varrho+\mathrm d\theta^B)
\] and define

\phantomsection\label{eq:online-extension-polar-variation} \[
\mathfrak V_{\widehat h}(\phi,\beta)
=a^2\,\mathrm d\varrho^2
+e^{2\phi}\widehat\gamma_{AB}
\bigl((b+\beta)^A\mathrm d\varrho+\mathrm d\theta^A\bigr)
\bigl((b+\beta)^B\mathrm d\varrho+\mathrm d\theta^B\bigr),
\] where \(\beta\) is tangent to the coordinate spheres. Set \[
P_2(\phi,\beta)
=D R_\delta\!\left[D\mathfrak V_\delta(0,0)(\phi,\beta)\right].
\] Fix the explicit spherical-harmonic inverses \(B_1,B_2\) of the flat
exterior constructions in \hyperref[ref-Czimek2018]{Czimek} (2018,
Lemmas 4.6 and 5.10), with the sign of the scalar residual chosen to
agree with this \(P_2\). The particular inverse formulas and their
ranges are retained. They are not replaced by arbitrary right inverses
of the same differential operators.

Choose a fixed smooth radial function \(\ell\) that vanishes near
\(\overline B_1\) and equals \(\log\varrho\) for \(\varrho\ge2\), and
put \(w_{\varepsilon_\gamma}=e^{-\varepsilon_\gamma\ell}\). Lemma B.2
supplies \(\varepsilon_*>0\) such that, after fixing

\phantomsection\label{eq:online-extension-selected-weight} \[
0<\varepsilon_\gamma<\min(\varepsilon_*,1/2),
\qquad \gamma=-\frac12-\varepsilon_\gamma,
\] the conjugated inverses

\phantomsection\label{eq:online-extension-fixed-conjugated-inverses} \[
B_{j,\gamma}
=w_{\varepsilon_\gamma} B_j(I+E_{j,\varepsilon_\gamma}B_j)^{-1}w_{\varepsilon_\gamma}^{-1},
\qquad
E_{j,\varepsilon_\gamma}=w_{\varepsilon_\gamma}^{-1}P_jw_{\varepsilon_\gamma}-P_j,
\quad j=1,2,
\] are bounded right inverses with the mappings

\phantomsection\label{eq:online-extension-inverse-domains} \[
\begin{aligned}
B_{1,\gamma}:\overline H^{s-2}_{\gamma-2}(T^*)
&\longrightarrow
\overline H^{s-1}_{\gamma-1}(\operatorname{Sym}^2_0T^*),\\
B_{2,\gamma}:\overline H^{s-2}_{\gamma-2}
&\longrightarrow
\overline H^s_\gamma\oplus
\overline H^s_\gamma(TS_\varrho),
\qquad s\ge2.
\end{aligned}
\] The same operators are used at every integer order. Their
residual-space Neumann inverses are selected at one base order; the fact
that \(E_{j,\varepsilon_\gamma}\) has one lower differential order gives
all higher estimates without reducing \(\varepsilon_\gamma\). On their
base spaces the exact range identities are

\phantomsection\label{eq:online-extension-exact-ranges} \[
\operatorname{Ran}B_{j,\gamma}
=w_{\varepsilon_\gamma}\operatorname{Ran}B_j,
\qquad j=1,2.
\] The inverses preserve real fields, zero-extension traces and proper
Euclidean rotations. The rotation statement uses the fixed radial
cutoffs and the same radial formulas throughout each harmonic
multiplicity space. Lemmas B.2 and B.5 prove these properties for the
chosen normalization. They are needed for the regularity of the selected
nonlinear branch.

\subsection{B.2. The relative-extension
theorem}\label{b.2.-the-relative-extension-theorem}\label{sec:online-extension-common-branch}

\pandocstatementheading{\leavevmode\phantomsection\label{thm:online-relative-extension-one-branch}Theorem B.1 (Relative maximal extension on one smooth branch).}

Fix the weight and inverse ranges above and an integer \(S_0\ge4\).
Increase \(S_0\) beforehand if subsequent finite-order chart estimates
require it. There is a base threshold \(\epsilon>0\), chosen with fixed
positivity and inverse margins, with the following property.

Let \((\widehat h,\widehat K)\) be smooth whole fields on
\(\mathbb R^3\) such that \[
\widehat h-\delta\in H^M_\gamma,
\qquad
\widehat K\in H^{M-1}_{\gamma-1}
\quad\text{for every integer }M\ge S_0,
\] and

\phantomsection\label{eq:online-extension-base-smallness} \[
\|\widehat h-\delta\|_{H^{S_0}_\gamma}
+\|\widehat K\|_{H^{S_0-1}_{\gamma-1}}<\epsilon.
\] Assume that \(\widehat h\) is positive, that
\(\operatorname{tr}_{\widehat h}\widehat K=0\) on the whole space, and
that on \(B_1\)

\phantomsection\label{eq:online-extension-local-constraint-input} \[
R(\widehat h)=|\widehat K|_{\widehat h}^2,
\qquad
\operatorname{div}_{\widehat h}\widehat K=0.
\] The initialized relative operation selects a whole pair \[
(h,K)=\mathfrak E(\widehat h,\widehat K)
\] with

\phantomsection\label{eq:online-extension-output-contract} \[
\begin{gathered}
(h,K)|_{B_1}=(\widehat h,\widehat K)|_{B_1},\\
R(h)=|K|_h^2,
\qquad \operatorname{div}_hK=0,
\qquad \operatorname{tr}_hK=0,\\
h-\delta\in H^M_\gamma,
\qquad K\in H^{M-1}_{\gamma-1}
\quad\text{for every integer }M\ge S_0.
\end{gathered}
\] The corrections \(h-\widehat h\) and \(K-\widehat K\) have the
corresponding zero-extension regularity. At the base order their whole
output norm is bounded by a fixed multiple of the whole input norm in
the smallness condition. At each higher finite order, the output norm is
bounded in terms of finitely many higher input norms and the fixed base
margins. No higher norm is required to be small. The operation fixes
every supplied whole pair in its domain that already satisfies the
maximal vacuum constraints everywhere.

Let the prefields form a smooth family on a finite-dimensional parameter
domain, smooth in all the indicated weighted seminorms and remaining in
the selected base neighborhood. Then the same output is smooth at every
finite weighted order. On a compact parameter set contained in the
interior of a slightly larger retained compact set, every finite number
of parameter derivatives in any specified output order is bounded by
finitely many weighted input seminorms and parameter seminorms on that
larger set, together with the fixed margins. The branch and base
neighborhood are chosen before the requested derivative orders.

The proof is completed in
\hyperref[sec:online-extension-proof-of-smooth-theorem]{§B.8}. The
intervening results specify the inverse ranges and establish the
estimates for this selected operation.

\subsection{B.3. Fixed inverses, nonlinear repairs and the base
branch}\label{b.3.-fixed-inverses-nonlinear-repairs-and-the-base-branch}\label{sec:online-extension-linear-and-base}

The sphere-tangent variable \(\beta\) is a vector field on the Euclidean
exterior, with its Cartesian tensor norm. Thus its norm is the norm of
the vector field itself, rather than of its components in an unscaled
angular coordinate basis. The fixed Euclidean metric on each sphere is
\(\gamma^\circ(\varrho)=\delta|_{TS_\varrho}=\varrho^2g_{S^2}\); in
particular its Laplacian includes the factor \(\varrho^{-2}\) relative
to the unit-sphere Laplacian. With these conventions the second flat
operator is

\phantomsection\label{eq:online-extension-scalar-flat-operator} \[
P_2(\phi,\beta)
=-4\partial_\varrho^2\phi
-\frac{12}{\varrho}\partial_\varrho\phi
-\frac{4}{\varrho^2}\phi
-2\Delta_{\gamma^\circ}\phi
+\frac{2}{\varrho^3}\partial_\varrho
  \bigl(\varrho^3\operatorname{div}_{\gamma^\circ}\beta\bigr).
\] The external linear input is the constructive part of
\hyperref[ref-Czimek2018]{Czimek} (2018, Lemmas 4.6 and 5.10 and their
proofs in §§4.3 and 5.3). Fix the transition function in those
constructions, using the same radial formula throughout each
spherical-harmonic multiplicity space. They give real linear right
inverses \[
\begin{aligned}
B_1:\overline H^{q}_{-5/2}(T^*)
&\longrightarrow
\overline H^{q+1}_{-3/2}(\operatorname{Sym}^2_0T^*),\\
B_2:\overline H^{q}_{-5/2}
&\longrightarrow
\overline H^{q+2}_{-1/2}\oplus
\overline H^{q+2}_{-1/2}(TS_\varrho),
\qquad q\ge0,
\end{aligned}
\] bounded at every integer order and defined by the same formulas at
those orders. Part (1) of each cited lemma gives the base estimate
\(q=0\). For \(q\ge1\), set the regularity index in part (2) to \(q+2\):
the tensor output has order \(q+1\) in Lemma 4.6, and both outputs have
order \(q+2\) in Lemma 5.10. These are estimates in the barred spaces
for the fields constructed in the respective proofs. In particular,
Proposition 5.17 gives the boundary regularity of the sphere-tangent
field \(\xi\) constructed by the angular Hodge equations, including its
membership in the barred space. The transition function and mode
formulas are fixed before the regularity index is chosen. With these
choices the normalization here is \[
P_jB_j=I,\qquad j=1,2.
\] The equation displayed in the cited Lemma 5.10 uses the negative of
the curvature derivative displayed in its Lemma 5.7. Accordingly, its
constructive solution is applied to the negative residual when defining
\(B_2\). This fixes a sign and leaves its range unchanged. The concrete
ranges of these constructive maps are retained; the Hilbert orthogonal
complements introduced separately in the cited implicit-function
arguments are not used.

The equation numbers in the next two paragraphs refer to
\href{https://arxiv.org/pdf/1609.08814v2}{Czimek, arXiv:1609.08814v2}.
On the momentum side, (4.39)--(4.56) construct the radial scalar
\(\delta\), the mixed field \(\epsilon\) and the tangential tracefree
field \(\widehat\eta\) of the \emph{same} tensor \(k\). Lemma 4.11
verifies the equations. Lemma 4.16, (4.93), bounds this \(k\) at order
\(w-1\) by the residual at orders \(w-2\) and zero, and puts it in the
barred space of weight \(-3/2\). The auxiliary curl tensor has order
\(w-2\) and weight \(-5/2\). Taking \(w=q+2\) gives the asserted
estimate on the explicit image of \(B_1\). The low degrees in \(\delta\)
and \(\epsilon\) remain in this inverse. Equations (4.53) and (4.56)
specify the exact component of the angular divergence of
\(\widehat\eta\) only in degrees \(l\ge2\). The half-residual terms
cancel on that projection. In the physical sphere frame the angular
derivatives acquire \(\varrho^{-1}\), giving the third equation in
\hyperref[eq:online-momentum-radial-divergence]{Lemma B.10}. There is no
tracefree tensor harmonic of degree zero or one.

On the scalar side, multiplying (5.31) by four identifies its left side
with \(-P_2\) and its right side with the cited residual \(h\).
Equations (5.33) and (5.34) select the radial and higher-mode scalar
solutions. Equations (5.35)--(5.36) fix the two cutoff moments,
including the coefficient \(-1/2\) multiplying the residual moment.
Equation (5.37) fixes the sphere-tangent vector by the radial divergence
equation, vanishing angular curl and zero boundary value. Lemma 5.13
checks the equation. Proposition 5.14 gives the two moment cancellations
(5.39)--(5.40), and the barred estimate (5.42) for \(u^{[\ge1]}\).
Proposition 5.17 gives (5.52) for \(u^{[0]}\) and (5.53) for \(\xi\),
both in the barred space of weight \(-1/2\) at order \(w\). Thus all
required zero-extension traces through order \(w-1\) for the scalar
outputs (through \(w-2\) for the momentum output) follow at that finite
order. The critical-weight estimates concern these fixed linear
solutions. Lemma B.2 passes to the stronger weight by conjugation and
identifies the image there. Lemmas B.8--B.11 use the resulting image
equations and angular control to raise the regularity of the one
nonlinear solution selected at base order.

\pandocstatementheading{\leavevmode\phantomsection\label{lem:online-extension-stronger-weight-inverse}Lemma B.2 (One stronger weight at every order).}

There is \(\varepsilon_*>0\) such that each choice
\(0<\varepsilon_\gamma<\min(\varepsilon_*,1/2)\) gives the operators
\(B_{j,\gamma}\) in the stated exterior spaces, with
\(\gamma=-1/2-\varepsilon_\gamma\). They are bounded right inverses at
every integer order, use one fixed formula, and preserve real fields,
zero-extension regularity and proper Euclidean rotations. Their base
ranges satisfy \[
\operatorname{Ran}B_{j,\gamma}
=w_{\varepsilon_\gamma}\operatorname{Ran}B_j.
\] The threshold \(\varepsilon_*\) is chosen at order zero for the
residual and is independent of every subsequent regularity order.

\emph{Proof.} Write \(m_1=1\) and \(m_2=2\). For this proof let
\(Y_j^q\) be the corresponding residual space at order \(q\) and weight
\(-5/2\), and let \(Z_j^q\) be the output space at order \(q+m_j\) and
weight \(-5/2+m_j\). Thus \(B_j:Y_j^q\to Z_j^q\). The scalar output has
two components.

Choose the fixed radial \(\ell\) to vanish on a neighborhood of
\(\overline B_1\) and equal \(\log\varrho\) for \(\varrho\ge2\). It
satisfies \(|\ell|\le C(1+\log\langle x\rangle)\), and its
positive-order derivatives satisfy
\(|\partial^\alpha\ell|\le C_\alpha\langle x\rangle^{-|\alpha|}\).
Consequently multiplication by \(w=e^{-\varepsilon_\gamma\ell}\) is an
isomorphism \(H^k_\omega\to H^k_{\omega-\varepsilon_\gamma}\) at every
finite \(k\), with inverse multiplication by \(w^{-1}\). This follows by
differentiating the product and using
\(|\partial^\alpha w|\le C_{\alpha,\varepsilon_\gamma}
\langle x\rangle^{-\varepsilon_\gamma-|\alpha|}\) and the corresponding
inverse estimate. Both multipliers preserve the zero-extension subspaces
and all algebraic tensor types.

In Cartesian components, a coefficient of differential order \(d\) in
\(P_j\) has size \(O(\varrho^{d-m_j})\), with the corresponding bound
after each derivative. For \(P_1\) this is immediate; for \(P_2\) it
follows either from its displayed polar expression or from the
degree-zero coefficients of the linearized polar metric variation.
Leibniz' rule therefore gives \[
w^{-1}P_jw=P_j+E_j,\qquad
E_j=\sum_{|\alpha|\le m_j-1}e_{j,\alpha}(x)\partial^\alpha,
\] where \(E_j=0\) near \(\overline B_1\) and, for
\(0<\varepsilon_\gamma\le1/2\), \[
|\partial^\nu e_{j,\alpha}(x)|
\le C_{\alpha,\nu}\varepsilon_\gamma
\langle x\rangle^{|\alpha|-m_j-|\nu|}.
\] In particular, \[
\|E_jB_j\|_{Y_j^0\to Y_j^0}
\le C\varepsilon_\gamma.
\] Choose \(\varepsilon_*\) so that the two norms are less than \(1/2\).
The inverse \(Q_j=(I+E_jB_j)^{-1}\) is then defined on \(Y_j^0\) by its
convergent Neumann series. Define \[
\widetilde B_j=B_jQ_j,\qquad
B_{j,\gamma}=w\widetilde B_jw^{-1}.
\] The identity \((P_j+E_j)\widetilde B_j=I\) proves the required
right-inverse equation after conjugation.

It remains to prove higher bounds for this same operator. The
coefficient estimate above has a useful precise consequence: \[
E_j:Z_j^q\longrightarrow Y_j^{q+1}
\] is bounded for each \(q\ge0\). Indeed, a derivative of order
\(|\alpha|\le m_j-1\) leaves at least \(q+1\) derivatives of the output
available. Its weight becomes \(-5/2+m_j-|\alpha|\); the coefficient has
the compensating weight \(|\alpha|-m_j\). Differentiating the
coefficient preserves this balance. Moreover, \(E_j\) vanishes near the
inner sphere, so the extra residual derivative does not require an extra
trace of the input there.

Let \(u=\widetilde B_j f\) be the solution already constructed in
\(Z_j^0\). It satisfies \[
u=B_jf-B_jE_ju.
\] If \(f\in Y_j^{q+1}\) and \(u\in Z_j^q\), both terms on the right
belong to \(Z_j^{q+1}\), with \[
\|u\|_{Z_j^{q+1}}
\le C_{q+1}\bigl(
\|f\|_{Y_j^{q+1}}+\|u\|_{Z_j^q}\bigr).
\] Induction proves all higher estimates. Only their constants depend on
the order; the Neumann inverse and the chosen \(\varepsilon_\gamma\)
remain fixed. Multiplication by \(w\) gives precisely the claimed
shifted weights.

The base range identity follows directly from the surjectivity of
\(Q_j:Y_j^0\to Y_j^0\): \[
\operatorname{Ran}\widetilde B_j=\operatorname{Ran}B_j,
\qquad
\operatorname{Ran}B_{j,\gamma}
=w\operatorname{Ran}B_j.
\] Every step preserves real fields and zero-extension regularity.
Proper rotations commute with \(P_j\), the fixed radial multipliers, and
the constructive flat inverses: the radial formulas depend on harmonic
degree, while the angular projections and oriented Hodge operators
commute with those rotations. They therefore commute with \(E_j\), its
Neumann inverse and \(B_{j,\gamma}\). Finally, at each fixed order
\(B_{j,\gamma}P_j\) is a bounded projection onto
\(\operatorname{Ran}B_{j,\gamma}\), because \(P_jB_{j,\gamma}=I\). The
selected range is therefore closed in the indicated output space; limits
of correction iterates remain in that range. \hfill\(\square\)

\textbf{Exterior repairs at one fixed Sobolev order.}

Fix a finite \(S\ge4\) and the weight chosen in Lemma B.2. Write \[
\begin{aligned}
\mathcal X_S&=H^S_\gamma(\operatorname{Sym}^2T^*),\\
\mathcal V_S&=\overline H^{S-1}_{\gamma-1}
(\operatorname{Sym}^2_0T^*),\\
\mathcal Z_S&=\overline H^S_\gamma
\oplus\overline H^S_\gamma(TS_\varrho),\\
\mathcal Y_S^1&=\overline H^{S-2}_{\gamma-2}(T^*),\qquad
\mathcal Y_S^0=\overline H^{S-2}_{\gamma-2}.
\end{aligned}
\] Metrics are written as \(\delta+\mathcal X_S\). The bar is omitted
when a tensor is defined on the whole space; norms of zero-extended
tensors agree with the corresponding exterior norms.

\pandocstatementheading{\leavevmode\phantomsection\label{lem:online-extension-normalized-repairs}Lemma B.3 (Normalized exterior repairs).}

On a sufficiently small neighborhood at this fixed order there are
smooth operations \[
\mathfrak D(h,\rho)=\operatorname{TF}_h v,\qquad
v\in\operatorname{Ran}B_{1,\gamma},
\] and \[
\mathfrak G(\widehat h,s)
=\mathfrak V_{\widehat h}(z),\qquad
z\in\operatorname{Ran}B_{2,\gamma},
\] such that \[
\begin{gathered}
\operatorname{div}_h\mathfrak D(h,\rho)=\rho,\qquad
\operatorname{tr}_h\mathfrak D(h,\rho)=0,\qquad
\mathfrak D(h,0)=0,\\
R(\mathfrak G(\widehat h,s))=R(\widehat h)+s,\qquad
\mathfrak G(\widehat h,0)=\widehat h.
\end{gathered}
\] Both corrections have zero-extension regularity and preserve the
supplied fields on \(B_1\). On a smaller retained neighborhood they
satisfy \[
\begin{aligned}
\|\mathfrak D(h,\rho)-\mathfrak D(h',\rho')\|_{H^{S-1}_{\gamma-1}}
&\le C_S\bigl(
\|\rho-\rho'\|_{\mathcal Y_S^1}
+\|h-h'\|_{\mathcal X_S}\|\rho'\|_{\mathcal Y_S^1}\bigr),\\
\|\mathfrak G(\widehat h,s)-\mathfrak G(\widehat h,s')\|_{\mathcal X_S}
&\le C_S\|s-s'\|_{\mathcal Y_S^0},\\
\|\mathfrak G(\widehat h,s)-\widehat h\|_{\mathcal X_S}
&\le C_S\|s\|_{\mathcal Y_S^0}.
\end{aligned}
\] The neighborhoods and constants here may depend on \(S\). These
operations agree at different finite orders whenever both are defined on
the same input.

\emph{Proof.} We first record the estimates that control the two
nonlinear equations. On a dyadic annulus \(\varrho\asymp 2^\nu\),
rescale to a fixed annulus. The weighted norm is equivalent to the
square sum of \(2^{-\nu\omega}\) times the ordinary Sobolev norms of the
rescaled Cartesian coefficients. The usual three-dimensional product
estimate on the fixed annulus, followed by this square summation, gives
\[
\|uv\|_{H^{S-2}_{\omega_1+\omega_2}}
\le C_S\|u\|_{H^S_{\omega_1}}
\|v\|_{H^{S-2}_{\omega_2}}.
\] The same argument applies when both factors have \(S-1\) derivatives
and the output has at most \(S-1\) derivatives. Here \(S-2>3/2\), so
each lowest undifferentiated Sobolev space needed in these products is
an algebra. Smooth composition with matrix inversion near the identity,
with the exponential near zero, and with the tracefree projection
follows by differentiating these finite-dimensional functions and
applying the same product estimate. Differentiation lowers the weight by
one. A factor of negative weight can also be included without weakening
the desired decay.

For clarity, these estimates also apply to the polar variation. The
Euclidean radial vector, radial one-form and tangential projection have
bounded coefficients with derivatives of order \(k\) bounded by
\(C_k\varrho^{-k}\) on \(\Omega\). The polar variables of \(h\) are
obtained algebraically from its radial and tangential blocks, with the
tangential block inverted on the tangent planes. Its inverse is
uniformly bounded after rescaling each sphere. In Cartesian tensor
components the map \((h,z)\mapsto\mathfrak V_h(z)-h\) is therefore
smooth \((\delta+\mathcal X_S)\times\mathcal Z_S\to\mathcal X_S\) near
\((\delta,0)\), vanishes when \(z=0\), and has the corresponding bounded
derivatives. This avoids assigning Cartesian weights to bare angular
coordinate components.

The Cartesian scalar curvature formula consists of terms
\(h^{-1}h^{-1}\partial^2h\) and
\(h^{-1}h^{-1}h^{-1}(\partial h)(\partial h)\) with fixed index
contractions. The preceding estimates make \(R\) a smooth map from the
metric neighborhood to \(H^{S-2}_{\gamma-2}\). They also make \[
L_hv:=\operatorname{div}_h\operatorname{TF}_h v
\] a bounded map \(\mathcal V_S\to\mathcal Y_S^1\), smooth in \(h\).
Since \(L_\delta=P_1\) on Euclidean tracefree tensors, one has \[
\begin{aligned}
\|(L_h-P_1)v\|_{\mathcal Y_S^1}
&\le C_S\|h-\delta\|_{\mathcal X_S}\|v\|_{\mathcal V_S},\\
\|(L_h-L_{h'})v\|_{\mathcal Y_S^1}
&\le C_S\|h-h'\|_{\mathcal X_S}\|v\|_{\mathcal V_S}.
\end{aligned}
\] For the scalar equation put \[
\mathcal S_h(z)=R(\mathfrak V_h(z))-R(h),\qquad
\mathcal N_h(z)=\mathcal S_h(z)-P_2z.
\] The smoothness just proved, \(\mathcal S_h(0)=0\), and
\(D_z\mathcal S_\delta(0)=P_2\) give, on a fixed small ball, \[
\|\mathcal N_h(z)-\mathcal N_h(z')\|_{\mathcal Y_S^0}
\le C_S\bigl(
\|h-\delta\|_{\mathcal X_S}
+\|z\|_{\mathcal Z_S}+\|z'\|_{\mathcal Z_S}\bigr)
\|z-z'\|_{\mathcal Z_S}.
\] Indeed, integrate \(D_z\mathcal S_h-P_2\) along the segment from
\(z'\) to \(z\) and use the bounded mixed and second derivatives. All
these residuals have the indicated barred regularity. A zero-extended
correction is a whole Sobolev tensor vanishing on \(B_1\); the local
differential expressions formed from it are whole Sobolev tensors
vanishing there as well.

For the momentum equation, restrict the unknown by writing
\(v=B_{1,\gamma}f\). Set \[
T_h=(L_h-P_1)B_{1,\gamma}:
\mathcal Y_S^1\longrightarrow\mathcal Y_S^1.
\] After shrinking the metric neighborhood, \(\|T_h\|<1/2\). The
required solution is explicitly \[
v=B_{1,\gamma}(I+T_h)^{-1}\rho,\qquad
\mathfrak D(h,\rho)=\operatorname{TF}_h v.
\] This proves its equations, range and zero normalization. The
resolvent identity \[
(I+T_h)^{-1}-(I+T_{h'})^{-1}
=(I+T_h)^{-1}(T_{h'}-T_h)(I+T_{h'})^{-1}
\] and the estimates for \(L_h\) and \(\operatorname{TF}_h\) prove the
first difference bound. They also prove smooth dependence on \(h\) and
\(\rho\).

For the scalar equation use the iteration \[
z_{n+1}
=B_{2,\gamma}\bigl(s-\mathcal N_{\widehat h}(z_n)\bigr),
\qquad z_0=0.
\] Choose the metric and residual neighborhoods so that, on
\(\|z\|_{\mathcal Z_S}\le2\|B_{2,\gamma}\|\|s\|_{\mathcal Y_S^0}\), the
displayed remainder estimate gives contraction constant less than
\(1/2\). Since \(\mathcal N_{\widehat h}(0)=0\), the same estimate
proves that this ball is invariant. The limit has
\(\|z\|\le2\|B_{2,\gamma}\|\|s\|\), remains in the fixed inverse range,
and satisfies \(\mathcal S_{\widehat h}(z)=s\) after applying \(P_2\).
For \(s=0\) every iterate is zero. Difference of two fixed-point
equations, followed by absorption of the contraction term, gives the
scalar residual difference estimate; the smooth metric variation then
gives both stated estimates for \(\mathfrak G\).

Smooth dependence follows in the residual variable as well. The equation
\(\mathcal S_{\widehat h}(B_{2,\gamma}f)=s\) has derivative \(I\) in
\(f\) at the flat zero input. On the selected small ball its derivative
remains invertible by a Neumann series, so the Banach implicit-function
theorem supplies all finite parameter derivatives. The initialized
solution above is that branch by uniqueness. The zero extensions of
\(z\) make \(\mathfrak G-\widehat h\) a whole tensor in \(H^S_\gamma\)
vanishing on \(B_1\).

For two finite orders the momentum solution is also obtained by the
iteration \[
v_{n+1}
=B_{1,\gamma}\bigl(\rho-(L_h-P_1)v_n\bigr),\qquad v_0=0.
\] Both this iteration and the scalar iteration use the same operators
and initial values, independently of the norm in which convergence is
proved. If their input lies in both selected neighborhoods, the
higher-order limit and the lower-order limit coincide by lower-order
uniqueness. This proves compatibility. \hfill\(\square\)

\pandocstatementheading{\leavevmode\phantomsection\label{prop:online-extension-initialized-base-branch}Proposition B.4 (The initialized relative extension at the base order).}

For the prefields of Theorem B.1, choose the threshold at the one fixed
order \(S=S_0\). The repairs of Lemma B.3 determine an initialized
operation \(\mathfrak E\) whose output is a whole maximal vacuum datum
in \(H^{S_0}_\gamma\oplus H^{S_0-1}_{\gamma-1}\), with the metric
measured relative to \(\delta\). It agrees with the prefields on
\(B_1\), its corrections have the corresponding zero-extension
regularity, and its base output norm is bounded by a fixed multiple of
the base input norm. It fixes already exact whole data. On every
supplied smooth finite-dimensional family in its domain it is smooth in
the base spaces, with finite parameter and difference estimates. These
assertions require only the fixed base threshold.

\emph{Proof.} The construction is carried out at \(S=S_0\). All later
regularity arguments concern this initialized output.

Let \(d=(\widehat h,\widehat K)\) satisfy the hypotheses of Theorem B.1,
and set

\phantomsection\label{eq:online-extension-fixed-premomentum} \[
\widehat a=\widehat K-
\frac{\operatorname{tr}_\delta\widehat K}
{\operatorname{tr}_\delta\widehat h}\widehat h.
\] The denominator is bounded away from zero by the base metric margin.
One has \(\operatorname{tr}_\delta \widehat a=0\) and
\(\operatorname{TF}_{\widehat h}\widehat a=\widehat K\): the second
identity follows by taking the \(\widehat h\) trace and using
\(\operatorname{tr}_{\widehat h}\widehat K=0\). In particular, \[
\|\widehat a\|_{H^{S-1}_{\gamma-1}}
\le C_S\|\widehat K\|_{H^{S-1}_{\gamma-1}}.
\] Consider candidate whole fields \((h,K)\) agreeing with \(d\) on
\(B_1\), and define

\phantomsection\label{eq:online-extension-iteration-blocks} \[
\begin{aligned}
A_d(h)&=\operatorname{TF}_h \widehat a,\\
M_d(h)&=A_d(h)+
\mathfrak D\bigl(h,-\operatorname{div}_hA_d(h)\bigr),\\
H_d(h,K)&=\mathfrak G\bigl(
\widehat h,|K|_h^2-R(\widehat h)\bigr).
\end{aligned}
\] The first momentum residual vanishes on \(B_1\) because
\(A_d(\widehat h)=\widehat K\) there. The scalar residual vanishes there
by the local Hamiltonian constraint. Both are whole tensors in
\(H^{S-2}_{\gamma-2}\): for the quadratic term the sharper product
weight is \(2\gamma-2<\gamma-2\). Their restrictions therefore belong to
\(\overline H^{S-2}_{\gamma-2}(\Omega)\). Whole-space Sobolev
regularity, together with vanishing throughout the ball, gives every
trace required for zero extension. Boundary values alone would control
only the zeroth trace.

Lemma B.3 gives \[
\operatorname{div}_hM_d(h)=0,\qquad
\operatorname{tr}_hM_d(h)=0,
\] and both blocks preserve the prescribed interior. The same estimates
imply, for candidate fields in a fixed small neighborhood, \[
\begin{aligned}
\|M_d(h)-M_d(h')\|_{H^{S-1}_{\gamma-1}}
&\le C_S\|\widehat a\|_{H^{S-1}_{\gamma-1}}
\|h-h'\|_{H^S_\gamma},\\
\|H_d(h,K)-H_d(h',K')\|_{H^S_\gamma}
&\le C_S\bigl(\|K\|_{H^{S-1}_{\gamma-1}}+
\|K'\|_{H^{S-1}_{\gamma-1}}\bigr)
\|K-K'\|_{H^{S-1}_{\gamma-1}}\\
&\quad+C_S\|K'\|_{H^{S-1}_{\gamma-1}}^2
\|h-h'\|_{H^S_\gamma}.
\end{aligned}
\] For the first estimate, \(A_d(h)\) and each of its metric variations
are linear in \(\widehat a\). Thus
\(\|\operatorname{div}_hA_d(h)\|\le C_S\|\widehat a\|\) and the
difference of these residuals is at most \(C_S\|\widehat a\|\|h-h'\|\).
Apply the momentum difference estimate in Lemma B.3, retaining its
residual factor. For the second, expand the difference of the two
quadratic norms, use smooth dependence of \(h^{-1}\), and apply the
scalar residual difference estimate with the same fixed premetric
\(\widehat h\).

Define \[
\mathcal T_d(h,K)
=\bigl(H_d(h,K),M_d(H_d(h,K))\bigr).
\] Here is an explicit way to choose its invariant ball. Put
\(e=\|\widehat h-\delta\|_{H^S_\gamma}
+\|\widehat K\|_{H^{S-1}_{\gamma-1}}\). The scalar repair estimate and
the scalar-curvature product estimate give \[
\|H_d(h,K)-\delta\|_{H^S_\gamma}
\le C_S e+C_S\|K\|_{H^{S-1}_{\gamma-1}}^2,
\] whereas \(\|M_d(h)\|_{H^{S-1}_{\gamma-1}}\le C_S e\) on the fixed
metric neighborhood. Choose a fixed \(C_0\) larger than twice the sum of
these linear constants. After choosing the base threshold for \(e\)
sufficiently small, \(\mathcal T_d\) maps the closed set \[
\left\{(h,K):
\begin{gathered}
(h,K)|_{B_1}=d|_{B_1},\\
\|h-\delta\|_{H^S_\gamma}
+\|K\|_{H^{S-1}_{\gamma-1}}\le C_0e
\end{gathered}
\right\}
\] to itself, and all metrics in that set have the chosen positivity
margin. The two difference estimates make its Lipschitz constant less
than \(1/2\) after a further base restriction. The restriction depends
on \(S\), the fixed inverses and the margins, with no reference to any
higher order.

Starting at \((h_0,K_0)=(\widehat h,\widehat K)\), the iteration

\phantomsection\label{eq:online-extension-base-iteration} \[
(h_{n+1},K_{n+1})
=\bigl(H_d(h_n,K_n),M_d(H_d(h_n,K_n))\bigr)
\] converges in the base space to a unique fixed point in this ball. At
the fixed point the defining equations of the repairs give all three
maximal constraints. The invariant-ball estimate gives the claimed base
norm bound. Every iterate agrees with \(d\) on \(B_1\), so the limit
does as well; its differences from \(d\) are whole Sobolev fields
vanishing there and hence have the claimed barred regularity.

If \(d\) already satisfies all maximal constraints on the whole space,
then at \((\widehat h,\widehat K)\) both residuals are zero. The
identities \(\mathfrak D(h,0)=0\) and \(\mathfrak G(h,0)=h\) show that
the first iterate equals the initial value. Every later iterate and the
selected output are therefore equal to \(d\).

For base parameter smoothness, use the fixed Banach space of corrections
\[
\mathcal C_S=
\overline H^S_\gamma(\Omega;\operatorname{Sym}^2T^*)
\oplus
\overline H^{S-1}_{\gamma-1}(\Omega;\operatorname{Sym}^2T^*).
\] If \(d=d(p)\) is a supplied smooth family satisfying the local
constraints, write the fixed-point equation in the form \[
\zeta=\mathcal T_{d(p)}(d(p)+\zeta)-d(p),
\qquad \zeta\in\mathcal C_S.
\] The right side is smooth in \(p\) and \(\zeta\). It takes values in
the same fixed correction space because the residuals vanish identically
on \(B_1\) for every \(p\) and every zero-extended correction; their
parameter derivatives do so as well. Its derivative in \(\zeta\) has
norm less than \(1/2\). Hence the implicit-function theorem gives a
smooth base solution. The inverse derivative is bounded by \(2\).
Differentiating the equation \(q\) times gives the same inverse applied
to finite sums of derivatives of the block maps, input derivatives up to
order \(q\), and lower derivatives of \(\zeta\). Induction gives all
specified finite parameter bounds at the base order on a retained
compact parameter set.

Differences of two locally constrained prefields can be estimated
without drawing a line segment inside the constrained input set. Write
their fixed-point corrections as \(\zeta_d,\zeta_{d'}\). The contraction
estimate gives \[
\begin{aligned}
\|\zeta_d-\zeta_{d'}\|_{\mathcal C_S}
&\le\frac12\|\zeta_d-\zeta_{d'}\|_{\mathcal C_S}\\
&\quad+
\|\mathcal T_d(d+\zeta_{d'})-d
-\mathcal T_{d'}(d'+\zeta_{d'})+d'\|_{\mathcal C_S}.
\end{aligned}
\] Both block evaluations are defined after restricting the prefields to
a smaller retained neighborhood inside the common small correction
neighborhood: the same correction vanishes on \(B_1\), leaving each
prefield's own exact interior unchanged. The estimates proved above
bound the second term by the finite input difference. Absorbing the
first term proves the base difference estimate. The construction uses
the same initialized block maps at every finite order where the
corresponding contraction is available, so those constructions agree by
base uniqueness. The higher regularity on one base domain requires the
range arguments that follow. \hfill\(\square\)

\subsection{B.4. Angular regularity and radial
profiles}\label{b.4.-angular-regularity-and-radial-profiles}\label{sec:online-extension-angular-regularity}

Fix the exterior inverses, the weight
\(\gamma=-1/2-\varepsilon_\gamma\), and the initialized nonlinear
branches described above. Take \(S_0\ge4\) and choose the base domains
using rotation-invariant weighted norms. One can use the Euclidean norms
of the full derivative tensors in the weighted integrals; these are
equivalent to the coordinate norms. The inverse formulas use fixed
radial cutoffs and apply the same radial operation to every member of a
spherical-harmonic multiplicity space.

For \(A\in\operatorname{SO}(3)\), the notation \(A^*\) includes the
pullback of the actual tensor indices. Let \(Z_1,Z_2,Z_3\) be the
standard infinitesimal rotations. For a finite ordered word
\(I=(i_1,\ldots,i_q)\), put
\(\mathcal L_Z^I=\mathcal L_{Z_{i_1}}\cdots\mathcal L_{Z_{i_q}}\).

\pandocstatementheading{\leavevmode\phantomsection\label{lem:online-extension-rotation-equivariance}Lemma B.5 (Rotation equivariance).}

The fixed inverses and initialized nonlinear operations satisfy

\phantomsection\label{eq:online-extension-rotation-identities} \[
\begin{aligned}
B_{j,\gamma}A^*f&=A^*B_{j,\gamma}f,\qquad j=1,2,\\
\mathfrak D(A^*h,A^*\rho)&=A^*\mathfrak D(h,\rho),\\
\mathfrak G(A^*\widehat h,A^*s)&=A^*\mathfrak G(\widehat h,s),\\
\mathfrak E(A^*\widehat h,A^*\widehat K)
&=A^*\mathfrak E(\widehat h,\widehat K).
\end{aligned}
\] These identities hold on the fixed base domains, before any
improvement of spatial regularity.

\emph{Proof.} The flat inverses are composed of radial differential
equations, round-sphere differential operators, spherical Hodge inverses
on specified harmonic subspaces, and scalar multipliers depending on the
harmonic degree. Proper rotations preserve the round metric, its
orientation, and the radial foliation. They therefore commute with the
gradient, divergence, curl, and the uniquely normalized Hodge inverses.
The excluded low modes are entire harmonic subspaces, so their
projections also commute with rotations. On each remaining multiplicity
space the radial equations and their boundary conditions are identical.
A change of orthonormal harmonic basis consequently leaves the
constructed operator unchanged. This proves the assertion for \(B_1\)
and \(B_2\).

The multiplier \(w_{\varepsilon_\gamma}\) is radial. Both \(P_j\) and
\(E_{j,\varepsilon_\gamma}=w_{\varepsilon_\gamma}^{-1}P_jw_{\varepsilon_\gamma}-P_j\)
commute with rotations. Every partial sum of the base Neumann inverse of
\(I+E_{j,\varepsilon_\gamma}B_j\) commutes with rotations, and the
identity passes to its operator-norm limit. The formula for
\(B_{j,\gamma}\) gives its asserted equivariance. In particular its
retained range is invariant.

Tracefree projection and divergence commute with metric pullback. The
polar metric variation does also, because a rotation preserves the
radial coordinate and its tangent spheres. Rotating either initialized
block solution therefore gives a solution of the rotated block equation
in the same retained range and small branch. The uniqueness supplied by
the base implicit equation proves the two block identities. The
premomentum \[
\widehat a=\widehat K-
\frac{\operatorname{tr}_\delta\widehat K}
{\operatorname{tr}_\delta\widehat h}\widehat h
\] is equivariant as well. Thus the initial pair, both iteration blocks,
and the invariant contraction ball for \(\mathfrak E\) are preserved by
rotations. Every iterate obeys the last identity, which passes to the
base-space limit. \hfill\(\square\)

\pandocstatementheading{\leavevmode\phantomsection\label{thm:online-extension-all-angular-derivatives}Theorem B.6 (Angular derivatives of the selected output).}

Let the smooth prefields \((\widehat h,\widehat K)\) have every finite
weighted norm and satisfy the local constraints and the fixed
base-smallness condition. For the single initialized output
\((h,K)=\mathfrak E(\widehat h,\widehat K)\) and every finite angular
word \(I\),

\phantomsection\label{eq:online-extension-angular-base-estimates} \[
\mathcal L_Z^I(h-\delta)\in H^{S_0}_\gamma,
\qquad
\mathcal L_Z^I K\in H^{S_0-1}_{\gamma-1}.
\] The scalar correction \(z=(\phi,\beta)\) and the Euclidean tracefree
momentum correction \(v\) have the corresponding angular regularity in
their fixed base spaces. Each specified norm is bounded using finitely
many higher weighted input norms and the fixed base margins. The
conclusions are uniform on supplied finite-dimensional smooth families
with these uniform bounds. No higher input norm is required to be small.

\emph{Proof.} For every tensor field \(u\) and finite word \(I\), the
coefficients of \(Z_i\) being linear and their index actions constant
give

\phantomsection\label{eq:online-extension-angular-orbit-bound} \[
\|\mathcal L_Z^Iu\|_{H^s_\omega}
\le C_{s,I,\omega}\|u\|_{H^{s+|I|}_\omega}.
\] Indeed a factor of \(x\) compensates the weight change of one spatial
derivative in the defining norm. Applying the product rule proves the
estimate successively for each letter of \(I\). Difference quotients,
first on smooth compactly supported fields and then by density in the
required weighted spaces, show that \(A\mapsto A^*u\) is smooth into a
fixed \(H^s_\omega\) when all higher weighted norms of \(u\) are finite.
Its derivatives are the corresponding tensor Lie derivatives. The same
argument gives locally uniform estimates for supplied smooth families.

The rotated prefields still satisfy the local constraints, and their
base norms stay in the same rotation-invariant domain. Apply the smooth
base parameter dependence of the initialized construction to this
rotation orbit and, when present, to the supplied passive parameters.
Lemma B.5 identifies the output as the orbit of the original \((h,K)\).
Its derivatives in the base output spaces, identified distributionally,
are precisely the angular Lie derivatives in the statement. Products of
one-parameter rotations give every ordered word.

The parameter estimates of the base implicit equations supply finite
bounds for these derivatives. More explicitly, at each order their
highest variation is acted on by the same uniformly invertible base
derivative. The other terms contain lower variations and finitely many
derivatives of the supplied orbit. The fixed-order multiplication and
composition estimates and induction on the variation order therefore
give bounds in finitely many of the norms in the right-hand side of the
orbit estimate. The branch and its base threshold are unchanged.

The scalar block has a uniquely selected correction pair in
\(\operatorname{Ran}B_{2,\gamma}\), and that correction pair itself
satisfies the same equivariance identity. Its arguments in the outer
fixed point are \(\widehat h\) and \(|K|_h^2-R(\widehat h)\), whose
angular derivatives have the required base regularity by the result just
proved and the fixed-order product estimates. The identical argument
applies to the selected \(v\in\operatorname{Ran}B_{1,\gamma}\) in the
momentum block. These are the internal variables of the original output,
so this argument does not introduce another spatial branch.
\hfill\(\square\)

The rotation fields span the tangent directions of each sphere. The
additional normal regularity comes from the ranges of the chosen
inverses. Their boundary profiles have a useful property that we record
first.

\pandocstatementheading{\leavevmode\phantomsection\label{lem:online-extension-radial-profile-regularity}Lemma B.7 (Radial profiles with rapidly decreasing angular
coefficients).}

Let \(Y_{l\mu}\), \(l\ge l_0\ge1\), be one of the fixed scalar, vector,
or tensor spherical-harmonic families on the unit sphere, radially
identified with the spheres \(\varrho\ge1\). Suppose

\phantomsection\label{eq:online-extension-profile-rapid-coefficients} \[
\sum_{l,\mu}(1+l)^{2N}|a_{l\mu}|^2<\infty
\quad\text{for every finite integer }N\ge0.
\] Let smooth radial functions \(p_l\) have support in
\([1+c/l,1+C/l]\), for fixed \(0<c<C\), and suppose that for every
finite \(j\)

\phantomsection\label{eq:online-extension-profile-radial-bounds} \[
\sup_\varrho|\partial_\varrho^jp_l(\varrho)|
\le C_j(1+l)^{n_j}
\] for some finite \(n_j\). Then \[
u(\varrho,\theta)=
\sum_{l,\mu}a_{l\mu}p_l(\varrho)Y_{l\mu}(\theta)
\] has a smooth compactly supported zero extension to \(\mathbb R^3\).
Its restriction to the exterior belongs to every
\(\overline H^m_\omega\). Each specified finite norm is bounded by one
sufficiently high finite sum in the coefficient condition. The assertion
holds for finite sums of such families and for their fixed radial
primitives when the corresponding total moment vanishes. For the latter
assertion the fixed radial factors and their derivatives must have
polynomial bounds in \(l\) on these collars, as do the factors occurring
in the exterior inverses.

\emph{Proof.} All supports lie in the bounded annulus
\(1\le\varrho\le1+C\), where Cartesian Sobolev norms are equivalent to
finite sums of radial and spherical Sobolev norms. Round-sphere elliptic
estimates and orthogonality give, for a finite harmonic sum,

\phantomsection\label{eq:online-extension-profile-sobolev-estimate} \[
\|u\|_{H^m}^2
\le C_m\sum_{j=0}^m\sum_{l,\mu}
(1+l)^{2(m-j)}|a_{l\mu}|^2
\|\partial_\varrho^jp_l\|_{L^2(1,1+C)}^2.
\] The same estimate holds in each of the indicated tensor bundles;
their fixed harmonic eigenvalues are comparable with \((1+l)^2\). Fixed
radial identifications contribute only bounded smooth factors on this
annulus. Choose an integer \(N_m\ge\max_{0\le j\le m}(m-j+n_j)\). The
radial bounds control the displayed expression by
\(C_m\sum(1+l)^{2N_m}|a_{l\mu}|^2\). Applying the estimate to tails
proves convergence at every prescribed Sobolev order.

Every finite harmonic sum vanishes in a neighborhood of the inner
sphere, so its zero extension is smooth. The just-proved convergence of
these zero extensions at every order gives a smooth limit on the whole
space. It vanishes on \(B_1\) and outside the fixed outer ball. This
also proves all zero-extension trace conditions. Weighted norms are
equivalent to unweighted ones on the common support.

For a radial primitive, vanishing of its total moment makes the
primitive zero on both sides of the supporting collar. Its zeroth-order
estimate follows by integrating over an interval of length at most
\(C/l\). Each positive derivative is estimated by the product rule from
the original profile and its fixed radial factors. All such estimates
remain polynomial in \(l\), with order-dependent constants. The
preceding argument then applies to the primitive. At no point is a
constant required to be uniform in the derivative order.
\hfill\(\square\)

\subsection{B.5. The scalar range and its elliptic
equation}\label{b.5.-the-scalar-range-and-its-elliptic-equation}\label{sec:online-extension-scalar-range}

The scalar right inverse has a sphere-tangent variable whose radial
dependence is fixed by moments. This range condition first makes that
variable smooth; a scalar elliptic estimate then improves the other
variable. The premetric in this argument is the supplied smooth
\(\widehat h\).

\pandocstatementheading{\leavevmode\phantomsection\label{lem:online-extension-scalar-tangent-smooth}Lemma B.8 (Smoothness of the sphere-tangent scalar correction).}

Let \(z=(\phi,\beta)\in\operatorname{Ran}B_{2,\gamma}\) have every
angular Lie derivative in its base \(\overline H^{S_0}_\gamma\) space.
Then the zero extension of \(\beta\) is smooth, with

\phantomsection\label{eq:online-extension-scalar-tangent-support} \[
\operatorname{supp}\beta\subset\{x:1\le|x|\le2\}.
\] For every finite \(m\) its spatial norm is bounded by finitely many
angular base norms of \(z\). These estimates are uniform on families
with the corresponding uniform bounds. The support assertion does not
require a common positive distance between the support and the inner
sphere.

\emph{Proof.} The exact range identity gives
\(z=w_{\varepsilon_\gamma}(u,\xi)\) with
\((u,\xi)\in\operatorname{Ran}B_2\). Use the curvature derivative
\hyperref[eq:online-extension-scalar-flat-operator]{defined above}. Use
\(f=-P_2(u,\xi)\) in the formulas of \hyperref[ref-Czimek2018]{Czimek}
(2018, (5.31) and (5.35)--(5.37)). This is the sign of the residual in
the displayed construction there. The normalized right inverse in this
paper satisfies \(P_2B_2=I\); changing the sign of the supplied residual
changes neither its range nor the following regularity argument. The
flat differential operator and radial multiplier show that
\(f\in\overline H^{S_0-2}_{-5/2}\) and that every angular derivative of
\(f\) belongs to that space, with bounds from the angular base norms of
\(z\).

Use harmonics
\(Y_{l\mu}(\varrho,\theta)=\varrho^{-1}Y_{l\mu}(1,\theta)\), orthonormal
on \(S_\varrho\), and let \(f_{l\mu}(\varrho)\) be the corresponding
coefficients. Put \[
\lambda_l=\sqrt{\frac{l(l+1)}2},
\qquad \chi_l(\varrho)=\chi(l(\varrho-1)),
\qquad l\ge1,
\] where \(0\le\chi\le1\), \(\chi=0\) on \((-\infty,1/10]\), and
\(\chi=1\) on \([1,\infty)\). The retained profiles are

\phantomsection\label{eq:online-extension-scalar-moment-profiles} \[
\begin{aligned}
c_{l\mu}
&=-\frac12\int_1^\infty
\varrho^{1-\lambda_l}f_{l\mu}(\varrho)\,\mathrm d\varrho,\\
A_l&=\int_1^\infty
\varrho^{\lambda_l+1}\chi_l'(\varrho)\,\mathrm d\varrho,
\qquad
D_l=\int_1^\infty
\varrho^{\lambda_l+1}\chi_l''(\varrho)\,\mathrm d\varrho,\\
\widetilde c_{l\mu}&=c_{l\mu}\frac{A_l}{D_l},\\
\zeta_{l\mu}(\varrho)
&=c_{l\mu}\varrho^{\lambda_l-1}\chi_l'(\varrho)
-\widetilde c_{l\mu}\varrho^{\lambda_l-1}\chi_l''(\varrho),\\
\zeta(\varrho,\theta)
&=\sum_{l\ge1,\mu}
\zeta_{l\mu}(\varrho)Y_{l\mu}(\varrho,\theta).
\end{aligned}
\] These formulas determine the tangent variable by

\phantomsection\label{eq:online-extension-scalar-tangent-hodge-equation}
\[
\operatorname{div}_{S_\varrho}\xi
=\varrho^{-3}\int_1^\varrho s^3\zeta(s,\theta)\,\mathrm ds,
\qquad
\operatorname{curl}_{S_\varrho}\xi=0.
\] We verify all properties of these expressions needed here. Since
\(\lambda_l\ge1\), Cauchy--Schwarz gives

\phantomsection\label{eq:online-extension-scalar-moment-bound} \[
|c_{l\mu}|\le
\frac{1}{2\sqrt{2\lambda_l-1}}
\|\varrho f_{l\mu}\|_{L^2(1,\infty)}.
\] The \(H^0_{-5/2}\) norm controls the square sum of the norms on the
right: orthonormality on \(S_\varrho\) leaves the radial measure
\(\mathrm d\varrho\), and the weight is comparable with \(\varrho^2\)
for \(\varrho\ge1\). The rotation Casimir multiplies the \(l\)th scalar
harmonic by \(l(l+1)\). Applying the same bound after finitely many
rotations therefore yields

\phantomsection\label{eq:online-extension-scalar-moments-rapid} \[
\sum_{l,\mu}(1+l)^{2N}|c_{l\mu}|^2<\infty
\quad\text{for every finite }N,
\] with finite bounds in the corresponding angular norms of \(f\).

The denominator in the profile definition cannot vanish. Integration by
parts gives \[
D_l=-(\lambda_l+1)
\int_1^\infty\varrho^{\lambda_l}\chi_l'(\varrho)\,\mathrm d\varrho.
\] Writing \(b_l=1+1/l\), another integration by parts and
\(0\le\chi_l\le1\) give \[
\int_1^{b_l}\varrho^{\lambda_l}\chi_l'(\varrho)\,\mathrm d\varrho
=b_l^{\lambda_l}
-\lambda_l\int_1^{b_l}
\varrho^{\lambda_l-1}\chi_l(\varrho)\,\mathrm d\varrho
\ge1.
\] Thus \(D_l\le-(\lambda_l+1)\). On the supporting collar,
\(\varrho^{\lambda_l+1}\) is bounded independently of \(l\), so
\(|A_l|\le C\|\chi'\|_{L^1}\). Consequently \[
|\widetilde c_{l\mu}|
\le C(1+l)^{-1}|c_{l\mu}|.
\] Each profile is supported in \([1+(10l)^{-1},1+l^{-1}]\). Every
radial derivative of \(\varrho^{\lambda_l-1}\) costs a finite power of
\(1+l\) on this collar, as does each derivative of \(\chi_l\). For
example, after factoring out \(c_{l\mu}\), the \(j\)th derivative of the
combined profile is bounded by \(C_j(1+l)^{j+1}\). The radial factor
\(\varrho^{-1}\) in the harmonic normalization preserves this property.
Lemma B.7 and the rapid coefficient estimates show that \(\zeta\) has a
smooth compactly supported zero extension.

The two moment identities can be checked directly:

\phantomsection\label{eq:online-extension-scalar-exact-moments} \[
\begin{aligned}
\int_1^\infty\varrho^{1-\lambda_l}
\left(\frac12 f_{l\mu}(\varrho)+\zeta_{l\mu}(\varrho)\right)
\,\mathrm d\varrho&=0,\\
\int_1^\infty\varrho^2\zeta_{l\mu}(\varrho)\,\mathrm d\varrho&=0.
\end{aligned}
\] For the first, \(\int\chi_l'=1\) and \(\int\chi_l''=0\), so the
contribution of \(\zeta_{l\mu}\) is \(c_{l\mu}\). The definition of
\(c_{l\mu}\) cancels the residual term. For the second, the integral is
\(c_{l\mu}A_l-\widetilde c_{l\mu}D_l\). This is zero by definition. The
second identity is the cancellation required for the tangent variable.
In unit-sphere components, \[
\int_1^\varrho s^3\zeta(s,\theta)\,\mathrm ds
=\sum_{l,\mu}
\left(\int_1^\varrho s^2\zeta_{l\mu}(s)\,\mathrm ds\right)
Y_{l\mu}(1,\theta).
\] Each coefficient on the right vanishes before and after its original
collar, by the second moment identity. Its derivatives have polynomial
harmonic-degree bounds. The scalar constant mode is absent. The unique
curl-free solution of the spherical divergence equation is therefore the
spherical gradient of the mean-zero Poisson solution. On a nonconstant
mode its inverse involves the factor \(1/[l(l+1)]\) and fixed powers of
\(\varrho\), so it preserves polynomial bounds on the fixed annulus.
Lemma B.7 applied to these primitives proves the asserted smoothness and
support for \(\xi\), hence for \(\beta=w_{\varepsilon_\gamma}\xi\).

For any requested spatial order, the profile estimate uses a finite
coefficient order, that coefficient order uses finitely many rotation
derivatives of \(f\), and the differential formula for \(f\) uses
finitely many angular base norms of \(z\). This proves the stated finite
dependence and its uniform family version. The collars shrink toward the
inner sphere as \(l\) increases; convergence of their zero extensions is
what proves smoothness at that sphere. \hfill\(\square\)

\pandocstatementheading{\leavevmode\phantomsection\label{prop:online-extension-scalar-one-derivative-gain}Proposition B.9 (Scalar regularity of the same base solution).}

Let \((h,K)\) be the initialized base fixed point for smooth supplied
prefields, and write \(h=\mathfrak V_{\widehat h}(\phi,\beta)\). Retain
the fixed base positivity margins and all angular estimates of Theorem
B.6. If, for an integer \(m\ge S_0\),

\phantomsection\label{eq:online-extension-scalar-induction-hypothesis}
\[
\phi\in\overline H^m_\gamma,
\qquad K\in H^{m-1}_{\gamma-1},
\] then \(\phi\in\overline H^{m+1}_\gamma\) and
\(h-\delta\in H^{m+1}_\gamma\). Their finite bounds depend on the
displayed norms, finitely many supplied premetric and angular norms, and
the fixed base margins. If the hypothesis holds after every finite
angular word, the conclusion holds after every finite angular word. For
each specified word only finitely many preceding norms are needed. These
conclusions hold uniformly on families with the indicated uniform
bounds, without another smallness condition.

\emph{Proof.} Lemma B.8 makes \(\beta\) smooth at every spatial order
for this very solution. We give the elliptic equation and its estimate,
including the behavior at the inner sphere and at infinity.

Write \[
\widehat h=a^2\,\mathrm d\varrho^2+
\widehat\gamma_{AB}
(b^A\mathrm d\varrho+\mathrm d\theta^A)
(b^B\mathrm d\varrho+\mathrm d\theta^B),
\qquad
N'=a^{-1}(\partial_\varrho-b-\beta).
\] For the following calculation use the positive second fundamental
form of the radial leaves, namely \(\tfrac12\mathcal L_{N'}\) of their
metric. Set \[
T_0=\frac1{2a}
\bigl(\partial_\varrho\widehat\gamma
-\mathcal L_{b+\beta}\widehat\gamma\bigr),
\qquad H_0=\operatorname{tr}_{\widehat\gamma}T_0,
\qquad Q_0=|T_0|_{\widehat\gamma}^2.
\] These fields depend only on \(\widehat h\) and \(\beta\). The leaf
metric of \(h\) is \(q=e^{2\phi}\widehat\gamma\), with second
fundamental form \(e^{2\phi}(T_0+(N'\phi)\widehat\gamma)\). Its mean
curvature and squared norm are \[
H=H_0+2N'\phi,
\qquad
|T|_q^2=Q_0+2H_0N'\phi+2(N'\phi)^2.
\] The contracted Gauss identity in this convention is \[
R(h)=R(q)-2N'H-H^2-|T|_q^2-2a^{-1}\Delta_q a.
\] Substituting the two-dimensional conformal identities
\(R(q)=e^{-2\phi}(R(\widehat\gamma)-2\Delta_{\widehat\gamma}\phi)\) and
\(\Delta_q a=e^{-2\phi}\Delta_{\widehat\gamma}a\) gives

\phantomsection\label{eq:online-extension-scalar-complete-equation} \[
\begin{aligned}
R(h)={}&-4N'(N'\phi)
-2e^{-2\phi}\Delta_{\widehat\gamma}\phi
+\mathcal Q(\phi,\mathrm d\phi;\widehat h,\beta),\\
\mathcal Q={}&
e^{-2\phi}\left(R(\widehat\gamma)
-2a^{-1}\Delta_{\widehat\gamma}a\right)
-2N'H_0-H_0^2-Q_0\\
&\quad-6H_0N'\phi-6(N'\phi)^2.
\end{aligned}
\] In particular the second-order part is

\phantomsection\label{eq:online-extension-scalar-principal-part} \[
-4N'(N'\phi)-2e^{-2\phi}\Delta_{\widehat\gamma}\phi.
\] This agrees with the scalar variation formula in
\hyperref[ref-Czimek2018]{Czimek} (2018, Lemma 5.5). The displayed
expression for \(\mathcal Q\) shows that every remaining term contains
at most one derivative of \(\phi\).

After changing the sign of the second-order operator, its principal form
is

\phantomsection\label{eq:online-extension-scalar-positive-form} \[
4\,\xi(N')^2+
2e^{-2\phi}|\xi|^2_{\widehat\gamma^{-1},\mathrm{tan}}.
\] It is uniformly positive relative to the Euclidean squared norm.
Indeed the tangential term controls the restriction of \(\xi\) to a
leaf, and the nonzero radial component \(a^{-1}\) of \(N'\) then
controls its remaining component. The base bounds give uniform upper and
lower bounds for \(a\), \(e^\phi\), the leaf metric relative to the
Euclidean sphere metric, and the shift. These bounds fix the ellipticity
margin independently of \(m\).

The actual fixed-point equation is \(R(h)=|K|_h^2\). At the induction
order the scalar inhomogeneity and all first-order terms have
\(H^{m-1}_{\gamma-2}\) regularity. We spell out the cancellation needed
for this weighted assertion. At
\((\widehat h,\beta,\phi)=(\delta,0,0)\), \[
R(\widehat\gamma)=\frac2{\varrho^2},
\quad H_0=\frac2\varrho,
\quad Q_0=\frac2{\varrho^2},
\quad N'H_0=-\frac2{\varrho^2},
\quad \Delta_{\widehat\gamma}a=0.
\] Their contribution to \(\mathcal Q\) is zero. Subtract these flat
values before estimating the prescribed terms. Their differences contain
\(\widehat h-\delta\), its derivatives through order two, or the
compactly supported smooth \(\beta\) and its derivatives. They
consequently have the required weight \(\gamma-2\). The term
\((e^{-2\phi}-1)\varrho^{-2}\) has the same weight. Terms linear in
\(\mathrm d\phi\) have coefficients of size \(O(\varrho^{-1})\), plus
faster-decaying prescribed terms. Quadratic terms in \(\mathrm d\phi\)
and in \(K\) have weight \(2\gamma-2\), which is at least the required
decay since \(\gamma<0\). Weighted product and composition estimates,
with \(m-1\ge3\), give the asserted \(H^{m-1}\) regularity. The
individual flat \(\varrho^{-2}\) terms are used only in their cancelling
combination.

On a bounded annulus, write the equation in Cartesian coordinates as

\phantomsection\label{eq:online-extension-scalar-local-coordinate-equation}
\[
A^{ij}(x,\phi)\partial_i\partial_j\phi
=F(x,\phi,\mathrm d\phi,K),
\] where \(A\) has the positive form above; the prescribed arguments
\(\widehat h\) and \(\beta\) are suppressed. Its coefficients have the
known \(H^m\) regularity and a fixed ellipticity margin. The right-hand
side is in \(H^{m-1}\). For a multi-index \(|\alpha|=m-1\),
\(w=\partial^\alpha\phi\) initially belongs to \(H^1\) and satisfies \[
A^{ij}\partial_i\partial_jw
=\partial^\alpha F-
[\partial^\alpha,A^{ij}]\partial_i\partial_j\phi.
\] The right-hand side is in \(L^2\). In the commutator the highest
derivative of \(\phi\) is \(m\), not \(m+1\). The standard product
estimate is bounded by a constant times \[
\|\mathrm dA\|_{L^\infty}\|\phi\|_{H^m}
+\|A\|_{H^{m-1}}\|\mathrm d^2\phi\|_{L^\infty},
\] with the usual lower-order terms on a slightly larger annulus. These
quantities are finite because \(m\ge4\) and \(H^4\) in three dimensions
controls the second derivatives in \(L^\infty\). The coefficients are
Lipschitz. Converting the equation for \(w\) to divergence form puts the
term \((\partial_iA^{ij})\partial_jw\) in \(L^2\). The interior
difference-quotient \(H^2\) estimate therefore applies to this existing
\(H^1\) solution. It gives \(w\in H^2\) on a smaller annulus, hence
\(\phi\in H^{m+1}\) there. The constants depend on norms already
estimated and on the ellipticity margin fixed at base order. No
smallness condition at the new order enters this step.

For an annulus crossing \(\partial B_1\), extend \(\phi\) and \(\beta\)
by zero. The initial zero-extension regularity makes
\(h=\mathfrak V_{\widehat h}(\phi,\beta)\) a whole \(H^{S_0}\) metric
there. Its curvature has no boundary distribution. On the inner side
\(h=\widehat h\) and \(K=\widehat K\), so the same equation is exactly
the supplied local constraint. Thus the scalar equation holds
distributionally across the sphere. The polar coefficients are smooth
away from the origin, and the preceding interior estimate applies across
the sphere. The gain is a gain for the zero extension itself, so it
preserves the required vanishing traces.

For the global weighted estimate, cover the remaining exterior by
overlapping annuli \(\{R<\varrho<2R\}\) with \(R=2^j\) and larger
parents \(\{R/2<\varrho<4R\}\). Rescale \(x=Ry\). On a fixed annulus in
\(y\), the known weighted \(H^m\) norms control the coefficient norms
uniformly in \(R\), because \(\gamma<0\); the rescaled Euclidean polar
coefficients are fixed smooth functions. The ellipticity margin is the
base margin. For \(R>2\) the tangent correction is zero. The local
estimate just proved has the form

\phantomsection\label{eq:online-extension-scalar-annular-estimate} \[
\|\phi(R\cdot)\|_{H^{m+1}(A)}
\le C_m\left(
\|R^2F(R\cdot)\|_{H^{m-1}(A^*)}
+\|\phi(R\cdot)\|_{H^m(A^*)}\right),
\] where \(C_m\) is bounded by the fixed margin and finitely many
already controlled weighted norms. Multiplication by \(R^{-\gamma}\),
squaring, and summation over the finite-overlap annuli gives the
\(H^{m+1}_\gamma\) estimate, using the weighted \(H^{m-1}_{\gamma-2}\)
norm of the inhomogeneity. A bounded annulus handles the inner sphere.
The zero extension remains zero on the ball, so its new regularity is
precisely \(\overline H^{m+1}_\gamma\). The algebraic formula for
\(\mathfrak V_{\widehat h}\) then gives \(h-\delta\in H^{m+1}_\gamma\).

Finally retain angular words. First perform the zero-word improvement
above, and then induct on word length. Applying an ordered word \(I\) to
the scalar equation gives a linear equation for \(\mathcal L_Z^I\phi\)
with the same second-order principal coefficients. All other second
derivatives of angular variations have smaller word length and have
already gained the spatial derivative. The highest angular variations of
\(K\) have the assumed \(H^{m-1}_{\gamma-1}\) regularity; the prescribed
fields and \(\beta\) have every required order. The zeroth-order
linearized coefficient may contain second derivatives of the zero-word
\(\phi\), which are now in \(H^{m-1}\). Products with the known \(H^m\)
angular unknown therefore have the required \(H^{m-1}\) regularity. The
same differentiated \(H^2\) estimate, the equation across the sphere,
and the annular summation give the improved norm for \(I\). Each finite
step uses finitely many lower words. All estimates remain uniform for
families with the stated bounds, and every step concerns the original
initialized solution. \hfill\(\square\)

\subsection{B.6. The normal equations in the momentum
range}\label{b.6.-the-normal-equations-in-the-momentum-range}\label{sec:online-momentum-range-normal-equations}

The momentum equation has three components, whereas a symmetric
tracefree tensor in three dimensions has five. The selected exterior
range supplies the two remaining equations. We derive those equations
before estimating them, since a right inverse with the same mapping
properties need not select this range.

Throughout this section, \(\varrho=|x|\) is the exterior radius. Tensor
components are taken in Euclidean orthonormal radial frames parallel
along each radial ray. A tangent tensor in such a frame is identified
with a tensor on the unit sphere. Write \(D\), \(d\) and \(c\) for the
unit-sphere gradient of a scalar, divergence of a symmetric tracefree
tensor, and scalar curl of a one-form, respectively. Thus a physical
angular derivative contains the factor \(\varrho^{-1}\).

For a Euclidean tracefree symmetric tensor \(u\), put

\phantomsection\label{eq:online-momentum-physical-decomposition} \[
u_{NN}=a_u,\qquad u_{NA}=e_A,\qquad
u_{AB}=T_{AB}-\tfrac12a_u(g_{S^2})_{AB},
\qquad \operatorname{tr}_{S^2}T=0.
\] The round-sphere Hodge decomposition splits one-forms into exact and
coexact parts, denoted by the subscripts \(E\) and \(H\). On harmonic
degrees \(l\ge2\), the map \(d\) is an isomorphism from tracefree
symmetric tensor harmonics to one-form harmonics. Let \(L_E\) and
\(L_H\) be its inverse on the exact and coexact parts, extended by zero
on degrees below two. Both have angular order \(-1\), and \(L_ED\) has
order zero. Write \(T=T_E+T_H\) for the corresponding decomposition of
\(T\). There are no tracefree symmetric tensor harmonics of degrees zero
or one. The operators and projections just defined commute with proper
rotations.

On coexact one-forms define \[
\Delta_H=c^{-1}\Delta_{S^2}c.
\] Here \(c\) is invertible from coexact one-forms to mean-zero scalars.
In degree \(l\), the eigenvalue of \(\Delta_H\) is \(-l(l+1)\). This
definition fixes the convention; \(\Delta_H\) is not the rough Laplacian
on one-forms.

\pandocstatementheading{\leavevmode\phantomsection\label{lem:online-momentum-explicit-range-identities}Lemma B.10 (Normal equations of the selected flat momentum range).}

Let \(u=B_1f\) lie in the explicit flat momentum range used to define
the exterior normalization. Suppose \(u\in\overline H^{S_0-1}_{-3/2}\),
with \(S_0\ge4\), after every finite angular Lie derivative. There are
smooth tangent fields \(\zeta_E,\zeta_H\), whose zero extensions are
supported in \(1\le\varrho\le3/2\), such that

\phantomsection\label{eq:online-momentum-electric-range} \[
T_E=-L_E\bigl(\tfrac12Da_u+e_E+\varrho\zeta_E\bigr)
\] and

\phantomsection\label{eq:online-momentum-magnetic-range} \[
\begin{aligned}
\partial_\varrho T_H-\frac2\varrho T_H
={}&-L_H\partial_\varrho e_H
 +\frac1\varrho L_H(\Delta_H+2)e_H\\
&-\varrho L_H\partial_\varrho\zeta_H+L_H\zeta_H.
\end{aligned}
\] Projection to degrees \(l\ge2\) is understood wherever \(L_E\) or
\(L_H\) occurs. Each finite spatial norm of the profiles is bounded by
finitely many angular base norms of \(u\). Both identities hold
distributionally across \(\varrho=1\) after zero extension.

\emph{Proof.} The inverse \(B_1\) is the construction of
\hyperref[ref-Czimek2018]{Czimek} (2018, §4.3), with the fixed
degree-dependent radial cutoffs stipulated above. Its residual is
\(f=\operatorname{div}_\delta u\). Since Euclidean divergence commutes
with rotations, \(f\) has every angular derivative in
\(\overline H^{S_0-2}_{-5/2}\).

First consider the boundary profiles. Use the harmonic normalization of
that construction: the scalar and one-form harmonics on \(S_\varrho\)
have \(L^2(S_\varrho)\) norm one. Write \(f_{N,l\mu}\), \(f_{E,l\mu}\)
and \(f_{H,l\mu}\) for the resulting coefficient functions. Equations
(4.43)--(4.44) and (4.48)--(4.49) of the cited construction define

\phantomsection\label{eq:online-momentum-profile-moments} \[
\begin{aligned}
b_{E,l\mu}
 &=\int_1^\infty \varrho^{-l+1}
       \left(\frac{l}{\sqrt{l(l+1)}}f_{N,l\mu}
                          -f_{E,l\mu}\right)\,\mathrm d\varrho,\\
b_{H,l\mu}
 &=-\int_1^\infty \varrho^{-1-\nu_l}f_{H,l\mu}\,\mathrm d\varrho,
\qquad \nu_l=\sqrt{l(l+1)+4},\qquad l\ge2.
\end{aligned}
\] The corresponding normalized profile coefficients are

\phantomsection\label{eq:online-momentum-profile-coefficients} \[
\begin{aligned}
(\zeta_E)_{l\mu}
 &=b_{E,l\mu}\varrho^{l-1}
       \partial_\varrho\chi(l(\varrho-1)),\\
(\zeta_H)_{l\mu}
 &=b_{H,l\mu}\varrho^{1+\nu_l}
       \partial_\varrho\chi(l(\varrho-1)).
\end{aligned}
\] The fixed smooth transition function \(\chi\) is zero on
\((-\infty,1/10]\) and one on \([1,\infty)\). The normalized harmonics
become \(\varrho^{-1}\) times the unit-sphere harmonics when expressed
in the physical frames used here. Consequently the radial profiles in
physical components contain the powers \(\varrho^{l-2}\) and
\(\varrho^{\nu_l}\), respectively. This conversion does not change the
moment coefficients.

Cauchy--Schwarz gives

\phantomsection\label{eq:online-momentum-moment-bounds} \[
\begin{aligned}
|b_{E,l\mu}|
 &\le(2l-1)^{-1/2}
       \left(\|\varrho f_{N,l\mu}\|_{L^2(1,\infty)}
             +\|\varrho f_{E,l\mu}\|_{L^2(1,\infty)}\right),\\
|b_{H,l\mu}|
 &\le(2\nu_l+3)^{-1/2}
                  \|\varrho f_{H,l\mu}\|_{L^2(1,\infty)}.
\end{aligned}
\] For example, in the first line write
\(\varrho^{-l+1}f_{N,l\mu}=\varrho^{-l}(\varrho f_{N,l\mu})\) and use
\(\int_1^\infty\varrho^{-2l}\,\mathrm d\varrho=(2l-1)^{-1}\). With the
stated harmonic normalization, the square sum of the source norms on the
right is equivalent to the critical weighted \(L^2_{-5/2}\) norm of
\(f\). Applying powers of the rotation Casimir operator supplies every
polynomial harmonic-degree weight. Hence, for every finite \(N\),

\phantomsection\label{eq:online-momentum-rapid-moments} \[
\sum_{l,\mu}(1+l)^{2N}
  \bigl(|b_{E,l\mu}|^2+|b_{H,l\mu}|^2\bigr)<\infty,
\] with a bound in finitely many angular base norms of \(u\).

Each nonzero profile in
\hyperref[eq:online-momentum-profile-coefficients]{the coefficient
formulas} is supported in \[
1+(10l)^{-1}\le\varrho\le1+l^{-1}.
\] On this collar, \(\varrho^{l-2}\) and \(\varrho^{\nu_l}\) are
uniformly bounded. Every specified number of radial derivatives costs a
finite polynomial in \(1+l\); the constant can depend on that number.
The rapidly decreasing moment coefficients therefore give convergence in
every Sobolev norm of the zero-extended harmonic sums. Indeed, any
finite sum vanishes on a neighborhood of the inner sphere, and the
polynomial derivative bounds make the tails converge in each prescribed
norm. Sobolev embedding then proves that the full zero extensions are
smooth. Their supports lie in \(1\le\varrho\le3/2\), since \(l\ge2\).
This also proves the asserted finite dependence of each profile norm.

We next derive the tensor equations. Put
\(\rho=\operatorname{div}_\delta u\) and let \(\sigma\) be the symmetric
curl tensor in the explicit construction. The radial divergence formula
and the electric Hodge equation in \hyperref[ref-Czimek2018]{Czimek}
(2018, equations (4.38), (4.53) and (4.56)) give, in physical
components,

\phantomsection\label{eq:online-momentum-radial-divergence} \[
\begin{aligned}
\rho_N
 &=\partial_\varrho a_u+3\varrho^{-1}a_u
                         +\varrho^{-1}\operatorname{div}_{S^2}e,\\
\rho_{\mathrm{tan}}
 &=\partial_\varrho e+3\varrho^{-1}e
          +\varrho^{-1}dT-\tfrac1{2\varrho}Da_u,\\
\varrho^{-1}dT_E
 &=-\zeta_E-\varrho^{-1}\Pi_{\ge2}\bigl(\tfrac12Da_u+e_E\bigr).
\end{aligned}
\] In the third equation, \(\Pi_{\ge2}\) projects one-forms on the unit
sphere onto harmonic degrees \(l\ge2\). Since \(L_E\Pi_{\ge2}=L_E\),
applying \(L_E\) to that equation gives
\hyperref[eq:online-momentum-electric-range]{the electric identity}.

On the coexact subspace of degree at least two, equations (4.47) and
(4.50) of that construction become

\phantomsection\label{eq:online-momentum-curl-equation} \[
\sigma_{NN}=\varrho^{-1}ce_H,
\qquad
\left(\partial_\varrho^2+\frac3\varrho\partial_\varrho
           +\frac{\Delta_{S^2}-3}{\varrho^2}\right)\sigma_{NN}
 =\partial_\varrho
       \bigl(\varrho^{-1}c(\rho_H+\zeta_H)\bigr).
\] Substitution of the first identity and cancellation of \(c/\varrho\)
give \[
\partial_\varrho^2e_H+\varrho^{-1}\partial_\varrho e_H
 +\varrho^{-2}(\Delta_H-4)e_H
 =\partial_\varrho(\rho_H+\zeta_H)
                         -\varrho^{-1}(\rho_H+\zeta_H).
\] The coexact part of
\hyperref[eq:online-momentum-radial-divergence]{the divergence formula}
is \(\rho_H=\partial_\varrho e_H+3\varrho^{-1}e_H+\varrho^{-1}dT_H\).
Insert it into the last equation. The second normal derivatives of
\(e_H\) cancel, and multiplication by \(\varrho\) leaves \[
d\partial_\varrho T_H-2\varrho^{-1}dT_H
 =-\partial_\varrho e_H+\varrho^{-1}(\Delta_H+2)e_H
                    -\varrho\partial_\varrho\zeta_H+\zeta_H.
\] Application of \(L_H\) proves
\hyperref[eq:online-momentum-magnetic-range]{the magnetic identity}.

These calculations are valid for the selected base solution. The range
formulas and their derivatives are continuous as distributional
operations on its base spaces, so one may first use the defining smooth
approximations of the flat inverse. The zero-extended field \(u\)
belongs to \(H^{S_0-1}\) locally and its angular derivatives have the
same property. Since \(S_0-1\ge3\), the traces needed in the
second-order calculation vanish at \(\varrho=1\). The zero-extended
residual has the corresponding barred regularity and the profiles are
smooth zero extensions. Differentiation therefore produces no boundary
delta terms. The identities consequently hold on an annulus crossing the
inner sphere. \hfill\(\square\)

The stronger-weight normalization preserves the normal principal part.
Write \[
v=w u\in\operatorname{Ran}B_{1,\gamma},
\qquad w=e^{-\varepsilon_\gamma\ell}.
\] The exact range equality established for the conjugated inverse
identifies \(u=w^{-1}v\) with an element of the original flat range. The
electric equation is

\phantomsection\label{eq:online-momentum-weighted-electric} \[
(T_v)_E
 =-L_E\bigl(\tfrac12Da_v+(e_v)_E
                              +\varrho w\zeta_E\bigr).
\] For the magnetic part set

\phantomsection\label{eq:online-momentum-weighted-W} \[
W_H=(T_v)_H+L_H(e_v)_H+\varrho wL_H\zeta_H.
\] Multiplication of the flat identity by \(w\) gives exactly

\phantomsection\label{eq:online-momentum-weighted-magnetic} \[
\partial_\varrho W_H
 -\left(\frac2\varrho+\frac{w'}w\right)W_H
 =\frac1\varrho L_H\Delta_H(e_v)_H.
\] Thus the weight changes only known coefficients of order zero. All
derivatives of \(w'/w\) satisfy
\(|\partial_\varrho^j(w'/w)|\le C_j\varrho^{-1-j}\); it also vanishes
near the inner sphere.

\pandocstatementheading{\leavevmode\phantomsection\label{prop:online-momentum-one-normal-gain}Proposition B.11 (One normal regularity gain in the fixed momentum
range).}

Let \(m\ge S_0\ge4\). Let \(h\) be a positive whole metric satisfying
one fixed sufficiently small \(L^\infty\) bound on \(h-\delta\). Suppose
\(v\in\operatorname{Ran}B_{1,\gamma}\) is a base solution of

\phantomsection\label{eq:online-momentum-variable-metric-equation} \[
\operatorname{div}_h\operatorname{TF}_h v=\rho.
\] Assume that after every finite angular word,

\phantomsection\label{eq:online-momentum-gain-hypotheses} \[
h-\delta\in H^m_\gamma,\qquad
v\in\overline H^{m-1}_{\gamma-1},\qquad
\rho\in\overline H^{m-1}_{\gamma-2}.
\] Then \(v\in\overline H^m_{\gamma-1}\) after every finite angular
word. Each specified finite output norm is bounded by finitely many of
the displayed norms and the fixed base inverse margin. The smallness
threshold is independent of \(m\) and of the angular order. Uniform
finite input bounds give uniform output bounds for families.

\emph{Proof.} Lemma B.10 applies to \(u=w^{-1}v\) and supplies smooth
boundary profiles, with bounds from finitely many angular base norms of
\(v\). Differentiate \hyperref[eq:online-momentum-weighted-electric]{the
weighted electric equation} and expand
\hyperref[eq:online-momentum-weighted-magnetic]{the weighted magnetic
equation}. They parameterize the complete normal derivative as

\phantomsection\label{eq:online-extension-momentum-normal-system} \[
\partial_\varrho v=JU+F_\varrho,
\qquad U=(\partial_\varrho a_v,\partial_\varrho e_v).
\] The fixed unit-sphere operator \(J\) has components

\phantomsection\label{eq:online-momentum-normal-J} \[
\begin{aligned}
(JU)_{NN}&=U_0,\qquad (JU)_{NA}=(U_1)_A,\\
(JU)_{\mathrm{tan},0}
 &=-L_E\bigl(\tfrac12DU_0+(U_1)_E\bigr)-L_H(U_1)_H,\\
(JU)_{AB}
 &=(JU)_{\mathrm{tan},0,AB}-\tfrac12U_0(g_{S^2})_{AB}.
\end{aligned}
\] It is bounded from the \(L^2\) scalar--one-form pair to \(L^2\)
tracefree tensors: \(L_ED\) has order zero and \(L_E,L_H\) have order
\(-1\). It commutes with rotations. No low harmonic is discarded from
\(U\); the zero action of \(L_E,L_H\) below degree two says only that
the tangential tracefree part has no such modes.

The \(NN\) and \(NA\) components of \(F_\varrho\) vanish. Its tangential
tracefree component is explicitly

\phantomsection\label{eq:online-momentum-explicit-remainder} \[
\begin{aligned}
(F_\varrho)_{\mathrm{tan},0}
={}&-L_E\partial_\varrho(\varrho w\zeta_E)
 -\partial_\varrho(\varrho wL_H\zeta_H)\\
&+\left(\frac2\varrho+\frac{w'}w\right)W_H
 +\frac1\varrho L_H\Delta_H(e_v)_H.
\end{aligned}
\] Its nonsmooth-input terms have at most one angular derivative of
\(v\) and no normal derivative of \(v\). All other terms are the already
smooth profiles. Hence, using finitely many extra angular words,
\(F_\varrho\) has radial Sobolev order \(m-1\) at weight \(\gamma-2\).

Write the normal principal coefficient of
\(\operatorname{div}_h\operatorname{TF}_h\) as \(M_h\). In the radial
orthonormal frame its action on a Euclidean tracefree tensor \(t\) is

\phantomsection\label{eq:online-momentum-principal-matrix} \[
(M_h t)_j=h^{Ni}t_{ij}
                   -\tfrac13\delta_{jN}h^{ab}t_{ab}.
\] This follows by differentiating
\(\operatorname{TF}_h t=t-\tfrac13(h^{ab}t_{ab})h\) and retaining the
derivatives of \(t\) in the normal direction. Derivatives of \(h\) and
derivatives of the radial frame belong to the remaining terms. At the
Euclidean metric, \(M_\delta t=(t_{NN},t_{NA})\), because
\(\operatorname{tr}_\delta t=0\). Therefore

\phantomsection\label{eq:online-momentum-fixed-normal-margin} \[
M_\delta J=I,\qquad
\|(M_h-M_\delta)J\|_{L^2(S^2)\to L^2(S^2)}
 \le C\|h-\delta\|_{L^\infty(S^2)}.
\] Choose the base metric neighborhood once so that the last expression
is at most \(1/2\) at every radius. The operator \[
\mathcal A_h=I+(M_h-M_\delta)J
\] then has an inverse on \(L^2(S^2)\) of norm at most two. This is the
only smallness used in the normal regularity argument.

Substitute \hyperref[eq:online-extension-momentum-normal-system]{the
parameterization} into
\hyperref[eq:online-momentum-variable-metric-equation]{the momentum
equation}. It becomes

\phantomsection\label{eq:online-momentum-invertible-normal-equation} \[
\mathcal A_hU=G_\varrho,
\qquad
G_\varrho=\rho-M_hF_\varrho-Q(h,v),
\] where \(Q(h,v)\) contains only the angular derivatives of \(v\), the
radial-frame terms of size \(\varrho^{-1}v\), and products of first
derivatives of \(h\) with \(v\). We verify the required regularity of
this equation on a fixed annulus before summing over annuli.

All-angular \(H^m\) regularity of \(h\) implies \(H^m\) regularity in
the radial variable with values in any prescribed finite sphere Sobolev
space. The analogous statement holds at radial order \(m-1\) for \(v\)
and \(\rho\). Conversely, a finite number of angular words controls each
specified sphere Sobolev norm; the rotation fields span the sphere
tangent bundle and their Casimir is elliptic. Choose a sphere Sobolev
index large enough for the finite products under consideration. The
one-dimensional radial product estimates and the sphere Sobolev algebra
estimates then show that \(G_\varrho\) has radial order \(m-1\) with
every specified angular derivative. The highest radial derivative of
\(h\) appearing here has order \(m\), from the differentiated
first-derivative coefficient in \(Q(h,v)\). It is an \(L^2\) factor
multiplied by lower derivatives of \(v\), which are bounded by the
preceding radial regularity. No derivative of \(v\) of radial order
\(m\) occurs in \(G_\varrho\).

The coefficients of \(M_h-M_\delta\), after any fixed angular word, have
bounded radial derivatives through order \(m-1\) on this annulus.
Indeed, the additional radial derivative of \(h\) gives \(H^1\) control
of its \((m-1)\)st radial derivative, and the one-dimensional Sobolev
embedding gives the bound with values in a sufficiently high sphere
Sobolev space. Smooth dependence of inverse metric coefficients
preserves these bounds on the positive metric neighborhood.

The preceding regularity of \(v\) already gives \(U\in L^2\) locally.
Thus \hyperref[eq:online-momentum-invertible-normal-equation]{the normal
equation} identifies it almost everywhere with
\(\mathcal A_h^{-1}G_\varrho\). To justify differentiating this identity
without assuming the desired derivatives, use difference quotients. For
a radial increment \(t\), put
\(\delta_tU(\varrho)=(U(\varrho+t)-U(\varrho))/t\), and use the same
convention for the other fields. Subtraction of the equation at two
radii gives \[
\mathcal A_h(\varrho+t)\,\delta_tU
 =\delta_tG_\varrho-(\delta_t\mathcal A_h)U.
\] The same inverse bound applies at the shifted radius. For a
one-parameter rotation, pull back the equation and subtract it; \(J\)
commutes with the pullback, so the resulting difference equation has the
same form and the same inverse margin. Repeating this procedure gives,
for every finite mixed radial and rotation derivative \(D^\alpha\), an
equation whose highest derivative is

\phantomsection\label{eq:online-momentum-differentiated-normal-equation}
\[
\mathcal A_hD^\alpha U
 =D^\alpha G_\varrho
  -\sum_{0<\beta\le\alpha}
     C_{\alpha,\beta}
       (D^\beta\mathcal A_h)D^{\alpha-\beta}U.
\] For ordered rotation words, the sum means the corresponding finite
Leibniz sum, with their order retained. Every term on the right uses a
strictly lower total derivative of \(U\). Induction, the coefficient
bounds just proved, and the \(L^2\) inverse bound therefore control all
angular derivatives of \(U\) with radial order at most \(m-1\). Bounded
difference quotients identify the weak derivatives at each step. The
constants involve finitely many coefficient and input norms, while the
highest derivative is always inverted at the same base margin.

The \hyperref[eq:online-extension-momentum-normal-system]{normal system}
now gives radial order \(m-1\) for \(\partial_\varrho v\). Together with
the available angular derivatives, this yields \(H^m\) regularity of
\(v\) on the annulus. The argument also applies to an annulus crossing
\(\varrho=1\): the range equations hold there for the zero extension,
\(\rho\) is a barred field, and \(h\) is the given whole metric.
Consequently it improves the zero extension itself, including its
boundary traces.

For completeness, the weighted summation uses the same margin. Cover the
exterior by overlapping dyadic annuli \(A_R\) with \(R=2^j\), together
with one annulus crossing the inner sphere, and rescale \(x=Ry\). Use
physical tensor components for \(h_R(y)=h(Ry)\) and \(v_R(y)=v(Ry)\).
The rescaled normal equation has right-hand side \(R\rho(Ry)\); \(J\) is
unchanged, and the factors \(\varrho^{-1}\) and \(w'/w\) have uniformly
bounded rescaled derivatives. Weighted norms are equivalent to sums with
factors \(R^{-2(\gamma-1)}\) for \(v_R\) and \(R^{-2(\gamma-2)}\) for
\(\rho(Ry)\). The additional factor \(R\) in the latter equation gives
exactly these matching weights.

Every fixed rescaled coefficient norm of \(h_R-\delta\) is bounded by
the corresponding global weighted norm, since its local bound carries
\(R^\gamma\) and \(\gamma<0\). Product terms thus have uniform finite
annular constants. Summing the squared local estimates proves
\(v\in H^m_{\gamma-1}\) after each retained angular word. Since the
improved field is still the zero extension, it belongs to
\(\overline H^m_{\gamma-1}\). The entire argument applies with uniform
constants on a family having the stipulated finite bounds.
\hfill\(\square\)

\pandocstatementheading{\leavevmode\phantomsection\label{cor:online-momentum-gain-in-relative-fixed-point}Corollary B.12 (Momentum regularity in the relative extension).}

Let \((h,K)\) be the initialized relative-extension fixed point for
smooth supplied prefields. Suppose its all-angular regularity has
reached \(h-\delta\in H^m_\gamma\) and \(K\in H^{m-1}_{\gamma-1}\),
where \(m\ge S_0\). Then \(K\in H^m_{\gamma-1}\) after every finite
angular word, on the same base domain and with finite bounds of the
stated kind.

\emph{Proof.} The actual momentum block has the form

\phantomsection\label{eq:online-momentum-relative-range-equation} \[
K=\operatorname{TF}_h(\widehat a+v),\qquad
v\in\operatorname{Ran}B_{1,\gamma},\qquad
\operatorname{div}_h\operatorname{TF}_h v
              =-\operatorname{div}_h\operatorname{TF}_h\widehat a,
\] where \(\widehat a\) is the fixed smooth Euclidean tracefree
premomentum. On Euclidean tracefree tensors, the inverse of
\(\operatorname{TF}_h\) is explicit. If \(\operatorname{tr}_hK=0\), it
is

\phantomsection\label{eq:online-momentum-algebraic-inverse} \[
t=K-\frac{\operatorname{tr}_\delta K}
                  {\operatorname{tr}_\delta h}\,h,
\qquad \operatorname{tr}_\delta t=0,
\qquad \operatorname{TF}_h t=K.
\] Thus \(v=t-\widehat a\) has the preceding spatial and angular
regularity of \(K\). The denominator has a fixed positive lower bound in
the metric neighborhood. The residual
\(\rho=-\operatorname{div}_h\operatorname{TF}_h\widehat a\) has
all-angular \(H^{m-1}_{\gamma-2}\) regularity: it contains at most one
derivative of \(h\), and \(\widehat a\) has every required finite
weighted norm. On the inner ball \(h=\widehat h\) and
\(\operatorname{TF}_{\widehat h}\widehat a=\widehat K\), so the supplied
local momentum constraint makes \(\rho\) zero there. Since it is already
a whole \(H^{m-1}_{\gamma-2}\) field, its exterior restriction belongs
to the barred space at that order. Apply Proposition B.11 to this
existing \(v\). The formula for \(K\) then gives its additional
derivative and finite bound. No new fixed point is selected.
\hfill\(\square\)

\subsection{B.7. Parameter regularity on the same
branch}\label{b.7.-parameter-regularity-on-the-same-branch}\label{sec:online-parameter-promotion}

The spatial estimates concern one selected extension at all finite
orders. To obtain smooth dependence on passive parameters in each of
those spaces, it is enough to combine their locally uniform bounds with
the base parameter theorem. The following interpolation argument
supplies the precise implication and the finite dependence of its
estimates.

\pandocstatementheading{\leavevmode\phantomsection\label{lem:online-parameter-promotion}Lemma B.13 (Parameter regularity from uniform higher spatial bounds).}

Let \(P\) be a finite-dimensional open parameter domain, let \(s_0\) be
the fixed base order, and put

\phantomsection\label{eq:online-parameter-spatial-scale} \[
X_s=H^s_\gamma\oplus H^{s-1}_{\gamma-1},\qquad s\ge s_0.
\] Suppose \(F:P\to X_{s_0}\) is smooth. Suppose also that for every
finite integer \(M\ge s_0\) and every compact \(\mathcal K'\Subset P\),

\phantomsection\label{eq:online-parameter-higher-bounds} \[
F(p)\in X_M,\qquad
\sup_{p\in\mathcal K'}\|F(p)\|_{X_M}<\infty.
\] Then \(F:P\to X_s\) is smooth for every finite \(s\ge s_0\). On
\(\mathcal K\Subset\operatorname{int}\mathcal K'\), every specified
finite number of parameter derivatives in \(X_s\) is bounded by one
sufficiently high spatial bound in
\hyperref[eq:online-parameter-higher-bounds]{the hypothesis}, finitely
many base-space parameter derivative bounds on \(\mathcal K'\), and the
parameter-chart and cutoff margins. The assertion applies to fixed
finite products and to tensor components.

\emph{Proof.} We first record the interpolation inequality in the
precise weighted scale. Choose a smooth dyadic partition of unity in the
spatial variable, including one cutoff near the origin. Rescale the
\(j\)th exterior annulus by \(x=2^jy\). For each fixed weight
\(\omega\), the squared weighted norm of a field \(u\) is equivalent to
\[
\sum_{j\ge0}2^{-2j\omega}
  \bigl\|(\chi_j u)(2^j\,\cdot)\bigr\|_{H^s}^2.
\] The same annular factor occurs at every derivative order because the
derivative weight in the original norm compensates the rescaling. Extend
each cut-off rescaled field by zero to a fixed bounded region and use
ordinary Sobolev interpolation there. Hölder's inequality in the annular
sum gives \[
\|u\|_{H^s_\omega}
 \le C\|u\|_{H^{s_0}_\omega}^{1-\theta}
          \|u\|_{H^M_\omega}^{\theta},
\qquad
\theta=\frac{s-s_0}{M-s_0},\qquad s_0<s<M.
\] Apply this to both components, with orders shifted by one in the
second component, and apply the finite-dimensional version of Hölder's
inequality to their sum. The result is

\phantomsection\label{eq:online-parameter-weighted-interpolation} \[
\|u\|_{X_s}
 \le C\|u\|_{X_{s_0}}^{1-\theta}\|u\|_{X_M}^{\theta}.
\] All orders here are fixed and finite. The constants need not be
uniform as \(M\) increases. Noninteger intermediate orders may
equivalently be defined by the same rescaled Sobolev construction.

Work in one parameter chart, write its coordinates as
\(p\in\mathbb R^d\), and multiply \(F\) by a smooth parameter cutoff
supported in \(\operatorname{int}\mathcal K'\) and equal to one on a
neighborhood of \(\mathcal K\). Extend this product by zero outside the
chart. We continue to write \(F\) for the product. A finite number of
cutoff derivatives enters the bounds below.

The parameter convolutions used next are well defined in each \(X_M\),
even before high-space continuity has been proved. One direct
verification uses the Hilbert structure. In the integer weighted Sobolev
norm, choose a countable orthonormal basis by Gram--Schmidt from a dense
family of smooth compactly supported fields. For a fixed such basis
element \(e\), the functional \(u\mapsto\langle u,e\rangle_{X_M}\) is
continuous in the base distributional topology, by integration by parts
against the smooth compactly supported coefficients. Its value on
\(F(p)\) is therefore measurable. The finite basis expansions of
\(F(p)\) are measurable \(X_M\)-valued maps and converge pointwise in
\(X_M\). This proves strong measurability. The local high-space bound,
together with the parameter cutoff, then gives Bochner integrability.
These integrals agree with their base-space versions because the
inclusion \(X_M\hookrightarrow X_{s_0}\) is continuous and injective.

Fix the desired derivative order \(q\). Choose a finite integer \(M>s\)
and a finite integer \(L>q\) such that

\phantomsection\label{eq:online-parameter-finite-order-choice} \[
L(1-\theta)>q,
\qquad \theta=\frac{s-s_0}{M-s_0}.
\] For example, one may take an integer \(M>2s-s_0\) and \(L=2(q+1)\).
Choose a compactly supported smooth signed kernel \(\kappa\) of integral
one, with all moments of total degrees \(1,\ldots,L-1\) equal to zero.
Such a kernel is obtained from a finite linear combination of distinct
dilations of any fixed integral-one smooth kernel: solve the finite
Vandermonde system for their coefficients so that the dilation factors
cancel each positive degree below \(L\). Put \[
\kappa_t(p)=t^{-d}\kappa(p/t),\qquad F_t=\kappa_t*F.
\] For \(0<t\le t_0\), where \(t_0\) is chosen below the parameter
cutoff margin, \(F_t\) is smooth with values in every fixed \(X_M\).
Differentiating the kernels gives, for \(|\alpha|\le q\),

\phantomsection\label{eq:online-parameter-high-mollifier-bound} \[
\|\partial_p^\alpha(F_{t/2}-F_t)\|_{X_M}
 \le C t^{-|\alpha|}
                \sup_{\mathcal K'}\|F\|_{X_M}.
\] Taylor's formula for the base-space-valued map, with remainder of
order \(L\), and the vanishing moments give

\phantomsection\label{eq:online-parameter-base-mollifier-bound} \[
\|\partial_p^\alpha(F_{t/2}-F_t)\|_{X_{s_0}}
 \le C t^{L-|\alpha|}
       \sum_{j=0}^L\sup_{\mathcal K'}\|D_p^jF\|_{X_{s_0}}.
\] One can obtain this bound directly by differentiating the difference
kernel and subtracting the degree-\((L-1)\) Taylor polynomial of \(F\).
Integration by parts shows that every polynomial term has zero integral;
the differentiated kernel contributes \(t^{-|\alpha|}\) and the Taylor
remainder contributes \(t^L\). Thus no high-space parameter derivative
of \(F\) is used.

Interpolation of \hyperref[eq:online-parameter-high-mollifier-bound]{the
high-space bound} and
\hyperref[eq:online-parameter-base-mollifier-bound]{the base-space
bound} gives

\phantomsection\label{eq:online-parameter-summable-differences} \[
\begin{aligned}
\|\partial_p^\alpha(F_{t/2}-F_t)\|_{X_s}
 &\le C t^{L(1-\theta)-|\alpha|}
       A_L^{1-\theta}B_M^\theta,\\
A_L&=\sum_{j=0}^L\sup_{\mathcal K'}
                              \|D_p^jF\|_{X_{s_0}},\qquad
B_M=\sup_{\mathcal K'}\|F\|_{X_M}.
\end{aligned}
\] The exponent is positive for every \(|\alpha|\le q\). Summation at
\(t=2^{-n}t_0\) therefore proves uniform convergence in \(X_s\) of
\(F_t\) and all its parameter derivatives through order \(q\) on a
neighborhood of \(\mathcal K\) with compact closure in
\(\operatorname{int}\mathcal K'\). The undifferentiated limit is the
given map \(F\), since convergence also holds in \(X_{s_0}\). The
fundamental theorem of calculus along coordinate segments identifies the
limits of the first derivatives with derivatives of \(F\) in \(X_s\);
iteration identifies all derivatives through order \(q\). Their
continuity follows from uniform convergence of continuous maps.

The first mollified term and the summable differences bound these
derivatives using only \(A_L\), \(B_M\), the finitely chosen orders and
the cutoff margin. For example, the bound can be written as
\(C A_L^{1-\theta}B_M^\theta\), with a larger fixed constant accounting
for the initial scale, and hence also as \(C(A_L+B_M)\). Since \(s\) and
\(q\) were arbitrary finite orders, this proves smoothness in every
\(X_s\). A finite parameter-chart cover proves the assertion on a
general finite-dimensional parameter domain. \hfill\(\square\)

Applying Lemma B.13 to the initialized extension now combines its
base-order smooth dependence with its locally uniform higher spatial
bounds. The parameter neighborhood chosen at base order therefore serves
at each finite derivative order.

\subsection{B.8. Proof of the relative-extension
theorem}\label{b.8.-proof-of-the-relative-extension-theorem}\label{sec:online-extension-proof-of-smooth-theorem}

\emph{Proof of Theorem B.1.} Choose the weight offset by Lemma B.2 and
the finite base order \(S_0\) before the higher orders to be estimated.
Choose the base neighborhood in Proposition B.4 small enough to retain
its contraction, the scalar positivity margin of Proposition B.9, and
the normal inverse margin of Proposition B.11. These are finitely many
conditions at the one base order. Sobolev embedding supplies the
required uniform metric and scalar bounds. Rotation-invariant equivalent
norms make the neighborhood invariant under proper rotations.

Proposition B.4 gives the initialized pair
\((h,K)=\mathfrak E(\widehat h,\widehat K)\), its exact protected
agreement, its base bound and its fixed-point property on exact whole
data. Let \(z=(\phi,\beta)\) be its actual scalar correction and \(v\)
its Euclidean tracefree momentum correction. They belong to the
prescribed ranges. Theorem B.6 gives all their angular derivatives in
the base spaces. Lemma B.8 therefore makes \(\beta\) smooth with a
compactly supported zero extension, with each finite norm controlled by
finitely many input norms.

We improve the same pair by induction on the spatial order. At
\(m=S_0\), and after every finite angular word, \[
\phi\in\overline H^{m}_\gamma,
\qquad K\in H^{m-1}_{\gamma-1}.
\] Assume this at one finite \(m\ge S_0\). Proposition B.9 gives
\(\phi\in\overline H^{m+1}_\gamma\) and \(h-\delta\in H^{m+1}_\gamma\),
with every angular word retained. The momentum block is \[
K=\operatorname{TF}_h(\widehat a+v),
\qquad
\operatorname{div}_h\operatorname{TF}_h v
=-\operatorname{div}_h\operatorname{TF}_h\widehat a,
\qquad v\in\operatorname{Ran}B_{1,\gamma}.
\] The algebraic inverse in Corollary B.12 gives the already known
\(H^{m-1}_{\gamma-1}\) regularity of \(v\) after every angular word. The
fixed prefield \(\widehat a\) is smooth at all weighted orders. Together
with the improved metric, it gives the required \(H^{m-1}_{\gamma-2}\)
regularity of the residual. That residual vanishes on \(B_1\) as a whole
Sobolev field, hence belongs to the required barred space. Proposition
B.11 and Corollary B.12 give \(v,K\in H^m_{\gamma-1}\) with every
angular word. This closes the induction at \(m+1\).

Each finite induction uses only finitely many input norms and angular
words. The scalar step uses a positive elliptic principal form. The
momentum step inverts the fixed angular operator at its base \(L^2\)
margin and then differentiates. Both steps use the smallness fixed at
base order. The zero-extension estimates hold across the inner sphere,
so the induction gives all whole weighted orders and the stated
zero-extension regularity of the corrections.

For a supplied finite-dimensional smooth parameter family, Proposition
B.4 first gives smooth dependence in \[
X_{S_0}=H^{S_0}_\gamma\oplus H^{S_0-1}_{\gamma-1}
\] for \(F=(h-\delta,K)\). On a retained compact parameter set, Theorem
B.6 gives locally uniform angular bounds. The finite spatial induction
then gives locally uniform bounds in every \(X_M\). Apply Lemma B.13 to
this same map \(F\). Every specified output derivative on a compact
\(\mathcal K\Subset\operatorname{int}\mathcal K'\) is bounded by one
sufficiently high spatial bound and finitely many base parameter bounds
on \(\mathcal K'\). The preceding construction expresses both in
finitely many supplied input seminorms on that larger set, with the
fixed margins. This proves the parameter assertion in Theorem B.1 with
its stated order of quantifiers.

Spatial induction and parameter differentiation apply to the solution
initialized in Proposition B.4. Hence its equality with the prescribed
data on the protected ball and its fixed-point identity on exact whole
data still hold. This proves all assertions of Theorem B.1.
\hfill\(\square\)

For the supplied extension the initialization and the chosen ranges
determine the normalized correction uniquely. The prescribed interior
data can admit other whole maximal vacuum extensions.

\section{AI and tool use disclosure}\label{ai-and-tool-use-disclosure}\label{sec:ai-and-tool-use-disclosure}

\begingroup\hyphenpenalty=10000\relax
AI systems were used extensively throughout this research and the
preparation of this manuscript. The models used were GPT-6 Pro,
gpt-6.2-astra (high), Claude Fable 5.2 (xhigh), Claude Opus 5.5 (xhigh),
and Gemini 4 Carbon (high). The computational supplement records the
symbolic checks performed with SymPy in the SageMath environment.\par\endgroup

\section{References}\label{references}\label{sec:online-coverage-references}

\leavevmode\phantomsection\label{ref-AdamoMason2014}Tim Adamo, Lionel Mason (2014). Conformal and Einstein gravity from
twistor actions. \emph{Classical and Quantum Gravity} \textbf{31},
045014. \href{https://doi.org/10.1088/0264-9381/31/4/045014}{DOI}.

\leavevmode\phantomsection\label{ref-BarGinouxPfaffle2007}Christian Bär, Nicolas Ginoux, Frank Pfäffle (2007). \emph{Wave
Equations on Lorentzian Manifolds and Quantization}. European
Mathematical Society. \href{https://arxiv.org/pdf/0806.1036}{Author
text}.

\leavevmode\phantomsection\label{ref-BautistaIbortLafuente2014}A. Bautista, A. Ibort, J. Lafuente (2014). On the space of light rays of
a spacetime and a reconstruction theorem by Low. \emph{Classical and
Quantum Gravity} \textbf{31}, 075020.
\href{https://doi.org/10.1088/0264-9381/31/7/075020}{DOI}.

\leavevmode\phantomsection\label{ref-BernalSanchez2005}Antonio N. Bernal, Miguel Sánchez (2005). Smoothness of time functions
and the metric splitting of globally hyperbolic spacetimes.
\emph{Communications in Mathematical Physics} \textbf{257}, 43--50.
\href{https://arxiv.org/pdf/gr-qc/0401112}{Author text}.

\leavevmode\phantomsection\label{ref-BernalSanchez2006}Antonio N. Bernal, Miguel Sánchez (2006). Further results on the
smoothability of Cauchy hypersurfaces and Cauchy time functions.
\emph{Letters in Mathematical Physics} \textbf{77}, 183--197.
\href{https://arxiv.org/pdf/gr-qc/0512095}{Author text}.

\leavevmode\phantomsection\label{ref-ChruscielCong2023}Piotr T. Chruściel, Wan Cong (2023). Gluing variations. \emph{Classical
and Quantum Gravity} \textbf{40}, 165009.
\href{https://doi.org/10.1088/1361-6382/ace494}{Published article}.

\leavevmode\phantomsection\label{ref-Czimek2018}Stefan Czimek (2018). An extension procedure for the constraint
equations. \emph{Annals of PDE} \textbf{4}, article 2.
\href{https://doi.org/10.1007/s40818-017-0039-3}{DOI}.
\href{https://arxiv.org/abs/1609.08814v2}{arXiv version}.

\leavevmode\phantomsection\label{ref-DouglisNirenberg1955}Avron Douglis, Louis Nirenberg (1955). Interior estimates for elliptic
systems of partial differential equations. \emph{Communications on Pure
and Applied Mathematics} \textbf{8}(4), 503--538.
\href{https://doi.org/10.1002/cpa.3160080406}{DOI}.

\leavevmode\phantomsection\label{ref-Duistermaat1972FIO}J. J. Duistermaat (1972). Applications of Fourier integral operators.
\emph{Séminaire Équations aux dérivées partielles (Polytechnique)}
\textbf{1971--1972}, exposé 27, 1--24.
\href{https://www.numdam.org/article/SEDP_1971-1972____A27_0.pdf}{Online
text}.

\leavevmode\phantomsection\label{ref-DuistermaatHormander1972}J. J. Duistermaat, L. Hörmander (1972). Fourier integral operators. II.
\emph{Acta Mathematica} \textbf{128}, 183--269.
\href{https://doi.org/10.1007/BF02392165}{DOI}.

\leavevmode\phantomsection\label{ref-EastwoodPenroseWells1981}Michael G. Eastwood, Roger Penrose, R. O. Wells, Jr.~(1981). Cohomology
and massless fields. \emph{Communications in Mathematical Physics}
\textbf{78}, 305--351. \href{https://doi.org/10.1007/BF01942327}{DOI}.

\leavevmode\phantomsection\label{ref-FriedrichNagy1999}Helmut Friedrich, Gabriel Nagy (1999). The Initial Boundary Value
Problem for Einstein's Vacuum Field Equation. \emph{Communications in
Mathematical Physics} \textbf{201}, 619--655.
\href{https://users.math.msu.edu/users/gnagy/papers/hFgN99.pdf}{Online
text}.

\leavevmode\phantomsection\label{ref-GiaquintaMartinazzi2012}Mariano Giaquinta, Luca Martinazzi (2012). An Introduction to the
Regularity Theory for Elliptic Systems, Harmonic Maps and Minimal
Graphs. Edizioni della Normale. Second edition.
\href{https://doi.org/10.1007/978-88-7642-443-4}{DOI}.

\leavevmode\phantomsection\label{ref-gerard-wrochna2014}Christian Gérard, Michał Wrochna (2014). Construction of Hadamard states
by pseudo-differential calculus. \emph{Communications in Mathematical
Physics} \textbf{325}, 713--755.
\href{https://doi.org/10.1007/s00220-013-1824-9}{DOI}.
\href{https://arxiv.org/abs/1209.2604v4}{arXiv version}.

\leavevmode\phantomsection\label{ref-HughesKatoMarsden1977}Thomas J. R. Hughes, Tosio Kato, Jerrold E. Marsden (1977). Well-posed
quasi-linear second-order hyperbolic systems with applications to
nonlinear elastodynamics and general relativity. \emph{Archive for
Rational Mechanics and Analysis} \textbf{63}(3), 273--294.
\href{https://doi.org/10.1007/BF00251584}{DOI}.
\href{https://www.cds.caltech.edu/~marsden/bib/1977/02-HuKaMa1977/HuKaMa1977.pdf}{Online
text}.

\leavevmode\phantomsection\label{ref-IsthmusLocal2026}Lior Isthmus (2026a). A projective twistor correspondence for local
analytic Einstein germs. Version 1.1.
\href{https://doi.org/10.5281/zenodo.24197319}{DOI}.
\href{https://hexagonmath.org/2611.00123v2}{Online text}.

\leavevmode\phantomsection\label{ref-IsthmusGlobal2026}Lior Isthmus (2026b). Global twistor sources for analytic Einstein
spacetimes. Version 1.0.1.
\href{https://doi.org/10.5281/zenodo.24239028}{DOI}.
\href{https://hexagonmath.org/2611.00425v2}{Online text}.

\leavevmode\phantomsection\label{ref-LeBrun1984}Claude R. LeBrun (1984). Twistor CR manifolds and three-dimensional
conformal geometry. \emph{Transactions of the American Mathematical
Society} \textbf{284}(2), 601--616.
\href{https://doi.org/10.1090/S0002-9947-1984-0743735-9}{DOI}.

\leavevmode\phantomsection\label{ref-LeFlochNguyen2023}Philippe G. LeFloch, The-Cang Nguyen (2023). The seed-to-solution method
for the Einstein constraints and the asymptotic localization problem.
\emph{Journal of Functional Analysis} \textbf{285}, 110106.
\href{https://doi.org/10.1016/j.jfa.2023.110106}{DOI}.

\leavevmode\phantomsection\label{ref-Low2001}Robert J. Low (2001). The space of null geodesics. \emph{Nonlinear
Analysis} \textbf{47}, 3005--3017.
\href{https://www.sciencedirect.com/science/article/abs/pii/S0362546X01004217}{Publisher}.

\leavevmode\phantomsection\label{ref-MaoOhTao2023}Yuchen Mao, Sung-Jin Oh, Zhongkai Tao (2023). Initial data gluing in the
asymptotically flat regime via solution operators with prescribed
support properties.
\href{https://arxiv.org/abs/2308.13031}{arXiv:2308.13031}.

\leavevmode\phantomsection\label{ref-Mason1989}Lionel J. Mason (1989). Hypersurface twistors. \emph{Twistor Newsletter}
\textbf{28}, 20--23.
\href{https://people.maths.ox.ac.uk/lmason/Tn/28/TN28-06.pdf}{Online
text}.

\leavevmode\phantomsection\label{ref-MasonPenrose1989}Lionel J. Mason, Roger Penrose (1989). A twistorial approach to the full
vacuum equations. \emph{Twistor Newsletter} \textbf{29}, 1--5.
\href{https://people.maths.ox.ac.uk/lmason/Tn/29/TN29-02.pdf}{Online
text}.

\leavevmode\phantomsection\label{ref-Melrose2006}Richard Melrose (2006). Lectures on Pseudodifferential operators.
Massachusetts Institute of Technology.
\href{https://math.mit.edu/~rbm/18.157-F05.pdf}{Online text}. MIT
18.157, version 0.7E.

\leavevmode\phantomsection\label{ref-Penrose1976}Roger Penrose (1976). Nonlinear gravitons and curved twistor theory.
\emph{General Relativity and Gravitation} \textbf{7}(1), 31--52.
\href{https://doi.org/10.1007/BF00762011}{DOI}.

\leavevmode\phantomsection\label{ref-Penrose2015}Roger Penrose (2015). Palatial twistor theory and the twistor googly
problem. \emph{Philosophical Transactions of the Royal Society A}
\textbf{373}, 20140237.
\href{https://doi.org/10.1098/rsta.2014.0237}{DOI}.

\leavevmode\phantomsection\label{ref-Radzikowski1996}M. J. Radzikowski (1996). Micro-local approach to the Hadamard condition
in quantum field theory on curved space-time. \emph{Communications in
Mathematical Physics} \textbf{179}, 529--553.
\href{https://doi.org/10.1007/BF02100096}{DOI}.

\leavevmode\phantomsection\label{ref-Sharma2021}Atul Sharma (2021). Twistor action for general relativity.
\href{https://arxiv.org/abs/2104.07031v2}{arXiv:2104.07031v2}.

\leavevmode\phantomsection\label{ref-Sbierski2016}Jan Sbierski (2016). On the existence of a maximal Cauchy development
for the Einstein equations: a dezornification. \emph{Annales Henri
Poincaré} \textbf{17}(2), 301--329.
\href{https://doi.org/10.1007/s00023-015-0401-5}{DOI}.
\href{https://arxiv.org/abs/1309.7591v3}{arXiv version}.

\leavevmode\phantomsection\label{ref-York1974}James W. York, Jr.~(1974). Covariant decompositions of symmetric tensors
in the theory of gravitation. \emph{Annales de l'Institut Henri
Poincaré, Section A} \textbf{21}, 319--332.
\href{https://www.numdam.org/item/AIHPA_1974__21_4_319_0/}{Original
article}.

\end{document}
